\documentclass[11pt]{amsart}
\usepackage[margin=1.1in]{geometry}
\usepackage{amsmath,amssymb,amsthm}
\usepackage{enumitem}
\usepackage{booktabs}
\usepackage[dvipsnames]{xcolor}
\usepackage[colorlinks=true,linkcolor=MidnightBlue,citecolor=MidnightBlue,urlcolor=MidnightBlue]{hyperref}

\newtheorem{theorem}{Theorem}[section]
\newtheorem{proposition}[theorem]{Proposition}
\newtheorem{lemma}[theorem]{Lemma}
\newtheorem{corollary}[theorem]{Corollary}
\theoremstyle{definition}
\newtheorem{definition}[theorem]{Definition}
\newtheorem{example}[theorem]{Example}
\newtheorem{remark}[theorem]{Remark}
\newtheorem{question}[theorem]{Question}
\newtheorem{fact}[theorem]{Fact}
\newtheorem{conjecture}[theorem]{Conjecture}

\newcommand{\R}{\mathcal{R}}
\newcommand{\Lr}{\mathcal{L}}
\newcommand{\D}{\mathcal{D}}
\newcommand{\Hh}{\mathcal{H}}
\renewcommand{\P}{\mathcal{P}}
\newcommand{\PP}{\mathcal{P}}
\newcommand{\Part}{\Pi}
\newcommand{\supp}{\operatorname{supp}}
\newcommand{\T}{\mathsf{T}}
\newcommand{\N}{\mathsf{N}}
\newcommand{\Sym}{\operatorname{Sym}}
\newcommand{\rk}{\operatorname{rank}}
\newcommand{\Z}{\mathbb{Z}}
\newcommand{\Inv}{\operatorname{Inv}}

\title[Embeddability in partition monoids]{Necessary conditions for $*$-embeddability in partition monoids:\\ notes on the projection condition and a new obstruction}
\date{\today}

\begin{document}
\maketitle

\begin{abstract}
These notes concern regular $*$-semigroups that $*$-embed in a partition monoid $\P_X$. We (1) reformulate the known necessary condition ``projections with a common lower bound commute'' as ``for every projection $p$, the set $\{x : px=xp=p\}$ is an inverse semigroup'' and record some of its consequences; (2) show that the local (completely simple) structure is not constrained at all: every finite completely simple regular $*$-semigroup embeds in some $\P_n$, by an explicit construction, and $\P_4$ already contains non-orthodox completely simple $*$-subsemigroups; (3) prove a family of new necessary conditions indexed by digraphs, the simplest being the \emph{triangle condition} and the \emph{hexagon condition}: quasi-identities in six projection variables which hold in every $\P_X$, are equivalent for finite projection algebras, and fail in a $61$-element regular $*$-semigroup that satisfies the projection condition; the witness is a quotient of $\P_3$, and the obstruction is visible in the projection algebra alone; (4) show that the triangle condition is a \emph{master condition}: in every projection algebra it implies all the other digraph conditions, which therefore hold in every $\P_X$; (5) show that the projection condition and the triangle condition are nevertheless \emph{not sufficient}: a $20$-element regular $*$-semigroup (a Rees matrix semigroup over the trivial group on the path $a-c-b-d$, with its zero split into a $2\times2$ rectangular band) satisfies both, and all the digraph conditions, but embeds in no $\P_X$; the obstruction is a new quasi-identity of ``wiring'' type, valid in every $\P_X$, which mixes matched pairs and commuting pairs, and, from a census of the projection algebras of order $7$, a fourth quasi-identity, closing a matched path $a\bowtie c\bowtie b$ to $a\bowtie b$ whenever some projection commutes with $a$ and $c$ and sees $b$ as it sees $c$, with a $19$-element counterexample to the sufficiency of the other three, and, from a census of order $8$ (by one-point extensions, which makes orders $8$ and $9$ accessible), seven further candidate obstructions of order $8$ (projection algebras with no embedding into $\P(\P_m)$ for $m\le8$) and quasi-identities of three new shapes, valid in $\P(\P_m)$ for $m\le7$, that they violate, six of which, the four-variable condition $bcb=b,\ ata\le b,\ ctc\le a\Rightarrow ata=btb$, the path-closing condition (R$'$) with two auxiliary projections, which subsumes the fourth quasi-identity, the path-closing conditions (P$_1$) and (Q) of length $3$, the latter subsuming the wiring condition, and the path-closing conditions K4 and K0 with two auxiliaries, which yield (E$_6$) and (E$_1$), we prove in every $\P_X$, so that all seven witnesses of order $8$ are minimal non-embeddable projection algebras, and a tenth, of order $9$, follows; for the family of graph semigroups $Q(G,\bot)$ to which the example belongs, and for a layered generalisation of it, the projection condition and the wiring condition are also sufficient, and they say that $\bot$ is a proper colouring whose bichromatic subgraphs are disjoint unions of complete bipartite graphs; and, as the main new result of these notes, we prove an infinite family of \emph{chain conditions} (W$_n$), valid in every $\P_X$, by an exact counting invariant of the identification graph, exhibit minimal non-embeddable projection algebras of order $2n+5$ violating them, and deduce that the quasivariety generated by the $\P(\P_X)$ is not finitely based (Section~\ref{subsec:count}); (6) report computations: every regular $*$-semigroup of order $\le 8$ satisfying the projection condition embeds in a partition monoid (in $\P_5$ if it is not inverse), and every projection algebra of order $\le5$ satisfying the projection and triangle conditions embeds in $\P(\P_4)$; and (7) make some general remarks (for finite $S$ the quasi-identities of $\{\P_X\}$ decide embeddability; $\P_X$ embeds in the involuted semigroup of binary relations on $2^X$).
\end{abstract}

\tableofcontents

\section{Setting and the basic facts about $\P_X$}\label{sec:setup}

Throughout, a \emph{regular $*$-semigroup} is a semigroup with a unary operation satisfying $x^{**}=x$, $(xy)^*=y^*x^*$, $xx^*x=x$; its \emph{projections} are the elements $p=p^2=p^*$, and $\theta_p(q)=pqp$ is the projection algebra operation. For projections $p,q$ we write $p\le q$ for the natural order ($p=pq=qp$, equivalently $p=qpq$). The set of projections of $S$ is $\P(S)$.

\subsection{Partitions}
For a set $X$ let $X'=\{x':x\in X\}$ be a disjoint copy, and let $\P_X$ be the partition monoid on $X$: its elements are set partitions $\alpha$ of $X\sqcup X'$, the product $\alpha\beta$ is obtained by identifying the lower points of $\alpha$ with the upper points of $\beta$, taking connected components, and forgetting the middle layer; $\alpha^*$ interchanges $X$ and $X'$. This is a regular $*$-semigroup.

A block of $\alpha$ is \emph{transversal} if it meets both $X$ and $X'$, and \emph{upper} (resp.\ \emph{lower}) \emph{nontransversal} if it lies in $X$ (resp.\ $X'$). The \emph{rank} of $\alpha$ is the number of transversal blocks. The \emph{upper partition} $\ker\alpha$ of $\alpha$ is the partition of $X$ into the upper parts of the blocks of $\alpha$; dually the lower partition $\operatorname{coker}\alpha$.

A projection $p$ of $\P_X$ is determined by the pair $(\varepsilon_p,\T(p))$, where $\varepsilon_p$ is a partition of $X$ and $\T(p)$ is a set of $\varepsilon_p$-classes: the blocks of $p$ are $A\cup A'$ for $A\in\T(p)$ (the \emph{$\T$-blocks}) and $B$, $B'$ for the remaining classes $B$ (the \emph{$\N$-blocks}); we write $\N(p)$ for the set of these classes and
\[
\supp(p)=\textstyle\bigcup\T(p)\subseteq X .
\]
Thus $\rk(p)=|\T(p)|$, and $p$ has rank $0$ iff $\T(p)=\emptyset$.

\subsection{Products}
Let $\alpha,\beta\in\P_X$. The classes of the join $\operatorname{coker}\alpha\vee\ker\beta$ (a partition of the middle copy of $X$) are the connected components of the middle layer; for such a class $K$ let $U_K\subseteq X$ be the union of the upper parts of those transversal blocks of $\alpha$ whose lower parts lie in $K$, and let $L_K\subseteq X'$ be the union of the lower parts of those transversal blocks of $\beta$ whose upper parts lie in $K$. Then the blocks of $\alpha\beta$ are
\begin{itemize}[nosep]
\item $U_K\cup L_K$ for the classes $K$ with $U_K\ne\emptyset\ne L_K$ (the transversal blocks);
\item the upper nontransversal blocks of $\alpha$, and the sets $U_K$ with $U_K\ne\emptyset=L_K$ (upper nontransversal);
\item the lower nontransversal blocks of $\beta$, and the sets $L_K$ with $L_K\ne\emptyset=U_K$ (lower nontransversal).
\end{itemize}
We shall use this only for projections, in the following forms. For projections $p,q$ the middle join is $\varepsilon_p\vee\varepsilon_q$, and $U_K$ (resp.\ $L_K$) is the union of the $\T(p)$-blocks (resp.\ $\T(q)$-blocks) contained in $K$.

\begin{lemma}\label{lem:PX}
Let $p,q$ be projections of $\P_X$.
\begin{enumerate}[label=\textup{(\alph*)}]
\item The blocks of $\theta_p(q)=pqp$ are: all $\N(p)$-blocks, and the sets $U_K$ for the classes $K$ of $\varepsilon_p\vee\varepsilon_q$ meeting $\supp(p)$; the latter is a $\T$-block of $\theta_p(q)$ iff $K$ meets $\supp(q)$. In particular the blocks of $\theta_p(q)$ not contained in $\supp(p)$ are exactly the $\N(p)$-blocks. Consequently, if $e\le p$ then $\supp(e)\subseteq\supp(p)$ and every block of $e$ not contained in $\supp(p)$ is an $\N(p)$-block.
\item Suppose $pq=qp$. Then every class $K$ of $\varepsilon_p\vee\varepsilon_q$ is of one of the following types: \textup{(i)} $K=U_K=L_K$; \textup{(ii)} $K=U_K$ is a single $\N(q)$-block; \textup{(iii)} $K=L_K$ is a single $\N(p)$-block; \textup{(iv)} $K$ is a single block of both $p$ and $q$. The $\T$-blocks of $pq$ are the classes of type \textup{(i)}; hence $\supp(p)\cap\supp(q)=\supp(pq)$, and every $\N(q)$-block contained in $\supp(p)$ is a union of $\T(p)$-blocks.
\item Suppose $pqp=p$ and $qpq=q$. Then every class of $\varepsilon_p\vee\varepsilon_q$ meeting $\supp(p)\cup\supp(q)$ contains exactly one $\T(p)$-block and exactly one $\T(q)$-block. In particular $\rk p=\rk q$, and no $\T(p)$-block is an $\N(q)$-block.
\end{enumerate}
\end{lemma}

\begin{proof}
(a) $\theta_p(q)=(pq)(pq)^*$ is the projection whose $\T$-blocks are the upper parts of the transversal blocks of $pq$ and whose $\N$-blocks are the upper nontransversal blocks of $pq$; now read off the description of $pq$. For the last sentence use $e=\theta_p(e)$.

(b) Since $pq$ is symmetric, $U_K\cup L_K$ is a block only if $U_K=L_K$ (when both are nonempty), and the family of upper nontransversal blocks $\N(p)\cup\{U_K:L_K=\emptyset\ne U_K\}$ equals the family of lower ones $\N(q)\cup\{L_K:U_K=\emptyset\ne L_K\}$. Let $K$ be a class. If $U_K\ne\emptyset\ne L_K$, an $\N(p)$-block inside $K$ would be an $\N$-block of $pq$, hence (as $U_K\neq\emptyset$) an $\N(q)$-block; but a set that is a block of both $p$ and $q$ is a class by itself, so $K$ contains no $\N(p)$-blocks and dually no $\N(q)$-blocks: $K=U_K=L_K$, type (i). If $U_K\ne\emptyset=L_K$ then $U_K$ is an $\N(q)$-block; $K$ is a union of $q$-blocks containing the $q$-block $U_K$ and $p$-blocks; as before the remaining $p$-blocks in $K$ would be common blocks, so $K=U_K$, type (ii). Type (iii) is dual, and if $U_K=L_K=\emptyset$ then $K$ consists of $\N(p)$- and $\N(q)$-blocks, which must coincide as above, type (iv). Finally, an $\N(q)$-block $B\subseteq\supp(p)$ lies in a class $K$; $K$ is not of type (i) or (iii) (these consist of $\T(q)$-blocks only), nor of type (iv) ($B$ meets $\supp(p)$), so $K=B$ is of type (ii).

(c) $pqp=p$ says that the upper projection of $pq$ is $p$: its $\T$-blocks $\{U_K: U_K\ne\emptyset\ne L_K\}$ are exactly the $\T(p)$-blocks, and its $\N$-blocks $\N(p)\cup\{U_K:U_K\ne\emptyset=L_K\}$ are exactly the $\N(p)$-blocks. So every class meeting $\supp(p)$ contains exactly one $\T(p)$-block and at least one $\T(q)$-block; by symmetry, exactly one of each. A $\T(p)$-block which is an $\N(q)$-block would be a class by itself, containing no $\T(q)$-block.
\end{proof}

We say that projections $p,q$ form a \emph{matched pair}, $p\bowtie q$, if $pqp=p$ and $qpq=q$; by (c) this means that $p$ and $q$ have the same rank and their product has that rank (in the trace of the $\D$-class, $pq\ne0$), and it gives a bijection $\T(p)\to\T(q)$ (``the block in the same class''). This is exactly the \emph{friendly} relation $p\,\mathcal F\,q$ of East, Gray, Azeef Muhammed and Ru\v skuc (2025), defined by $pqp=p$ and $qpq=q$; we keep the symbol $\bowtie$ and the word ``matched'' to emphasise the bijection of $\T$-blocks, but the reader may read $\bowtie$ as $\mathcal F$ throughout. Note that $\mathcal F$ is reflexive and symmetric but not transitive, even in $\P(\P_X)$: in $\P_2$ the projection $p$ with the single $\T$-block $\{1,2\}$ is matched with $q=(\{1\}\ \T,\ \{2\}\ \N)$ and with $r=(\{2\}\ \T,\ \{1\}\ \N)$, but $qrq$ has rank $0$ (a $\T$-block and its partner need not meet, they need only lie in the same class of the join).

\section{The projection condition}\label{sec:Q1}

The known necessary condition is
\begin{equation}\label{eq:Q1}\tag{Q1}
r=prp\ \&\ r=qrq\ \Longrightarrow\ pq=qp\qquad(p,q,r\ \text{projections}),
\end{equation}
equivalently: an idempotent lying above a projection in the natural order is a projection. We record some reformulations and consequences.

\begin{proposition}\label{prop:Q1}
For a regular $*$-semigroup $S$ the following are equivalent.
\begin{enumerate}[label=\textup{(\arabic*)}]
\item $S$ satisfies \eqref{eq:Q1}.
\item For every projection $p$, the set $S_p=\{x\in S: px=xp=p\}$ (the elements above $p$ in the natural order) is an inverse semigroup.
\item Two projections commute if and only if they have a common lower bound in $\P(S)$.
\item The projection algebra satisfies the quasi-identity $\theta_p(r)=r=\theta_q(r)\Rightarrow\theta_p(q)=\theta_q(p)$.
\end{enumerate}
\end{proposition}

\begin{proof}
$S_p$ is closed under multiplication and under $*$ (if $px=xp=p$ then $px^*=(xp)^*=p$), so it is a regular $*$-semigroup with zero $p$; it is inverse iff all its idempotents are projections, and its idempotents are exactly the idempotents $e\ge p$. This gives (1)$\Leftrightarrow$(2). For (3): if $pq=qp$ then $pq$ is a projection below $p$ and $q$; the converse is \eqref{eq:Q1}. For (4): for projections $p,q$, $\theta_p(q)=\theta_q(p)$ iff $pq=qp$ (if $pqp=qpq=:r$ then $pq=(pq)^2=pr=r$ and $qp=rp=r$).
\end{proof}

\begin{remark}[The condition in $\P_X$]\label{rem:Sp}
In $\P_X$, an element $x$ lies above a projection $p$ iff $x$ maps every $\T(p)$-block onto itself by a ``block bijection'' (a partition of $A\sqcup A'$ into transversal blocks) and permutes the $\N(p)$-blocks among themselves in the same way, so that $(\P_X)_p\cong\prod_{A\in\T(p)}\mathcal I^*_A\times\mathcal J_p$, with $\mathcal I^*_A$ the dual symmetric inverse monoid on $A$ and $\mathcal J_p$ an inverse monoid of partial block bijections between $\N(p)$-blocks. For $p$ the projection with all blocks singletons and nontransversal, $(\P_X)_p$ is the symmetric inverse monoid $\mathcal I_X$. So \eqref{eq:Q1} is exactly the statement that these ``unitary'' subsemigroups are inverse, which is why it resembles $E$-unitarity.
\end{remark}

Some consequences of \eqref{eq:Q1} which may be useful in the search for further conditions:
\begin{itemize}[nosep]
\item For every projection $r$, the projections $\ge r$ pairwise commute; so every principal up-set of $\P(S)$ is a semilattice under multiplication.
\item Two distinct projections in the same $\D$-class with nonzero product ($p\bowtie q$, $p\ne q$) have no common lower bound (they would commute and then $p=pqp=pq=qp=q$). In particular a completely $0$-simple $*$-subsemigroup of $S$ that contains its own zero is a $0$-direct union of groups (a Brandt semigroup), and if $S$ has a zero then \eqref{eq:Q1} forces $S$ to be inverse.
\item If $x\ge p$ then $xx^*$ and $x^*x$ commute, being projections above $p$.
\end{itemize}

\section{The local structure is unconstrained}\label{sec:local}

A natural first guess for a second condition is a restriction on the ``local'' pieces of $S$, i.e.\ on its principal factors (the traces of its $\D$-classes), which are completely $0$-simple regular $*$-semigroups. Some care is needed in saying what such a restriction could be. The zero of a principal factor is not in general an element of $S$, and the second consequence listed above shows that when it is -- that is, when $S$ contains a completely $0$-simple $*$-subsemigroup together with its zero -- \eqref{eq:Q1} already forces that subsemigroup to be a Brandt semigroup. So the zero pattern of a principal factor (which pairs $p,q$ of its projections have $pq=0$) is not local information: it is decided by how the $\D$-class sits inside $S$, and it does interact with the rest of $S$, as the hexagon condition of Section~\ref{sec:hex} shows. What remains genuinely local is the nonzero part of a principal factor -- a set of projections of a $\D$-class with pairwise nonzero products, together with the corresponding sandwich entries and the maximal subgroup. These are exactly the completely simple $*$-subsemigroups of $S$ (conversely, every completely simple $*$-subsemigroup of $S$ lies in a single $\D$-class and arises this way), and any condition on them would have to be a condition on finite completely simple regular $*$-semigroups. The next theorem shows there is none: every one of them embeds.

\begin{theorem}\label{thm:cs}
Every finite completely simple regular $*$-semigroup embeds in a finite partition monoid. Explicitly, let $S=\mathcal M[G;I,I;P]$ be in normal form ($p_{ii}=1$, $p_{ji}=p_{ij}^{-1}$, $(i,g,j)^*=(j,g^{-1},i)$), let $G\le\Sym(k)$ act on $\{1,\dots,k\}$ on the right, and let
\[
X=\{a_{i,s}: i\in I,\ s\le k\}\ \sqcup\ \{c_{ij,s}: i\ne j\in I,\ s\le k\},\qquad |X|=k|I|^2 .
\]
Send $(i,g,j)$ to the partition with transversal blocks $\{a_{i,s},a'_{j,sg}\}$ ($s\le k$), upper nontransversal blocks $\{c_{il,s},\,a_{l,sp_{il}},\,c_{li,sp_{il}}\}$ ($l\ne i$, $s\le k$) together with singletons $\{c_{lm,s}\}$ for $i\notin\{l,m\}$, and lower nontransversal blocks defined in the same way from $j$. This is an injective $*$-homomorphism $S\to\P_X$.
\end{theorem}

\begin{proof}[Proof sketch]
The transversal blocks of the image of $e_i=(i,1,i)$ are the singletons $a_{i,s}$, so the images of the $e_i$ are projections of rank $k$. In the middle layer of $e_ie_j$ the point $a_{i,s}$ lies in the $e_j$-block $\{c_{ji,t},a_{i,s},c_{ij,s}\}$ (with $s=tp_{ji}$), and $c_{ij,s}$ lies in the $e_i$-block $\{c_{ij,s},a_{j,sp_{ij}},c_{ji,sp_{ij}}\}$; no other blocks are attached (all remaining connectors are singletons for $e_i$ or $e_j$), so the class of $a_{i,s}$ is $\{a_{i,s},c_{ij,s},c_{ji,sp_{ij}},a_{j,sp_{ij}}\}$. Hence $e_ie_j$ has rank $k$ and its matching $\T(e_i)\to\T(e_j)$ is $s\mapsto sp_{ij}$. Classes containing no $a_{i,\cdot}$ or $a_{j,\cdot}$ contribute nothing. The images of the group elements are the elements of the rank-$k$ $\Hh$-classes of $\P_X$ determined by the permutations, and the multiplication $(i,g,j)(l,h,m)=(i,gp_{jl}h,m)$ is then read off block by block. Injectivity is clear. (All of this was also verified by computer for several nonabelian and non-flat sandwich matrices.)
\end{proof}

Since every finite group embeds in some $\Sym(k)$, every finite completely simple regular $*$-semigroup embeds. Two remarks on what this rules out:

\begin{remark}\label{rem:orthodox}
Orthodoxy of the completely simple $*$-subsemigroups is not necessary either, even though $\P_n$ for small $n$ looks ``flat'': the projections
\begin{align*}
p&: \{0\},\{1\}\ \text{transversal},\ \{2,3\}\ \text{nontransversal};\\
q&: \{1\},\{3\}\ \text{transversal},\ \{0,2\}\ \text{nontransversal};\\
r&: \{0\},\{3\}\ \text{transversal},\ \{1,2\}\ \text{nontransversal}
\end{align*}
of $\P_4$ satisfy $pqp=p$, $qrq=q$, $rpr=r$ (and the symmetric identities), but $pqr\ne pr$: the sandwich matrix of the rank-$2$ $\D$-class of $\P_4$ has a non-flat triangle, and the $*$-subsemigroup generated by $p,q,r$ is a non-orthodox completely simple regular $*$-semigroup $\mathcal M[\Z_2;3,3;P]$ with $p_{12}\ne1$.
\end{remark}

\begin{remark}[Normal $*$-bands]\label{rem:normal}
For a strong semilattice $Y$ of square $*$-bands $A_\alpha\times A_\alpha$ with structure maps $f_{\alpha\beta}\colon A_\alpha\to A_\beta$ ($\alpha\ge\beta$), \eqref{eq:Q1} holds iff all $f_{\alpha\beta}$ are injective (a projection $(c,c)$ of a lower component lies below $(a,b)$ iff $f(a)=f(b)=c$). If $Y$ has a least element $0$, injectivity lets us regard every $A_\alpha$ as a subset of $A=A_0$, and then the band is the $*$-subsemigroup $\{(\alpha,a,b): a,b\in A_\alpha\}$ of $Y\times(A\times A)$. Since semilattices and square bands embed in partition monoids and $\P_X\times\P_Y\le\P_{X\sqcup Y}$, we get: \emph{a finite normal $*$-band embeds in a partition monoid iff it satisfies \eqref{eq:Q1}}. Several strong semilattices of completely simple $*$-semigroups with non-injective group maps were also found to embed by computer.
\end{remark}

\section{The triangle condition and its relatives}\label{sec:hex}

We now give a family of new necessary conditions. Each is a quasi-identity in $2k$ projection variables ($k\ge3$), and they are independent of \eqref{eq:Q1}. The family is indexed by a pattern of matched pairs, which we encode as a digraph.

\begin{definition}
Let $k\ge3$. A set $E$ of ordered pairs $(i,j)$ of distinct elements of $[k]=\{1,\dots,k\}$ (a loopless digraph on $[k]$) is \emph{admissible} if
\begin{enumerate}[label=\textup{(\roman*)},nosep]
\item every vertex has an out-edge and an in-edge in $E$, and
\item for every $m\in[k]$ there is an edge $(i,j)\in E$ with $m\notin\{i,j\}$ (no vertex is incident to all edges).
\end{enumerate}
\end{definition}

\begin{theorem}[Matching-pattern condition]\label{thm:hex}
Let $E$ be an admissible digraph on $[k]$, and let $a_1,\dots,a_k,b_1,\dots,b_k$ be projections of a partition monoid $\P_X$ such that
\begin{enumerate}[label=\textup{(H\arabic*)},leftmargin=3em]
\item $a_ia_j=a_ja_i=:w$ and $b_ib_j=b_jb_i=:z$ for all $i\ne j$, the products being independent of the pair $\{i,j\}$;
\item $a_ib_i=b_ia_i$ for all $i$;
\item $a_ib_ja_i=a_i$ and $b_ja_ib_j=b_j$ for all $(i,j)\in E$.
\end{enumerate}
Then $w=z$. Consequently $a_i=b_j$ for all $(i,j)\in E$, and if the bipartite graph $\{a_i\text{---}b_j:(i,j)\in E\}$ is connected then all $2k$ projections coincide.
\end{theorem}

The hypotheses say: the $a_i$ are $k$ pairwise commuting projections with common meet $w$, the $b_i$ are $k$ pairwise commuting projections with common meet $z$; $a_i$ and $b_i$ commute; and $a_i\bowtie b_j$ for $(i,j)\in E$. All of them are equations of the projection algebra (commutation of projections is expressible by $\theta_p(q)=\theta_q(p)$: if $pqp=qpq=:w$ then $pq\in H_w$ and $w\cdot pq=(pq)^2$, so $pq$ is an idempotent of the group $H_w$, i.e.\ $pq=w=qp$), so each instance of the theorem is a quasi-identity of the projection algebras of partition monoids. The instances of most interest are:
\begin{itemize}[nosep]
\item $k=3$, $E=$ all six pairs: the \emph{hexagon condition} (the six matched pairs form the hexagon $a_1\,b_2\,a_3\,b_1\,a_2\,b_3$). This was the original form of the theorem.
\item $k=3$, $E=\{(1,2),(2,3),(3,1)\}$, a directed $3$-cycle: the \emph{triangle condition}. Only three matched pairs are assumed, $a_1\bowtie b_2$, $a_2\bowtie b_3$, $a_3\bowtie b_1$, so this instance has strictly weaker hypotheses than the hexagon condition and implies it. It is the unique admissible digraph on three vertices with three edges (a derangement of $[3]$ is a $3$-cycle), hence the strongest instance with $k=3$.
\item $k=4$, $E$ a directed $4$-cycle $\{(1,2),(2,3),(3,4),(4,1)\}$; and $k=4$, $E$ two disjoint $2$-cycles $\{(1,2),(2,1),(3,4),(4,3)\}$ (the ``double square'').
\end{itemize}
Restricting to a subset of at least three indices only removes hypotheses, so an instance of the theorem follows from every instance obtained by deleting indices. Every admissible digraph contains, on some subset of its vertices, a directed cycle of length $\ge3$ or two vertex-disjoint $2$-cycles (if all directed cycles are $2$-cycles and any two of them share a vertex, they all share one vertex $v$, or three of them form a directed triangle; in the former case take an edge $(p,q)$ avoiding $v$, follow out-edges from $q$ until $v$ is reached and in-edges from $p$ back to $v$: the resulting closed walk through $v$ contains no vertex twice, since a repetition would produce a cycle avoiding $v$, so it is a directed cycle of length $\ge3$). Both of these are admissible. Hence every instance of the theorem is implied by one of the \emph{basic} instances: a directed $\ell$-cycle $C_\ell$ ($\ell\ge3$), or the double square $D_2$. In fact the triangle condition alone implies every non-star instance, in every projection algebra (Theorem~\ref{thm:master}), so only the case $E=C_3$ of the theorem needs a combinatorial proof. For finite projection algebras the hexagon condition implies all these instances as well (Theorem~\ref{thm:equiv} below), although the implications are not equational ones; in general projection algebras it does not even imply the triangle condition (Theorem~\ref{thm:hinf}).

Write $m_i:=a_ib_i=b_ia_i$. The proof uses the following ``transport'' relations, which are formal consequences of (H1)--(H3); here $(i,j)\in E$:
\begin{equation}\label{eq:H4}
a_im_ja_i=w,\qquad b_jm_ib_j=z,\qquad a_iza_i=m_i,\qquad b_jwb_j=m_j .
\end{equation}
Indeed $a_im_ja_i=a_ia_jb_ja_i=a_j(a_ib_ja_i)=a_ja_i=w$; $b_jm_ib_j=b_jb_ia_ib_j=b_i(b_ja_ib_j)=b_ib_j=z$; $a_iza_i=a_ib_jb_ia_i=(a_ib_ja_i)b_i=m_i$; and $b_jwb_j=b_ja_ia_jb_j=(b_ja_ib_j)a_j=m_j$. By admissibility (i), every index $i$ has an out-edge and an in-edge, so for every $i$ there are $j,j'$ with $a_im_ja_i=w$, $a_iza_i=m_i$, $b_im_{j'}b_i=z$ and $b_iwb_i=m_i$. (In an earlier version of these notes \eqref{eq:H4} was listed as a fourth hypothesis of the hexagon condition; that it is redundant was pointed out by Prover9, working from the axioms of projection algebras alone.)

For the hexagon condition Prover9 also shows that the hypotheses can be pared down to $21$ equations: in (H1) it suffices to assume the six ``cyclic'' equations $a_2a_1a_2=a_3a_2a_3=a_1a_3a_1$ and $b_2b_1b_2=b_3b_2b_3=b_1b_3b_1$ (if $w$ denotes the common value of the first three, then $a_1a_2a_1=(a_1a_2a_1)^2=a_1(a_2a_1a_2)a_1=a_1wa_1=a_1(a_3a_1a_3)a_1=(a_1a_3a_1)^2=w$, and similarly $a_3a_2a_3=a_1a_3a_1=w$, so the $a_i$ pairwise commute with product $w$ by the remark above); (H2) consists of the three equations $a_ib_ia_i=b_ia_ib_i$; and (H3) consists of twelve equations. The same reduction of (H1) applies to every instance; for the triangle condition the hypotheses are then $6+3+6=15$ equations. The Prover9 inputs are the files \texttt{hexagon.in}, \texttt{triangle.in} and \texttt{octagon.in}.

\begin{proof}
Write $A_i=\supp(a_i)$, $B_i=\supp(b_i)$, $W=\supp(w)$, $Z=\supp(z)$, and $E_p$ for the partition of $X$ underlying a projection $p$ (its blocks are the $\T$-blocks and $\N$-blocks of $p$). We use Lemma~\ref{lem:PX} throughout.

\emph{Step 0: supports.} By (H1) and Lemma~\ref{lem:PX}(b), $A_i\cap A_j=W$ and $B_i\cap B_j=Z$ for $i\ne j$; by (H2), $A_i\cap B_i=\supp(m_i)$. In particular a point of $X$ lies in at most one $A_i$ unless it lies in $W$ (and then in all of them), and similarly for the $B_i$ and $Z$. Since $k\ge3$,
\begin{equation}\label{eq:cover}
\text{every } x\notin W\cup Z \text{ lies outside } A_i\cup B_i \text{ for some } i .
\end{equation}

\emph{Step 1: $E_w$ and $E_z$ agree outside $W\cup Z$.} By \eqref{eq:H4} and Lemma~\ref{lem:PX}(a) (applied to $z=\theta_{b_i}(m_{j'})$, $w=\theta_{a_i}(m_j)$, $m_i=\theta_{a_i}(z)$, $m_i=\theta_{b_i}(w)$), for every $i$: the blocks of $E_z$ not contained in $B_i$ are exactly the $\N(b_i)$-blocks; the blocks of $E_w$ not contained in $A_i$ are exactly the $\N(a_i)$-blocks; and the blocks of $E_{m_i}$ not contained in $A_i$ (resp.\ $B_i$) are exactly the $\N(a_i)$-blocks (resp.\ $\N(b_i)$-blocks). Consequently, for $x\notin A_i\cup B_i$, the block of $E_z$ containing $x$ is the $\N(b_i)$-block of $x$, which is the block of $E_{m_i}$ containing $x$, which is the $\N(a_i)$-block of $x$, which is the block of $E_w$ containing $x$. By \eqref{eq:cover}, $E_w$ and $E_z$ have the same block at every point of $X\setminus(W\cup Z)$; moreover such a common block is an $\N(a_i)$-block and an $\N(b_i)$-block, hence contained in $X\setminus(A_i\cup B_i)\subseteq X\setminus(W\cup Z)$. It follows that a block of $E_w$ or of $E_z$ is either a common block contained in $X\setminus(W\cup Z)$, or is contained in $W\cup Z$.

\emph{Step 2: matching of $\T(w)$ and $\T(z)$.} Fix an edge $(i,j)\in E$. From (H1) and \eqref{eq:H4}, $wzw=a_i(a_jza_j)a_i=a_im_ja_i=w$ and $zwz=b_j(b_iwb_i)b_j=b_jm_ib_j=z$, so $w\bowtie z$. By Lemma~\ref{lem:PX}(c), every class of $E_w\vee E_z$ meeting $W\cup Z$ contains exactly one $\T(w)$-block and exactly one $\T(z)$-block. Common blocks are classes by themselves, so the class of a $\T(w)$-block $\Omega$ consists of blocks contained in $W\cup Z$: these are $\T(w)$-blocks, $\N(w)$-blocks contained in $Z\setminus W$, $\T(z)$-blocks and $\N(z)$-blocks contained in $W\setminus Z$. An $\N(z)$-block contained in $W$ meets only $\T(w)$-blocks, so it meets only one of them (two would lie in a common class); hence every $\N(z)$-block contained in $W$ lies inside a single $\T(w)$-block, and dually. Therefore the class of $\Omega$ is
\[
\{\Omega\}\ \cup\ \{\N(z)\text{-blocks inside }\Omega\}\ \cup\ \{\T(z)\text{-blocks meeting }\Omega\}\ \cup\ \{\N(w)\text{-blocks inside those}\},
\]
and since it contains exactly one $\T(z)$-block, \emph{there is exactly one $\T(z)$-block $Z_\Omega$ meeting $\Omega$}, and dually exactly one $\T(w)$-block $W_\Theta$ meeting a given $\T(z)$-block $\Theta$.

\emph{Step 3: $\Omega\subseteq Z_\Omega$.} Suppose $y\in\Omega\setminus Z_\Omega$. Then $y\notin Z$ (a $\T(z)$-block containing $y$ would meet $\Omega$, hence be $Z_\Omega$), so the block $N$ of $E_z$ containing $y$ is an $\N(z)$-block, and by Step~2, $N\subseteq\Omega\subseteq W\subseteq A_1\cap\dots\cap A_k$. For each $i$, $N$ is a block of $E_z=\theta_{b_i}(m_{j'})$ (with $(j',i)\in E$), so by Lemma~\ref{lem:PX}(a) either $N\subseteq B_i$ or $N$ is an $\N(b_i)$-block. As $N\cap Z=\emptyset$ and $B_i\cap B_l=Z$, $N\subseteq B_i$ holds for at most one index $i=m$; by admissibility (ii) there is an edge $(i,j)\in E$ with $m\notin\{i,j\}$, so that $N$ is both an $\N(b_i)$-block and an $\N(b_j)$-block. Since $N\subseteq A_i$ and $a_ib_i=b_ia_i$, Lemma~\ref{lem:PX}(b) shows that $N$ is a union of $\T(a_i)$-blocks; let $\tau$ be one of them. Now $a_i\bowtie b_j$ by (H3). The class of $\varepsilon_{a_i}\vee\varepsilon_{b_j}$ containing $\tau$ contains the $b_j$-block $N\supseteq\tau$, hence all $\T(a_i)$-blocks inside $N$; by Lemma~\ref{lem:PX}(c) there is only one, so $N=\tau$. But then $N$ is simultaneously a $\T(a_i)$-block and an $\N(b_j)$-block, which is impossible for a matched pair by Lemma~\ref{lem:PX}(c). Hence $\Omega\subseteq Z_\Omega$.

\emph{Step 4: conclusion.} Dually $Z_\Omega\subseteq W_{Z_\Omega}$ (the dual argument uses an edge $(i,j)$ chosen by (ii) as in Step~3, with the roles of $a$ and $b$ exchanged: an $\N(w)$-block $N\subseteq Z\subseteq B_j$ is a union of $\T(b_j)$-blocks because $b_j$ commutes with $a_j$, and $b_j\bowtie a_i$); since $\Omega$ meets $Z_\Omega$ and $W_{Z_\Omega}$ is the only $\T(w)$-block meeting $Z_\Omega$, we get $W_{Z_\Omega}=\Omega$ and so $\Omega=Z_\Omega$. Thus $\T(w)=\T(z)$ and $W=Z$. Inside $W=Z$ both $E_w$ and $E_z$ consist of these $\T$-blocks, and outside $W\cup Z$ they agree by Step~1. Hence $E_w=E_z$ and $\T(w)=\T(z)$, i.e.\ $w=z$.

Finally, $w=z$ is a common lower bound of $a_i$ and $b_j$, so by \eqref{eq:Q1} (which holds in $\P_X$) $a_ib_j=b_ja_i$ for all $i,j$; for $(i,j)\in E$ this gives $a_i=a_ib_ja_i=a_ib_j$, i.e.\ $a_i\le b_j$, and dually $b_j\le a_i$, so $a_i=b_j$. If the bipartite graph of $E$ is connected, all $2k$ projections coincide.
\end{proof}

\begin{remark}\label{rem:2k}
(1) The hypothesis $k\ge3$ is needed: for $k=2$ (two matched pairs $a_1\bowtie b_2$, $a_2\bowtie b_1$, the ``square'') the statement is false, and $\P_2$ already contains counterexamples (\texttt{digraph.py}). The admissibility conditions, on the other hand, are conditions of the proof rather than of the statement: for every $X$ the exact condition is that $E$ is not a star (Theorems~\ref{prop:path} and~\ref{thm:master} below; for instance the $k=3$ ``star'' $\{(1,2),(2,1),(1,3),(3,1)\}$ -- a star in the graph-theoretic but not in our sense, since it has two tails and two heads -- the transitive tournament $\{(1,2),(2,3),(1,3)\}$, the path $\{(1,2),(2,3)\}$ and the two disjoint pairs $\{(1,2),(3,4)\}$ are all valid in $\P_n$ for finite $n$; an earlier version of this remark wrongly reported counterexamples to the first two, due to a bug in the search program).

(2) Exhaustive search (allowing coincidences among the variables) confirms that the hexagon, triangle, directed $4$-cycle, double-square and cyclic-octagon instances have no counterexamples among the projections of $\P_n$ for $n\le4$ (\texttt{digraph.py}, \texttt{octcheck.py}).

(3) For finite algebras the hexagon condition implies all the others (Theorem~\ref{thm:equiv}) and is equivalent to the conditions on three indices; in general it does not even imply the triangle condition (Theorem~\ref{thm:hinf} below), while the conditions on four or more indices do not imply the hexagon condition even for finite algebras ($\P(Q_0)$). Whether the hexagon condition implies $C_4$, $D_2$ or the cyclic octagon condition in general is open. A separating model for ``hexagon $\not\Rightarrow$ octagon'' (or ``hexagon $\not\Rightarrow$ triangle'') cannot be small: in any regular $*$-semigroup with the hexagon hypotheses, a pair $a_i,b_j$ ($i\ne j$) that is \emph{not} assumed matched can be neither matched (else the configuration contains a hexagon) nor commuting: in the cyclic octagon, if $a_i$ and $b_{i+2}$ commuted with product $u$, then $a_{i+2}ua_{i+2}=(a_{i+2}a_i)(b_{i+2}a_{i+2})=wm_{i+2}$ would be a projection, so $w$ and $m_{i+2}$ would commute; but $w\bowtie m_l$ always ($wm_lw=wzw=w$ and $m_lwm_l=a_lzwza_l=m_l$ by \eqref{eq:H4}), whence $w=m_{i+2}$, a degenerate case. Mace4 experiments by M.~Kinyon find that the smallest projection algebra violating the cyclic octagon condition has $14$ elements -- the $8$ vertices, $w$, $z$ and the four $m_i$ -- and that it violates the hexagon condition as well; indeed, of the $3^8$ ways to complete such a $14$-element table, exactly one is a projection algebra, and in it every antipodal pair $a_i,b_{i+2}$ is matched (\texttt{oct14.py}). So a separating example needs further projections $a_ib_ja_i$ in ``general position''. On the partition side, the quotients $\P_4/\mu_1$ (57 projections) and $\P_5/\mu_1$ (303 projections) contain hexagon configurations with $w\ne z$ but no octagon configurations; $\P_5$ contains no octagon configuration with eight distinct vertices even when the products in (H1) are only required to commute, while $\P_6$ does.
\end{remark}

\subsection{The hexagon condition implies the triangle condition (finite case)}\label{subsec:equiv}

Since the triangle condition assumes only three of the six matched pairs of the hexagon condition, it implies the hexagon condition trivially. The converse is true for finite projection algebras, but for a rather subtle reason: no fixed sextuple of terms in $a_1,\dots,b_3$ satisfies the hexagon hypotheses in every triangle configuration (see Remark~\ref{rem:equiv}), and the six projections to which the hexagon condition has to be applied are found by an iteration. Throughout this subsection $a_1,a_2,a_3,b_1,b_2,b_3$ is a triangle configuration ($a_1\bowtie b_2$, $a_2\bowtie b_3$, $a_3\bowtie b_1$, and (H1), (H2)) in a projection algebra $P=\P(S)$, $w$ and $z$ are the products, and $m_i=a_ib_i$. Indices are mod~$3$. We call a projection $x$ \emph{bounded} if $x\le a_i$ and $w,m_i\le x$ for some $i$, or $x\le b_j$ and $z,m_j\le x$ for some $j$; e.g.\ the six vertices are bounded.

We use freely the transport relations \eqref{eq:H4}, which by the argument of Lemma~\ref{lem:rev} hold for all $i\ne j$.

\begin{lemma}\label{lem:rev}
$b_jm_ib_j=z$ and $a_jm_ia_j=w$ for all $i\ne j$.
\end{lemma}
\begin{proof}
Put $e=b_jm_ib_j$. Writing $m_i=b_ia_i$ and $b_jb_i=z$ gives $e=za_ib_j$; writing $m_i=a_ib_i$ gives $e=b_ja_iz$. Hence $e=ze=ez$, so $e\le z$. Conversely $zez=zm_iz=za_iz=z$ because $z\le b_{i+1}\bowtie a_i$ (here index $i$ has an out-edge); so $z\le e$ and $e=z$. Dually for $w$.
\end{proof}

\begin{lemma}[bounded elements]\label{lem:bounded}
Let $e\le a_i$, $e'\le a_j$ with $i\ne j$, and $w\le e,e'$. Then $ee'=e'e=w$. Dually two bounded elements under distinct $b$'s commute with product $z$. If $e\le a_i$, $f\le b_i$ and $m_i\le e,f$, then $ef=fe=m_i$.
\end{lemma}
\begin{proof}
$ee'e=e(a_ia_j)e'(a_ja_i)e=ewe'we=ewe=w$, and symmetrically $e'ee'=w$; two projections with $pqp=qpq=w$ commute with product $w$ (Section~\ref{sec:hex}). Likewise $efe=e(a_ib_i)f(b_ia_i)e=em_ifm_ie=m_i=fef$.
\end{proof}

Thus any six bounded projections, one under each vertex, satisfy the commutation hypotheses (H1), (H2) of the hexagon condition with the same $w$ and $z$; only the six cross matchings are at issue. They will come from the following \emph{walk}: for a projection $x\le a_1$ put
\[
x\ \xrightarrow{\ \theta_{b_3}\ }\ x_1\ \xrightarrow{\ \theta_{a_2}\ }\ x_2\ \xrightarrow{\ \theta_{b_1}\ }\ x_3\ \xrightarrow{\ \theta_{a_3}\ }\ x_4\ \xrightarrow{\ \theta_{b_2}\ }\ x_5\ \xrightarrow{\ \theta_{a_1}\ }\ x_6=:\Phi(x)\le a_1 ,
\]
alternately projecting onto the free partner ($a_i\mapsto b_{i-1}$) and transporting through a matched pair ($b_j\mapsto a_{j-1}$, where $a_{j-1}\bowtie b_j$). The six elements $x_1,\dots,x_6$ lie under $b_3,a_2,b_1,a_3,b_2,a_1$ respectively.

\begin{lemma}[invariance]\label{lem:inv}
The walk preserves boundedness: if $x\le a_i$ is bounded then so is $\theta_{b_{i-1}}(x)$, and if $y\le b_j$ is bounded then so is $\theta_{a_{j-1}}(y)$.
\end{lemma}
\begin{proof}
Let $y=\theta_b(x)$ with $b=b_{i-1}$. Since $z\le b$, $yzy=b\,x(bzb)x\,b=b(xzx)b$, and $xzx=x(a_iza_i)x=xm_ix=m_i$, so $yzy=bm_ib=z$, i.e.\ $z\le y$. Similarly $ym_{i-1}y=b(xm_{i-1}x)b=b(x(a_im_{i-1}a_i)x)b=b(xwx)b=bwb=m_{i-1}$. The transport step is the same computation with the roles of $a$ and $b$ exchanged: for $x'=\theta_a(y)$, $a=a_{j-1}$, one has $x'wx'=a(ywy)a=a(y(b_jwb_j)y)a=am_ja=w$ and $x'm_{j-1}x'=a(y(b_jm_{j-1}b_j)y)a=aza=m_{j-1}$.
\end{proof}

\begin{lemma}[monotonicity]\label{lem:mono}
In every projection algebra, $x\le y$ implies $\theta_p(x)\le\theta_p(y)$.
\end{lemma}
\begin{proof}
In the binary notation $u*v=u\theta_v$, the axioms imply $(x*y)*(z*y)=x*(z*y)$ and $(x*y)*(y*z)=(x*y)*z$ (both are found by Prover9 in a few steps from (P2)--(P5)). If $x\le y$, i.e.\ $x*y=x$, then $(x*p)*(y*p)=x*(y*p)=(x*y)*(y*p)=(x*y)*p=x*p$, i.e.\ $\theta_p(x)\le\theta_p(y)$.
\end{proof}

\begin{lemma}[stabilisation]\label{lem:stab}
Suppose $\P(S)$ is finite. If $x_0\to x_1=\theta_{p_1}(x_0)\to\dots\to x_n=\theta_{p_n}(x_{n-1})=x_0$ is a closed walk of projections, then $x_{j-1}\bowtie x_j$ for all $j$.
\end{lemma}
\begin{proof}
Put $s_j=x_{j-1}p_j$, so that $s_j^*s_j=x_j$ and $s_js_j^*=\theta_{x_{j-1}}(p_j)\le x_{j-1}$, and $s=s_1s_2\cdots s_n$. Since $s_j^*s_js_{j+1}=x_js_{j+1}=s_{j+1}$, we get $s^*s=x_n=x_0$, and $ss^*\le s_1s_1^*=\theta_{x_0}(p_1)\le x_0$ (using $s_j(\cdot)s_j^*$ maps projections $\le s_j^*s_j$ to projections $\le s_js_j^*$). Now $t\mapsto s^*ts$ is an order isomorphism from the projections below $ss^*$ onto the projections below $s^*s$; as $ss^*\le s^*s$ and $\P(S)$ is finite, $ss^*=s^*s$. Hence $\theta_{x_0}(p_1)=x_0$, so $x_0x_1x_0=(x_0p_1x_0)(p_1x_0)=x_0$, while $x_1x_0x_1=x_1$ always ($\theta_{\theta_p(q)}(q)=\theta_p(q)$ is an axiom consequence). Rotating the cycle gives the other pairs.
\end{proof}

\begin{theorem}\label{thm:equiv}
Let $P$ be a finite projection algebra (more generally, $P=\P(S)$ with $S$ a regular $*$-semigroup having finitely many projections) satisfying the hexagon condition. Let $k\ge3$, let $E$ be a digraph on $[k]$ in which every vertex has an in-edge and an out-edge, and let $a_1,\dots,a_k,b_1,\dots,b_k$ satisfy (H1)--(H3) of Theorem~\ref{thm:hex} for $E$. Then $w=z$. In particular, for finite projection algebras the hexagon condition implies the triangle condition and the matching-pattern conditions of all digraphs with in- and out-edges at every vertex (on any number $k\ge3$ of vertices, e.g.\ $C_\ell$, $D_2$ and the cyclic octagon); on three vertices the converse holds too, since every $E$ on $[3]$ has fewer hypotheses than the hexagon, so for finite projection algebras the hexagon condition, the triangle condition and all the other non-star conditions on three indices (Theorem~\ref{prop:path}) are equivalent. The conditions on $k\ge4$ indices are strictly weaker: $\P(Q_0)$ satisfies the $C_4$, $D_2$ and cyclic-octagon conditions (\texttt{mace4parse.py}), as no fourth pair can be added to its hexagon.
\end{theorem}
\begin{proof}
By a theorem of Imaoka (see East and Azeef Muhammed, arXiv:2301.04845, Section~4) every projection algebra is $\P(S)$ for some regular $*$-semigroup $S$. We first treat a triangle configuration $a_1,\dots,b_3$ in $P$; at the end we explain why only the transport relations were used. The operator $\Phi$ is a composite of maps $\theta_p$, hence monotone by Lemma~\ref{lem:mono}, and $\Phi(a_1)\le a_1$. Therefore $a_1\ge\Phi(a_1)\ge\Phi^2(a_1)\ge\cdots$ is a descending chain, which is eventually constant since $P$ is finite: there is $x^*=\Phi^n(a_1)$ with $\Phi(x^*)=x^*$. By Lemma~\ref{lem:inv} (and induction from the bounded element $a_1$) all elements of the walk from $x^*$ are bounded, and this walk $x^*\to x_1\to\dots\to x_6=x^*$ is closed, so by Lemma~\ref{lem:stab} consecutive elements are matched. Put
\[
A_3=x_1\le b_3,\ \ A_1=x_3\le b_1,\ \ A_2=x_5\le b_2,\qquad B_2=x_2\le a_2,\ \ B_3=x_4\le a_3,\ \ B_1=x_6\le a_1 .
\]
By Lemma~\ref{lem:bounded} the $A_j$ pairwise commute with product $z$, the $B_i$ with product $w$, and $A_iB_i=B_iA_i$; the six consecutive matchings $x_1\bowtie x_2\bowtie\dots\bowtie x_6\bowtie x_1$ are exactly the matchings $A_i\bowtie B_j$, $i\ne j$. So $(B_1,B_2,B_3\,|\,A_1,A_2,A_3)$ satisfies the hexagon hypotheses with products $w$ and $z$, and the hexagon condition gives $w=z$.

For a general $E$ with in- and out-edges at every vertex, fix three indices $i,j,l$ and run the same walk $a_i\to b_j\to a_l\to b_i\to a_j\to b_l\to a_i$. Lemmas~\ref{lem:rev}, \ref{lem:bounded}, \ref{lem:mono} and \ref{lem:stab} are general, and the proof of Lemma~\ref{lem:inv} used, at a step from a vertex of index $i$ to a vertex of index $j$, only the transport relations $a_iza_i=m_i$, $b_jm_ib_j=z$ (an out-edge at $i$) and $a_im_ja_i=w$, $b_jwb_j=m_j$ (an in-edge at $j$), and dually. So the walk from $a_i$ stays bounded, and the rest of the argument is unchanged.
\end{proof}

\begin{remark}\label{rem:equiv}
(1) The proof was found with substantial help from Prover9 and Mace4, run by M.~Kinyon. Mace4 models with $w\ne z$ showed that the derived projections forming the hexagon vary from model to model (the walk from $a_1$ stabilises after $0$, $1$ or $2$ turns in the models found), and Prover9 proved the ground equations behind Lemmas \ref{lem:bounded}, \ref{lem:inv} and \ref{lem:mono}, from whose output the hand proofs were extracted.

(2) The implication cannot be obtained by a single application of the hexagon condition to a fixed sextuple of terms in $a_1,\dots,b_3$: for any such sextuple, one of its hexagon hypotheses $t=t'$ fails in some model of the triangle hypotheses in which $w=z$ (Mace4 finds them), while $w=z$ fails in others; since the theory ``axioms $+$ triangle hypotheses'' is Horn, its initial model satisfies neither $t=t'$ nor $w=z$. For infinite projection algebras the hexagon condition does not imply the triangle condition (Theorem~\ref{thm:hinf}); in the example the walk closes after one turn but Lemma~\ref{lem:stab} fails, and Mace4 also finds finite models of the triangle hypotheses in which the chain $a_1>\Phi(a_1)>\Phi^2(a_1)$ descends strictly for two turns before closing (sizes $13$, $21$, and at least $30$ for three turns; \texttt{chain.py}).

(3) Applied to $\P(\P_X)$ for finite $X$, which satisfies the hexagon condition by Theorem~\ref{thm:hex}, the theorem shows that condition (ii) of admissibility is not needed in Theorem~\ref{thm:hex} for finite $X$: in- and out-edges at every vertex suffice. (Indeed Steps 1 and 2 of the proof of Theorem~\ref{thm:hex} only use (i) -- in Step 1 the $\N(b_i)$-block and the $\N(a_i)$-block of a point outside $A_i\cup B_i$ coincide by Lemma~\ref{lem:PX}(b), since $a_ib_i=b_ia_i$ -- and only Step 3 used (ii).) Theorem~\ref{prop:path} below goes further.

(4) The argument shows more generally that in a finite projection algebra, for any bounded $x\le a_1$, the walk from $x$ ends in a closed $6$-cycle which is a hexagon configuration with products $w$ and $z$. Two further closed cycles always exist, through $z$ and through $w$ ($z\to m_2\to z\to m_3\to z\to m_1\to z$ and dually); they are not hexagons, which is why the bounds in Lemma~\ref{lem:inv} involve $m_i$ as well as $w$ or $z$.
\end{remark}


\begin{theorem}[non-star patterns]\label{prop:path}
Let $E$ be a digraph on $[k]$, $k\ge3$. Call $E$ a \emph{star} if it is empty, or all its edges have the same tail (an out-star), or all have the same head (an in-star); equivalently, $E$ is not a star iff it has two distinct tails and two distinct heads.
\begin{enumerate}[label=\textup{(\alph*)}]
\item If $E$ is not a star, then in every projection algebra the hypotheses (H1)--(H3) for $E$ imply $w\bowtie z$ and the four transport relations
\begin{equation}\label{eq:transport}
a_iza_i=m_i,\qquad b_iwb_i=m_i,\qquad za_iz=z,\qquad wb_iw=w
\end{equation}
at every non-isolated index $i$ (hence also $b_jm_ib_j=z$ and $a_jm_ia_j=w$ for all $j\ne i$); moreover each isolated index $i$ can be replaced by the pair $a_i':=\theta_{a_i}(b_{j_0})$, $b_i':=\theta_{b_i}(a_{j_0})$, where $j_0$ is any non-isolated index, after which (H1)--(H3) still hold with the same $w,z$ and \eqref{eq:transport} holds at every index. Consequently, in a finite projection algebra satisfying the hexagon condition, and in particular in $\P(\P_X)$ for finite $X$, (H1)--(H3) for $E$ force $w=z$.
\item If $E$ is a star, then $\P_2$ contains projections satisfying (H1)--(H3) for $E$ with $w\ne z$.
\end{enumerate}
So for finite $X$ the matching-pattern condition of Theorem~\ref{thm:hex} holds in $\P_X$ for exactly the non-star digraphs $E$. The smallest instances beyond the hexagon are the \emph{path} $a_1\bowtie b_2$, $a_2\bowtie b_3$, the \emph{$2$-cycle-plus-spare} $a_1\bowtie b_2$, $a_2\bowtie b_1$ (both on three indices, with (H1), (H2)), and the \emph{two disjoint pairs} $a_1\bowtie b_2$, $a_3\bowtie b_4$ on four indices.
\end{theorem}
\begin{proof}
(a) For an index $i$ with an out-edge $(i,l)$ we have $a_iza_i=a_ib_lb_ia_i=(a_ib_la_i)b_i=m_i$ and $za_iz=zb_la_ib_lz=z$; dually an in-edge $(l,i)$ gives $b_iwb_i=m_i$ and $wb_iw=w$. By the argument of Lemma~\ref{lem:rev}, $za_iz=z$ gives $b_jm_ib_j=z$ for all $j\ne i$, and $wb_iw=w$ gives $a_jm_ia_j=w$ for all $j\ne i$.

\emph{Step 1: $b_iwb_i=m_i$ at every tail $i$.} Let $i\ne i'$ be tails. Since $w\le a_i$, $wb_iw=wa_ib_ia_iw=wm_iw=w(a_iza_i)w=wzw$, and likewise $wb_{i'}w=wzw$. Hence, $b_iwb_i$ being a projection,
\[
b_iwb_i=b_i(wb_iw)b_i=b_i(wb_{i'}w)b_i=b_ia_ia_{i'}b_{i'}a_{i'}a_ib_i=m_i\,m_{i'}\,m_i=a_ib_i\,(b_{i'}a_{i'}b_{i'})\,b_ia_i=a_iza_{i'}za_i=a_iza_i=m_i ,
\]
using $b_ib_{i'}=z$, $za_{i'}z=z$ and $a_iza_i=m_i$. Dually $a_jza_j=m_j$ at every head $j$, since there are two distinct heads: $za_jz=zm_jz=z(b_jwb_j)z=zwz=zb_{j'}wb_{j'}z=zm_{j'}z$, so $a_jza_j=a_j(zm_{j'}z)a_j=m_jm_{j'}m_j=b_ja_j(a_{j'}b_{j'}a_{j'})a_jb_j=b_jwb_{j'}wb_j=b_jwb_j=m_j$.

\emph{Step 2: $w\bowtie z$.} Choose a tail $i$ and a head $j$ with $i\ne j$ (possible since there are two of each). Then $wzw=a_ia_jza_ja_i=a_im_ja_i=a_ia_jb_ja_ja_i=wb_jw=w$ by Step~1 and the in-edge at $j$, and $zwz=b_jb_iwb_ib_j=b_jm_ib_j=b_jb_ia_ib_ib_j=za_iz=z$ by Step~1 and the out-edge at $i$.

\emph{Step 3: the remaining relations.} At a tail $i$ we now have $a_iza_i=m_i$, $za_iz=z$ (out-edge) and $b_iwb_i=m_i$ (Step~1), and $wb_iw=w$ follows from $w\bowtie z$ by monotonicity (Lemma~\ref{lem:mono}): $z\le b_i$ gives $w=wzw\le wb_iw\le w$. Dually at a head. So \eqref{eq:transport} holds at every non-isolated index, and Lemma~\ref{lem:rev} gives $b_jm_ib_j=z$, $a_jm_ia_j=w$ for all $j\ne i$.

\emph{Step 4: isolated indices.} Let $i$ be isolated and $j_0$ non-isolated. Since $w\le a_i$, monotonicity gives $z=zwz\le za_iz\le z$, so $za_iz=z$, and dually $wb_iw=w$; by Lemma~\ref{lem:rev} again, $b_lm_ib_l=z$ and $a_lm_ia_l=w$ for $l\ne i$. The other two relations of \eqref{eq:transport} may fail at $i$ (Mace4 finds models with $a_iza_i<m_i$, and so does $\P_2$), but the two sandwiches are equal:
\[
u_i:=a_iza_i=b_iwb_i=m_im_{j_0}m_i .
\]
Indeed $z=b_im_{j_0}b_i$ (Lemma~\ref{lem:rev} at $j_0$), and writing $m_{j_0}=a_{j_0}b_{j_0}$, $b_{j_0}b_i=z$ in one copy and $m_{j_0}=b_{j_0}a_{j_0}$ in the other, $z=b_ia_{j_0}z=za_{j_0}b_i$; hence $a_iza_i=a_i(b_ia_{j_0}z)(za_{j_0}b_i)a_i=m_i(a_{j_0}za_{j_0})m_i=m_im_{j_0}m_i$, and dually $b_iwb_i=m_i(b_{j_0}wb_{j_0})m_i=m_im_{j_0}m_i$. Thus $u_i=\theta_{m_i}(m_{j_0})\le m_i$, and for $l\ne i$ we get $b_lu_ib_l=z$: it is $\le b_lm_ib_l=z$ by monotonicity, and $z(b_lu_ib_l)z=zu_iz=za_iza_iz=z$; dually $a_lu_ia_l=w$.

Put $a_i'=a_ib_{j_0}a_i$ and $b_i'=b_ia_{j_0}b_i$. Then $a_i'\le a_i$ and $w\le a_i'$, since $a_i'wa_i'=a_i(b_{j_0}wb_{j_0})a_i=a_im_{j_0}a_i=w$; dually $z\le b_i'\le b_i$. By Lemma~\ref{lem:bounded}, $a_i'$ commutes with every $a_l$ ($l\ne i$), and with $a_l'$ if $l$ is also isolated, with product $w$, and dually for $b_i'$: (H1) holds for the new family, with the same $w$ and $z$. For (H2),
\[
a_i'b_i'a_i'=a_ib_{j_0}(m_ia_{j_0}m_i)b_{j_0}a_i=a_ib_{j_0}\bigl(b_i(a_ia_{j_0}a_i)b_i\bigr)b_{j_0}a_i=a_i(b_{j_0}u_ib_{j_0})a_i=a_iza_i=u_i ,
\]
and dually $b_i'a_i'b_i'=b_iwb_i=u_i$; so $a_i'b_i'=b_i'a_i'=u_i=:m_i'$. (H3) is unaffected, as $i$ lies on no edge. Finally \eqref{eq:transport} holds at $i$ for the new pair: $a_i'za_i'=a_ib_{j_0}(a_iza_i)b_{j_0}a_i=a_i(b_{j_0}u_ib_{j_0})a_i=u_i=m_i'$, dually $b_i'wb_i'=m_i'$, and $za_i'z=z$, $wb_i'w=w$ as before since $w\le a_i'$, $z\le b_i'$. The relations at the other indices only involve $a_l,b_l,m_l,w,z$ for $l\ne i$ and are unchanged.

\emph{Step 5.} After the replacement all four relations hold at every index, and with them the proof of Lemma~\ref{lem:inv}, hence of Theorem~\ref{thm:equiv}, goes through for any three indices of $E$: the walk from $a_i$ ends in a closed $6$-cycle which is a hexagon configuration with products $w$ and $z$, and the hexagon condition gives $w=z$.

(b) In $\P_2$ let $F$ be the projection with the single block $\{1,2,1',2'\}$ (rank $1$) and $e$ the projection with blocks $\{1,2\}$, $\{1',2'\}$ (rank $0$); then $e\le F$, $eF=Fe=e$, $F\bowtie F$. For an out-star with tail $v$ (or for $E=\emptyset$) take $a_v=F$, $a_i=e$ for $i\ne v$, and $b_j=F$ for all $j$: then $w=e$, $z=F$, all pairs commute, $a_v\bowtie b_j$ for all $j$, and $w\ne z$. For an in-star with head $v$ take $a_i=F$ for all $i$, $b_v=F$ and $b_j=e$ for $j\ne v$.
\end{proof}

The transport relations are all that the matchings of $E$ contribute; once they are available, matchings can be manufactured at will, because of the following elementary fact.

\begin{lemma}\label{lem:friends}
For any two projections $p,q$ of a regular $*$-semigroup, $\theta_p(q)\bowtie\theta_q(p)$.
\end{lemma}
\begin{proof}
$pq$ is idempotent, being $(pq)(pq)^*(pq)$. Hence $(pqp)(qpq)(pqp)=p(qp)^4=p(qp)=pqp$ and symmetrically $(qpq)(pqp)(qpq)=qpq$.
\end{proof}

Call a family $a_1,\dots,a_k,b_1,\dots,b_k$ satisfying (H1), (H2) \emph{transported} if the four relations \eqref{eq:transport} hold at every index. By Theorem~\ref{prop:path}(a) and its proof, a configuration for a non-star digraph becomes transported after the isolated indices are replaced, and in a transported family $w\bowtie z$, $b_jm_ib_j=z$ and $a_jm_ia_j=w$ for all $i\ne j$ (Lemma~\ref{lem:rev}), and the bounded-elements Lemma~\ref{lem:bounded} applies.

\begin{lemma}[edge creation]\label{lem:create}
Let $a_1,\dots,a_k,b_1,\dots,b_k$ be transported, $k\ge2$, and $i\ne j$. Put $a_i'=\theta_{a_i}(b_j)$, $b_j'=\theta_{b_j}(a_i)$, and replace $a_i$ by $a_i'$ and $b_j$ by $b_j'$. The new family is again transported, with the same $w$, $z$ and $m_1,\dots,m_k$; $a_i'\bowtie b_j'$; and every pair $a_l,b_m$ with $l\ne i$, $m\ne j$ is unchanged.
\end{lemma}
\begin{proof}
$a_i'\le a_i$, and $w,m_i\le a_i'$: indeed $a_i'wa_i'=a_i(b_jwb_j)a_i=a_im_ja_i=w$ and $a_i'm_ia_i'=a_i(b_jm_ib_j)a_i=a_iza_i=m_i$. Dually $b_j'\le b_j$ and $z,m_j\le b_j'$. By Lemma~\ref{lem:bounded}, $a_i'$ commutes with every $a_l$ ($l\ne i$) with product $w$, and with $b_i$ with product $m_i$; dually for $b_j'$; so (H1), (H2) hold with the same $w,z,m_l$. The transport relations at $i$: $a_i'za_i'=a_i(b_jm_ib_j)a_i=m_i$, $b_iwb_i=m_i$ is unchanged, $za_i'z=z$ since $z=zwz\le za_i'z\le z$ by monotonicity ($w\le a_i'$), and $wb_iw=w$ is unchanged; dually at $j$; the other indices are untouched. Finally $a_i'\bowtie b_j'$ by Lemma~\ref{lem:friends}.
\end{proof}

\begin{theorem}[the triangle condition is a master condition]\label{thm:master}
\begin{enumerate}[label=\textup{(\alph*)}]
\item Every transported family on $k\ge3$ indices yields, on any three of its indices, a triangle configuration with the same $w$ and $z$. Consequently, in \emph{every} projection algebra, the triangle condition implies the matching-pattern condition of every non-star digraph.
\item For every set $X$, the matching-pattern condition of $E$ holds in $\P_X$ iff $E$ is not a star ($|X|\ge2$ for the ``only if'').
\end{enumerate}
\end{theorem}
\begin{proof}
(a) Apply Lemma~\ref{lem:create} three times on the indices $1,2,3$: create $(1,2)$, replacing $a_1,b_2$; then $(2,3)$, replacing $a_2,b_3$; then $(3,1)$, replacing $a_3,b_1$. Each step leaves the pairs created before untouched, since they involve different elements: after the three steps $a_1'\bowtie b_2'$, $a_2'\bowtie b_3'$, $a_3'\bowtie b_1'$, and (H1), (H2) hold with the original $w,z$ by the lemma. For the consequence, take a configuration for a non-star $E$, replace its isolated indices (Theorem~\ref{prop:path}), obtaining a transported family with the same $w,z$, and build the triangle configuration; the triangle condition gives $w=z$.

(b) The triangle condition holds in $\P_X$ for every $X$ by Theorem~\ref{thm:hex}, so the non-star conditions hold by (a); the stars fail already in $\P_2\subseteq\P_X$ (Theorem~\ref{prop:path}(b)).
\end{proof}

Thus the matchings of a pattern serve only to establish the transport relations; and among all matching-pattern conditions the triangle condition is the strongest one valid in partition monoids, in the sense that it implies all the others in every projection algebra, while the path and the $2$-cycle conditions are equivalent to it (they imply it trivially, having fewer hypotheses). In the proof of Theorem~\ref{thm:hex} only the triangle is now needed; the other admissible digraphs follow from it algebraically. The hexagon condition is strictly weaker than the triangle condition in general (Theorem~\ref{thm:hinf} below), although the two are equivalent on finite projection algebras (Theorem~\ref{thm:equiv}).


\begin{remark}\label{rem:path}
(1) Theorem~\ref{prop:path} was found in three stages, each driven by Mace4 and Prover9 experiments of M.~Kinyon. First, on three indices, exhaustive search in $\P_n$, $n\le4$, showed that two matched pairs force $w=z$ iff they form a directed path $\{(1,2),(2,3)\}$ (up to relabelling) or the $2$-cycle $\{(1,2),(2,1)\}$, while two pairs with a common tail $\{(1,2),(1,3)\}$ or a common head $\{(2,1),(3,1)\}$, the pair $\{(1,2),(3,2)\}$, and a single pair all admit counterexamples in $\P_2$ (\texttt{digraph3.py}, \texttt{digraph4.py}); the path is the case of a \emph{full} vertex (in- and out-edge), where $w\bowtie z$ is immediate ($wzw=a_i(a_{j_0}za_{j_0})a_i=a_im_{j_0}a_i=w$), and for such $E$ with every vertex on an edge the argument of Steps 3 and 5 was the first version of the theorem.

(2) The $2$-cycle-plus-spare required the treatment of isolated indices (Step~4). There is no projection algebra of size $\le12$ with $a_1\bowtie b_2$, $a_2\bowtie b_1$, (H1), (H2), $w\ne z$ and none of the pairs $(a_1,b_3),(a_2,b_3),(a_3,b_1),(a_3,b_2)$ matched; there is exactly one of size $13$, $45$ of size $14$ and $2054$ of size $15$ (\texttt{square.in}, isofiltered). All of them violate the hexagon condition in the same way: the walk from $a_1$ closes after a single turn, $a_1\to b_3a_1b_3\to a_2\to b_1\to a_3b_1a_3\to b_2\to a_1$, and $\{a_1,a_2,a_3b_1a_3\mid b_1,b_2,b_3a_1b_3\}$ is a hexagon with $w\ne z$ generating a copy of $\P(Q_0)$ (\texttt{sqfilter2.py}; the size-$13$ algebra is $\P(Q_0)$ with $a_3,b_3$ adjoined above it, and $1934$ of the size-$15$ algebras contain it). The elements $a_3b_1a_3=\theta_{a_3}(b_1)$ and $b_3a_1b_3=\theta_{b_3}(a_1)$ appearing here are exactly the replacements $a_3'$, $b_3'$ of Step~4, and in $49$ of the models $a_3za_3<m_3$, which showed that the transport relations themselves cannot be recovered at an isolated index; the identity $a_3za_3=b_3wb_3=m_3m_1m_3$ and the computation of $a_3'b_3'a_3'$ were first checked on these models (\texttt{sqstar.py}) and then proved by hand.

(3) With four indices the full-vertex condition is no longer the boundary: the pattern $\{(1,2),(3,4)\}$ of two \emph{disjoint} matched pairs has no full vertex, yet admits no counterexample in $\P_n$ for $n\le4$ nor in the $57$-element projection algebra of $\P_4/\mu_1$, which violates the hexagon condition (\texttt{nofull.py}, \texttt{nofull3.py}, \texttt{p4mu\_disjoint.py}), and the walk argument is unavailable for it as it stands, since a closed walk alternating between out-vertices $\{1,3\}$ and in-vertices $\{2,4\}$ only produces the $k=2$ square. Mace4 finds a projection algebra with $w\ne z$ for this pattern (the $14$-element algebra of Remark~\ref{rem:2k}(3), in which all twelve cross pairs are matched, and, excluding the extra pairs that would create a full vertex, a unique one of size $18$ with $a_1\bowtie b_4$, $a_3\bowtie b_2$ as well; \texttt{disjoint.in}, \texttt{disjoint2.in}), so the pattern is not a quasi-identity of all projection algebras, but both models violate the hexagon condition and in both $w\bowtie z$ and all sixteen transport relations hold. Prover9 then proved $wzw=w$ from the two disjoint pairs in $0.12$ seconds (\texttt{disjoint3.in}, a $165$-step proof); Steps 1 and 2 of the proof above are the hand version, and they apply verbatim whenever $E$ has two tails and two heads, which is what unifies the full-vertex and the disjoint-pair cases. Patterns with fewer than two tails or two heads are the stars, which fail in $\P_2$ by (b).

(4) The non-star conditions do not follow from the hexagon condition in infinite projection algebras (Theorem~\ref{thm:hinf}), but they do follow from the triangle condition, in every projection algebra (Theorem~\ref{thm:master}), and hence hold in $\P_X$ for every $X$. A direct combinatorial proof in $\P_X$ would have to handle the case of Step~3 of the proof of Theorem~\ref{thm:hex} in which the bad index is the middle vertex of a path; the algebraic reduction to the triangle avoids it.
\end{remark}

\subsection{The infinite case: the unrolled hexagon}\label{subsec:hinf}

The finiteness hypothesis of Theorems~\ref{thm:equiv} and~\ref{prop:path} cannot be dropped. The counterexample is obtained by ``unrolling'' the hexagon of $\P(Q_0)$ along its cycle of matchings $a_1\bowtie b_2\bowtie a_3\bowtie b_1\bowtie a_2\bowtie b_3\bowtie a_1$: each vertex is replaced by a descending chain, all matchings are taken level by level except the last one, which shifts the level by one.

Recall the structure of $\P(Q_0)$ (Section~\ref{subsec:witness}): it consists of six projections $a_1,a_2,a_3,b_1,b_2,b_3$ and five \emph{low} elements $w,z,m_1,m_2,m_3$ (the images of the five partitions of $\{1,2,3\}$), with $a_ia_j=w$, $b_ib_j=z$ ($i\ne j$), $a_ib_i=b_ia_i=m_i$, $a_i\bowtie b_j$ ($i\ne j$); $\theta_\ell(x)=\ell$ for every low $\ell$ and every $x$ (so the low elements form a right-zero subalgebra, $\ell*\ell'=\ell'$, and any two of them are matched); and $\theta_{a_i}(w)=w$, $\theta_{a_i}(z)=m_i$, $\theta_{a_i}(m_i)=m_i$, $\theta_{a_i}(m_j)=w$ ($j\ne i$), dually $\theta_{b_i}(z)=z$, $\theta_{b_i}(w)=m_i$, $\theta_{b_i}(m_i)=m_i$, $\theta_{b_i}(m_j)=z$.

\begin{definition}\label{def:hinf}
The \emph{unrolled hexagon} is the set
\[
H_\infty=\{w,z,m_1,m_2,m_3\}\ \cup\ \{e_{i,n},f_{i,n}: i\in\{1,2,3\},\ n\in\N\}
\]
with the operation $x*y=\theta_y(x)$ defined as follows, where $\ell$ denotes a low element, $i,j\in\{1,2,3\}$, and $M=\{(1,2),(3,2),(3,1),(2,1),(2,3)\}$ is the set of \emph{level-preserving} pairs (all $(i,j)$ with $i\ne j$ except the \emph{crossing} pair $(1,3)$):
\begin{enumerate}[label=\textup{(\roman*)},nosep]
\item $\theta_\ell(x)=\ell$ for all $x$; $\theta_{e_{i,n}}(\ell)$ and $\theta_{f_{i,n}}(\ell)$ are $\theta_{a_i}(\ell)$ and $\theta_{b_i}(\ell)$ as in $\P(Q_0)$ (independent of $n$);
\item $\theta_{e_{i,m}}(e_{i,n})=e_{i,\max(n,m)}$, $\theta_{f_{i,m}}(f_{i,n})=f_{i,\max(n,m)}$; $\theta_{e_{j,m}}(e_{i,n})=w$ and $\theta_{f_{j,m}}(f_{i,n})=z$ for $i\ne j$;
\item $\theta_{f_{i,m}}(e_{i,n})=\theta_{e_{i,m}}(f_{i,n})=m_i$;
\item for $(i,j)\in M$: $\theta_{f_{j,m}}(e_{i,n})=f_{j,\max(n,m)}$ and $\theta_{e_{i,m}}(f_{j,n})=e_{i,\max(n,m)}$;
\item for the crossing pair: $\theta_{f_{3,m}}(e_{1,n})=f_{3,\max(m,n-1)}$ and $\theta_{e_{1,m}}(f_{3,n})=e_{1,\max(m,n+1)}$.
\end{enumerate}
\end{definition}

Thus $e_{i,0}>e_{i,1}>\cdots$ and $f_{i,0}>f_{i,1}>\cdots$ are chains (``level'' $n$), $w,m_i\le e_{i,n}$ and $z,m_i\le f_{i,n}$, no low element is below another one and no chain element is below a low one. Two chain elements are matched iff they are equal, or of the form $e_{i,n},f_{j,n}$ with $(i,j)\in M$, or $e_{1,n+1},f_{3,n}$; a low element is matched with every low element and with no chain element. Two elements commute iff they are on the same chain, or are $e_{i,n},e_{j,m}$ (product $w$), $f_{i,n},f_{j,m}$ ($z$), $e_{i,n},f_{i,m}$ ($m_i$), or a low element and an element above it. Level $0$ is a copy of $\P(Q_0)$ with the matching $a_1\bowtie b_3$ broken in one direction: $\theta_{f_{3,0}}(e_{1,0})=f_{3,0}$ but $\theta_{e_{1,0}}(f_{3,0})=e_{1,1}<e_{1,0}$.

\begin{theorem}\label{thm:hinf}
\begin{enumerate}[label=\textup{(\alph*)}]
\item $H_\infty$ is a projection algebra satisfying \eqref{eq:Q1}.
\item $H_\infty$ satisfies the hexagon condition and the matching-pattern condition of every non-star digraph on $k\ge4$ vertices.
\item $(e_{1,0},e_{2,0},e_{3,0}\mid f_{1,0},f_{2,0},f_{3,0})$ is a triangle configuration with $w\ne z$. More generally, $H_\infty$ violates the matching-pattern condition of every digraph on three vertices other than the hexagon, and of every star.
\end{enumerate}
Consequently the hexagon condition does not imply the triangle condition, nor any other non-star condition on three indices, in the variety of projection algebras; and a regular $*$-semigroup $S$ with $\P(S)\cong H_\infty$ (which exists by Imaoka's theorem) satisfies \eqref{eq:Q1} and the hexagon condition but embeds in no partition monoid.
\end{theorem}

\begin{proof}
(a) By construction, the \emph{type} of a product (low, or which chain) depends only on the types of the factors, and the level of a product of two chain elements of levels $n,m$ is $\max(n+c,m+c')$ with $c,c'\in\{-1,0,1\}$, while a product involving a low element is low. Hence for a term $t(x_1,\dots,x_r)$ of depth $d$ and an assignment of types to the variables, either $t$ is low for all levels, or its value lies on a fixed chain at level $\max_{v\in S}(n_v+c_v)$ for some set $S$ of the variables and constants $|c_v|\le d$. Two such expressions of the same type agree on $\N^r$ as soon as they agree on $\{0,\dots,2d+1\}^r$, since evaluating at $n_v=2d+1$, $n_u=0$ ($u\ne v$) tells whether $v\in S$ and recovers $c_v$. The axioms (P1)$'$--(P5)$'$ have depth $\le4$, so it suffices to check them on all triples of level $\le9$; this was done by computer (\texttt{hinf2.py}, levels $\le12$). For \eqref{eq:Q1}: by the description of the order, two elements with a common lower bound are comparable, or of the form $e_{i,n},f_{i,m}$, or $e_{i,n},e_{j,m}$, or $f_{i,n},f_{j,m}$, or a low element and an element above it, and all these pairs commute.

(b) We first classify the tuples $(x_1,\dots,x_k)$, $k\ge2$, in which all products $x_i*x_j$ ($i\ne j$, both orders) are equal to one element $w'$, using the description of commuting pairs: ($\alpha$) all $x_i$ on one chain, $w'$ on that chain; ($\beta$) each $x_i$ is $w$ or an $e$-element, the $e$-elements on pairwise distinct chains, $w'=w$; ($\beta'$) dually with $z$ and $f$-elements; ($\gamma$) each $x_i$ is $m_p$ or on $e_p$ or on $f_p$, at most one on $e_p$ and at most one on $f_p$, $w'=m_p$. (Two low elements have equal products in both orders only if they are equal; a low $\ell$ and a chain element $c$ only if $\ell\le c$; and $e_{i,n},f_{j,m}$ with $i\ne j$ never.)

\emph{The hexagon.} Let $(x_1,x_2,x_3\mid y_1,y_2,y_3)$ satisfy the hexagon hypotheses with products $w',z'$, and suppose $w'\ne z'$. If some $x_i$ is low, then $w'=x_i$, the $y_j$ ($j\ne i$) are low, being matched with $x_i$, so $z'=y_j$, and $x_l$ ($l\ne i$) is low, being matched with some $y_j$, $j\notin\{i,l\}$; so $x_1=x_2=x_3=w'$ and $y_j=z'$ for $j\ne i$; finally $y_i$ is matched with the low element $x_l$, so $y_i$ is low, $y_i=z'$, and (H2) at $i$ gives $w'=z'$. So all six are chain elements. Case ($\gamma$) is then impossible (three chain elements in $e_p\cup f_p$ with at most one on each). In case ($\alpha$), say all $x_i\in e_p$: the elements matched with an element of $e_p$ are that element itself and the $f_{q,m}$ with $(p,q)\in M\cup\{(1,3)\}$, $q\ne p$; if $y_j=f_{q,m}$, then by (H2) $y_j$ commutes with $x_j\in e_p$, which requires $q=p$; so every $y_j$ lies on $e_p$ and equals both $x_i$, $i\ne j$, whence all six elements coincide. In case ($\beta$), $x_i=e_{\sigma(i),n_i}$ with $\sigma$ injective and $w'=w$. A $y_j$ equal to some $x_i$ would be matched with the other $x_{i'}$, on a different $e$-chain, impossible; so $y_j=f_{\tau(j),m_j}$ with $(\sigma(i),\tau(j))\in M\cup\{(1,3)\}$ for all $i\ne j$. If $\tau$ is not injective, $z'$ is on an $f$-chain, all $\tau(j)$ are equal, $=q$, and the three values $\sigma(i)$ lie in the two-element set $\{i:(i,q)\in M\cup\{(1,3)\}\}$, impossible; so $\tau$ is a permutation and $z'=z$. By (H2), $x_i$ and $y_i$ commute, which fails if $\sigma(i)\ne\tau(i)$; so $\sigma=\tau$. The six matchings $x_i\bowtie y_j$ then involve each ordered pair $(\sigma(i),\sigma(j))$, $i\ne j$, exactly once; the five level-preserving ones give $n_i=m_j$ and the crossing one gives $n_{i_0}=m_{j_0}+1$. The five equalities connect all six levels ($K_{3,3}$ minus a perfect matching is a $6$-cycle, and deleting one more edge leaves a path), so all levels are equal, a contradiction. Case ($\beta'$) is dual.

\emph{Non-star patterns on $k\ge4$ vertices.} Every non-star digraph contains two disjoint edges, a directed path of length two, or a $2$-cycle, and its condition on $[k]$ follows from the condition of $\{(1,2),(3,4)\}$, $\{(1,2),(2,3)\}$ or $\{(1,2),(2,1)\}$ on $[4]$ (restrict to four indices containing the two edges). In each of these patterns the two edges have distinct tails and distinct heads. Suppose $(x_1,\dots,x_4\mid y_1,\dots,y_4)$ satisfies (H1)--(H3) for one of them with $w'\ne z'$. Since $k=4>3$, in case ($\beta$) some $x_i=w$, and in case ($\gamma$) at least two $x_i$ equal $m_p$. Recall that the elements matched with a chain element $c$ are $c$ itself and chain elements on a chain of the other letter, and that the elements commuting with $e_{p,n}$ are those of $e_p$, $f_p$, of the other $e$-chains, and $w,m_p$.

\emph{Case ($\alpha$), all $x_i\in e_p$.} The head $y_b$ of each edge is matched with its tail $x_a\in e_p$ and commutes with $x_b\in e_p$, hence $y_b=x_a$: the $y$'s contain two elements of $e_p$ (the two tails, of distinct indices). This rules out cases ($\beta$), ($\gamma$) (at most one element on $e_p$) and ($\beta'$) for the $y$'s, and case ($\alpha$) on a chain other than $e_p$. So all $y_j\in e_p$, $w'=e_{p,N}$, $z'=e_{p,N'}$, at most one $x$ has level $<N$ and at most one $y$ has level $<N'$. One of the two heads has level $N'$ and equals a tail, whose level is $\le N$; so $N'\le N$, and symmetrically $N\le N'$; so $w'=z'$.

\emph{Case ($\beta$), some $x_i=w$.} Then $y_i$ is $w$ or an $e$-element. So the $y$'s are not in case ($\beta'$); in case ($\beta$), $z'=w=w'$. In case ($\alpha$) they lie on one chain $e_q$; each tail is matched with its head, an element of $e_q$, hence equals it, and two tails on one $e$-chain contradict ($\beta$). In case ($\gamma$) they lie in $e_q\cup f_q\cup\{m_q\}$ with at most one on $e_q$; every $x_i=w$ needs $y_i\in e_q$, so exactly one $x_i$ is $w$ and the others lie on the three distinct $e$-chains, and the two of them not on $e_q$ need $y$'s in $e_q\cup f_q\cup\{m_q\}$ commuting with them, i.e.\ again on $e_q$; impossible.

\emph{Case ($\gamma$), at least two $x_i=m_p$.} An element commutes with $m_p$ iff it is above $m_p$, i.e.\ lies in $e_p\cup f_p\cup\{m_p\}$. So the $y$'s are not in case ($\gamma$) for $q\ne p$ (and for $q=p$, $z'=m_p=w'$), nor in cases ($\beta$), ($\beta'$) (two of them would lie on $e_p$, resp.\ $f_p$). In case ($\alpha$) they lie on $e_p$ or on $f_p$; the tails equal their heads, so two $x$'s lie on that chain, contradicting ($\gamma$).

The cases ($\alpha$) on $f_p$ and ($\beta'$) are dual. Hence $w'=z'$.

(c) The level-$0$ vertices satisfy (H1), (H2), $e_{1,0}\bowtie f_{2,0}$, $e_{2,0}\bowtie f_{3,0}$, $e_{3,0}\bowtie f_{1,0}$ (all in $M$), and $w\ne z$. For $E$ on $[3]$ missing the ordered pair $(p,q)$, choose a permutation $\sigma$ with $(\sigma(p),\sigma(q))=(1,3)$ and put $x_i=e_{\sigma(i),0}$, $y_j=f_{\sigma(j),0}$: all matchings required by $E$ are in $M$, hence hold at level $0$. For an out-star with tail $1$ on $[k]$ take $x_1=e_{1,0}$, $x_i=w$ ($i\ge2$), and all $y_j=e_{1,0}$: then $w'=w$, $z'=e_{1,0}$, all required pairs commute, and $x_1\bowtie y_j$ trivially. In-stars are dual, and the empty pattern is contained in a star.

The final assertion follows because $\P(\P_X)$ satisfies the triangle condition for every $X$ (Theorem~\ref{thm:hex}), and the projection algebra of a $*$-subsemigroup of $\P_X$ is a subalgebra of $\P(\P_X)$.
\end{proof}

\begin{remark}\label{rem:hinf}
(1) Together with Theorems~\ref{thm:equiv}, \ref{prop:path} and $\P(Q_0)$ this gives the complete picture of the implications between the matching-pattern conditions: on finite projection algebras the hexagon condition is equivalent to every non-star condition on three indices and implies every non-star condition on more indices, but not conversely ($\P(Q_0)$ satisfies all the conditions on $k\ge4$ indices); on arbitrary projection algebras the non-star conditions on three indices other than the hexagon are strictly stronger than the hexagon condition ($H_\infty$). On arbitrary projection algebras the triangle, path and $2$-cycle conditions are equivalent and imply all the others (Theorem~\ref{thm:master}). Whether the hexagon condition implies the conditions on $k\ge4$ indices in general is open; an ``unrolled octagon'' would separate them, but there the products of non-adjacent vertices have to be new elements (Remark~\ref{rem:2k}(3)) and we have not constructed it. All these conditions were verified on $H_\infty$ by computer as well (\texttt{hinf2\_run.py}, \texttt{hinf2\_pat.py}).

(2) In $H_\infty$ the walk of Section~\ref{subsec:equiv} from $e_{1,0}$ closes after one turn, $e_{1,0}\to f_{3,0}\to e_{2,0}\to f_{1,0}\to e_{3,0}\to f_{2,0}\to e_{1,0}$, but its consecutive elements $e_{1,0},f_{3,0}$ are not matched: in the semigroup $S$ the element $s=e_{1,0}f_{3,0}e_{2,0}f_{1,0}e_{3,0}f_{2,0}$ has $s^*s=e_{1,0}$ and $ss^*=e_{1,1}<e_{1,0}$. So it is the Dedekind-finiteness used in Lemma~\ref{lem:stab}, not the descending chain condition, that fails here. Mace4 finds finite models of the triangle hypotheses with $w\ne z$ in which the chain $a_1>\Phi(a_1)>\Phi^2(a_1)$ descends strictly for two turns before closing (sizes $13$ and $21$; three turns need at least $30$ elements; \texttt{chain.py}), which suggests that an example of the other kind, with an infinite strictly descending walk, exists as well.

(3) The non-star conditions all hold in $\P_X$ for every $X$ (Theorem~\ref{thm:master}); $H_\infty$ is not a partition monoid. Before this was known, the natural place to look for counterexamples was among \emph{periodic} projections of $\P_{\Z\times[m]}$, invariant under the level shift $(n,i)\mapsto(n+1,i)$, since these can carry the level-shifting matchings of $H_\infty$ inside an honest partition monoid. Such a projection is described by a partition of $[m]$ into types $c$, each with a modulus $d_c\ge0$ and a potential $\phi_c:c\to\Z$ (mod $d_c$) up to a constant -- its blocks are $\{(n+\phi(i),i):i\in c\}$ for $d_c=0$ and $\{(n,i):i\in c,\ n\equiv\phi(i)+r\ (\mathrm{mod}\ d_c)\}$ for $d_c>0$ -- together with a $\T$/$\N$ label per type, and the products $\theta_p(q)$ are computed exactly by a union-find with potentials on $[m]$ (\texttt{periodic.py}; the description of $\theta_p(q)$ in Lemma~\ref{lem:PX}(a) makes this a finite computation). Exhaustive search over all periodic projections with $m\le3$ points per level, potentials in $\{-1,0,1\}$ and moduli $\le2$ ($460$ projections for $m=3$), and with $m=4$, potentials in $\{-1,0,1\}$ and moduli $\le2$ ($5064$ projections), finds no counterexample to the path condition, the $2$-cycle condition, or (for $m\le3$) the two-disjoint-pairs, $K_{2,2}$ and $C_4$ conditions, while reproducing the counterexamples to the star patterns and, as a check, none to the triangle and hexagon conditions (\texttt{periodic3.py}). These searches were run before Theorem~\ref{thm:master} was found; that theorem explains their outcome, and the direct combinatorial attack on Step~3 of the proof of Theorem~\ref{thm:hex} at the middle vertex of a path (where one only obtains that the offending $\N(z)$-block $N\subseteq B_2$ is a union of blocks that are simultaneously $\T(a_1)$- and $\T(b_2)$-blocks, a union of $\T(a_3)$-blocks, and contained in a single $\T(a_2)$-block) is no longer needed.
\end{remark}

\subsection{The witness}\label{subsec:witness}

\begin{theorem}\label{thm:witness}
Let $\rho$ be the relation on $\P_3$ defined by $\alpha\,\rho\,\beta$ iff $\alpha=\beta$, or $\rk\alpha,\rk\beta\le1$ and $\ker\alpha=\ker\beta$, $\operatorname{coker}\alpha=\operatorname{coker}\beta$, or $\rk\alpha=\rk\beta=2$ and $\alpha\,\Hh\,\beta$. Then $\rho$ is a $*$-congruence, $Q:=\P_3/\rho$ has $67$ elements and satisfies \eqref{eq:Q1}, and $Q$ embeds in no partition monoid. The same holds for the $61$-element $*$-subsemigroup $Q_0$ of $Q$ generated by the six projections of rank $2$, and for the $103$-element quotient $\P_3/\mu$ where $\mu$ identifies only the elements of rank $\le1$ with equal kernel and cokernel.
\end{theorem}

\begin{proof}
That $\rho$ and $\mu$ are $*$-congruences, that the quotients satisfy \eqref{eq:Q1}, and that $Q_0$ has $61$ elements were checked by direct computation (see Section~\ref{sec:comp}; $\rho$ and $\mu$ are among the ten $*$-congruences of $\P_3$). Write $a_i$ for the image of the rank-$2$ projection of $\P_3$ whose transversal blocks are $\{i\}$ and $\{1,2,3\}\setminus\{i\}$, and $b_i$ for the image of the rank-$2$ projection whose transversal blocks are the two singletons $\{j\}$, $\{k\}$ ($\{i,j,k\}=\{1,2,3\}$) and whose nontransversal block is $\{i\}$. In $Q$ (and already in $\P_3/\mu$) the hypotheses (H1)--(H3) hold with $E$ the set of all six pairs (so in particular for the triangle): $w$ is the class of the partition of $\{1,2,3\}$ into one block, $z$ that of the discrete partition, $m_i$ that of the partition $\{i\}\mid\{j,k\}$, and the products $a_ib_j$ ($i\ne j$) have rank $2$; all this is a direct verification in $\P_3$ (also done by computer). But $w\ne z$, since the two partitions are different. By Theorem~\ref{thm:hex} no $*$-homomorphism $Q\to\P_X$ is injective on $\{w,z\}$.
\end{proof}

\begin{remark}
(1) The obstruction lives in the projection algebra: the eleven-element projection algebra of $Q_0$ (six projections of rank $2$ and the five ``partition'' projections $w,z,m_1,m_2,m_3$) satisfies the hypotheses of Theorem~\ref{thm:hex} with $w\ne z$, so it is not a subalgebra of the projection algebra of any $\P_X$. Both known obstructions, \eqref{eq:Q1} and the hexagon condition, are therefore statements about projection algebras. \emph{(Added 11 Oct 2026.)} $\P(Q_0)$ is a minimal non-embeddable projection algebra, and of smallest possible order: its six maximal proper subalgebras (omit one of the $a_i$, $b_i$) all embed in $\P(\P_4)$ (SAT, \texttt{pa8/w3/subsQ0.py}), and a Mace4 search by M.K.\ shows that no projection algebra of order $\le10$ satisfies \eqref{eq:Q1} and violates the triangle condition, the smallest one of order $11$ being isomorphic to $\P(Q_0)$ (\texttt{pa8/w3/model11b.json}). This also settles the entry ``minimality not checked'' of the status table.

(2) Concretely, $Q_0=D\cup(\Part_3\times\Part_3)$ where $\Part_3$ is the set of five partitions of $\{1,2,3\}$, $\Part_3\times\Part_3$ is a $5\times5$ square band (the image of the ideal of rank $\le1$, via $\alpha\mapsto(\ker\alpha,\operatorname{coker}\alpha)$), and $D$ is the rank-$2$ $\D$-class of $\P_3$ with its $\Hh$-classes collapsed, a $6\times6$ ``trace'' with trivial groups whose zero pattern is the $6$-cycle $a_1b_2a_3b_1a_2b_3$ of matched pairs. The elements of $D$ act on $\Part_3$ by $\varepsilon\mapsto\ker(x\alpha)$ for any $\alpha$ with $\ker\alpha=\varepsilon$. Every proper $*$-subsemigroup of $Q_0$ generated by five of the six rank-$2$ projections embeds in $\P_3$; all six are needed.

(3) Computer searches confirm that $\P_3/\mu$ has no $*$-embedding into $\P_n$ for $n\le7$, that $Q_0$ has none for $n\le5$, and that the hypotheses of Theorem~\ref{thm:hex} have no solution with $w\ne z$ among the projections of $\P_n$ for $n\le4$ (a check of the theorem, not needed for it).

(4) The congruence $\mu$ is the $*$-invariant congruence of $\P_3$ that collapses the ideal of rank $\le1$ onto the square band $\Part_3\times\Part_3$; in the classification of the congruences of partition monoids (East--Ruškuc; East--Mitchell--Ruškuc--Torpey) it is one of the congruences below the Rees congruence of that ideal, and $\rho$ is its join with the ``$\Hh$-relation on rank $2$''. For $n=4$ the analogous quotients contain those of $\P_3$ and give nothing new (Section~\ref{sec:comp}); the further obstruction of Section~\ref{sec:wiring} was found instead by enumerating small projection algebras.
\end{remark}

\section{The conditions are not sufficient: a wiring obstruction}\label{sec:wiring}

In this section we show that \eqref{eq:Q1} together with the triangle condition, and indeed together with all the matching-pattern conditions of Section~\ref{sec:hex}, is not sufficient for embeddability, even for finite regular $*$-semigroups. The obstruction is a quasi-identity of a new kind: its hypotheses involve matched pairs and commuting pairs at the same time, and its validity in $\P_X$ rests on the fact that commuting projections share their $\N$-blocks (Lemma~\ref{lem:PX}(b)), so that matched pairs to a commuting pair are ``wired'' through the same nontransversal blocks.

\subsection{Graph semigroups}\label{subsec:graphsg}

Let $G=(V,E)$ be a graph, $B$ a set and $\bot\colon V\to B$ a map; extend $\bot$ to $V\sqcup B$ by $\bot(m)=m$ for $m\in B$, and write $v\bowtie v'$ for $v=v'$ or $vv'\in E$. On the set
\[
Q(G,\bot)=(V\times V)\sqcup(B\times B)
\]
define $(x,y)^*=(y,x)$ and
\[
(p,q)(r,s)=\begin{cases}(p,s)&\text{if }q,r\in V\text{ and }q\bowtie r,\\ (\bot p,\bot s)&\text{otherwise.}\end{cases}
\]

\begin{lemma}\label{lem:graphsg}
$Q(G,\bot)$ is a regular $*$-semigroup with $|V|^2+|B|^2$ elements. Its projections are the pairs $(v,v)$, $v\in V$, and $(m,m)$, $m\in B$; identifying them with $V\sqcup B$, the projection algebra is
\[
\theta_v(v')=\begin{cases}v&v\bowtie v'\\ \bot v&\text{otherwise,}\end{cases}\qquad \theta_v(m)=\bot v,\qquad \theta_m(x)=m\qquad(v,v'\in V,\ m\in B,\ x\in V\sqcup B).
\]
Its matched pairs of top projections are the edges of $G$, and it satisfies \eqref{eq:Q1} iff $\bot$ is a proper colouring of $G$ (adjacent vertices have distinct values).
\end{lemma}
\begin{proof}
For $x=(p,q)$, $y=(r,s)$, $z=(t,u)$ one checks the four cases according to whether $q\bowtie r$ and $s\bowtie t$: if both hold, $(xy)z=(p,u)=x(yz)$; in the other three cases both products equal $(\bot p,\bot u)$, because a pair with a coordinate in $B$ is never matched. The identities $(xy)^*=y^*x^*$ and $xx^*x=x$ are immediate ($q\bowtie q$ for all $q$). The projections are the symmetric pairs, and $\theta_{(v,v)}(v',v')=(v,v')(v,v)=(v,v)$ if $v\bowtie v'$, while otherwise it is $(\bot v,\bot v')(v,v)=(\bot v,\bot v)$; the other values are computed in the same way. The projections below $(v,v)$ are $(v,v)$ and $(\bot v,\bot v)$, and $(m,m)$ has no other projection below it; so two distinct top projections $v,v'$ have a common lower bound iff $\bot v=\bot v'$, and then $(v,v)(v',v')=(v,v')\ne(v',v)=(v',v')(v,v)$ if $vv'\in E$, while the two products are both $(\bot v,\bot v)$ if $vv'\notin E$.
\end{proof}

Thus $Q(G,\bot)$ is the Rees matrix semigroup over the trivial group with sandwich matrix ``adjacency plus identity'' of $G$, in which the zero has been replaced by the rectangular band $B\times B$ via $\bot$. For $|V|\le3$ and every proper colouring $\bot$ the semigroup $Q(G,\bot)$ embeds in $\P_4$ (computer search); our counterexample is the smallest graph for which this fails, and Theorem~\ref{thm:graphsg} in Section~\ref{subsec:graphthm} determines exactly which $Q(G,\bot)$ embed.

\begin{definition}\label{def:Qpath}
Let $\mathfrak Q=Q(G,\bot)$ where $G$ is the path $a-c-b-d$ on $V=\{a,b,c,d\}$, $B=\{w,z\}$, $\bot(a)=\bot(b)=w$ and $\bot(c)=\bot(d)=z$; so $|\mathfrak Q|=20$. Its projection algebra $\mathfrak P=\{a,b,c,d,w,z\}$ has the matched pairs
\[
a\bowtie c,\qquad b\bowtie c,\qquad b\bowtie d,\qquad w\bowtie z,
\]
and every other value of $\theta$ is ``the bottom of the subscript'': $\theta_a(b)=\theta_b(a)=w$, $\theta_c(d)=\theta_d(c)=z$, $\theta_a(d)=w$, $\theta_d(a)=z$, $\theta_a(z)=\theta_b(z)=w$, $\theta_c(w)=\theta_d(w)=z$, and $\theta_w\equiv w$, $\theta_z\equiv z$.
\end{definition}

The colouring is proper, so $\mathfrak Q$ satisfies \eqref{eq:Q1}. The triangle condition holds in $\mathfrak P$ (a direct search over the $6^6$ assignments), hence by Theorem~\ref{thm:master} so do all the non-star matching-pattern conditions; in fact the same search shows that the hexagon, cyclic octagon, directed $4$-cycle and double $2$-cycle conditions hold in $\mathfrak P$, and so do the $2$-index square and $3$-index star conditions, which fail in $\P_2$.

\subsection{The obstruction}

\begin{theorem}\label{thm:wiring}
Let $a,b,c,d$ be projections of a partition monoid $\P_X$ such that
\[
a\bowtie c,\qquad b\bowtie c,\qquad b\bowtie d,\qquad cd=dc=:z,\qquad \theta_a(z)=\theta_b(z)=\theta_a(d).
\]
Then $z=c$ and $a=b$.
\end{theorem}

This is a quasi-identity of projection algebras: $cd=dc$ is equivalent to $\theta_c(d)=\theta_d(c)$ in every regular $*$-semigroup (if $cdc=dcd=:e$ then $cd\in\Hh_e$ and $(cd)^2=cdcd=ed=e$, so $dc=(cd)^*=(cd)^{-1}=cd$), and $p\bowtie q$ means $\theta_p(q)=p$, $\theta_q(p)=q$. In the binary notation $x*y=\theta_y(x)$ of Lemma~\ref{lem:mono}, the theorem reads
\[
c*a=a,\ a*c=c,\ c*b=b,\ b*c=c,\ d*b=b,\ b*d=d,\ d*c=c*d,\ (d*c)*a=(d*c)*b=d*a\ \Longrightarrow\ a=b,
\]
and it is independent of the conditions found so far: it is not implied by \eqref{eq:Q1} and the triangle condition ($\mathfrak P$), and together with \eqref{eq:Q1} it does not imply the triangle condition, since it holds in the projection algebra $\P(Q_0)$ of Theorem~\ref{thm:witness} and in the Mace4 models of Remark~\ref{rem:path} (a direct check of all quadruples). It does not imply \eqref{eq:Q1} either (it holds in the $5$-element $Q(K_2,\bot)$ with constant $\bot$).

Before the proof we note the consequence.

\begin{corollary}\label{cor:wiring}
The projection algebra $\mathfrak P$ does not embed in $\P(\P_X)$, and the $20$-element regular $*$-semigroup $\mathfrak Q$ does not $*$-embed in $\P_X$, for any set $X$. Consequently \eqref{eq:Q1} together with the triangle condition (or with all the matching-pattern conditions of Section~\ref{sec:hex}) is not sufficient for a finite regular $*$-semigroup to embed in a partition monoid; the smallest counterexample has order between $9$ and $20$ (between $9$ and $19$ by Theorem~\ref{thm:R}).
\end{corollary}
\begin{proof}
In $\mathfrak P$ the hypotheses of Theorem~\ref{thm:wiring} hold, but $a\ne b$. The bound $9$ is Section~\ref{sec:comp}.
\end{proof}

The rest of the section proves Theorem~\ref{thm:wiring}. Recall from Lemma~\ref{lem:PX}(c) that a matched pair $p\bowtie q$ comes with a bijection $\T(p)\to\T(q)$, sending a $\T(p)$-block to the unique $\T(q)$-block in the same class of $\varepsilon_p\vee\varepsilon_q$; we write $\varphi\colon\T(a)\to\T(c)$, $\psi\colon\T(b)\to\T(c)$, $\chi\colon\T(b)\to\T(d)$ for the three bijections, and $w=\theta_a(z)$. For projections $p,q$ let $\Gamma(p,q)$ be the bipartite graph whose vertices are the blocks of $p$ and the blocks of $q$ (a set which is a block of both appears twice) and whose edges join a $p$-block to a $q$-block meeting it; the classes of $\varepsilon_p\vee\varepsilon_q$ are the unions of the connected components of $\Gamma(p,q)$, and Lemma~\ref{lem:PX}(c) says that for $p\bowtie q$ every component contains at most one $\T(p)$-block and at most one $\T(q)$-block, one of each if any. We say that a subgraph \emph{touches} a block if one of its vertices meets that block.

\smallskip
\emph{Step 1: the structure of $z$.} By Lemma~\ref{lem:PX}(b) applied to $cd=dc$, every class $K$ of $\varepsilon_c\vee\varepsilon_d$ is of one of the types (i) $K$ is a union of $\T(c)$-blocks and a union of $\T(d)$-blocks; (ii) $K$ is a union of $\T(c)$-blocks and a single $\N(d)$-block; (iii) $K$ is a union of $\T(d)$-blocks and a single $\N(c)$-block; (iv) $K$ is a single block of both. By Lemma~\ref{lem:PX}(a) the blocks of $z=\theta_c(d)$ are exactly these classes, the $\T$-blocks being those of type (i). In particular $\N(z)$ consists of the $\N(c)$-blocks (types (iii), (iv)) and the type (ii) classes, and symmetrically of the $\N(d)$-blocks and the type (iii) classes.

\smallskip
\emph{Step 2: transport.} Let $p\bowtie q$ with bijection $\varphi\colon\T(p)\to\T(q)$, and let $x\le q$. For a set $Y$ which is a union of $\T(q)$-blocks put $\varphi^{-1}[Y]=\bigcup\{\varphi^{-1}(C):C\in\T(q),\ C\subseteq Y\}$. We claim: the blocks of $\theta_p(x)$ are the $\N(p)$-blocks and the sets $\varphi^{-1}[Y]$ for the blocks $Y\subseteq\supp(q)$ of $x$, with the label of $Y$; and the class of $\varepsilon_p\vee\varepsilon_x$ containing $\varphi^{-1}[Y]$ is
\[
\Omega_Y=\varphi^{-1}[Y]\cup\textstyle\bigcup\{\text{$\N(p)$-blocks}\}=Y\cup\bigcup\{\text{$\N(q)$-blocks}\},
\]
these being all the classes meeting $\supp(p)$. Indeed, by Lemma~\ref{lem:PX}(a) the blocks of $x$ outside $\supp(q)$ are $\N(q)$-blocks and those inside are unions of $\T(q)$-blocks. The classes of $\varepsilon_p\vee\varepsilon_q$ meeting $\supp(p)$ are the sets $A\cup(\N(p)\text{-blocks})=\varphi(A)\cup(\N(q)\text{-blocks})$, $A\in\T(p)$, each containing exactly one $\T(q)$-block. Passing to $\varepsilon_p\vee\varepsilon_x\ge\varepsilon_p\vee\varepsilon_q$, a block $Y\subseteq\supp(q)$ of $x$ merges the classes of the $\T(q)$-blocks it contains, and the other blocks of $x$ lie inside single classes; since each class contains one $\T(q)$-block, which lies in one $Y$, the classes of $\varepsilon_p\vee\varepsilon_x$ meeting $\supp(p)$ are the $\Omega_Y$. Now Lemma~\ref{lem:PX}(a) gives the blocks $\Omega_Y\cap\supp(p)=\varphi^{-1}[Y]$, transversal iff $\Omega_Y$ meets $\supp(x)$, i.e.\ (as $\Omega_Y\cap\supp(q)=Y$) iff $Y$ is a $\T$-block of $x$.

Applied to $(a,c,z)$, $(b,c,z)$ and $(b,d,z)$, this describes $w$ in three ways: its blocks are the $\N(a)$-blocks and the sets $\varphi^{-1}[K]$, $K$ of type (i) or (ii); the $\N(b)$-blocks and the sets $\psi^{-1}[K]$, $K$ of type (i) or (ii); the $\N(b)$-blocks and the sets $\chi^{-1}[K]$, $K$ of type (i) or (iii); in each case the block is transversal iff $K$ is of type (i).

\smallskip
\emph{Step 3: the common join.} Since $w\le a$, every block of $w$ is a union of $a$-blocks, and by Step~2 every block of $w$ lies in a class of $\varepsilon_a\vee\varepsilon_z$; hence $\varepsilon_a\vee\varepsilon_z=\varepsilon_w\vee\varepsilon_z$, and likewise for $b$. Write $\Pi=\varepsilon_a\vee\varepsilon_z=\varepsilon_b\vee\varepsilon_z$. By Step~2 the classes of $\Pi$ meeting $\supp(a)$, equivalently meeting $\supp(c)$ (a class containing $\varphi(A)$ contains $A$), are the sets $\Omega_K=K\cup(\N(c)\text{-blocks})=\varphi^{-1}[K]\cup(\N(a)\text{-blocks})=\psi^{-1}[K]\cup(\N(b)\text{-blocks})$ for $K$ of type (i), (ii); and the classes meeting $\supp(d)$ are $\Omega'_K=K\cup(\N(d)\text{-blocks})=\chi^{-1}[K]\cup(\N(b)\text{-blocks})$ for $K$ of type (i), (iii). Fix a class $\Omega$ of $\Pi$ meeting one of the four supports and write $S_p=\Omega\cap\supp(p)$ and $\N(p)|_\Omega$ for the set of $\N(p)$-blocks contained in $\Omega$. All classes of $\varepsilon_a\vee\varepsilon_c$, $\varepsilon_b\vee\varepsilon_c$, $\varepsilon_b\vee\varepsilon_d$, $\varepsilon_a\vee\varepsilon_d$ are contained in classes of $\Pi$, so all the joins may be computed inside $\Omega$, using only the blocks contained in $\Omega$.

\emph{(F1) If $\Omega=\Omega_K$ with $K$ of type \textup{(i)}}, then $\Omega=\Omega'_K$, so $S_c=S_d=K$ and $S_a=S_b=\varphi^{-1}[K]=\psi^{-1}[K]=\chi^{-1}[K]=:W$; moreover $\N(a)|_\Omega=\N(b)|_\Omega$ and $\N(c)|_\Omega=\N(d)|_\Omega$. Indeed, the blocks of $w$ inside $\Omega$ are, by the first description, the $\N(a)$-blocks in $\Omega$ and the single $\T$-block $\varphi^{-1}[K]$ (the sets $\varphi^{-1}[K']$, $K'\ne K$, lie in other classes); by the second and third descriptions they are the $\N(b)$-blocks in $\Omega$ and $\psi^{-1}[K]$, resp.\ $\chi^{-1}[K]$. Comparing gives the claims about $a$ and $b$. An $\N(c)$-block $E\subseteq\Omega$ is a block of $z$ of type (iii) or (iv); type (iii) is impossible, since $E\subseteq\Omega\setminus K$, which is a union of $\N(d)$-blocks, whereas a type (iii) class is a union of $\T(d)$-blocks. So $E$ is of type (iv), an $\N(d)$-block; symmetrically.

\emph{(F2) If $\Omega=\Omega_K$ with $K$ of type \textup{(ii)}}, then $\Omega$ contains exactly one class $K'$ of type (iii), and $\Omega=\Omega'_{K'}$; $S_c=K$ is an $\N(d)$-block, $S_d=K'$ is an $\N(c)$-block, and the remaining $\N(c)$-blocks in $\Omega$ are exactly the remaining $\N(d)$-blocks in $\Omega$ (the type (iv) classes inside $\Omega$). The sets $S_a=\varphi^{-1}[K]$ and $S_b=\psi^{-1}[K]=\chi^{-1}[K']$ are $\N$-blocks of $w$, and either $S_a=S_b$ and $\N(a)|_\Omega=\N(b)|_\Omega$, or $S_a\cap S_b=\emptyset$, $S_a\in\N(b)$, $S_b\in\N(a)$ and $\N(a)|_\Omega\setminus\{S_b\}=\N(b)|_\Omega\setminus\{S_a\}$. Indeed, $\Omega$ contains a $\T(b)$-block $B$ and hence $\chi(B)$, so it meets $\supp(d)$ and $\Omega=\Omega'_{K'}$ for some $K'$ of type (i) or (iii); as $K'\subseteq\Omega\setminus K\subseteq\bigcup\N(c)$, $K'$ is of type (iii), and it is the only type (iii) class in $\Omega$ because $\Omega\setminus K'$ is a union of $\N(d)$-blocks; $\Omega\cap\supp(c)=K$ excludes further classes of types (i), (ii). The $\N(c)$-blocks in $\Omega$ are the blocks of $z$ of types (iii), (iv) in $\Omega$, i.e.\ $K'$ and the type (iv) classes, and symmetrically for $d$. Finally the blocks of $w$ inside $\Omega$ are the $\N(a)$-blocks in $\Omega$ together with $\varphi^{-1}[K]$, and also the $\N(b)$-blocks together with $\psi^{-1}[K]$, and also the $\N(b)$-blocks together with $\chi^{-1}[K']$; a union of $\T(b)$-blocks is not an $\N(b)$-block, so $\psi^{-1}[K]=\chi^{-1}[K']$, and comparing the first two descriptions gives the alternative.

\smallskip
\emph{Step 4: the hypothesis $\theta_a(d)=w$.} By Lemma~\ref{lem:PX}(a), the blocks of $\theta_a(d)$ inside $\supp(a)$ are the sets $L\cap\supp(a)$, $L$ a class of $\varepsilon_a\vee\varepsilon_d$, transversal iff $L$ meets $\supp(d)$; and each $L$ lies in a class of $\Pi$. Since $\theta_a(d)=w$, we get from Step~2: in a class $\Omega$ of type (F1) all $\T(a)$-blocks contained in $W$ lie in one class of $\varepsilon_a\vee\varepsilon_d$; in a class of type (F2) all $\T(a)$-blocks contained in $S_a$ lie in one class $L$ of $\varepsilon_a\vee\varepsilon_d$, and $L$ does not meet $\supp(d)$. We show that both are impossible unless they are vacuous, i.e.\ that $|\T(c)\cap K|=1$ for $K$ of type (i) (Step~5) and that there is no $K$ of type (ii) (Step~6). Then $\T(z)=\T(c)$ and $\N(z)=\N(c)$ by Step~1, i.e.\ $z=c$; and then $w=\theta_a(c)=a$ and $w=\theta_b(c)=b$, so $a=b$, proving the theorem.

\smallskip
\emph{Step 5: in a class $\Omega$ of type \textup{(F1)}, distinct $\T(a)$-blocks lie in distinct classes of $\varepsilon_a\vee\varepsilon_d$.} Let $\mathcal N'=\N(a)|_\Omega=\N(b)|_\Omega$ and $\mathcal N=\N(c)|_\Omega=\N(d)|_\Omega$, and let $\mathcal G$ be the graph on $\mathcal N'\sqcup\mathcal N$ joining blocks that meet. Since the blocks of $\mathcal N'$ are $a$-blocks and $b$-blocks and those of $\mathcal N$ are $c$-blocks and $d$-blocks, $\mathcal G$ is a subgraph of each of $\Gamma(a,c)$, $\Gamma(b,c)$, $\Gamma(b,d)$, $\Gamma(a,d)$, so each component $\kappa$ of $\mathcal G$ lies in a single component of each of them. By Lemma~\ref{lem:PX}(c) for the three matched pairs, $\kappa$ touches at most one $\T(a)$-block, at most one $\T(b)$-block, at most one $\T(c)$-block and at most one $\T(d)$-block, and if it touches $A$ and $C$ then $C=\varphi(A)$; if $B$ and $C$ then $C=\psi(B)$; if $B$ and $D$ then $D=\chi(B)$. Moreover, if $\kappa$ touches a $\T(a)$-block $A$ then the point of contact lies in $W=\bigcup\T(b)$, so $\kappa$ touches a $\T(b)$-block, and conversely; likewise for $c$ and $d$, as $K=\bigcup\T(c)=\bigcup\T(d)$. Consequently, if $\kappa$ touches $A$ and $D$ then $D=\sigma(A)$, where $\sigma=\chi\psi^{-1}\varphi$: $\kappa$ touches some $B$ and some $C$, and $C=\varphi(A)=\psi(B)$, $D=\chi(B)$. The same holds for a point $x\in W\cap K$ and the blocks $A,B,C,D$ containing it, since meeting blocks lie in the same component.

Now let $A_1\ne A_2$ be $\T(a)$-blocks in the same class of $\varepsilon_a\vee\varepsilon_d$, and take a path in $\Gamma(a,d)$ from $A_1$ to $A_2$ with no $\T(a)$-block in its interior. The interior vertices are blocks of $\mathcal N'$, of $\mathcal N$ and $\T(d)$-blocks. A maximal interior segment without $\T(d)$-blocks is a connected subgraph of $\mathcal G$, hence lies in a component $\kappa$. If the path contains no $\T(d)$-block, a single $\kappa$ touches $A_1$ and $A_2$, which is impossible. Otherwise let $D_1,\dots,D_r$ be the $\T(d)$-blocks on the path in order. The segment from $A_1$ to $D_1$ is empty ($A_1$ meets $D_1$) or a $\kappa$ touching $A_1$ and $D_1$; either way $D_1=\sigma(A_1)$, and likewise $D_r=\sigma(A_2)$. The segment between $D_i$ and $D_{i+1}$ is a $\kappa$ touching both, so $D_i=D_{i+1}$. Hence $\sigma(A_1)=\sigma(A_2)$ and $A_1=A_2$, a contradiction.

\smallskip
\emph{Step 6: in a class $\Omega$ of type \textup{(F2)}, some $\T(a)$-block lies in a class of $\varepsilon_a\vee\varepsilon_d$ meeting $\supp(d)$.} Write $\sim$ for ``in the same class of $\varepsilon_a\vee\varepsilon_d$'' (for blocks of $a$ and of $d$ inside $\Omega$), $\mathcal N$ for the common $\N(c)$/$\N(d)$-blocks in $\Omega$, and recall that the $c$-blocks in $\Omega$ are the $\T(c)$-blocks, $K'$ and $\mathcal N$, the $d$-blocks are the $\T(d)$-blocks, $K$ and $\mathcal N$, every $\N(b)$-block in $\Omega$ is an $\N(a)$-block or $S_a$, and every $\N(a)$-block in $\Omega$ is an $\N(b)$-block or $S_b$. Suppose, for a contradiction, that no $\T(a)$-block is $\sim$ to a $\T(d)$-block.

($\alpha$) \emph{Every $\T(a)$-block $A$ satisfies $A\sim K$.} Take a path in $\Gamma(a,c)$ from $A$ to $\varphi(A)$ with no $\T$-block in its interior; its interior $c$-blocks are blocks of $\mathcal N$ or $K'$, and its interior $a$-blocks are $\N(a)$-blocks. If $K'$ occurs, consider its first occurrence and the $a$-block $u$ preceding it ($u=A$ if $K'$ comes first): $u$ meets $K'=\bigcup\T(d)$, hence some $\T(d)$-block $D$, so $u\sim D$; the path from $A$ to $u$ is a path in $\Gamma(a,d)$ (its $c$-blocks are in $\mathcal N$), so $A\sim D$, contrary to assumption. So $K'$ does not occur, and the $a$-block $u$ preceding $\varphi(A)$ ($u=A$ if the path has length one) meets $\varphi(A)\subseteq K$; as $K$ is a $d$-block, $u\sim K$ and $A\sim K$.

($\beta$) Hence $K\not\sim D$ for every $\T(d)$-block $D$.

($\gamma$) \emph{Every $\T(d)$-block $D$ satisfies $S_b\sim D$, and $S_a\ne S_b$.} Let $D=\chi(B)$ and take a path in $\Gamma(b,d)$ from $B$ to $D$ with no $\T$-block in its interior; its interior $d$-blocks are in $\mathcal N\cup\{K\}$ and its interior $b$-blocks are $\N(a)$-blocks or $S_a$. If $S_a$ occurs, look at its last occurrence and the $d$-block $v$ following it: $v$ meets $S_a=\bigcup\T(a)$, hence some $\T(a)$-block $A$, and the path from $v$ to $D$ is a path in $\Gamma(a,d)$, so $A\sim D$, contrary to assumption. If $K$ occurs, the path from $K$ to $D$ is a path in $\Gamma(a,d)$, contradicting ($\beta$). So the interior consists of blocks of $\mathcal N$ and $\N(a)$-blocks, and the first $d$-block $v_1$ meets $B$. If $S_a=S_b$ then $B\cap v_1$ contains a point of some $\T(a)$-block $A$, and $A\sim D$ along the path, a contradiction; so $S_a\ne S_b$, $S_b$ is an $a$-block meeting $v_1$, and $S_b\sim D$ along the path (if the path is $B,D$ alone, $S_b$ meets $D$). As $\chi$ is a bijection onto the $\T(d)$-blocks in $\Omega$, this holds for every $D$.

($\delta$) Now take any $\T(b)$-block $B$ and a path $B=u_0,v_1,u_1,v_2,\dots,v_m=\psi(B)$ in $\Gamma(b,c)$ with no $\T$-block in its interior: the $c$-blocks $v_j$ ($j<m$) are in $\mathcal N\cup\{K'\}$ and the $b$-blocks $u_j$ ($1\le j<m$) are $\N(a)$-blocks or $S_a$. We show by induction on $j$ that $u_j$ is an $\N(a)$-block with $u_j\sim S_b$ (where $u_0=S_b\supseteq B$). For $u_{j-1}$ (an $a$-block $\sim S_b$) and $v_j$: if $v_j\in\mathcal N$ then $v_j$ is a $d$-block meeting $u_{j-1}$, so $v_j\sim S_b$, and if $v_j=K'$ then $u_{j-1}$ meets some $\T(d)$-block $D\sim S_b$ (by ($\gamma$)). If now $u_j=S_a$: in the first case $v_j$ meets $S_a$ in a point of some $\T(a)$-block $A$, so $A\sim v_j\sim S_b\sim D$ for any $D$; in the second case a point of $S_a\cap K'$ lies in some $\T(a)$-block $A$ and some $\T(d)$-block $D$, so $A\sim D$; both contradict the assumption. So $u_j$ is an $\N(a)$-block, and $u_j\sim v_j\sim S_b$ resp.\ $u_j\sim D\sim S_b$ (it meets $K'$ in a point of some $D$). Finally $u_{m-1}$ (or $S_b$ itself if $m=1$) meets $\psi(B)\subseteq K$, so $K\sim S_b\sim D$, contradicting ($\beta$). This contradiction proves Step~6 and the theorem. \qed

\begin{remark}\label{rem:wiring}
(1) The hypothesis $\theta_a(d)=\theta_a(z)$ was used only in Step~4; Theorem~\ref{thm:S} below removes it. (2) The mechanism is different from that of the matching-pattern conditions. There the matchings alone carry the information (Theorem~\ref{thm:master}); here the commuting pair $c,d$ forces $c$ and $d$ to share their $\N$-blocks (Lemma~\ref{lem:PX}(b)), and the equality $\theta_a(z)=\theta_b(z)$ forces $a$ and $b$ to share theirs (Step~3); through these shared blocks the matchings $a\bowtie c$, $b\bowtie c$, $b\bowtie d$ transport the pairing of $\T$-blocks to a pairing between $a$ and $d$. The hubs $K$, $K'$ of Step~6 (a union of $\T(c)$-blocks which is a single $\N(d)$-block, and vice versa) are the only place where this is delicate.
\end{remark}

\subsection{The full wiring condition}\label{subsec:S}

\begin{theorem}\label{thm:S}
In every partition monoid $\P_X$, for projections $a,b,c,d$,
\[
a\bowtie c,\quad b\bowtie c,\quad b\bowtie d,\quad cd=dc,\quad \theta_a(cd)=\theta_b(cd)\ \Longrightarrow\ a\bowtie d .
\tag{S}
\]
In the binary notation: $c*a=a$, $a*c=c$, $c*b=b$, $b*c=c$, $d*b=b$, $b*d=d$, $d*c=c*d$, $(d*c)*a=(d*c)*b$ imply $d*a=a$ and $a*d=d$.
\end{theorem}

Condition (S) implies Theorem~\ref{thm:wiring} in every projection algebra (Prover9 finds the derivation in a fraction of a second): $a\bowtie d$ gives $\theta_a(d)=a$, so $\theta_a(z)=a=\theta_a(c)$ with $z\le c$, and $\theta_a$ is injective on $c^{\downarrow}$ (its inverse there is $\theta_c$), so $z=c$; then $w=\theta_a(c)=a$ and $w=\theta_b(c)=b$. The converse fails, even in the presence of \eqref{eq:Q1} and the triangle condition: see Example~\ref{ex:P9}.

\begin{proof}
We keep the notation of the proof of Theorem~\ref{thm:wiring}; Steps 1--3 there did not use the extra hypothesis. Since $\varepsilon_a\vee\varepsilon_d\le\Pi$, it suffices to show that inside each class $\Omega$ of $\Pi$ meeting $\supp(a)\cup\supp(d)$ every class of $\varepsilon_a\vee\varepsilon_d$ contains exactly one $\T(a)$-block and one $\T(d)$-block. Classes of type (F1) are handled by Step~5 and its mirror image: Step~5 shows that distinct $\T(a)$-blocks are not connected in $\Gamma(a,d)$; the same argument with the roles of $a$ and $d$ exchanged (a path between distinct $\T(d)$-blocks with no interior $\T(d)$-block; the segments between consecutive $\T(a)$-blocks are components $\kappa$ touching two $\T(a)$-blocks, which is impossible, so there is at most one $\T(a)$-block on the path, and it is $\sigma^{-1}$ of both ends) shows that distinct $\T(d)$-blocks are not connected; and every $\T(a)$-block $A$ is connected to a $\T(d)$-block, because a path in $\Gamma(a,c)$ from $A$ to $\varphi(A)$ ends with an $a$-block meeting $\varphi(A)\subseteq K=\bigcup\T(d)$, hence meeting some $\T(d)$-block, and its other vertices are $\Gamma(a,d)$-vertices; dually for $\T(d)$-blocks. It remains to treat a class $\Omega$ of type (F2), which we do in three steps: a bound on $p=|\T(a)\cap\Omega|$, a reduction to a finite combinatorial statement, and a verification of that statement by computer.

\smallskip
\emph{The bound.} We use the notation of Step~6: $K=S_c$, $K'=S_d$, $\mathcal N$ the common $\N(c)/\N(d)$-blocks in $\Omega$, and the alternative (B1) $S_a=S_b$ with $\N(a)|_\Omega=\N(b)|_\Omega$, or (B2) $S_a\cap S_b=\emptyset$ with $S_a\in\N(b)$, $S_b\in\N(a)$ and the remaining $\N(a)$- and $\N(b)$-blocks in $\Omega$ equal; in both cases write $\mathcal N'$ for the common ones. Let $\mathcal G$ be the graph on $\mathcal N'\sqcup\mathcal N$ joining blocks that meet, with components $\kappa$; as in Step~5, every $\kappa$ lies in a single component of each of the four graphs $\Gamma(a,c)$, $\Gamma(b,c)$, $\Gamma(b,d)$, $\Gamma(a,d)$, and (by the three hypotheses) touches at most one $\T(a)$-, one $\T(b)$-, one $\T(c)$- and one $\T(d)$-block. Now let $A\in\T(a)$ and take a path in $\Gamma(a,c)$ from $A$ to $\varphi(A)$ with no $\T$-block in its interior. Its interior vertices are plain blocks, $K'$ and (in case (B2)) $S_b$. If it contains no hub, it is $A$ followed by a single $\kappa$ touching $A$ and $\varphi(A)$, or $A$ meets $\varphi(A)$; in case (B1) the point of contact of $\kappa$ (or of $A$) with $A$ lies in some $B\in\T(b)$ and $\kappa$ (or that point) meets $K$, so $B$ lies in the component of $K$ in $\Gamma(b,d)$ and in the component of $\varphi(A)$ in $\Gamma(b,c)$, whence $A=\varphi^{-1}\psi(B_K)$ where $B_K$ is the unique $\T(b)$-block in the $\Gamma(b,d)$-component of $K$; in case (B2) the point of contact lies in $S_a$, so $\varphi(A)$ lies in the component of $S_a$ in $\Gamma(b,c)$, whence $A=\varphi^{-1}(C_a)$ where $C_a$ is the unique $\T(c)$-block of that component. If the path contains $K'$ then $A$ lies in the component of $K'$ in $\Gamma(a,c)$, which contains a unique $\T(a)$-block; and if it contains $S_b$ then $A$ lies in the component of $S_b$ in $\Gamma(a,c)$, again with a unique $\T(a)$-block. Hence $p\le2$ in case (B1) and $p\le3$ in case (B2).

\smallskip
\emph{The reduction.} Every point of $\Omega$ lies in one block of each of $a,b,c,d$. On the $a$-side its blocks are either a pair $(A,B)$ of $\T$-blocks (case (B1)), or $(A,S_a)$ or $(S_b,B)$ (case (B2)), or a common block $F\in\mathcal N'$; on the $c$-side they are $(C,K)$, $(K',D)$ or a common $E\in\mathcal N$. The four graphs $\Gamma(\cdot,\cdot)$, and hence the truth of the hypotheses and of the conclusion, depend only on which pairs (an $a$-side datum, a $c$-side datum) occur as points. Points whose two data are both plain lie in a component $\kappa$ of $\mathcal G$, and in all four graphs the blocks of $\kappa$ are connected to each other and to exactly the $\T$-blocks (and, through them, hubs) that $\kappa$ touches. Consequently the component structure of the four graphs is determined by (i) the set of \emph{direct overlaps}, i.e.\ pairs of a special $a$-side datum and a special $c$-side datum occurring as a point, and (ii) the set of \emph{types} of the components $\kappa$, the type of $\kappa$ being the set of $\T$-blocks it touches (at most one of each of the four kinds, and in case (B1) a $\T(a)$-block iff a $\T(b)$-block, since the contact points lie in cells $A\cap B$, which must be nonempty). Two components of the same type can be merged without changing anything, and a component touching no $\T$-block is irrelevant. So a class of type (F2) is described, up to the equivalence that matters, by a subset of the at most $(p+1)^4$ types and a set of direct overlaps, subject to: every $\T$-block occurs in some type or overlap (blocks are nonempty), and in case (B1) a type or overlap involving the cell $(A,B)$ requires that cell to be nonempty, i.e.\ to occur.

\smallskip
\emph{The verification.} For each $p\le6$ and each of the cases (B1), (B2) we encoded the finitely many such descriptions as a propositional formula (presence variables for types, overlaps and cells; reachability in each of the four finite graphs by iterated squaring), added the fineness of $\Gamma(a,c)$, $\Gamma(b,c)$, $\Gamma(b,d)$ and the negation of the fineness of $\Gamma(a,d)$, and found the formula unsatisfiable (three different SAT solvers, well under a second each; \texttt{hubtypes.py}). Together with the bound $p\le3$ this proves the theorem in classes of type (F2). (The same encoding shows that classes of type (F2) with $p\ge3$ do not exist at all, so the actual bound is $p\le2$; and a randomised test with an independent implementation of the reachability found no counterexample either, \texttt{hubcheck\_rand.py}.) (A computer-free proof of the hub case is a special case of Theorem~\ref{thm:Q}, see Corollary~\ref{cor:Q} and Remark~\ref{rem:Q}(2).)
\end{proof}

\begin{example}\label{ex:P9}
Let $\mathfrak P_9=\{a,b,c,d,m,n,k,w,z\}$ with $\theta$ given by the table (rows $p$, columns $q$, entry $\theta_p(q)$):
\[
\begin{array}{c|ccccccccc}
 & a & b & c & d & m & n & k & w & z\\ \hline
a & a & m & a & m & m & m & m & w & w\\
b & m & b & b & b & m & m & m & w & w\\
c & c & c & c & z & n & n & z & z & z\\
d & k & d & z & d & k & z & k & z & z\\
m & m & m & m & m & m & m & m & w & w\\
n & n & n & n & z & n & n & z & z & z\\
k & k & k & z & k & k & z & k & z & z\\
w & w & w & w & w & w & w & w & w & w\\
z & z & z & z & z & z & z & z & z & z
\end{array}
\]
It is $\mathfrak P$ with a middle layer inserted: the matched pairs are $a\bowtie c$, $b\bowtie c$, $b\bowtie d$, $m\bowtie n$, $m\bowtie k$, $w\bowtie z$; $m\le a,b$, $n\le c$, $k\le d$, $w\le m$, $z\le n,k$; and $\theta_a(d)=\theta_a(b)=\theta_b(a)=m$, $\theta_d(a)=k$, $\theta_c(d)=\theta_d(c)=z$, $\theta_a(z)=\theta_b(z)=w$. Mace4 produced it as a model of the projection algebra axioms, \eqref{eq:Q1}, the triangle clause and the quasi-identity of Theorem~\ref{thm:wiring}, with the hypotheses of (S) as ground equations and $\theta_a(d)\ne a$ (there is none of order $\le8$; the run needs the memory limit raised because of the six-variable triangle clause). So $\mathfrak P_9$ satisfies \eqref{eq:Q1}, the triangle condition and Theorem~\ref{thm:wiring}, but violates (S): the pair $a,d$ is only ``half wired'', $\theta_a(d)=m$ lying strictly between $w$ and $a$. By Theorem~\ref{thm:S} it does not embed in any $\P(\P_X)$ (the SAT search confirms this for $\P(\P_m)$, $m\le8$), and any regular $*$-semigroup with projection algebra $\mathfrak P_9$ (it has at least $16+9+4=29$ elements) is a further finite counterexample to sufficiency which Theorem~\ref{thm:wiring} does not detect. Thus (S) is the condition to keep.
\end{example}

\subsection{Graph semigroups: the complete answer}\label{subsec:graphthm}

For the family of Section~\ref{subsec:graphsg} the two conditions \eqref{eq:Q1} and (S) are also sufficient, and both have a transparent graph-theoretic meaning. For a colouring $\bot\colon V\to B$ write $V_m=\bot^{-1}(m)$, and for $m\ne n$ let $G[V_m,V_n]$ be the bipartite graph of edges between $V_m$ and $V_n$.

\begin{theorem}\label{thm:graphsg}
Let $G=(V,E)$ be a graph, $B$ a set and $\bot\colon V\to B$ a map. The following are equivalent:
\begin{enumerate}[label=\textup{(\roman*)},nosep]
\item $Q(G,\bot)$ $*$-embeds in some partition monoid;
\item the projection algebra $V\sqcup B$ of $Q(G,\bot)$ embeds in some $\P(\P_X)$;
\item $V\sqcup B$ satisfies \eqref{eq:Q1} and (S);
\item $\bot$ is a proper colouring, and for any two colours $m\ne n$ the bipartite graph $G[V_m,V_n]$ contains no induced path $a-c-b-d$ with $\bot a=\bot b=m$ and $\bot c=\bot d=n$; equivalently, every $G[V_m,V_n]$ is a disjoint union of complete bipartite graphs and isolated vertices.
\end{enumerate}
When they hold, $Q(G,\bot)$ embeds in $\P_X$ with $|X|=|V|+\beta+2|B|$, where $\beta$ is the number of connected components with at least one edge of the graphs $G[V_m,V_n]$, $m\ne n$, and the image consists of elements of rank $\le1$.
\end{theorem}

\begin{proof}
(i)$\Rightarrow$(ii) is clear and (ii)$\Rightarrow$(iii) is \eqref{eq:Q1} together with Theorem~\ref{thm:S}.

(iii)$\Rightarrow$(iv). By Lemma~\ref{lem:graphsg}, \eqref{eq:Q1} says that $\bot$ is proper. We work out what (S) says in $V\sqcup B$. A vertex and a bottom element are never matched ($\theta_v(m)=\bot v\ne v$), so in an instance of (S) either all of $a,b,c,d$ lie in $B$, where everything is matched and the conclusion holds, or all lie in $V$, with $a\bowtie c$, $b\bowtie c$, $b\bowtie d$ meaning equality or adjacency. Two distinct vertices $c,d$ commute iff $\theta_c(d)=\theta_d(c)$, i.e.\ iff they are non-adjacent with $\bot c=\bot d$ (if $c\bowtie d$ the two values are $c\ne d$). If $c=d$ then $cd=c$, and $\theta_a(c)=a$, $\theta_b(c)=b$, so the last hypothesis forces $a=b$ and the conclusion $a\bowtie d=b\bowtie d$ holds. If $c\ne d$ then $cd=\bot c\in B$, $\theta_a(cd)=\bot a$ and $\theta_b(cd)=\bot b$. Hence, for a proper colouring, (S) is exactly the statement
\[
a\bowtie c,\quad b\bowtie c,\quad b\bowtie d,\quad c\ne d,\quad \bot c=\bot d,\quad \bot a=\bot b\ \Longrightarrow\ a\bowtie d\qquad(a,b,c,d\in V).
\]
An induced path $a-c-b-d$ in some $G[V_m,V_n]$ with $a,b\in V_m$, $c,d\in V_n$ satisfies the hypotheses ($c,d$ are distinct and, having the same colour, non-adjacent), so (S) gives $a\bowtie d$, i.e.\ $ad\in E$, contradicting ``induced''. Finally, a bipartite graph without induced $4$-vertex paths is a disjoint union of complete bipartite graphs: if a connected component is not complete bipartite, choose non-adjacent $p\in V_m$, $q\in V_n$ in it at minimal distance; the distance is odd and $\ge3$, and the first four vertices of a shortest path from $p$ to $q$ form an induced path $p-c-b-d$ with $b\in V_m$, $c,d\in V_n$ ($p\not\sim d$ by minimality). The converse is obvious.

(iv)$\Rightarrow$(i). Let $\mathcal C$ be the set of connected components with at least one edge of the graphs $G[V_m,V_n]$, $m\ne n$; by (iv) each $C\in\mathcal C$ is complete bipartite between its two nonempty sides $C\cap V_m$ and $C\cap V_n$. Let
\[
X=\{x_v:v\in V\}\sqcup\{y_C:C\in\mathcal C\}\sqcup\{z_m,z_m':m\in B\},
\]
and for $m\in B$ let $\varepsilon_m$ be the partition of $X$ whose non-singleton blocks are $\{z_m,z_m'\}$ and, for every $C\in\mathcal C$ one of whose colours is $m$, the block $\{y_C\}\cup\{x_u:u\in C\setminus V_m\}$ (the point $y_C$ together with the points of the side of $C$ \emph{not} coloured $m$). These blocks are pairwise disjoint: a vertex $u\in V_n$ lies in at most one component of $G[V_m,V_n]$, and components with different second colours involve points $x_u$ of different colours. Note that $x_v$ is a singleton of $\varepsilon_{\bot v}$.

\emph{Claim.} Let $v\in V_m$ and $v'\in V_n$. Then $x_v$ and $x_{v'}$ lie in the same class of $\varepsilon_m\vee\varepsilon_n$ iff $v\bowtie v'$ (i.e.\ $v=v'$ or $vv'\in E$). Indeed, if $m=n$ the join is $\varepsilon_m$, in which both points are singletons, and $v\bowtie v'$ iff $v=v'$ since the colouring is proper. If $m\ne n$ and $v$ lies in a component $C\in\mathcal C$ of $G[V_m,V_n]$, then the set $\{y_C\}\cup\{x_u:u\in C\}$ is a union of blocks of $\varepsilon_m$ (the block of $y_C$, which is $\{y_C\}\cup\{x_u:u\in C\cap V_n\}$, and the singletons $x_u$, $u\in C\cap V_m$) and likewise of $\varepsilon_n$, so it is a class of the join; if $v$ lies in no such component, $\{x_v\}$ is a class. So $x_v\sim x_{v'}$ iff $v,v'$ lie in a common component of $G[V_m,V_n]$ iff $vv'\in E$, by completeness of the components.

Now define $\varphi\colon Q(G,\bot)\to\P_X$: $\varphi(v,v')$ is the partition of $X\sqcup X'$ whose blocks are the blocks of $\varepsilon_{\bot v}$ on $X$ and of $\varepsilon_{\bot v'}$ on $X'$, except that the two singletons $\{x_v\}$ and $\{x'_{v'}\}$ are merged into one transversal block; and $\varphi(m,n)$ consists of the blocks of $\varepsilon_m$ on $X$ and of $\varepsilon_n$ on $X'$ (rank $0$). Clearly $\varphi(x^*)=\varphi(x)^*$, and $\varphi$ is injective: $\varphi(x)$ determines its rank, its upper and lower partitions (which determine the colours, since $\{z_m,z_m'\}$ is a block of $\varepsilon_m$ only) and, in rank $1$, the transversal pair $(x_v,x'_{v'})$. For the products, consider $\varphi(v,q)\varphi(r,s)$ with $v,q,r,s\in V$. The only point of the upper row of $\varphi(v,q)$ connected to the middle row is $x_v$, and its partner there is $x_q$; the middle row is glued by $\varepsilon_{\bot q}\vee\varepsilon_{\bot r}$, so by the Claim $x_q$ reaches the point $x_r$, and through it the lower point $x'_s$, iff $q\bowtie r$. Hence the product has the upper blocks of $\varepsilon_{\bot v}$, the lower blocks of $\varepsilon_{\bot s}$, and the transversal block $\{x_v,x'_s\}$ if $q\bowtie r$ and no transversal block otherwise: it is $\varphi(v,s)$ or $\varphi(\bot v,\bot s)$, as required. If one of the factors has rank $0$ there is no transversal block in the product, whose upper and lower partitions are those of the first and the last factor: $\varphi(m,n)\varphi(r,s)=\varphi(m,\bot s)$, $\varphi(v,q)\varphi(m,n)=\varphi(\bot v,n)$, $\varphi(m,n)\varphi(m',n')=\varphi(m,n')$. So $\varphi$ is an injective $*$-homomorphism.
\end{proof}

\begin{remark}\label{rem:graphthm}
(1) For an induced subgraph $G'=G[V']$ the semigroup $Q(G',\bot|_{V'})$ is a $*$-subsemigroup of $Q(G,\bot)$ (products of pairs over $V'\sqcup\bot(V')$ stay there). So a non-embeddable $Q(G,\bot)$ with $\bot$ proper contains a copy of $\mathfrak Q$, and $\mathfrak Q$ is the smallest non-embeddable graph semigroup satisfying \eqref{eq:Q1}; the ``finite list of forbidden coloured subgraphs'' consists of the monochromatic edge and the bichromatic path $a-c-b-d$.

(2) The theorem was found by a computer survey (\texttt{graphsg\_survey.py}): for the $108$ coloured graphs with at most $4$ vertices (up to isomorphism and renaming of colours) the projection algebra embeds in some $\P(\P_m)$, $m\le5$, exactly when (iv) holds (the search was run for the $48$ proper colourings; the others fail \eqref{eq:Q1}), and then already for $m\le4$; the triangle condition holds in all of them. The explicit embedding of the proof was then verified by computer for all $788$ coloured graphs with $5$ vertices ($190$ embeddable) and all $13\,712$ with $6$ vertices ($1\,824$ embeddable) (\texttt{graphsg\_thm.py}). The construction is not economical: for $\mathfrak Q$ with $\bot$ injective it uses $|X|=4+3+8=15$ while $\P_4$ suffices.

(3) The triangle condition holds in every $Q(G,\bot)$ with $\bot$ proper, and for a trivial reason. Three pairwise commuting projections with a common pairwise product $w$ are either all equal or three distinct pairwise non-adjacent vertices of one colour $w\in B$ (or elements of $B$, all equal to $w$); a triangle configuration with $w\ne z$ therefore has $a_i\in V_w$ and $b_i\in V_z$ with $w\ne z$, and then $a_i$ and $b_i$ cannot commute. Thus the family $Q(G,\bot)$ sees only \eqref{eq:Q1} and (S), and by the theorem it can witness no further obstruction: an embedding, when it exists, uses only projections of rank $\le1$, whose $\T$-blocks are single points. The layered graph semigroups of Section~\ref{subsec:layered}, which contain $\mathfrak P_9$, allow the commuting products to lie several levels below the matched vertices; they see no further obstruction either (Theorem~\ref{thm:layered}).
\end{remark}

\subsection{Layered graph semigroups}\label{subsec:layered}

Theorem~\ref{thm:graphsg} shows that graph semigroups cannot witness obstructions beyond \eqref{eq:Q1} and (S). The algebra $\mathfrak P_9$ of Example~\ref{ex:P9} suggests a generalisation in which the bottom band of $Q(G,\bot)$ is replaced by another graph semigroup, and so on, so that the product of two unmatched projections may lie several levels below them. It turns out that this family, too, sees only \eqref{eq:Q1} and (S).

\begin{definition}\label{def:layered}
A \emph{layered graph} is a finite set $V=V_1\sqcup\dots\sqcup V_L$ ($L\ge1$) with a reflexive symmetric relation $\bowtie$ on each level $V_i$, complete on the bottom level $V_L$, and maps $\bot\colon V_i\to V_{i+1}$ ($1\le i<L$) which preserve $\bowtie$. For $p\in V_i$ and $\ell\ge i$ write $p_\ell=\bot^{\ell-i}p$ ($p$ descended to level $\ell$), call $p_L$ the \emph{root} of $p$ and $T_r=\{p:p_L=r\}$ the \emph{tree} of the root $r\in V_L$; the trees partition $V$. For $q,r\in V$ let
\[
f(q,r)=\min\{\ell\ge\max(\operatorname{lv}q,\operatorname{lv}r): q_\ell\bowtie r_\ell\},
\]
which exists because the bottom level is complete. The \emph{layered graph semigroup} $Q(V)$ is the set of pairs $(p,q)$ with $p,q$ on a common level, with $(p,q)^*=(q,p)$ and
\[
(p,q)(r,s)=(p_\ell,s_\ell),\qquad \ell=f(q,r).
\]
\end{definition}

For $L=2$ this is $Q(G,\bot)$ ($V_1=V$, $V_2=B$); every map into a complete graph preserves $\bowtie$. The layered graph with $V_1=\{a,b,c,d\}$ ($a\bowtie c$, $b\bowtie c$, $b\bowtie d$), $V_2=\{m,n,k\}$ ($m\bowtie n$, $m\bowtie k$), $V_3=\{w,z\}$, $\bot\colon a,b\mapsto m$, $c\mapsto n$, $d\mapsto k$, $m\mapsto w$, $n,k\mapsto z$ gives a $29$-element regular $*$-semigroup whose projection algebra is $\mathfrak P_9$ (\texttt{layered.py}).

\begin{lemma}\label{lem:layered}
$Q(V)$ is a regular $*$-semigroup; its projections are the pairs $(p,p)$, identified with $V$, and
\[
\theta_p(q)=p_{f(p,q)},\qquad q\le p\iff q=p_\ell\text{ for some }\ell .
\]
(If $\bot$ does not preserve $\bowtie$, the multiplication is not associative.) Moreover:
\begin{enumerate}[label=\textup{(\alph*)},nosep]
\item \eqref{eq:Q1} holds in $Q(V)$ iff distinct matched elements have different roots;
\item two elements $p,q$ commute iff $p_f=q_f$ for $f=f(p,q)$, and then $pq=p_f$; under \eqref{eq:Q1} this says that $p,q$ lie in a common tree, and $p_f$ is their least common ancestor;
\item assuming \eqref{eq:Q1}, (S) holds in $Q(V)$ iff for every level $i$ and every two roots $r\ne r'$ the bipartite graph of matched pairs between $T_r\cap V_i$ and $T_{r'}\cap V_i$ contains no induced path on four vertices, i.e.\ is a disjoint union of complete bipartite graphs and isolated vertices;
\item assuming \eqref{eq:Q1}, the triangle condition holds in $Q(V)$.
\end{enumerate}
\end{lemma}

\begin{proof}
Preservation of $\bowtie$ gives $q_\ell\bowtie r_\ell\Rightarrow q_{\ell'}\bowtie r_{\ell'}$ for $\ell'\ge\ell$, hence $f(q_k,r)=\max(k,f(q,r))$ for $k\ge\operatorname{lv}q$. So $((p,q)(r,s))(t,u)=(p_{\ell_2},u_{\ell_2})$ with $\ell_2=f(s_{\ell_1},t)=\max(f(q,r),f(s,t))$, $\ell_1=f(q,r)$, and symmetrically $(p,q)((r,s)(t,u))=(p_{\ell_4},u_{\ell_4})$ with $\ell_4=f(q,r_{\ell_3})=\max(f(s,t),f(q,r))$, $\ell_3=f(s,t)$. The identities $(xy)^*=y^*x^*$ and $xx^*x=x$ are clear ($f$ is symmetric and $f(q,q)=\operatorname{lv}q$). If $p\bowtie q$ but $\bot p\not\bowtie\bot q$ (so $\operatorname{lv}p+1<L$), then $((p,p)(q,q))(\bot q,\bot q)=(p,q)(\bot q,\bot q)=(\bot p,\bot q)$, while $(p,p)((q,q)(\bot q,\bot q))=(p,p)(\bot q,\bot q)=(p_\ell,q_\ell)$ with $\ell\ge\operatorname{lv}p+2$. The projections are the symmetric pairs, and $\theta_p(q)=(p,p)(q,q)(p,p)=(p_f,q_f)(p,p)=(p_f,p_f)$ with $f=f(p,q)$, since $f(q_f,p)=f$; so $\theta_p(q)=q$ iff $q=p_f$.

(b) is the formula for $\theta$ together with the remark after Theorem~\ref{thm:hex}. Note that within a tree, $p_f\bowtie q_f$ with $p_f\ne q_f$ would be two distinct matched elements with a common descendant, so under \eqref{eq:Q1} ``matched'' means ``equal'' inside a tree and $f(p,q)$ is the level of the least common ancestor of $p$ and $q$.

(a) Matched elements lie on a common level, so ``no common descendant'' means ``different roots''. If $p\bowtie q$ are distinct with a common root then $r=p_L=q_L$ satisfies $r\le p,q$ but $\theta_p(q)=p\ne q=\theta_q(p)$. Conversely let $r\le p$, $r\le q$ with $p\ne q$, say $r=p_\ell=q_\ell$; at the level $f=f(p,q)\le\ell$ the matched elements $p_f,q_f$ have the common descendant $r$, so $p_f=q_f$ and $\theta_p(q)=p_f=q_f=\theta_q(p)$.

(c) Matched elements lie on a common level, so in an instance of (S) all of $a,b,c,d$ lie on one level $i$. If $c=d$ then $cd=c$, $\theta_a(c)=a$, $\theta_b(c)=b$ and the hypothesis $\theta_a(cd)=\theta_b(cd)$ gives $a=b$, so the instance is trivial. If $c\ne d$ then $cd=dc$ says by (b) that $c,d$ lie in a common tree, with $z:=c_f=d_f$, $f=f(c,d)$; since $a_f\bowtie c_f=z$ we get $f(a,z)=f$ and $\theta_a(z)=a_f$, similarly $\theta_b(z)=b_f$, so the last hypothesis says that $a,b$ lie in a common tree and their least common ancestor is at level $\le f$. By \eqref{eq:Q1} the tree of $a,b$ differs from that of $c,d$. Thus, writing $\mu(p,q)$ for the level of the least common ancestor of two elements of a tree, (S) says: for every path $a-c-b-d$ of matchings with $a,b\in T_r$, $c,d\in T_{r'}$, $r\ne r'$, if $\mu(a,b)\le\mu(c,d)$ then $a\bowtie d$. Applied to the reversed path $d-b-c-a$ (with $c,d$ in the role of $a,b$) it says the same under $\mu(c,d)\le\mu(a,b)$; one of the two always holds. So (S) is equivalent to: every path $a-c-b-d$ with $a,b\in T_r\cap V_i$ and $c,d\in T_{r'}\cap V_i$ closes up, $a\bowtie d$; that is, the bipartite graph in (c) has no induced $4$-vertex path, which (as in the proof of Theorem~\ref{thm:graphsg}) means that it is a disjoint union of complete bipartite graphs.

(d) Let $a_1,a_2,a_3,b_1,b_2,b_3$ be a triangle configuration. By (b) the $a_i$ pairwise commute, so they lie in one tree; likewise the $b_i$ lie in one tree; and $a_i$ commutes with $b_i$, so all six lie in one tree. Then $a_1\bowtie b_2$ forces $a_1=b_2$ by (a), and similarly $a_2=b_3$, $a_3=b_1$. Hence $\{a_1,a_2,a_3\}=\{b_1,b_2,b_3\}$ and $w=a_1a_2=b_2b_3=z$.
\end{proof}

\begin{theorem}\label{thm:layered}
For a layered graph $V$ the following are equivalent:
\begin{enumerate}[label=\textup{(\roman*)},nosep]
\item $Q(V)$ $*$-embeds in some partition monoid;
\item the projection algebra $V$ of $Q(V)$ embeds in some $\P(\P_X)$;
\item $V$ satisfies \eqref{eq:Q1} and (S);
\item distinct matched elements have different roots, and for every level $i$ and every two roots $r\ne r'$ the bipartite graph of matched pairs between $T_r\cap V_i$ and $T_{r'}\cap V_i$ is a disjoint union of complete bipartite graphs and isolated vertices.
\end{enumerate}
When they hold, $Q(V)$ embeds in $\P_X$ with $|X|=|V|+|V_L|+\beta$, where $\beta$ is the number of complete bipartite components in (iv), and every $\T$-block of the image is a single point.
\end{theorem}

\begin{proof}
(i)$\Rightarrow$(ii) is clear, (ii)$\Rightarrow$(iii) is \eqref{eq:Q1} and Theorem~\ref{thm:S}, and (iii)$\Leftrightarrow$(iv) is Lemma~\ref{lem:layered}(a),(c). For (iv)$\Rightarrow$(i) we generalise the construction of Theorem~\ref{thm:graphsg}, with trees in the role of colours and the non-root elements in the role of vertices. Let $\mathcal C$ be the set of complete bipartite components in (iv), over all levels $i<L$ and pairs of roots; each $C\in\mathcal C$ lies on one level and meets exactly two trees. Let
\[
X=\{x_v: v\in V\setminus V_L\}\sqcup\{y_C:C\in\mathcal C\}\sqcup\{z_r,z_r': r\in V_L\},
\]
and for a root $r$ let $\varepsilon_r$ be the partition of $X$ whose non-singleton blocks are $\{z_r,z_r'\}$ and, for each $C\in\mathcal C$ meeting $T_r$, the block $\{y_C\}\cup\{x_u: u\in C\setminus T_r\}$. These blocks are disjoint (an element lies in at most one component per level and pair of trees, and for a fixed $r$ the elements $x_u$ in the blocks belong to the other tree of $C$). As in the Claim in the proof of Theorem~\ref{thm:graphsg}, for $u\in T_r$ and $v\in T_{r'}$, $r\ne r'$, the points $x_u,x_v$ lie in a common class of $\varepsilon_r\vee\varepsilon_{r'}$ iff $u,v$ lie in a common component $C\in\mathcal C$, iff $u\bowtie v$; the class of $x_u$ is $\{y_C\}\cup\{x_w:w\in C\}$ or $\{x_u\}$, and in particular contains no other point $x_{u'}$ with $u'\in T_r$ on a different level. Note also that every $x_v$, $v\in T_r$, is a singleton of $\varepsilon_r$.

For $p\in T_r$ let $\hat p$ be the projection with partition $\varepsilon_r$ whose $\T$-blocks are the singletons $\{x_{p_j}\}$, $\operatorname{lv}p\le j<L$ (the chain from $p$ down to, but excluding, the root); so $\hat r$ has rank $0$, $\rk\hat p=L-\operatorname{lv}p$, and $\supp\hat p$ determines $p$. We check $\theta_{\hat p}(\hat q)=\widehat{p_f}$, $f=f(p,q)$, using Lemma~\ref{lem:PX}(a). If $p,q\in T_r$, the join is $\varepsilon_r$ itself, whose classes meeting $\supp\hat p$ are the singletons $\{x_{p_j}\}$, transversal in $\theta_{\hat p}(\hat q)$ iff $x_{p_j}\in\supp\hat q$, i.e.\ iff $p_j=q_k$ for some $k\ge\operatorname{lv}q$, i.e.\ iff $j\ge f$ (the level of the least common ancestor); so $\theta_{\hat p}(\hat q)$ is $\varepsilon_r$ with the $\T$-blocks $\{x_{p_j}\}$, $j\ge f$, which is $\widehat{p_f}$. If $p\in T_r$, $q\in T_{r'}$ with $r\ne r'$, the classes of $\varepsilon_r\vee\varepsilon_{r'}$ meeting $\supp\hat p$ are the classes $K_j$ of the points $x_{p_j}$, each containing only this one $\T$-block of $\hat p$; $K_j$ meets $\supp\hat q=\{x_{q_k}:k\ge\operatorname{lv}q\}$ iff $p_j\bowtie q_k$ for some $k\ge\operatorname{lv}q$, which forces $k=j$, iff $j\ge\operatorname{lv}q$ and $p_j\bowtie q_j$, iff $j\ge f$ (by preservation of $\bowtie$ and the definition of $f$). So again $\theta_{\hat p}(\hat q)=\widehat{p_f}$. (For $q=r'$ a root, $\supp\hat q=\emptyset$ and $\theta_{\hat p}(\hat q)=\hat r=\widehat{p_L}$, as required; and $\theta_{\hat r}(\hat q)=\hat r$ for every $q$ since $\hat r$ has rank $0$.) Injectivity: $\varepsilon_r$ determines $r$ (the block $\{z_r,z'_r\}$), and the $\T$-blocks determine $p$. So $p\mapsto\hat p$ embeds the projection algebra.

For the semigroup, send $(p,q)$ (with $p,q$ on level $i$) to the element $\varphi(p,q)$ whose upper and lower partitions are those of $\hat p$ and $\hat q$, with the transversal blocks $\{x_{p_j},x'_{q_j}\}$, $i\le j<L$. This is injective and $*$-compatible, and $\varphi(p,q)\varphi(r,s)=\varphi(p_\ell,s_\ell)$, $\ell=f(q,r)$: in the product, the upper point $x_{p_j}$ is joined to the middle point $x_{q_j}$, whose class in the middle join $\varepsilon_{T(q)}\vee\varepsilon_{T(r)}$ contains a point $x_{r_k}$, $k\ge\operatorname{lv}r$ (necessarily $k=j$), iff $j\ge f(q,r)$, by the two computations above, and then $x_{p_j}$ is joined to $x'_{s_j}$; for $j<\ell$ the block of $x_{p_j}$ in the product is the nontransversal singleton $\{x_{p_j}\}$, and the other blocks are unchanged. This is $\varphi(p_\ell,s_\ell)$.
\end{proof}

\begin{remark}\label{rem:layered}
(1) The search behind the theorem (\texttt{layered\_search.py}, \texttt{layered\_search2.py}) enumerated the layered graphs with surjective descent maps and level sizes $(n_1,n_2,n_3)$, $n_1\le4$, $n_2\le3$, $n_3\le2$ (the largest shape $(4,3,2)$ only partly, the run having been interrupted), and the four-level shapes $(2,2,2,1)$, $(3,2,2,1)$, $(2,2,2,2)$, $(3,2,2,2)$; it computed their projection algebras, kept those satisfying \eqref{eq:Q1} (up to isomorphism) and tested the triangle condition, (S) and embeddability into $\P(\P_m)$, $m\le8$ (backtracking, then SAT). The triangle condition never failed, and every algebra satisfying (S) embedded, with $m\le7$; the trees (no matchings above the bottom level) need the largest $m$. The construction of the theorem was then verified as an injective $*$-homomorphism on all these structures and on several hundred random layered graphs with up to four levels, and the predicate (iv) was checked against \eqref{eq:Q1} and (S) computed from the table (\texttt{layered\_thm.py}).

(2) In a layered graph every principal down-set $p^{\downarrow}$ is a chain, every value $\theta_p(q)$ lies in that chain, and (under \eqref{eq:Q1}) the commuting pairs are the pairs in a common tree. This is the reason why the triangle condition is invisible here: it needs two commuting triples with a common meet, which in a tree forces the triples to coincide. Any semigroup separating \eqref{eq:Q1}$+$(S) from embeddability must therefore have projections with non-chain down-sets (as $\P(Q_0)$ does, where the triangle condition fails), or must force $\T$-blocks of size $>1$; both constructions of this section use only singleton $\T$-blocks.
\end{remark}

\subsection{A fourth obstruction}\label{subsec:fourth}

Theorems~\ref{thm:graphsg} and \ref{thm:layered} say where \emph{not} to look for further obstructions. A census of the projection algebras of order $7$ finds one.

\emph{The census.} Mace4 enumerated the projection algebras of order $7$ satisfying \eqref{eq:Q1} and (S) (the axioms, the clause for \eqref{eq:Q1} and the two four-variable clauses for (S); the six-variable triangle clause was omitted and the triangle condition tested afterwards). To make the run resumable it was split into $49$ chunks by the values of $0*1$ and $1*0$, of which $21$ are redundant by the symmetry exchanging $0$ and $1$ (\texttt{pa7\_driver.sh}); isofilter within the chunks and a canonical form across them (\texttt{pa7\_post.py}) leave $1\,663$ algebras up to isomorphism (there are $295$ of order $6$). All of them satisfy the triangle condition. For $1\,662$ of them the SAT search finds an embedding into $\P(\P_m)$ with $m\le6$. The remaining one has no embedding into $\P(\P_m)$ for $m\le8$ (the case $m=8$ takes a quarter of an hour). With the elements renamed it is the following algebra.

\begin{definition}\label{def:P7}
$\mathfrak P_7=\{a,c,b,x,u,v,w\}$ with $\theta$ given by (rows $p$, columns $q$, entry $\theta_p(q)$)
\[
\begin{array}{c|ccccccc}
 & a & c & b & x & u & v & w\\ \hline
a & a & a & u & u & u & u & u\\
c & c & c & c & v & v & v & v\\
b & w & b & b & w & w & w & w\\
x & u & v & v & x & u & v & v\\
u & u & u & u & u & u & u & u\\
v & v & v & v & v & v & v & v\\
w & w & w & w & w & w & w & w
\end{array}
\]
\end{definition}

In words: $a\bowtie c\bowtie b$ is a path; $u,v,w$ form a bottom rectangular band ($\theta_u\equiv u$ etc., all pairs matched); $u\le a$, $v\le c$, $w\le b$, and every value not listed so far is the bottom of the subscript---so far this is the projection algebra of the graph semigroup $Q(a-c-b,\bot)$ with the injective colouring $\bot a=u$, $\bot c=v$, $\bot b=w$, which embeds in $\P_4$. The new element is $x$, with $u,v\le x$: it commutes with $a$ and with $c$, $xa=u$, $xc=v$, and it sees $b$ as it sees $c$, $\theta_x(b)=v$, while $\theta_b(x)=w$. Deleting any element leaves a subalgebra (of order $6$, hence embeddable), so $\mathfrak P_7$ is a minimal non-embeddable algebra. There is a $19$-element regular $*$-semigroup $\mathfrak T_{19}$ with projection algebra $\mathfrak P_7$: adjoin to $Q(a-c-b,\bot)$ an element $x$ with $x\cdot(p,q)=(\lambda p,\bot q)$, $(p,q)\cdot x=(\bot p,\lambda q)$ for the pairs $(p,q)$ of top elements, $x\cdot(m,n)=(\lambda m,n)$, $(m,n)\cdot x=(m,\lambda n)$ for the pairs of bottom elements, $x^2=x$, $x^*=x$, where $\lambda a=\lambda u=u$ and $\lambda c=\lambda b=\lambda v=\lambda w=v$ (\texttt{fourth.py} checks associativity and computes the projection algebra).

$\mathfrak P_7$ satisfies \eqref{eq:Q1}, the triangle condition and (S) (by the census, and directly by computer), hence all the matching-pattern conditions. The obstruction is the following quasi-identity.

\begin{proposition}\label{prop:R}
Let
\[
a\bowtie c,\quad b\bowtie c,\quad xa=ax,\quad xc=cx,\quad \theta_x(b)=\theta_x(c)\ \Longrightarrow\ a\bowtie b .
\tag{R}
\]
\begin{enumerate}[label=\textup{(\alph*)},nosep]
\item (R) holds in $\P(\P_m)$ for $m\le5$ (all $580\,902$ instances of the hypotheses in $\P(\P_5)$ were checked), and for $m\le6$ in the following stronger sense: a SAT search for four projections of $\P_6$ satisfying the hypotheses with $\theta_a(b)\ne a$, or with $\theta_b(a)\ne b$, is unsatisfiable (\texttt{fourth.py}, \texttt{pa\_sat.py} with a partial table).
\item (R) fails in $\mathfrak P_7$, with $(a,b,c,x)$ as named; it is the only instance. Theorem~\ref{thm:R} below proves (R) in every $\P_X$; the computations in (a) were the evidence that led to it.
\item (R) holds in every projection algebra of order $\le7$ satisfying \eqref{eq:Q1} and (S) other than $\mathfrak P_7$, and in every projection algebra of order $\le6$ satisfying \eqref{eq:Q1} and the triangle condition, in particular in $\mathfrak P$; and it holds in $\mathfrak P_9$. Hence (R) is independent of \eqref{eq:Q1}, the triangle condition and (S): $\mathfrak P$ satisfies (R) but not (S), $\mathfrak P_7$ satisfies (S) but not (R), and $\P(Q_0)$ and $Q(K_2,\text{const})$ show the rest as before.
\item In $\P(\P_5)$ the hypothesis $\theta_a(b)=\theta_a(x)$ (which also holds in $\mathfrak P_7$) does not strengthen (R): adding it changes nothing, and each of the five hypotheses listed is needed.
\end{enumerate}
\end{proposition}

\begin{theorem}\label{thm:R}
(R) holds in every partition monoid $\P_X$. Consequently $\mathfrak P_7$ embeds in no $\P(\P_X)$, $\mathfrak T_{19}$ embeds in no $\P_X$, and the smallest regular $*$-semigroup satisfying \eqref{eq:Q1}, the triangle condition and (S) that embeds in no partition monoid has order between $9$ and $19$.
\end{theorem}

\begin{proof}
Let $a,b,c,x$ be projections of $\P_X$ with $a\bowtie c$, $b\bowtie c$, $xa=ax=:u$, $xc=cx=:v$ and $\theta_x(b)=v$. We use the notation of Lemma~\ref{lem:PX}; ``block'' means a block of the relevant partition $\varepsilon_p$ of $X$, and two blocks of different partitions are \emph{adjacent} if they meet. A \emph{chain} in a join $\varepsilon_p\vee\varepsilon_q$ is a sequence of blocks, alternately of $p$ and of $q$, in which consecutive blocks are adjacent; two blocks lie in a common class of the join iff they are joined by a chain. Write $A_1,\dots,A_s$ for the $\T(a)$-blocks, $C_i$ for the $\T(c)$-block matched with $A_i$ by $a\bowtie c$ and $B_i$ for the $\T(b)$-block matched with $C_i$ by $b\bowtie c$ (Lemma~\ref{lem:PX}(c)); $\Gamma_i$ and $\Delta_i$ denote the classes of $\varepsilon_a\vee\varepsilon_c$ and $\varepsilon_b\vee\varepsilon_c$ containing $C_i$, so that $\Gamma_i$ contains exactly the $\T$-blocks $A_i,C_i$ and $\Delta_i$ exactly $B_i,C_i$. Write $H$ for a generic $\T(x)$-block and $Y=X\setminus\supp x$.

\smallskip
\emph{Step 1: the commuting pairs.} By Lemma~\ref{lem:PX}(b) the classes of $\varepsilon_x\vee\varepsilon_a$ are of four types: the $\T(u)$-blocks $U_1,\dots,U_t$ (each a union of $A$'s and a union of $H$'s), the $\N(a)$-blocks inside $\supp x$ (unions of $H$'s), the $\N(x)$-blocks inside $\supp a$ (unions of $A$'s), and the common $\N$-blocks; and $\varepsilon_u=\varepsilon_x\vee\varepsilon_a$. Likewise for $c$, with the $\T(v)$-blocks $V_1,\dots,V_{t'}$. In particular every block of $a$ and every block of $c$ lies inside $\supp x$ or inside $Y$; every $A_i$ lies in some $U_l$ or in an $\N(x)$-block, and every $C_i$ in some $V_k$ or in an $\N(x)$-block. An $\N(x)$-block $M$ containing a $\T(c)$-block $C_i$ is a union of $C$'s; it is not an $\N(a)$-block (an $\varepsilon_a$-block which is a union of $\varepsilon_c$-blocks would be a class of $\varepsilon_a\vee\varepsilon_c$ containing $C_i$ and no $A$), so it is a union of $A$'s as well, and the class of $C_i$ in $\varepsilon_a\vee\varepsilon_c$, being made of the $A$'s and $C$'s inside $M$, is $A_i=C_i$. Hence, with $I=\{i:A_i\subseteq\supp u\}$:
\[
A_i\subseteq\supp u\iff C_i\subseteq\supp v,\qquad\text{and } A_i=C_i\subseteq Y\text{ for } i\notin I
\]
(if $A_i\subseteq\supp u\subseteq\supp x$ but $C_i\subseteq Y$, then $A_i=C_i\subseteq Y$, a contradiction). Similarly an $\N(x)$-block is an $\N(a)$-block iff it is an $\N(c)$-block. So \emph{$a$ and $c$ have the same blocks inside $Y$}, with the same labels. For $i\in I$ let $U_{l(i)}\supseteq A_i$ and $V_{k(i)}\supseteq C_i$.

\emph{Step 2: two elementary facts.} (a) An $\N(a)$-block meets at most one $C$, an $\N(c)$-block meets at most one $A$ and at most one $B$, an $\N(b)$-block meets at most one $C$, a $C$ meets at most one $A$ and at most one $B$: for instance an $\N(a)$-block meeting $C_i$ and $C_j$ lies in $\Gamma_i\cap\Gamma_j$. (b) $u\bowtie v$, and $\theta_u(b)=u$, $\theta_v(b)=v$. Indeed $g=ac$ satisfies $gg^*=a$, $g^*g=c$, and $xg=xac=uc=uxc=uv=axc=acx=gx$; so $u=xgg^*x=(uv)(uv)^*=uvu$ and $v=xg^*gx=(uv)^*(uv)=vuv$. Since $u\le x$, $\theta_u(y)=\theta_u(\theta_x(y))$ for every projection $y$ (as $u=ux=xu$); so $\theta_u(b)=\theta_u(v)=u$ and $\theta_v(b)=\theta_v(v)=v$. By $u\bowtie v$ every class of $\varepsilon_u\vee\varepsilon_v=\varepsilon_x\vee\varepsilon_a\vee\varepsilon_c$ meeting $\supp u\cup\supp v$ contains exactly one $U$ and one $V$; since $\Gamma_i$ lies in such a class, $l(i)=l(j)$ iff $k(i)=k(j)$ for $i,j\in I$, i.e.\ the $A$'s inside $U_l$ are matched with the $C$'s inside one $V_{k(l)}$, and $U_l$ and $V_{k(l)}$ lie in one class of $\varepsilon_x\vee\varepsilon_a\vee\varepsilon_c$; two $\T(x)$-blocks lying in $U_l$ and $V_k$ respectively are connected in that join only if $k=k(l)$.

\emph{Step 3: the hypothesis $\theta_x(b)=v$.} By Lemma~\ref{lem:PX}(a) the blocks of $\theta_x(b)$ are the $\N(x)$-blocks and the sets $U_K=K\cap\supp x$ (a union of $H$'s) for the classes $K$ of $\varepsilon_x\vee\varepsilon_b$ meeting $\supp x$, $U_K$ being transversal iff $K$ meets $\supp b$; and the blocks of $v$ are the $\N(x)$-blocks, the $\T$-blocks $V_k$ and the $\N(c)$-blocks inside $\supp x$. So the partition $\mathcal Q$ of $\supp x$ into the $V_k$ and the $\N(c)$-blocks inside $\supp x$ is also the partition induced by $\varepsilon_x\vee\varepsilon_b$: there is a class $K_k$ with $K_k\cap\supp x=V_k$, which meets $\supp b$, and for each $\N(c)$-block $\tilde N\subseteq\supp x$ a class $K_{\tilde N}$ with $K_{\tilde N}\cap\supp x=\tilde N$, which does not meet $\supp b$. Consequently every $\varepsilon_b$-block $Z$ meets $\supp x$ inside a single part of $\mathcal Q$.

\emph{Step 4: the classes $\Delta_i$ avoid $\supp x$ except through $C_i$.} We claim that no $\N(c)$-block inside $\supp x$ lies in a $\Delta_i$. Suppose $\tilde N\subseteq\Delta_i$ and take a shortest chain $\tilde N=E_0,Z_1,E_1,Z_2,\dots,B_i$ in $\Delta_i$ ($Z$'s blocks of $b$, $E$'s blocks of $c$). $Z_1$ meets an $H\subseteq\tilde N$, so $Z_1\subseteq K_{\tilde N}$. Inductively all blocks of the chain lie in $K_{\tilde N}$: an $\varepsilon_b$-block meeting a block of $K_{\tilde N}$ lies in $K_{\tilde N}$; an $\N(c)$-block $E$ meeting such a $Z$ is an $\N(x)$-block (then $E\subseteq K_{\tilde N}$) or lies inside $\supp x$ (then the $H\subseteq E$ meeting $Z$ lies in $K_{\tilde N}\cap\supp x=\tilde N$, so $E=\tilde N$, excluded by minimality); and if $E=C_i$ then either $C_i\subseteq V_{k(i)}$, and the $H\subseteq V_{k(i)}$ meeting $Z$ would lie in $\tilde N$, or $C_i$ lies in an $\N(x)$-block, which is then in $K_{\tilde N}$. So $B_i\subseteq K_{\tilde N}$, contradicting $K_{\tilde N}\cap\supp b=\emptyset$. Consequently: every $\N(c)$-block in $\Delta_i$ lies in $Y$ and is an $\N(x)$-block and an $\N(a)$-block (Step~1); an $\varepsilon_b$-block $Z$ meeting $\supp x$ inside some $V_k$ meets a $C_p\subseteq V_k$, so $Z\subseteq\Delta_p$ and $Z\cap V_k\subseteq C_p$ by Step~2(a); and an $\varepsilon_b$-block meeting $\supp x$ inside an $\N(c)$-block $\tilde N$ lies in a class of $\varepsilon_b\vee\varepsilon_c$ containing $\tilde N$, hence containing no $B$ and no $C$.

\emph{Step 5: the indices outside $I$.} Let $i\notin I$, so $A_i=C_i\subseteq Y$, and let $D$ be the class of $A_i$ in $\varepsilon_a\vee\varepsilon_b$. We claim that every $\varepsilon_a$-block in $D$ lies in $Y$. Otherwise take a shortest chain $A_i,Z_1,Y_1,\dots,Z_t,Y_t$ in $D$ from $A_i$ to a block $Y_t\subseteq\supp x$ of $a$; then $Y_1,\dots,Y_{t-1}\subseteq Y$ are blocks of $c$ (Step~1), so the chain up to $Z_t$ lies in $\varepsilon_b\vee\varepsilon_c$ and $Z_t\subseteq\Delta_i$. Now $Z_t$ meets $Y_t\subseteq\supp x$ inside one part of $\mathcal Q$ (Step~3): inside a $V_k$ is impossible, since then $Z_t\subseteq\Delta_p$ for a $C_p\subseteq V_k$, so $p=i$ and $C_i\subseteq\supp v$; and inside an $\N(c)$-block $\tilde N$ is impossible since then $\tilde N\subseteq\Delta_i$, contradicting Step~4. So all blocks of $a$ in $D$ are blocks of $c$, $D$ is contained in $\Delta_i$, and $D$ contains no $A$ or $B$ other than $A_i$ and $B_i$. Moreover $B_i\in D$: a chain from $C_i=A_i$ to $B_i$ in $\Delta_i$ has interior $\varepsilon_c$-blocks which are $\N(a)$-blocks (Step~4), so it is a chain in $\varepsilon_a\vee\varepsilon_b$. Thus the classes of $\varepsilon_a\vee\varepsilon_b$ containing $A_i$ or $B_i$ with $i\notin I$ contain exactly one $\T(a)$-block and one $\T(b)$-block. In particular every class containing a block $A_i$ or $B_i$ with $i\in I$ contains no $A_j$, $B_j$ with $j\notin I$.

\emph{Step 6: consecutive blocks of $a$ inside $\supp x$.} Let $Y_1,Z_1,\dots,Z_r,Y_2$ be a chain in $\varepsilon_a\vee\varepsilon_b$ whose end blocks $Y_1,Y_2$ are blocks of $a$ inside $\supp x$ and whose interior blocks of $a$ lie in $Y$. The interior blocks of $a$ are $\N(x)$-blocks and $\N(c)$-blocks (Step~1, Step~5: no $A_j$ with $j\notin I$ occurs), so all $Z$'s lie in one class $K$ of $\varepsilon_x\vee\varepsilon_b$ and in one class of $\varepsilon_b\vee\varepsilon_c$; $K$ meets $\supp x$ (as $Z_1$ meets $Y_1$), and $Z_1\cap Y_1$, $Z_r\cap Y_2$ lie in the part $K\cap\supp x$ of $\mathcal Q$. If this part is $V_k$, then $Z_1\subseteq\Delta_p$, $Z_r\subseteq\Delta_{p'}$ with $Z_1\cap V_k\subseteq C_p$, $Z_r\cap V_k\subseteq C_{p'}$ (Step~4), $p=p'$ because $Z_1,Z_r$ lie in one class of $\varepsilon_b\vee\varepsilon_c$, and $Y_1,Y_2$ both meet $C_p$. If the part is an $\N(c)$-block $\tilde N$, then $Y_1,Y_2$ both meet $\tilde N$. In both cases \emph{$Y_1$ and $Y_2$ lie in a common class of $\varepsilon_a\vee\varepsilon_c$}.

\emph{Step 7: no class of $\varepsilon_a\vee\varepsilon_b$ contains two $\T(a)$-blocks or two $\T(b)$-blocks.} Let $A_i,A_j$ lie in a class $D$, $i\ne j$. By Step~5, $i,j\in I$, and $D$ contains no $A_p$ with $p\notin I$. Take a chain from $A_i$ to $A_j$ in $D$ and list its blocks of $a$ lying inside $\supp x$: the first is $A_i$, the last is $A_j$, and by Step~6 consecutive ones lie in a common class of $\varepsilon_a\vee\varepsilon_c$. Hence $\Gamma_i=\Gamma_j$, so $i=j$. Now let $B_i,B_j$ lie in a class $D$, $i\ne j$; again $i,j\in I$ by Step~5. Note that $B_i$ lies in the class $K_{k(i)}$ of $\varepsilon_x\vee\varepsilon_b$: a chain from $C_i$ to $B_i$ in $\Delta_i$ has interior $\varepsilon_c$-blocks which are $\N(x)$-blocks (Step~4), and its first $\varepsilon_b$-block meets $C_i\subseteq V_{k(i)}$. Take a chain from $B_i$ to $B_j$ in $D$. If it contains no block of $a$ inside $\supp x$, its blocks of $a$ are blocks of $c$ and $\Delta_i=\Delta_j$. Otherwise let $Y_1$ be the first block of $a$ inside $\supp x$ on the chain; the blocks between $B_i$ and $Y_1$ are $\varepsilon_b$-blocks and $\N(x)$-blocks, so they lie in $K_{k(i)}$ and in $\Delta_i$; the last of them, $Z$, meets $Y_1$ inside $K_{k(i)}\cap\supp x=V_{k(i)}$, and $Z\cap V_{k(i)}\subseteq C_i$ (Step~4, as $Z\subseteq\Delta_i$); so $Y_1$ meets $C_i$ and lies in $\Gamma_i$. Symmetrically the last block of $a$ inside $\supp x$ lies in $\Gamma_j$, and by Step~6 all of them lie in one class of $\varepsilon_a\vee\varepsilon_c$; so $\Gamma_i=\Gamma_j$ and $i=j$.

\emph{Step 8: $xb=xcb$.} For a class $K$ of $\varepsilon_x\vee\varepsilon_b$ let $B(K)=\{p:B_p\subseteq K\}$, and let $C(K)$ be the set of $p$ with $C_p\subseteq V_k=K\cap\supp x$ or $C_p\subseteq M\subseteq K$ for an $\N(x)$-block $M$. We claim $B(K)=C(K)$, that no class $K$ meeting $\supp x$ contains an $\N(x)$-block containing a $C$, and that a class not meeting $\supp x$ contains at most one such $\N(x)$-block. Take a shortest chain from $B_p$ to $C_p$ in $\Delta_p$; its interior $\varepsilon_c$-blocks are $\N(x)$-blocks (Step~4), so all its blocks except $C_p$ lie in the class $K$ of $B_p$, and the last $\varepsilon_b$-block $Z$ meets $C_p$. If $C_p\subseteq V_k$, the $\T(x)$-block containing a point of $Z\cap C_p$ lies in $K$, so $K=K_k$ and $p\in C(K)$; if $C_p\subseteq M$ then $M\subseteq K$ and again $p\in C(K)$. Reading the chain backwards gives the converse, so $B(K)=C(K)$. Next let $M\subseteq K$ be an $\N(x)$-block containing $C$'s (a union of $C$'s and of $A$'s), let $W$ be any block of $x$ in $K$ other than $M$, and take a shortest chain $M,Z_1,X_1,Z_2,\dots,Z_t,W$ in $K$. Let $C_p\subseteq M$ be a $\T(c)$-block met by $Z_1$; then $Z_1\subseteq\Delta_p$. We show inductively that all $Z_j$ lie in $\Delta_p$ and all $X_j$ are $\N(x)$-blocks which are $\N(c)$-blocks. If $X_j$ is a $\T(x)$-block $H$, then $Z_j$ meets $H$, which lies in a $V_k$ (then $Z_j$ meets a $C_q\subseteq V_k$, $q=p$, contradicting $C_p\subseteq M\subseteq Y$) or in an $\N(c)$-block $\tilde N$ (then $\tilde N\subseteq\Delta_p$, contradicting Step~4). If $X_j$ is an $\N(x)$-block containing $C$'s, it meets $Z_j\subseteq\Delta_p$ in a point of some $C_q$, so $q=p$ and $X_j\supseteq C_p$, i.e.\ $X_j=M$, contradicting minimality. So $X_j$ is an $\N(c)$-block, $X_j\subseteq\Delta_p$, and $Z_{j+1}\subseteq\Delta_p$. The same alternatives apply to $W$: $W$ is not a $\T(x)$-block, so $K$ does not meet $\supp x$; and if $W$ contains $C$'s then $W=M$. This proves the claim. Now compare the elements $xb$ and $xcb$ of $\P_X$. The upper blocks of $xb$ are the $\N(x)$-blocks and the sets $U_K$, transversal iff $K$ meets $\supp b$; the lower blocks are the $\N(b)$-blocks and, for each class $K$, $\bigcup_{p\in B(K)}B_p'$, attached to $U_K$ if $K$ meets $\supp x$. For $xcb=x(cb)$ the middle join is $\varepsilon_x\vee\varepsilon_c=\varepsilon_v$: the upper blocks are the $\N(x)$-blocks, the $V_k$ (transversal) and the $\N(c)$-blocks inside $\supp x$; the lower blocks are the $\N(b)$-blocks, $\bigcup_{C_p\subseteq V_k}B'_p$ attached to $V_k$, and $\bigcup_{C_p\subseteq M}B'_p$ for the $\N(x)$-blocks $M$ containing $C$'s. By Step~3 and the claim these coincide: $B(K_k)=C(K_k)=\{p:C_p\subseteq V_k\}$, and the classes $K$ not meeting $\supp x$ with $B(K)\ne\emptyset$ correspond bijectively to the $\N(x)$-blocks $M$ containing $C$'s, with $B(K)=\{p:C_p\subseteq M\}$. So $xb=xcb$, and multiplying by $a$ on the left, $ub=axb=axcb=ucb$.

\emph{Step 9: every $\T(a)$-block shares its class with a $\T(b)$-block.} By Step~7 the upper blocks of the element $ab$ are the $\N(a)$-blocks and the single blocks $A_p$, and $A_p$ is transversal in $ab$ iff its class $D(A_p)$ in $\varepsilon_a\vee\varepsilon_b$ contains a $\T(b)$-block, in which case the lower part of that block is the union of the $B$'s in $D(A_p)$. Multiply by $x$ on the left: the middle join is $\varepsilon_x\vee\varepsilon_a$, so the transversal block of $x(ab)$ at $U_l$ has lower part the union of the $\T(b)$-blocks lying in the classes $D(A_p)$, $p\in I_l=\{p:A_p\subseteq U_l\}$; by Step~7 each $D(A_p)$ contains at most one $\T(b)$-block, and distinct $p$ give distinct classes, so this is a union of at most $|I_l|$ distinct $\T(b)$-blocks, with equality iff every $D(A_p)$, $p\in I_l$, contains one. On the other hand $x(ab)=(xa)b=ub=ucb$ by Step~8, and in $u(cb)$ the middle join is $\varepsilon_u\vee\varepsilon_c=\varepsilon_x\vee\varepsilon_a\vee\varepsilon_c$, whose class containing $U_l$ contains the $C$'s of $V_{k(l)}$ and no other $C$ (Step~2(b); an $\N(x)$-block which is a union of $A$'s and $C$'s is a class by itself); so the transversal block of $ucb$ at $U_l$ has lower part $\bigcup_{p\in I_l}B_p$, a union of $|I_l|$ distinct $\T(b)$-blocks. Comparing the two lower parts, which are unions of pairwise disjoint nonempty sets, $D(A_p)$ contains a $\T(b)$-block for every $p\in I_l$ and every $l$, i.e.\ for every $p\in I$; for $p\notin I$ this is Step~5.

\emph{Conclusion.} The $s$ blocks $A_1,\dots,A_s$ lie in $s$ distinct classes of $\varepsilon_a\vee\varepsilon_b$ (Step~7), each of which contains at least one of the $s$ blocks $B_1,\dots,B_s$ (Step~9); so each contains exactly one, and no other class meets $\supp a\cup\supp b$. By Lemma~\ref{lem:PX}(c) (and its proof, which shows that this property characterises matched pairs), $a\bowtie b$.
\end{proof}

\begin{remark}\label{rem:Rproof}
The proof uses the hypothesis $\theta_x(b)=\theta_x(c)$ only through Step~3, i.e.\ through the partition $\mathcal Q$ of $\supp x$ common to $\varepsilon_x\vee\varepsilon_b$ and $\varepsilon_x\vee\varepsilon_c$; the commutation of $x$ with $a$ and with $c$ enters through the four types of Lemma~\ref{lem:PX}(b), which make every block of $a$ and of $c$ lie inside $\supp x$ or outside it and make $a$ and $c$ coincide outside $\supp x$ (Step~1). The mechanism is thus different from that of (S): there the commuting pair $c,d$ shares its nontransversal blocks and the matchings transport the transversal blocks across them; here the single projection $x$ splits $X$ into a part where $a$ and $c$ agree and a part where $b$ and $c$ induce the same coarse structure, and the two parts communicate only through the $\T(c)$-blocks (Step~4). The identity $xb=xcb$ of Step~8 holds in $\mathfrak P_7$ as well (for the elements of $\mathfrak T_{19}$), so it is not by itself an obstruction; what $\mathfrak T_{19}$ violates is Step~9, i.e.\ the way $ab$ and $ucb$ are compared through their transversal blocks.


Like \eqref{eq:Q1}, the triangle condition and (S), the condition (R) defines a proper quasivariety, i.e.\ the class it defines is not closed under $*$-homomorphic images: $\mathfrak T_{19}$ is a $*$-homomorphic image of a $*$-subsemigroup of $\P_5$. Indeed, let $\mathfrak T_{26}$ be defined like $\mathfrak T_{19}$ but with the bottom band $\{u,v,v',w\}\times\{u,v,v',w\}$ and $\lambda b=\lambda v'=\lambda w=v'$, $\lambda c=\lambda v=v$, $\lambda a=\lambda u=u$ (so $xbx=v'$ is no longer equal to $xcx=v$). Then $\mathfrak T_{26}$ is a regular $*$-semigroup, identifying $v'$ with $v$ is a $*$-homomorphism onto $\mathfrak T_{19}$, and $\mathfrak T_{26}$ $*$-embeds in $\P_5$ (its $8$-element projection algebra embeds in $\P(\P_5)$ by the SAT search, and since $\mathfrak T_{26}$ is generated by its projections the embedding extends; \texttt{fourth\_pre.py}). The same example shows that the class of regular $*$-semigroups embeddable in partition monoids is not closed under $*$-homomorphic images with a $19$-element image, the previous witnesses being the $32$-element $Q(a-c-b-d,\bot)$ with $\bot$ injective and its image $\mathfrak Q$.
\end{remark}

\subsection{A uniform view: path-closing conditions}\label{subsec:pathclose}

Both (S) and (R) have the same shape: a \emph{matched path} $p_0\bowtie p_1\bowtie\dots\bowtie p_n$, an auxiliary projection $x$ tied to some vertices of the path by commutation or order, two vertices which are ``indistinguishable through $x$'', and the conclusion that the path \emph{closes}, $p_0\bowtie p_n$:
\begin{align*}
\text{(R)}:\quad& p_0\bowtie p_1\bowtie p_2,\ xp_0=p_0x,\ xp_1=p_1x,\ \theta_x(p_1)=\theta_x(p_2)\ \Longrightarrow\ p_0\bowtie p_2,\\
\text{(S)}:\quad& p_0\bowtie p_1\bowtie p_2\bowtie p_3,\ x=p_1p_3=p_3p_1,\ \theta_{p_0}(x)=\theta_{p_2}(x)\ \Longrightarrow\ p_0\bowtie p_3 .
\end{align*}
In (R) the auxiliary projection \emph{sees} two consecutive vertices alike; in (S) it \emph{is seen} alike by the two even vertices, and it lies below the two odd ones. To find out whether this scheme contains further obstructions we searched it systematically (\texttt{pathclose\_np.py}): for a path of length $n$ and the atomic relations $xp_i=p_ix$, $x\le p_i$, $p_i\le x$, $\theta_x(p_i)=\theta_x(p_j)$, $\theta_{p_i}(x)=\theta_{p_j}(x)$, $x=\theta_{p_i}(p_j)$, we listed all sets of at most three atoms $H$ such that ``path $+$ $H$ $\Rightarrow$ $p_0\bowtie p_n$'' holds for all instances in $\P(\P_4)$ with distinct vertices and $x$ not a vertex, is not vacuous (some closed instance satisfies $H$), and is minimal; each survivor was then tested on more than a million random instances in $\P(\P_5)$ (\texttt{pathclose5.py}; none failed) and, by Mace4, for countermodels among the projection algebras of order $\le9$ satisfying \eqref{eq:Q1} alone, \eqref{eq:Q1}$+$(S), and \eqref{eq:Q1}$+$(S)$+$(R) (\texttt{pc\_mace.sh}). The outcome:
\begin{itemize}[nosep]
\item $n=2$: ten minimal conditions. Two of them are (R) and its mirror image; the other eight have no countermodel of order $\le9$ with \eqref{eq:Q1}, so they are, at least up to that order, consequences of \eqref{eq:Q1}.
\item $n=3$: $42$ minimal conditions. Fourteen are consequences of \eqref{eq:Q1} in the same sense; the other $28$ fail in $\mathfrak P$ (a countermodel of order $6$ with \eqref{eq:Q1}) and have no countermodel of order $\le9$ with \eqref{eq:Q1}$+$(S). Among them are (S), its mirror image, the variants in which $x=p_1p_3$ is weakened to $x\le p_1,\,x\le p_3$ or to $x\le p_1,\,xp_3=p_3x$, and the symmetric form
\[
\text{(S$'$)}:\quad p_0\bowtie p_1\bowtie p_2\bowtie p_3,\ \theta_{p_0}(x)=\theta_{p_2}(x),\ \theta_{p_1}(x)=\theta_{p_3}(x),\ x\le p_i\text{ or }xp_i=p_ix\text{ for some }i\ \Longrightarrow\ p_0\bowtie p_3,
\]
in which $x$ is seen alike by the two even and by the two odd vertices. (R) plays no role for $n=3$: adding it to \eqref{eq:Q1} produces no further consequences here.
\item $n=4$: $18$ minimal conditions, all without countermodels of order $\le9$ with \eqref{eq:Q1} alone.
\item Two auxiliary projections $x,y$ (\texttt{pathclose2.py}; the atoms above for each of $x,y$, together with $xy=yx$, $x\le y$, $y\le x$, $\theta_x(p_i)=\theta_y(p_i)$, $\theta_{p_i}(x)=\theta_{p_i}(y)$, $x=\theta_y(p_i)$, $y=\theta_x(p_i)$, $x=\theta_{p_i}(y)$, $y=\theta_{p_i}(x)$; $x,y$ distinct and not vertices). Length $2$: $60$ minimal conditions, $40$ consequences of \eqref{eq:Q1} and $20$ of \eqref{eq:Q1}$+$(R) (up to order $9$), all surviving the random test in $\P(\P_5)$; among the latter is the cleaner form
\[
p_0\bowtie p_1\bowtie p_2,\quad xp_0=p_0x,\quad y\le p_1,\quad \theta_x(p_2)=\theta_y(p_2)\ \Longrightarrow\ p_0\bowtie p_2
\]
(in (R) one has $y=xp_1$). Length $3$: $184$ minimal conditions in $\P(\P_4)$, of which $16$ fail in $\P(\P_5)$ (all of the form ``$x$ is seen alike by $p_0,p_2$ and by $p_1,p_3$, and $y=\theta_x(p_i)$''); of the remaining $168$, $52$ are consequences of \eqref{eq:Q1} and $116$ of \eqref{eq:Q1}$+$(S) up to order $9$, none needing (R). Typical members of the second group are
\[
p_0\bowtie p_1\bowtie p_2\bowtie p_3,\quad xp_0=p_0x,\quad \theta_{p_1}(x)=\theta_{p_3}(x),\quad \theta_{p_0}(y)=\theta_{p_2}(y)\ \Longrightarrow\ p_0\bowtie p_3 ,
\]
with two different auxiliary projections seen alike by the odd and by the even vertices. No condition with a countermodel of order $\le9$ satisfying \eqref{eq:Q1}, (S) and (R) was found.
\end{itemize}
So within this scheme (one or two auxiliary projections, at most three atoms, paths of length $\le4$ resp.\ $\le3$) the only obstructions beyond \eqref{eq:Q1} are (R), at length $2$, and (S), at length $3$, each with a cluster of equivalent forms; and, curiously, (R) is of ``sees'' type and (S) of ``seen'' type, the ``seen'' conditions of length $2$ and the ``sees'' conditions of length $3$ being consequences of \eqref{eq:Q1}. The triangle condition is of a different shape (two commuting triples, an equality as conclusion) and needs six variables; whether it, too, is the first member of a family, and whether two auxiliary projections or longer paths produce further path-closing obstructions, is open.


\subsection{The census of order $8$: seven more candidate obstructions}\label{subsec:census8}

The chunked Mace4 enumeration of Section~\ref{subsec:fourth} does not scale to order $8$ (the first of $4096$ chunks alone has millions of labelled models). The following observation replaces it by a census ``by extension''.

\begin{lemma}\label{lem:extension}
Let $P$ be a finite projection algebra with $|P|\ge2$ and let $e$ be a maximal element of $(P,\le)$. Then $P\setminus\{e\}$ is a subalgebra. Consequently, if $\mathcal K$ is a class of projection algebras defined by quasi-identities, the members of $\mathcal K$ of order $n+1$ are exactly the one-point extensions, by a new maximal element, of the members of $\mathcal K$ of order $n$.
\end{lemma}

\begin{proof}
If $\theta_q(p)=e$ then $e\le q$, so $q=e$. Subalgebras of members of $\mathcal K$ are in $\mathcal K$.
\end{proof}

For each algebra $B$ of order $n$ Mace4 is given the table of $B$ as ground equations and finds all values $\theta_e(q)$, $\theta_q(e)$ ($q\in B$) with $e$ maximal, subject to the axioms, \eqref{eq:Q1}, (S) and (R); the results are reduced by a canonical form (\texttt{pa8/extend.py}, a fraction of a second per base algebra). Starting from the one-element algebra, the numbers of projection algebras satisfying \eqref{eq:Q1}, (S) and (R) are
\[
1,\ 2,\ 5,\ 16,\ 62,\ 295,\ 1\,662,\ 10\,828,\ 80\,039,\ 658\,727\qquad(\text{orders }1,\dots,10),
\]
in agreement with the direct enumerations of orders $6$ and $7$ ($1\,662=1\,663-1$, the missing algebra being $\mathfrak P_7$, which violates (R)); the step from order $8$ to order $9$ takes about $40$ minutes on two cores, the step from order $9$ to order $10$ ($1\,849\,833$ labelled extensions) $8.2$ hours on one core; the analysis of the $658\,727$ algebras of order $10$ has not been done. Every algebra of order $8$ in the list satisfies the triangle condition, and $10\,821$ of them embed in $\P(\P_m)$ with $m\le7$ (\texttt{pa8/post8.py}: backtracking for $m\le4$, SAT for $5\le m\le7$). The remaining seven do not embed in $\P(\P_m)$ for $m\le8$ (SAT; the case $m=8$ takes from $25$ seconds for $\mathfrak P_8^2$ to $2.6$ hours for $\mathfrak P_8^7$; \texttt{pa8/sat8.py}). Since every algebra of order $7$ satisfying \eqref{eq:Q1}, (S) and (R) embeds, all proper subalgebras of the seven embed; so each of them, if it embeds in no $\P(\P_X)$ at all, is a minimal non-embeddable algebra. Rows $p$, columns $q$, entries $\theta_p(q)$:
\begin{gather*}
\scriptsize
\mathfrak P_8^{1}:\ \begin{array}{c|cccccccc} & 0 & 1 & 2 & 3 & 4 & 5 & 6 & 7\\ \hline 0 & 0 & 0 & 0 & 0 & 0 & 0 & 0 & 0\\ 1 & 0 & 1 & 0 & 0 & 0 & 0 & 6 & 6\\ 2 & 0 & 0 & 2 & 0 & 0 & 2 & 0 & 2\\ 3 & 0 & 0 & 0 & 3 & 0 & 3 & 0 & 3\\ 4 & 4 & 4 & 4 & 4 & 4 & 4 & 4 & 4\\ 5 & 4 & 4 & 5 & 5 & 4 & 5 & 4 & 5\\ 6 & 6 & 6 & 6 & 6 & 6 & 6 & 6 & 6\\ 7 & 6 & 6 & 7 & 7 & 6 & 7 & 6 & 7\\\end{array}\qquad \mathfrak P_8^{2}:\ \begin{array}{c|cccccccc} & 0 & 1 & 2 & 3 & 4 & 5 & 6 & 7\\ \hline 0 & 0 & 0 & 0 & 0 & 0 & 0 & 0 & 0\\ 1 & 0 & 1 & 0 & 0 & 6 & 6 & 6 & 6\\ 2 & 0 & 0 & 2 & 0 & 0 & 0 & 0 & 2\\ 3 & 0 & 0 & 0 & 3 & 0 & 3 & 0 & 3\\ 4 & 4 & 4 & 4 & 4 & 4 & 4 & 4 & 4\\ 5 & 4 & 4 & 4 & 5 & 4 & 5 & 4 & 4\\ 6 & 6 & 6 & 6 & 6 & 6 & 6 & 6 & 6\\ 7 & 6 & 6 & 7 & 7 & 6 & 6 & 6 & 7\\\end{array}\\[1ex]
\mathfrak P_8^{3}:\ \begin{array}{c|cccccccc} & 0 & 1 & 2 & 3 & 4 & 5 & 6 & 7\\ \hline 0 & 0 & 0 & 0 & 0 & 0 & 0 & 0 & 0\\ 1 & 0 & 1 & 0 & 0 & 0 & 1 & 0 & 1\\ 2 & 0 & 0 & 2 & 2 & 0 & 0 & 0 & 2\\ 3 & 0 & 0 & 2 & 3 & 4 & 4 & 0 & 2\\ 4 & 4 & 4 & 4 & 4 & 4 & 4 & 4 & 4\\ 5 & 4 & 5 & 4 & 4 & 4 & 5 & 4 & 4\\ 6 & 6 & 6 & 6 & 6 & 6 & 6 & 6 & 6\\ 7 & 6 & 7 & 7 & 7 & 6 & 6 & 6 & 7\\\end{array}\qquad \mathfrak P_8^{4}:\ \begin{array}{c|cccccccc} & 0 & 1 & 2 & 3 & 4 & 5 & 6 & 7\\ \hline 0 & 0 & 0 & 0 & 0 & 0 & 0 & 0 & 0\\ 1 & 0 & 1 & 0 & 6 & 6 & 6 & 6 & 6\\ 2 & 0 & 0 & 2 & 0 & 0 & 2 & 0 & 0\\ 3 & 3 & 3 & 3 & 3 & 3 & 3 & 3 & 3\\ 4 & 3 & 3 & 3 & 3 & 4 & 3 & 3 & 4\\ 5 & 3 & 3 & 5 & 3 & 3 & 5 & 3 & 5\\ 6 & 6 & 6 & 6 & 6 & 6 & 6 & 6 & 6\\ 7 & 6 & 6 & 6 & 6 & 7 & 7 & 6 & 7\\\end{array}\\[1ex]
\mathfrak P_8^{5}:\ \begin{array}{c|cccccccc} & 0 & 1 & 2 & 3 & 4 & 5 & 6 & 7\\ \hline 0 & 0 & 0 & 0 & 0 & 0 & 0 & 0 & 0\\ 1 & 0 & 1 & 0 & 0 & 0 & 1 & 0 & 0\\ 2 & 0 & 0 & 2 & 6 & 7 & 7 & 6 & 7\\ 3 & 3 & 3 & 3 & 3 & 3 & 3 & 3 & 3\\ 4 & 3 & 3 & 4 & 3 & 4 & 3 & 3 & 4\\ 5 & 3 & 5 & 5 & 3 & 3 & 5 & 3 & 5\\ 6 & 6 & 6 & 6 & 6 & 6 & 6 & 6 & 6\\ 7 & 6 & 6 & 7 & 6 & 7 & 7 & 6 & 7\\\end{array}\qquad \mathfrak P_8^{6}:\ \begin{array}{c|cccccccc} & 0 & 1 & 2 & 3 & 4 & 5 & 6 & 7\\ \hline 0 & 0 & 0 & 0 & 0 & 0 & 0 & 0 & 0\\ 1 & 0 & 1 & 0 & 0 & 0 & 0 & 0 & 1\\ 2 & 0 & 0 & 2 & 0 & 2 & 2 & 0 & 2\\ 3 & 3 & 3 & 3 & 3 & 3 & 3 & 3 & 3\\ 4 & 3 & 3 & 4 & 3 & 4 & 4 & 3 & 3\\ 5 & 3 & 3 & 4 & 3 & 4 & 5 & 6 & 6\\ 6 & 6 & 6 & 6 & 6 & 6 & 6 & 6 & 6\\ 7 & 6 & 7 & 7 & 6 & 6 & 6 & 6 & 7\\\end{array}\\[1ex]
\mathfrak P_8^{7}:\ \begin{array}{c|cccccccc} & 0 & 1 & 2 & 3 & 4 & 5 & 6 & 7\\ \hline 0 & 0 & 0 & 0 & 0 & 0 & 0 & 0 & 0\\ 1 & 0 & 1 & 0 & 3 & 3 & 0 & 3 & 0\\ 2 & 0 & 0 & 2 & 0 & 2 & 5 & 5 & 5\\ 3 & 3 & 3 & 3 & 3 & 3 & 3 & 3 & 3\\ 4 & 3 & 3 & 4 & 3 & 4 & 6 & 6 & 3\\ 5 & 5 & 5 & 5 & 5 & 5 & 5 & 5 & 5\\ 6 & 6 & 6 & 6 & 6 & 6 & 6 & 6 & 6\\ 7 & 7 & 7 & 7 & 7 & 7 & 7 & 7 & 7\\\end{array}
\end{gather*}

Six of the seven are one-point extensions of projection algebras of graph semigroups. Deleting the element $x$ indicated below leaves the projection algebra of $Q(G,\bot)$ with $|V|=4$ and three colours (\texttt{pa8/cand8b.py}):
\begin{itemize}[nosep]
\item $\mathfrak P_8^1$, $x=1$: $V=\{2,3,5,7\}$, $E=\{25,27,35,37,57\}$ ($K_4$ minus the edge $23$), $\bot\colon2,3\mapsto0$, $5\mapsto4$, $7\mapsto6$; $x$ lies above $0$ and $6$, commutes with $2,3$ (product $0$) and with $7$ (product $6$), and $\theta_x(5)=0$, $\theta_5(x)=4$.
\item $\mathfrak P_8^2$, $x=1$: same $V$ and $\bot$, $E=\{27,35,37\}$ (a path $2-7-3-5$); $x$ as in $\mathfrak P_8^1$ except $\theta_x(5)=6$ (and $\theta_x(4)=6$).
\item $\mathfrak P_8^3$, $x=3$: $V=\{1,2,5,7\}$, $E=\{15,17,27\}$ (a path $5-1-7-2$), $\bot\colon1,2\mapsto0$, $5\mapsto4$, $7\mapsto6$; $x$ lies above $0$, $4$ and the vertex $2$, commutes with $1$ (product $0$) and with $5$ (product $4$), and $\theta_x(7)=2$, $\theta_7(x)=7$.
\item $\mathfrak P_8^4$, $x=1$: $V=\{2,4,5,7\}$, $E=\{25,47,57\}$ (a path $2-5-7-4$), $\bot\colon2\mapsto0$, $4,5\mapsto3$, $7\mapsto6$; $x$ lies above $0$ and $6$, commutes with $2$ (product $0$) and with $7$ (product $6$), and $\theta_x(4)=\theta_x(5)=6$, $\theta_4(x)=\theta_5(x)=3$.
\item $\mathfrak P_8^5$, $x=2$: $V=\{1,4,5,7\}$, $E=\{15,47,57\}$ (a path $1-5-7-4$), $\bot\colon1\mapsto0$, $4,5\mapsto3$, $7\mapsto6$; $x$ lies above $0$, $6$ and the vertex $7$, commutes with $1$ (product $0$), and $\theta_x(4)=\theta_x(5)=7$, $\theta_4(x)=4$, $\theta_5(x)=5$.
\item $\mathfrak P_8^6$, $x=5$: $V=\{1,2,4,7\}$, $E=\{17,24,27\}$ (a path $1-7-2-4$), $\bot\colon1,2\mapsto0$, $4\mapsto3$, $7\mapsto6$; $x$ lies above $3$, $6$ and the vertex $4$, commutes with $7$ (product $6$), and $\theta_x(1)=3$, $\theta_1(x)=0$, $\theta_x(2)=4$, $\theta_2(x)=2$.
\end{itemize}
The seventh, $\mathfrak P_8^7$, is different: its bottom band has five elements $0,3,5,6,7$, and above it lie $1\ge0,3$, $2\ge0,5$, $4\ge3,6$ with $2\bowtie4$, $1$ commuting with $2$ (product $0$) and with $4$ (product $3$); the element $7$ lies below no other element, and $\theta_1(7)=0$, $\theta_2(7)=5$, $\theta_4(7)=3$.

Each of the seven violates a quasi-identity valid in $\P(\P_6)$, found by the methods of Section~\ref{subsec:pathclose} (\texttt{pa8/pc3\_hit.py}: for an instance in $\mathfrak P_8^k$ of a matched path with auxiliaries, the minimal subsets of the atoms true there that are contained in no counterexample signature of $\P(\P_4)$, then a SAT check of validity without distinctness assumptions, \texttt{pa8/pc3\_verify.py}) or by greedy minimisation of the diagram of $\mathfrak P_8^k$ under a SAT validity test (\texttt{pa8/qi\_extract.py}). Since $\mathfrak P_8^k$ satisfies \eqref{eq:Q1}, the triangle condition, (S) and (R), each of these quasi-identities is independent of those four.

\emph{Caveat.} Apart from $\mathfrak P_8^7$, $\mathfrak P_8^3$, $\mathfrak P_8^2$, $\mathfrak P_8^4$ and $\mathfrak P_8^5$, whose conditions (V), (R$'$), (P$_1$), (P$_2$) and (Q) are proved in Theorems~\ref{thm:V}, \ref{thm:Rp}, \ref{thm:P1} and \ref{thm:Q} below, none of the seven has been \emph{proved} non-embeddable, and none of the other quasi-identities below has been proved for all $\P_X$: the evidence is the absence of embeddings into $\P(\P_m)$ for $m\le8$ and the validity of the quasi-identities in $\P(\P_m)$ for $m\le7$. That a finite algebra embeds in some $\P(\P_X)$ but in no $\P(\P_m)$ with $m$ small does happen: a semilattice with $k$ join-irreducible elements embeds in $\P(\P_k)$ (partial identities) and, e.g.\ for the $9$-element semilattice consisting of chains of lengths $1$, $4$ and $3$ glued at the bottom, in no $\P(\P_m)$ with $m\le7$ (SAT); so in the census of order $9$ an SAT search with $m\le7$ is not conclusive, and the SAT results below for the two algebras $\mathfrak P_8^1,\mathfrak P_8^6$ are evidence, not proof.

\begin{proposition}\label{prop:census8}
The following quasi-identities hold in $\P(\P_m)$ for $m\le7$ (SAT, variables not assumed distinct; \texttt{pa8/verify6.py}, the case $m=6$ taking between $20$ seconds and $12$ minutes each; \texttt{pa8/verify7b.py} for $m=7$ on M.K.'s machine: (Q), and hence (P$_2$), $11.3$ hours, (E$_6$) $21$ minutes, (E$_1$) $3$ minutes; (V), (R$'$), (P$_1$), (Q) and hence (P$_2$) are proved below, in Theorems~\ref{thm:V}, \ref{thm:Rp}, \ref{thm:P1}, \ref{thm:Q}).
\begin{enumerate}[label=\textup{(\alph*)},nosep]
\item \textup{(R$'$)}: $a\bowtie c,\ b\bowtie c,\ xc=cx,\ \theta_x(a)=\theta_x(c),\ ya=ay,\ \theta_x(b)=\theta_y(b)\ \Longrightarrow\ a\bowtie b$. It fails in $\mathfrak P_8^3$ ($(a,c,b,x,y)=(5,1,7,2,3)$) and in $\mathfrak P_7$ ($(a,c,b,x,y)=(a,c,b,v,x)$). (R$'$) holds in every $\P_X$ (Theorem~\ref{thm:Rp}). In every projection algebra (R$'$) implies (R): given the hypotheses of (R), put $x'=xc=cx$ and $y'=x$. Then $x'\le c$, so $x'c=cx'$; $\theta_{x'}(a)=cx\,a\,xc=c(xax)c=c(xa)c=x(cac)=xc=\theta_{x'}(c)$, using $xa=ax$, $xc=cx$ and $cac=c$; $y'=x$ commutes with $a$; and $\theta_{x'}(b)=c\,\theta_x(b)\,c=c\,\theta_x(c)\,c=cxc=xc=\theta_x(c)=\theta_x(b)=\theta_{y'}(b)$. So (R$'$) gives $a\bowtie b$. (This substitution was found by Prover9.) Each of the four hypotheses on $x,y$ is needed, and with $y=x$ the hypotheses force $a=c$ (since then $xa=xc\le a,c$ and matched projections with a common lower bound are equal by \eqref{eq:Q1}).
\item \textup{(V)}: $bcb=b,\ ata\le b,\ ctc\le a\ \Longrightarrow\ ata\le btb$. It fails in $\mathfrak P_8^7$ ($(a,b,c,t)=(1,2,4,7)$: $ata=0$, $btb=5$, $ctc=3$). (V) holds in every $\P_X$, with $ata=btb$ (Theorem~\ref{thm:V}). In $\P(\P_m)$, $m\le6$, the hypotheses of (V) also force $tat=tbt$ (SAT), and each of the three hypotheses is needed for $tat=tbt$ already in $\P(\P_4)$; but $tat=tbt$ together with $ata\le b$ does not give $ata\le btb$, so this is not a decomposition of (V) into two steps.
\item \textup{(P$_1$)}: $p_0\bowtie p_1\bowtie p_2\bowtie p_3,\ xp_0=p_0x,\ xp_1=p_1x,\ \theta_{p_0}(x)=\theta_{p_2}(x),\ \theta_x(p_1)=\theta_x(p_3)\ \Longrightarrow\ p_0\bowtie p_3$, and \textup{(P$_2$)}: $p_0\bowtie p_1\bowtie p_2\bowtie p_3,\ xp_0=p_0x,\ xp_2=p_2x,\ \theta_x(p_1)=\theta_x(p_2),\ \theta_{p_1}(x)=\theta_{p_3}(x)\ \Longrightarrow\ p_0\bowtie p_3$. (P$_1$) fails in $\mathfrak P_8^2$ ($(p_0,\dots,p_3;x)=(2,7,3,5;1)$) and (P$_2$) in $\mathfrak P_8^4$ ($(2,5,7,4;1)$); (P$_1$) holds in every $\P_X$ (Theorem~\ref{thm:P1}). More precisely (\texttt{pa8/pc1\_classify.py}): with one auxiliary projection, a path of length $3$, the atoms of Section~\ref{subsec:pathclose} and the chords $p_0\bowtie p_2$, $p_1\bowtie p_3$ as further atoms, there are $76$ minimal valid hypothesis sets of size $\le4$ and $8$ of size $5$, all valid in $\P(\P_5)$; exactly $8$ of size $4$ and $2$ of size $5$ have a countermodel among the $12\,870$ algebras of order $2\le n\le8$ satisfying \eqref{eq:Q1}, (S) and (R), always a single one, $\mathfrak P_8^2$ or $\mathfrak P_8^4$. Up to reversal of the path the eight of size $4$ are (P$_1$), (P$_2$), and
\begin{align*}
&xp_0=p_0x,\ xp_2=p_2x,\ \theta_{p_1}(x)=\theta_{p_3}(x),\ \theta_x(p_2)=\theta_x(p_3)\quad(\mathfrak P_8^4),\\
&xp_1=p_1x,\ xp_2=p_2x,\ \theta_x(p_0)=\theta_x(p_2),\ \theta_{p_1}(x)=\theta_{p_3}(x)\quad(\mathfrak P_8^2);
\end{align*}
the two of size $5$ are $xp_0=p_0x,\ xp_1=p_1x,\ xp_2=p_2x,\ \theta_x(p_0)=\theta_x(p_2),\ \theta_x(p_1)=\theta_x(p_3)$ and its reversal ($\mathfrak P_8^2$). All of them combine two commutation atoms, one ``seen alike'' atom for two vertices of the same parity and one ``sees alike'' atom.
\item \textup{(Q)}: $p_0\bowtie p_1\bowtie p_2\bowtie p_3,\ xp_0=p_0x,\ \theta_{p_1}(x)=\theta_{p_3}(x),\ yp_0=p_0y,\ yp_2=p_2y,\ \theta_y(p_1)=\theta_y(p_2)\ \Longrightarrow\ p_0\bowtie p_3$, which fails in $\mathfrak P_8^5$ ($(p_0,\dots,p_3;\,x,y)=(1,5,7,4;\,0,2)$) and also in $\mathfrak P_8^4$ ($(2,5,7,4;\,0,1)$). No condition of the scheme of Section~\ref{subsec:pathclose} with two auxiliaries, a path of length $\le3$, chords and at most $5$ atoms fails in $\mathfrak P_8^1$ or $\mathfrak P_8^6$.
\item For $\mathfrak P_8^6$: $dbd=d,\ aca=a,\ ada=a,\ u\le c,\ u\le d,\ bcb=aua\ \Longrightarrow\ aua=a$ (fails with $(a,b,c,d,u)=(2,1,5,7,6)$; the conclusion was found in the form $aua\bowtie a$, which is equivalent because $\theta_a(aua)=aua$). For $\mathfrak P_8^1$: $z\le b,\ aba=z,\ dbd=d,\ xcx=z,\ aca=a,\ bcb=b,\ e\le x,\ e\le d,\ ada=a\ \Longrightarrow\ aba=a$ (fails with $(a,b,c,d,x,e,z)=(2,3,5,7,1,6,0)$; here $e$ only serves, via \eqref{eq:Q1}, to make $x$ and $d$ commute). Both are greedy-minimal subsets of the diagram of the algebra, and need not be the most natural formulations.
\end{enumerate}
\end{proposition}


\begin{theorem}\label{thm:V}
(V) holds in every partition monoid $\P_X$: if $b,c,a,t$ are projections with $bcb=b$, $ata\le b$ and $ctc\le a$, then $ata=btb$; in particular $ata\le btb$. Consequently $\mathfrak P_8^7$ embeds in no $\P(\P_X)$, and it is a minimal non-embeddable projection algebra.
\end{theorem}

\begin{proof}
Write $z=ata$, $w=ctc$, $s=btb$. We first reduce the claim, by computations valid in every projection algebra satisfying \eqref{eq:Q1}, to a single statement about $\P_X$.

\smallskip
\emph{Step 1: three identities.} (i) \emph{If $k\le a$ and $ab=ba$ then $bkb\le a$}: indeed $a(bkb)a=(ab)k(ba)=(ba)k(ab)=b(aka)b=bkb$. (ii) \emph{If $ata\le b$ then $ab=ba$ and $ata=ete$, where $e=ab$}: $z=ata$ lies below $a$ and below $b$, so $a$ and $b$ commute by \eqref{eq:Q1}; $e=ab$ is a projection with $e\le a$, $e\le b$, and $eze=a(bzb)a=aza=z$, so $z\le e$ and $z=eze=e(ata)e=(ea)t(ae)=ete$ because $ea=ae=e$. (iii) \emph{If $bcb=b$ then $s=b(csc)b$}: $s=btb=(bcb)t(bcb)=b\,c(btb)c\,b$.

\emph{Step 2: reduction.} Suppose we know that
\begin{equation}\label{eq:lemD}
csc=w,\qquad\text{i.e.}\qquad c(btb)c=ctc .
\end{equation}
Then $s=b(csc)b=bwb$ by (iii), and $bwb\le a$ by (i), since $w\le a$ and $ab=ba$ (by (ii), as $z\le b$). So $s\le a$ and $s\le b$, hence $s\le e=ab$ and $s=ese=e(btb)e=(eb)t(be)=ete$; and $z=ete$ by (ii). Thus $s=z$. It remains to prove \eqref{eq:lemD} in $\P_X$; in fact we show that $\theta_c(s)$ and $\theta_c(t)$ have the same blocks.

\smallskip
\emph{Step 3: the setting in $\P_X$.} We use Lemma~\ref{lem:PX}(a) and the notation of the proof of Theorem~\ref{thm:R}: a chain in a join $\varepsilon_p\vee\varepsilon_q$ is a sequence of blocks, alternately of $p$ and of $q$, consecutive ones meeting; two blocks lie in a common class iff they are joined by a chain. Recall also that for projections $p,q$:
\begin{itemize}[nosep]
\item[($\alpha$)] $q\le p$ iff every block of $p$ lies in a block of $q$ and every $\N(p)$-block is a block of $q$; then every block of $q$ is an $\N(p)$-block or a union of $\T(p)$-blocks, and the $\N(p)$-blocks are $\N(q)$-blocks. (If $q=pqp$, this is Lemma~\ref{lem:PX}(a); conversely, if the two conditions hold then $\varepsilon_p\vee\varepsilon_q=\varepsilon_q$ and every $\varepsilon_q$-class meeting $\supp p$ is a union of $\T(p)$-blocks, so $pqp=q$.)
\item[($\beta$)] $pqp=p$ iff every class of $\varepsilon_p\vee\varepsilon_q$ meeting $\supp p$ contains exactly one $\T(p)$-block and meets $\supp q$ (Lemma~\ref{lem:PX}(a)). In particular every block of $q$ meets at most one $\T(p)$-block.
\end{itemize}
Put $J_a=\varepsilon_a\vee\varepsilon_t$, $J_b=\varepsilon_b\vee\varepsilon_t$, $J_c=\varepsilon_c\vee\varepsilon_t$. By Lemma~\ref{lem:PX}(a) the blocks of $z$ are the $\N(a)$-blocks and the sets $Z_K=K\cap\supp a$ for the classes $K$ of $J_a$ meeting $\supp a$, $Z_K$ being transversal iff $K$ meets $\supp t$; similarly the blocks of $w$ are the $\N(c)$-blocks and the sets $W_L=L\cap\supp c$ ($L\in J_c$), and those of $s$ are the $\N(b)$-blocks and the sets $S_H=H\cap\supp b$ ($H\in J_b$). From $w\le a$ and ($\alpha$): every block of $a$ lies in a block of $w$; every block of $w$ is an $\N(a)$-block or a union of $\T(a)$-blocks; and every $\N(a)$-block is an $\N(w)$-block, i.e.\ an $\N(c)$-block or a set $W_L$ with $L\cap\supp t=\emptyset$. From $z\le b$: every block of $b$ lies in a block of $z$; every block of $z$ is an $\N(b)$-block or a union of $\T(b)$-blocks; and every $\N(b)$-block is a block of $z$, non-transversal. From $bcb=b$ and ($\beta$): every block of $c$ meets at most one $\T(b)$-block, and for every $\T(b)$-block $B$ the class $\Gamma_B$ of $\varepsilon_b\vee\varepsilon_c$ containing $B$ meets $\supp c$. Since blocks of $a$ lie in blocks of $w$, which lie in classes of $J_c$, we have $J_a\le J_c$ (every class of $J_a$ lies in a class of $J_c$); likewise $J_b\le J_a$. For a class $K$ of $J_a$ let $\lambda(K)$ be the class of $J_c$ containing it. Call a block of $c$ \emph{inner} if it meets $\supp a$ and \emph{outer} otherwise. An inner $\T(c)$-block lies in a set $W_L$ which is a union of $\T(a)$-blocks; an inner $\N(c)$-block is a union of $\T(a)$-blocks (being a block of $w$ that is not an $\N(a)$-block); an outer block of $c$ inside a class $L$ of $J_c$ with $W_L\ne\emptyset$ is either a $\T(c)$-block inside $W_L$, $W_L$ being then an $\N(a)$-block, or an $\N(c)$-block which is an $\N(a)$-block.

\emph{Step 4: two facts.} (1) \emph{No $\T(b)$-block lies in a block $M$ which is both an $\N(a)$-block and an $\N(c)$-block.} Otherwise let $B\subseteq M$ be one. The only block of $c$ meeting $B$ is $M$, and the blocks of $b$ meeting $M$ lie inside $M$ ($M$ is a block of $z$, and blocks of $b$ lie in blocks of $z$); so $\Gamma_B=M$ does not meet $\supp c$, contradicting ($\beta$). Consequently every block which is both an $\N(a)$-block and an $\N(c)$-block is an $\N(b)$-block (a block of $z$ containing no $\T(b)$-block).

(2) \emph{Let $K$ be a class of $J_a$ and $M\subseteq K$ an $\N(a)$-block containing a $\T(b)$-block. Then $M$ is the only block of $z$ inside $K$ that meets $\supp b$; in particular $Z_K$, if nonempty, is an $\N(b)$-block.} By (1) $M$ is not an $\N(c)$-block, so $M=W_L$ with $L=\lambda(K)$, $L\cap\supp t=\emptyset$; this already shows that no other $\N(a)$-block in $K$ contains a $\T(b)$-block (it would also equal $W_{\lambda(K)}$). Suppose $Z_K$ contains a $\T(b)$-block $B$. Every $\T(a)$-block in $K$ lies in $L$ but not in $W_L=M$ (which is disjoint from $\supp a$), hence in an $\N(c)$-block, which is then inner, a union of $\T(a)$-blocks, and contained in $L$. Take a chain $B=B_0,C_1,B_1,C_2,\dots$ from $B$ to a $\T(c)$-block in $\Gamma_B$ (it exists by ($\beta$)). $C_1$ meets $B\subseteq Z_K\subseteq\supp a$ in a point of a $\T(a)$-block of $K$, so $C_1$ is an inner $\N(c)$-block, $C_1\subseteq L\cap\supp a$; then $B_1$ meets $\supp a$ in a point of a $\T(a)$-block $A_1\subseteq C_1\subseteq L$, whose class $K_1$ satisfies $\lambda(K_1)=L$, and $B_1\subseteq Z_{K_1}$; so $C_2$ meets a $\T(a)$-block of $K_1$, which lies in $L$ but not in $W_L$, and $C_2$ is again an inner $\N(c)$-block. Inductively no block of the chain is a $\T(c)$-block, a contradiction.

\emph{Step 5: the blocks of $\theta_c(s)$ and $\theta_c(t)$ inside $\supp c$.} By Lemma~\ref{lem:PX}(a) the blocks of $\theta_c(s)$ are the $\N(c)$-blocks and the sets $\Theta\cap\supp c$ for the classes $\Theta$ of $\varepsilon_c\vee\varepsilon_s$ meeting $\supp c$, transversal iff $\Theta$ meets $\supp s$; the blocks of $\theta_c(t)=w$ are the $\N(c)$-blocks and the $W_L$, transversal iff $L$ meets $\supp t$. Every block of $s$ lies in a class of $J_c$: an $\N(b)$-block is a block of $z$, hence lies in a class of $J_a$, and $S_H\subseteq H$, a class of $J_b\le J_a\le J_c$. Hence $\varepsilon_c\vee\varepsilon_s\le J_c$: every class $\Theta$ lies in a class $L$ of $J_c$, $\Theta\cap\supp c\subseteq W_L$, and if $\Theta$ meets $\supp s$ then $L$ meets $\supp t$ (a block $S_H\subseteq\Theta$ is transversal only if $H\subseteq L$ meets $\supp t$). So \eqref{eq:lemD} follows once we show, for every class $L$ of $J_c$ with $W_L\ne\emptyset$: \emph{(A) any two $\T(c)$-blocks in $W_L$ lie in a common class of $\varepsilon_c\vee\varepsilon_s$, and (B) if $L$ meets $\supp t$ then this class meets $\supp s$.} Fix $L$; then $W_L$ is a union of $\T(a)$-blocks (\emph{type T}) or an $\N(a)$-block $M$ (\emph{type N}, and then $L\cap\supp t=\emptyset$, so (B) is vacuous).

\emph{Type T.} Let $C,C'$ be $\T(c)$-blocks in $W_L$ and take a chain $C=C_0,T_1,C_1,\dots,T_r,C_r=C'$ in $J_c$ (blocks $C_i$ of $c$, blocks $T_i$ of $t$, all inside $L$). We claim that consecutive inner blocks of $c$ on the chain, say $C_i$ and $C_j$ with $C_{i+1},\dots,C_{j-1}$ outer, lie in a common class of $\varepsilon_c\vee\varepsilon_s$; since $C_0$ and $C_r$ are inner ($W_L\subseteq\supp a$), (A) follows. The outer blocks $C_{i+1},\dots,C_{j-1}$ are $\N(c)$-blocks which are $\N(a)$-blocks (the $\T(c)$-blocks of $L$ lie in $W_L\subseteq\supp a$), hence $\N(b)$-blocks by (1); the blocks $T_{i+1},C_{i+1},\dots,T_j$ lie in a single class $K''$ of $J_a$ (each lies in a class, and consecutive ones meet). Let $x\in T_{i+1}\cap C_i$ and $y\in T_j\cap C_j$; both lie in $\supp a$ ($C_i,C_j$ are inner: a $\T(c)$-block inside $W_L$, or an $\N(c)$-block which is a union of $\T(a)$-blocks), so $x,y\in Z_{K''}$. If $Z_{K''}$ is an $\N(b)$-block, it is a block of $s$ meeting $C_i$ and $C_j$. Otherwise $Z_{K''}$ is a union of $\T(b)$-blocks; let $B_x\ni x$, $B_y\ni y$ be the $\T(b)$-blocks; then $B_x,T_{i+1},C_{i+1},T_{i+2},\dots,T_j,B_y$ is a chain in $J_b$, so $B_x,B_y\subseteq H$ for one class $H$ of $J_b$, and the block $S_H\supseteq B_x\cup B_y$ of $s$ meets $C_i$ and $C_j$. This proves the claim. For (B), let $L$ meet $\supp t$, and take a chain $C=C_0,T_1,C_1,\dots,C_q,T_{q+1}$ in $L$ from $C$ to a $\T(t)$-block $T_{q+1}$. Let $C_j$ ($j\le q$) be the last inner block of $c$ on it; then $C_{j+1},\dots,C_q$ are $\N(b)$-blocks as above, and $T_{j+1},C_{j+1},\dots,T_{q+1}$ lie in one class $K'$ of $J_a$, which meets $\supp t$. The point $x\in T_{j+1}\cap C_j$ lies in $Z_{K'}$, a transversal block of $z$, hence a union of $\T(b)$-blocks (an $\N(b)$-block is a non-transversal block of $z$); the $\T(b)$-block $B_x\ni x$ is joined to $T_{q+1}$ by the chain $B_x,T_{j+1},C_{j+1},\dots,C_q,T_{q+1}$ in $J_b$, so its class $H$ meets $\supp t$, $S_H$ is a transversal block of $s$, and $S_H$ meets $C_j$, which lies in the class of $C$ by the claim.

\emph{Type N.} Here $W_L=M$ is an $\N(a)$-block and a block of $z$. If $M$ is an $\N(b)$-block, it is a block of $s$ containing every $\T(c)$-block of $W_L$, and (A) holds. Otherwise $M$ is a union of $\T(b)$-blocks $B_1,\dots,B_k$; each $\T(c)$-block $C\subseteq M$ lies inside one $B_i$ (it is covered by them and meets at most one), so (A) follows if we show that $B_1,\dots,B_k$ lie in a single class $H$ of $J_b$ (then $S_H\supseteq M$). Let $K_0$ be the class of $J_a$ containing $M$. By (2), $M$ is the only block of $z$ inside $K_0$ meeting $\supp b$; so every $\N(a)$-block in $K_0$ other than $M$ is an $\N(b)$-block, and $Z_{K_0}$ is an $\N(b)$-block if nonempty. Hence every block of $b$ inside $K_0$ is one of: a $\T(b)$-block $B_i$; an $\N(a)$-block $M'\ne M$ in $K_0$; or $Z_{K_0}$; and for a point $p\in K_0$ we write $\beta(p)$ for the block of $b$ containing it, which is $Z_{K_0}$ whenever $p\in\supp a$. Let $C\subseteq B_1$ and $C'\subseteq B_2$ be $\T(c)$-blocks and take a chain $C=C_0,T_1,C_1,\dots,T_r,C_r=C'$ in $L$. Every block of $t$ meeting $M$, $Z_{K_0}$ or an $\N(a)$-block in $K_0$ lies in $K_0$; in particular $T_1,T_r\subseteq K_0$. Consider a block $C_l$ of the chain with $T_l\subseteq K_0$, and let $u\in T_l\cap C_l$. If $C_l$ is a $\T(c)$-block, then $C_l\subseteq M$ lies inside one $B_i$, and $\beta(u)=\beta(v)=B_i$ for the point $v\in T_{l+1}\cap C_l$. If $C_l$ is an outer $\N(c)$-block, it is an $\N(a)$-block inside $K_0$ and an $\N(b)$-block by (1), and $\beta(u)=\beta(v)=C_l$. If $C_l$ is an inner $\N(c)$-block, then $u\in\supp a\cap K_0$ and $\beta(u)=Z_{K_0}$; if also $T_{l+1}\subseteq K_0$ then $\beta(v)=Z_{K_0}$; if not, let $T_{l'}$ ($l'>l+1$) be the next block of $t$ on the chain contained in $K_0$ (it exists, as $T_r\subseteq K_0$): the block $C_{l'-1}$ meets $T_{l'-1}\not\subseteq K_0$ and $T_{l'}\subseteq K_0$, so it meets two classes of $J_a$ and is therefore an inner $\N(c)$-block, and the point $v'\in T_{l'}\cap C_{l'-1}$ lies in $\supp a\cap K_0$, so $\beta(v')=Z_{K_0}=\beta(u)$. Since consecutive points of the chain on a common block $T_l\subseteq K_0$ have blocks $\beta$ joined by $T_l$ in $J_b$, we obtain, walking along the chain and skipping the excursions outside $K_0$, a chain in $J_b$ from $\beta(x_0)=B_1$ ($x_0\in T_1\cap C_0$) to $\beta(y_r)=B_2$ ($y_r\in T_r\cap C_r$). So $B_1,B_2$ lie in one class of $J_b$, as required.

\emph{Conclusion.} By Step 5 the blocks of $\theta_c(s)$ inside $\supp c$ are the sets $W_L$, transversal iff $L$ meets $\supp t$, and its other blocks are the $\N(c)$-blocks; so $\theta_c(s)=\theta_c(t)=w$, which is \eqref{eq:lemD}, and $ata=btb$ by Step 2. The statements about $\mathfrak P_8^7$ follow from Proposition~\ref{prop:census8}(b) and the minimality noted above.
\end{proof}

\begin{remark}\label{rem:V}
(1) The hypothesis $bcb=b$ enters the proof of \eqref{eq:lemD} only through the facts (1) and (2) of Step~4; SAT finds no counterexample in $\P(\P_m)$, $m\le6$, to the quasi-identity $ata\le b,\ ctc\le a\Rightarrow c(btb)c=ctc$ without it. (2) The algebraic part (Steps 1--2) shows more generally: in any projection algebra satisfying \eqref{eq:Q1}, if $bcb=b$, $ata\le b$, $ctc\le a$ and $c(btb)c=ctc$, then $ata=btb$; and $\theta_b\circ\theta_c$ is the identity on $\{q:q\le b\}$ whenever $bcb=b$ (as $q=bqb=bcb\,q\,bcb=b(cqc)b$), so $\theta_c$ separates the elements below $b$.
 (3) Prover9 (M.K.) finds that (V) is equivalent, in every projection algebra (without \eqref{eq:Q1}), to the three-atom clause $\theta_a(t)\le\theta_d(c),\ \theta_c(t)\le a\Rightarrow\theta_a(t)\le\theta_{\theta_d(c)}(t)$ (in the binary notation: $(A*B)*(C*D)\ne A*B\ \vee\ (A*C)*B\ne A*C\ \vee\ (A*B)*(A*(C*D))=A*B$, with $(A,B,C,D)=(t,a,c,d)$). This is (V) with $b$ replaced by $b'=\theta_d(c)$: the fixing hypothesis $bcb=b$ has been absorbed, because $\theta_d(c)\,c\,\theta_d(c)=\theta_d(c)$ is an identity of projection algebras ($c*(c*d)=c*d$ follows from the axioms), and conversely every $b$ with $bcb=b$ is of this form ($b=\theta_b(c)$). So (V) says: the statement $\theta_a(t)\le q,\ \theta_c(t)\le a\Rightarrow\theta_a(t)\le\theta_q(t)$, which for $q=c$ is a triviality under \eqref{eq:Q1} ($\theta_a(t)$ is then a common lower bound of $a$ and $c$, so $ac=ca=:e$ and $\theta_a(t)=\theta_e(t)=\theta_c(t)$), remains true for every restriction $q=\theta_d(c)$ of $c$; and its conclusion reads $\theta_a(t)\le\theta_b\theta_c\theta_b(t)$, which is the form $s=b(csc)b$ in which it enters Step~1(iii) of the proof. The case $q=c$ itself (the clause with $d=c$) is thus a consequence of \eqref{eq:Q1}, and $\mathfrak P_8^7$ satisfies it while violating (V).
\end{remark}

\begin{remark}\label{rem:census8}
(1) The path-closing scheme now contains, besides (R) and (S), the one-auxiliary conditions of size $4$ and $5$ of (c) and the two-auxiliary conditions (R$'$) (length $2$) and (Q) (length $3$); the searches of Section~\ref{subsec:pathclose} stopped at three atoms, which is why these were missed. Whether the conditions of (c) are the complete list of one-auxiliary conditions of length $3$ (no bound on the size) is a finite question for $\P(\P_4)$, but its answer must then be checked in larger $\P(\P_m)$, as all conditions above were only verified for $m\le7$. (2) (V) is of a new shape: four variables, one ``fixing'' hypothesis $bcb=b$ (half of a matched pair), two order hypotheses on products, and an order conclusion; Theorem~\ref{thm:V} shows that its content is the statement $\theta_c(\theta_b(t))=\theta_c(t)$ about restrictions, the rest being equational. (3) \emph{Order $9$.} The $80\,039$ algebras of order $9$ satisfying \eqref{eq:Q1}, (S) and (R) were analysed in the same way (\texttt{pa8/post8.py} with the quasi-identities of Proposition~\ref{prop:census8} as a first filter, \texttt{pa8/newconds.py}; then \texttt{pa8/pass2\_9.py}). All satisfy the triangle condition. Exactly $163$ contain one of $\mathfrak P_8^1,\dots,\mathfrak P_8^7$ as a subalgebra, and these are exactly the algebras violating one of (R$'$), (V), (P$_1$), (P$_2$), (Q) or the two conditions of (e), with one exception: one algebra violates (R$'$) and contains none of the seven (so (R$'$), which is valid by Theorem~\ref{thm:Rp}, has at least two minimal witnesses, and this algebra is a minimal non-embeddable algebra of order $9$). Of the remaining $79\,875$, $79\,783$ embed in $\P(\P_m)$ with $m\le7$, $39$ are semilattices (which embed in $\P(\P_8)$ but, by SAT, in no $\P(\P_7)$: they have $8$ join-irreducibles), $46$ embed in $\P(\P_8)$ but in no $\P(\P_7)$, and $7$ embed in no $\P(\P_m)$ with $m\le8$ (SAT, $1$ to $30$ minutes each; \texttt{pa8/cand9.json}, \texttt{pa8/cand9.out}). These seven are the candidates of order $9$: all their proper subalgebras embed. Three of them are one-point extensions of graph-semigroup projection algebras ($|V|=4$ with four colours, and $|V|=5$ twice), one has a five-element bottom band like $\mathfrak P_8^7$, and none violates a two-auxiliary path-closing condition of length $\le3$ with at most four atoms. One of them (the second in \texttt{cand9.json}) violates the five-atom condition $p_0\bowtie p_1\bowtie p_2\bowtie p_3,\ xp_0=p_0x,\ \theta_{p_0}(x)=\theta_{p_2}(x),\ yp_1=p_1y,\ \theta_y(p_0)=\theta_y(p_1),\ \theta_x(p_3)=\theta_y(p_3)\Rightarrow p_0\bowtie p_3$ (valid in $\P(\P_5)$), a length-$3$ relative of (R$'$); for all seven, greedy minimisation of the diagram (\texttt{pa8/qi\_all9.py}, results in \texttt{pa8/qi\_all9.jsonl}) gives quasi-identities with $9$ to $14$ atoms, valid in $\P(\P_5)$, which we have not tried to interpret. So the count of candidates for minimal non-embeddable algebras is $1, 7, 8$ for orders $7, 8, 9$ (counting the algebra violating (R$'$)), of which $1,5,1$ are proved (Theorems~\ref{thm:V}, \ref{thm:Rp}, \ref{thm:P1}, \ref{thm:Q}), and the number of points needed by embeddable algebras of order $n$ has reached $n-1$ at $n=9$.
 (4) \emph{Absorbing the normalising hypotheses.} As with (V) (Remark~\ref{rem:V}(3)), the hypotheses of (e) which merely say that a variable is a restriction of another can be absorbed by substitution: $z\le b$ by $z=\theta_b(z')$, $dbd=d$ by $d=\theta_{d'}(b)$, $e\le x$ by $e=\theta_x(e')$. Prover9 (M.K.) found in this way a seven-literal clause equivalent to (E$_1$) in every projection algebra, $\theta_B(A)=\theta_A(C),\ \theta_E(D)=\theta_A(C),\ \theta_B(D)=B,\ \theta_A(D)=A,\ \theta_E(F)\le\theta_G(A),\ \theta_B(\theta_G(A))=B\Rightarrow\theta_B(A)=B$, which is (E$_1$) under $(a,b,c,d,x,e,z)=(B,A,D,\theta_G(A),E,\theta_E(F),\theta_A(C))$; the converse direction is the instance $(A,B,C,D,E,F,G)=(b,a,z,c,x,e,d)$, which Prover9 has difficulty finding (instantiating a seven-variable clause by paramodulation), but which is immediate by hand. More useful is that under \eqref{eq:Q1} the variable $e$ can be eliminated altogether, since it only serves to make $x$ and $d$ commute: (E$_1$) is equivalent to
\[
\text{(E$_1'$)}\qquad aca=a,\ bcb=b,\ dbd=d,\ ada=a,\ xd=dx,\ xcx=aba\le b\ \Longrightarrow\ aba=a
\]
(given (E$_1$), put $e=xd$, which lies below $x$ and $d$; given (E$_1'$), $x$ and $d$ commute by \eqref{eq:Q1}, and Prover9 verifies this direction), and (E$_1'$) fails in $\mathfrak P_8^1$ with $(a,b,c,d,x)=(2,3,5,7,1)$. Similarly $u$ can be specialised to $cd$ in (E$_6$): the consequence
\[
\text{(E$_6'$)}\qquad aca=a,\ ada=a,\ dbd=d,\ cd=dc,\ bcb=a(cd)a\ \Longrightarrow\ a(cd)a=a
\]
still fails in $\mathfrak P_8^6$ (with $(a,b,c,d)=(2,1,5,7)$, $cd=6$), so it is an obstruction of five variables and five atoms, though possibly weaker than (E$_6$). The two share the pattern $aca=a$, $ada=a$, $dbd=d$: $a$ is covered by $c$ and by $d$ (in the sense $\theta_a(c)=a$), and $d$ by $b$. On the $12\,046$ instances of (E$_6'$) in $\P(\P_4)$ with $a,b,c,d$ distinct (\texttt{pa8/e6/}), the conclusion comes with much more: $bcb=aba=bab=a$, so $a\le b$; $a\bowtie cd$; $\theta_d(a)=cd$; and $\theta_c(a)=\theta_c(b)$. Of these only $\theta_{cd}(a)=cd$ is a consequence of the axioms, \eqref{eq:Q1} and the hypotheses (Prover9); $\theta_d(a)=cd$ is not a consequence of $aca=a$, $ada=a$, $cd=dc$ alone (it fails in $\P(\P_4)$ without the hypotheses on $b$), and adding it to the hypotheses does not make (E$_6'$) equational either. So a proof of (E$_6'$) will need a block argument of a new kind, without an auxiliary projection commuting with two of the others; we have not attempted it.
 (5) \emph{Towards (Q).} Write the path of (Q) as $a\bowtie c\bowtie b\bowtie d$, so that the hypotheses are $xa=ax$, $\theta_c(x)=\theta_d(x)$, $ya=ay$, $yb=by$, $\theta_y(c)=\theta_y(b)$ and the conclusion is $a\bowtie d$. On all $349\,343$ instances in $\P(\P_4)$ (\texttt{pa8/q/}): $cd=dc$ (by \eqref{eq:Q1}, as $\theta_c(x)=\theta_d(x)$ is a common lower bound), $xc=xd$, $cx=dx$, $\theta_y(c)=\theta_y(b)=\theta_y(d)$, $y\,ad=y\,abd=y\,acbd$ and $x\,ad=x\,acd=x\,acbd$; but $ad=acbd$ fails in $420$ instances with distinct path elements, and no join $\varepsilon_a\vee\varepsilon_d\vee\varepsilon_p\vee\varepsilon_q$ with $p,q$ among $c,b,cd,x,y,ya,yb,\theta_c(x),xa$ separates the $\T(a)$-blocks in all instances. So, unlike (R$'$), the conclusion of (Q) is not the coherence of the three matchings, and a proof cannot go through a separation statement for a coarser join. Note also that a proof of (Q) contains a proof of (S) (Proposition~\ref{prop:indep}(a)), hence of the hub case of Theorem~\ref{thm:S}, which until Section~\ref{subsec:Q} rested on a computer check. A search for a common generalisation of (R$'$), (P$_1$) and (Q) in the two-auxiliary path-closing language (all $627\,392$ paths of length $3$ in $\P(\P_4)$, $90$ atoms, all instantiations of the auxiliaries by restrictions of the given ones; \texttt{pa8/gen/}) found none, not even for any two of the three. Section~\ref{subsec:Q} proves (Q) (Theorem~\ref{thm:Q}).
\end{remark}


\subsection{Proof of (R$'$)}\label{subsec:Rp}

\begin{theorem}\label{thm:Rp}
(R$'$) holds in every partition monoid $\P_X$: if $a,b,c,x,y$ are projections with
\[
a\bowtie c,\quad b\bowtie c,\quad xc=cx,\quad \theta_x(a)=\theta_x(c),\quad ya=ay,\quad \theta_x(b)=\theta_y(b),
\]
then $a\bowtie b$. Consequently $\mathfrak P_8^3$ embeds in no $\P(\P_X)$ and is a minimal non-embeddable projection algebra, and so is the algebra of order $9$ of Remark~\ref{rem:census8}(3) that violates (R$'$) but contains none of $\mathfrak P_8^1,\dots,\mathfrak P_8^7$.
\end{theorem}

The proof has the same shape as that of Theorem~\ref{thm:V}: an algebraic part, valid in every projection algebra satisfying \eqref{eq:Q1}, reduces the theorem to a combinatorial statement about blocks, and the combinatorial part is a chain argument in $\P_X$. The algebraic part is longer here, and we found most of it with Prover9 (\texttt{pa8/rp/}); the combinatorial part was found by testing candidate statements against all instances of the hypotheses in $\P(\P_4)$ ($91\,172$ quintuples, \texttt{pa8/rp/harness*.py}) and, for the statements finally used, against random instances in $\P(\P_5)$.

\begin{proof}
Throughout, $r=\theta_x(b)=\theta_y(b)$, $v=xc=cx$ and $u=ya=ay$.

\smallskip
\emph{Step 1: algebraic consequences.} All identities in this step hold in every regular $*$-semigroup whose projection algebra satisfies \eqref{eq:Q1}, under the hypotheses of the theorem.

(i) $r\le x$ and $r\le y$ (as $r=\theta_x(b)$ and $r=\theta_y(b)$), so $xy=yx=:m$ by \eqref{eq:Q1}, and $r\le m$.

(ii) $\theta_x(c)=xcx=xc=v$, so the hypothesis $\theta_x(a)=\theta_x(c)$ reads $xax=v$. Hence $xa=xa\,ax\,xa=(xax)a=va$ and, applying $*$, $ax=av$.

(iii) \emph{$v\le y$ and $cy=yc=v$.} First, $\theta_a(r)=a(yby)a=(ay)b(ya)=ubu$. Next, for every element $s$, $\theta_x(\theta_u(s))=x\,ya\,s\,ay\,x=m\,a s a\,m=\theta_m(\theta_a(s))$, since $xya=ma$ and $ayx=am$. Applying this to $s=b$ and using $\theta_a(r)=ubu$ together with $\theta_a(b)=aba$: $\theta_m(aba)=\theta_x(ubu)=\theta_x(\theta_a(r))=xa\,xbx\,ax=(xax)b(xax)=vbv$, and $vbv=v$ because $v\le c$ and $cbc=c$ (for $q\le c$, $qbq=qc\,b\,cq=q(cbc)q=q$). So $v=\theta_m(aba)\le m\le y$. Now $v$ is a common lower bound of $c$ and $y$, so $cy=yc$ by \eqref{eq:Q1}; the projection $q=cy$ lies below $c$, so $qbq=q$ by the identity just used, and $qbq=c(yby)c=crc=c(xbx)c=(cx)b(xc)=vbv=v$. Thus $cy=yc=v$.

(iv) \emph{$\theta_a(x)=\theta_a(y)=u$ and $u\bowtie v$.} By (ii) and (iii), $axa=ava=a\,cy\,a=ac\,ya=aca\,y=ay=u$. Put $g=ac$, so that $gg^*=aca=a$ and $g^*g=cac=c$; then $yg=yac=ayc=acy=gy$, hence $u=ya=y\,gg^*=(yg)(yg)^*$ and $v=yc=y\,g^*g=(yg)^*(yg)$. So $u\bowtie v$: $uvu=ya\,yc\,ya=y\,ac\,ya=y\,aca\,y=u$ and $vuv=yc\,ya\,yc=y\,ca\,yc=y\,cac\,y=v$, using $ay=ya$, $cy=yc$, $aca=a$ and $cac=c$. (Added 15 Oct: an earlier version justified this by ``$hh^*\bowtie h^*h$ for every $h$'', which is false: for the partial bijection $h=\{0\mapsto1\}$ in $\P_2$ the projections $hh^*=[1](0)$ and $h^*h=[0](1)$ have product of rank $0$. What is true is that $hh^*$ and $h^*h$ are $\D$-related; matchedness needs the commutations above.)

(v) \emph{$yax=ax$, $xya=xa$, $axb=yab$ and $xab=xcb$.} We have $cbv=cb\,cx=(cbc)x=cx=v$ and, by (ii), $ax=av=acx$; hence $axbv=avbv=a(vbv)=av=ax$. Since $r\le y$, $y\,xbx=yr=r=xbx$; therefore $yax=y\,axbv=ay\,xbx\,c=a\,xbx\,c=axbv=ax$ (we used $v=xc$ and $ay=ya$). Applying $*$ gives $xay=xa$, i.e.\ $xya=xa$. Next $xb=(xbx)b=rb=(yby)b=yb$, so $axb=ayb=yab$. Finally $xcb=(xax)b=x(axb)=x(yab)=(xya)b=xab$.

(vi) \emph{$y\,ab=y\,acb$.} Multiplying $xab=xcb$ on the left by $a$ and using $axa=u$ and $axc=avc=av$ (as $vc=v$): $u\,b=av\,b$; and $y\,acb=(ya)(cb)=a(yc)b=avb$. So $ub=y\,ab=y\,acb=ucb$.

(vii) \emph{Reduction to $x\le y$.} The projection $m=xy$ satisfies the hypotheses in place of $x$: $mc=xyc=xv=v=vy=cxy=cm$ (as $v\le x$, $v\le y$); $\theta_m(c)=mcm=v$ and $\theta_m(a)=xy\,a\,yx=xux=x(ava)x=(xa)v(ax)=(va)v(av)=v(ava)v=vuv=v$ by (ii), (iv); and $\theta_m(b)=\theta_m(\theta_x(b))=\theta_m(r)=r=\theta_y(b)$, since $m\le x$ and $r\le m$. As the conclusion does not involve $x$, we may and do assume from now on that $x\le y$.

\smallskip
\emph{Step 2: the setting in $\P_X$.} We use Lemma~\ref{lem:PX}, the facts ($\alpha$), ($\beta$) recalled in the proof of Theorem~\ref{thm:V}, and chains as in the proof of Theorem~\ref{thm:R}. Write $Y=X\setminus\supp y$. Let $A_1,\dots,A_s$ be the $\T(a)$-blocks, $C_i$ the $\T(c)$-block matched with $A_i$ and $B_i$ the $\T(b)$-block matched with $C_i$ (Lemma~\ref{lem:PX}(c)); $\Gamma_i$ and $\Delta_i$ denote the classes of $\varepsilon_a\vee\varepsilon_c$ and of $\varepsilon_b\vee\varepsilon_c$ containing $C_i$, so that $\Gamma_i$ contains exactly the $\T$-blocks $A_i,C_i$ and $\Delta_i$ exactly $B_i,C_i$; for a block $E$ of $c$ we write $\Gamma(E)$, $\Delta(E)$ for the classes containing it. We shall show that
\begin{equation}\label{eq:Rp-goal}
\text{no class of }\varepsilon_a\vee\varepsilon_b\vee\varepsilon_c\text{ contains two }\T(c)\text{-blocks,}
\end{equation}
and then conclude by a counting argument. Since $A_i,B_i,C_i$ lie in one class of $\varepsilon_a\vee\varepsilon_b\vee\varepsilon_c$, \eqref{eq:Rp-goal} implies that every class of $\varepsilon_a\vee\varepsilon_b$ contains at most one $\T(a)$-block and at most one $\T(b)$-block, and that a class containing $A_i$ contains no $\T(b)$-block other than $B_i$.

By Lemma~\ref{lem:PX}(b) applied to the commuting pairs $y,a$ and $y,c$ and to $x,c$: every block of $a$, of $c$ and of $x$ lies inside $\supp y$ or is an $\N(y)$-block; an $\N(y)$-block is an $\N(a)$-block or a union of $\T(a)$-blocks, and an $\N(c)$-block or a union of $\T(c)$-blocks; every $\N(c)$-block lies inside $\supp x$ (and is then a union of $\T(x)$-blocks) or is an $\N(x)$-block. Since $x\le y$, every block of $x$ is an $\N(y)$-block or a union of $\T(y)$-blocks (($\alpha$)). The classes of $\varepsilon_y\vee\varepsilon_a$ meeting $\supp y$ are the $\T(u)$-blocks $U_1,\dots,U_t$ (unions of $\T(a)$-blocks and of $\T(y)$-blocks) and the $\N(a)$-blocks inside $\supp y$; similarly the classes of $\varepsilon_y\vee\varepsilon_c$ meeting $\supp y$ are the $\T(v)$-blocks $V_1,\dots,V_{t'}$ and the $\N(c)$-blocks inside $\supp y$.

\emph{Step 3: a second reduction.} Let $M$ be an $\N(y)$-block which is a union of $\T(a)$-blocks. If $M$ were an $\N(c)$-block, it would be a class of $\varepsilon_a\vee\varepsilon_c$ by itself, containing $\T(a)$-blocks but no $\T(c)$-block, contradicting $aca=a$ (($\beta$)); so $M$ is a union of $\T(c)$-blocks, and the class of $\varepsilon_a\vee\varepsilon_c$ of any $A_i\subseteq M$ lies inside $M$ and consists of $\T$-blocks only, whence $\Gamma_i=A_i=C_i$. Symmetrically, an $\N(y)$-block which is a union of $\T(c)$-blocks is a union of $\T(a)$-blocks. Let $x',y'$ be obtained from $x,y$ by declaring all these blocks $M$ transversal. We claim that $a,b,c,x',y'$ still satisfy the hypotheses (with $x'\le y'$); as the conclusion does not change, we may then assume that \emph{every $\N(y)$-block is an $\N(a)$-block and an $\N(c)$-block}, so that all $A_i$ and all $C_i$ lie inside $\supp y$. Indeed, the partitions $\varepsilon_{x'}=\varepsilon_x$, $\varepsilon_{y'}=\varepsilon_y$ are unchanged, and each $M$ is a class by itself in $\varepsilon_x\vee\varepsilon_a$, in $\varepsilon_x\vee\varepsilon_c$, in $\varepsilon_y\vee\varepsilon_a$ and in $\varepsilon_y\vee\varepsilon_c$, being a block of $x$ and of $y$ which is a union of blocks of $a$ and of $c$. Hence $x'c=cx'$ and $y'a=ay'$, $y'c=cy'$ by the converse of Lemma~\ref{lem:PX}(b) (if every class of $\varepsilon_p\vee\varepsilon_q$ is of one of the four types, the description of $pq$ shows that $pq$ is symmetric, so $pq=qp$): the class $M$ is now of type (i), the others are unchanged; $\theta_{x'}(a)=\theta_{x'}(c)$ because the blocks of these projections are those of $\theta_x(a)=\theta_x(c)$ except that each $M$ becomes a transversal block of both; and $\theta_{x'}(b)=\theta_{y'}(b)$ will follow from Step~4(c) below, whose proof uses only $x\le y$ and $\theta_x(b)=\theta_y(b)$.

\emph{Step 4: the classes of $\varepsilon_y\vee\varepsilon_b$.} Let $\sigma=\varepsilon_y\vee\varepsilon_b$. (a) \emph{$\varepsilon_x\vee\varepsilon_b=\sigma$.} Let $W$ be a block of $x$ which is a union of $\T(y)$-blocks. If $W$ is a $\T(x)$-block, let $\tilde K$ be its class in $\varepsilon_x\vee\varepsilon_b$; then $\tilde K\cap\supp x$ is a block of $\theta_x(b)=\theta_y(b)$ contained in $\supp y$, hence equal to $K\cap\supp y$ for a class $K$ of $\sigma$, and $W\subseteq K$. If $W$ is an $\N(x)$-block, it is itself a block of $\theta_x(b)=\theta_y(b)$ contained in $\supp y$, so again $W=K\cap\supp y$ for a class $K$ of $\sigma$. The remaining blocks of $x$ are $\N(y)$-blocks. So every block of $x$ lies in a single class of $\sigma$, i.e.\ $\varepsilon_x\le\sigma$, and $\varepsilon_x\vee\varepsilon_b=\sigma$. (b) \emph{Two kinds of classes.} By (a) and the two descriptions of $r$, for every class $K$ of $\sigma$ meeting $\supp y$ either $K\cap\supp y=\tilde N$ is an $\N(x)$-block (a union of $\T(y)$-blocks), and then $K\cap\supp b=\emptyset$ since $\tilde N$ is a non-transversal block of $\theta_x(b)$; or $K\cap\supp y=K\cap\supp x$. We call the $\N(x)$-blocks inside $\supp y$ the \emph{$\tilde N$-blocks}; each $\tilde N$-block is $K\cap\supp y$ for a class $K$ of $\sigma$ with $K\cap\supp b=\emptyset$, and every other class meeting $\supp y$ has $K\cap\supp y\subseteq\supp x$. (c) In the situation of Step~3, the blocks of $\theta_{x'}(b)$ and $\theta_{y'}(b)$ are indexed by the same classes $K$ of $\sigma$ (by (a)), and they can differ only for a class $K\supseteq M$, where they are $(K\cap\supp x)\cup M$ and $(K\cap\supp y)\cup M$, both transversal iff $K$ meets $\supp b$. If these differ, then by (b) $K\cap\supp y$ is an $\tilde N$-block and $K\cap\supp b=\emptyset$; but then, for $A_i=C_i\subseteq M$, every block of $b$ meeting $K$ lies in $K$, and every block of $c$ meeting a block $Z$ of $b$ inside $K$ lies in $K$ (a block of $c$ in $Y$ is an $\N(y)$-block, hence lies in a single class of $\sigma$; a block of $c$ inside $\supp y$ meets $Z\cap\supp y\subseteq\tilde N$, so it is $\tilde N$ or contained in it), so a chain from $C_i\subseteq K$ to $B_i$ in $\Delta_i$ stays inside $K$, contradicting $K\cap\supp b=\emptyset$. This completes Step~3. (d) \emph{Every $\tilde N$-block is an $\N(c)$-block, and its class $\Delta(\tilde N)$ contains no $\T(b)$-block and satisfies $\Delta(\tilde N)\cap\supp y=\tilde N$.} Let $K$ be the class of $\sigma$ with $K\cap\supp y=\tilde N$. As in (c), every block of $b$ or of $c$ reachable by a chain in $\varepsilon_b\vee\varepsilon_c$ from a block of $b$ inside $K$ lies inside $K$. If $\tilde N$ were a union of $\T(c)$-blocks (the alternative, $x$ and $c$ commuting), a block of $b$ meeting some $C_j\subseteq\tilde N$ would lie in $K$, and then $B_j\subseteq\Delta_j\subseteq K$, contradicting $K\cap\supp b=\emptyset$. So $\tilde N$ is an $\N(c)$-block, $\Delta(\tilde N)\subseteq K$, and the claims follow.

\emph{Step 5: the classes of $\varepsilon_x\vee\varepsilon_a$.} By Lemma~\ref{lem:PX}(a), the blocks of $\theta_x(a)$ inside $\supp x$ are the sets $K\cap\supp x$ for the classes $K$ of $\varepsilon_x\vee\varepsilon_a$ meeting $\supp x$, transversal iff $K$ meets $\supp a$; those of $\theta_x(c)$ are the $\T(v)$-blocks $V_k$ (transversal) and the $\N(c)$-blocks $\tilde N'$ inside $\supp x$ (non-transversal; they are the classes of $\varepsilon_x\vee\varepsilon_c$ of type (ii)). Since $\theta_x(a)=\theta_x(c)$: for every $k$ there is a class $K_k$ of $\varepsilon_x\vee\varepsilon_a$ with $K_k\cap\supp x=V_k$, and for every $\N(c)$-block $\tilde N'\subseteq\supp x$ a class $K_{\tilde N'}$ with $K_{\tilde N'}\cap\supp x=\tilde N'$ and $K_{\tilde N'}\cap\supp a=\emptyset$. Consequently \emph{every block of $c$ lies inside a single class of $\varepsilon_x\vee\varepsilon_a$}: a $\T(c)$-block lies in some $V_k\subseteq K_k$; an $\N(c)$-block inside $\supp x$ is some $\tilde N'\subseteq K_{\tilde N'}$; and the other $\N(c)$-blocks are $\N(x)$-blocks. Hence $\varepsilon_a\vee\varepsilon_c\le\varepsilon_x\vee\varepsilon_a$, and for an $\N(c)$-block $\tilde N'\subseteq\supp x$ we get $\Gamma(\tilde N')\subseteq K_{\tilde N'}$; so \emph{$\Gamma(\tilde N')$ contains no $\T(a)$-block, hence no $\T(c)$-block, and $\tilde N'$ is the only block of $c$ inside $\supp x$ that it contains} (a block of $c$ in $\Gamma(\tilde N')$ inside $\supp x$ lies in $K_{\tilde N'}\cap\supp x=\tilde N'$).

\emph{Step 6: the four kinds of blocks of $c$ and their classes.} By Steps~2--4, every block of $c$ is of exactly one of the following kinds: a $\T(c)$-block $C_i$ (inside $\supp y$, by Step~3); an $\N(c)$-block inside $\supp x$; an $\tilde N$-block; or an $\N(y)$-block inside $Y$, which is an $\N(a)$-block by Step~3 and is therefore a class of $\varepsilon_a\vee\varepsilon_c$ by itself. Collecting Steps~4(d) and 5:
\begin{itemize}[nosep]
\item[(1)] the blocks of $c$ in $\Gamma_i$ are $C_i$ and $\tilde N$-blocks (an $\N(c)$-block inside $\supp x$ lies in a class without $\T(a)$-blocks, and an $\N(y)$-block in $Y$ is a class by itself);
\item[(2)] the blocks of $c$ in $\Delta_i$ are $C_i$, $\N(c)$-blocks inside $\supp x$ and $\N(y)$-blocks in $Y$ (an $\tilde N$-block lies in a class without $\T(b)$-blocks);
\item[(3)] for an $\N(c)$-block $\tilde N'\subseteq\supp x$, the blocks of $c$ in $\Gamma(\tilde N')$ are $\tilde N'$ and $\tilde N$-blocks, and $\Gamma(\tilde N')$ contains no $\T(a)$-block;
\item[(4)] for an $\tilde N$-block $\tilde N$, the blocks of $c$ in $\Delta(\tilde N)$ are $\tilde N$ and $\N(y)$-blocks in $Y$ (as $\Delta(\tilde N)\cap\supp y=\tilde N$), and $\Delta(\tilde N)$ contains no $\T(b)$-block;
\item[(5)] an $\N(y)$-block $E\subseteq Y$ is a class of $\varepsilon_a\vee\varepsilon_c$ by itself.
\end{itemize}

\emph{Step 7: proof of \eqref{eq:Rp-goal}.} Two blocks $E,E'$ of $c$ lie in one class of $\varepsilon_a\vee\varepsilon_b\vee\varepsilon_c$ iff they are joined by a sequence $E=E_0,E_1,\dots,E_n=E'$ of blocks of $c$ in which consecutive blocks lie in a common class of $\varepsilon_a\vee\varepsilon_c$ or in a common class of $\varepsilon_b\vee\varepsilon_c$: indeed, if two blocks $W,W'$ (of $a$, $b$ or $c$) meet at a point $p$, and $E(p)$ is the block of $c$ containing $p$, then $W$ and $E(p)$ lie in a common class of $\varepsilon_a\vee\varepsilon_c$ or of $\varepsilon_b\vee\varepsilon_c$ (according to whether $W$ is a block of $a$ or of $b$; if $W$ is a block of $c$ then $W=E(p)$), so a chain of blocks meeting consecutively can be replaced by the sequence of the blocks of $c$ at its meeting points. Now start such a sequence at $E_0=C_p$ and use (1)--(5) to see which blocks of $c$ can occur. From $C_p$ one can move, within $\Gamma_p$, only to $\tilde N$-blocks of $\Gamma_p$, and within $\Delta_p$ only to $\N(c)$-blocks $\tilde N'\subseteq\supp x$ of $\Delta_p$ and to $\N(y)$-blocks $E\subseteq Y$ of $\Delta_p$. From an $\tilde N$-block $\tilde N$ of $\Gamma_p$ one can move within $\Gamma(\tilde N)=\Gamma_p$ (back) or within $\Delta(\tilde N)$ to $\N(y)$-blocks $E\subseteq Y$; from such an $E$ one can move within $\Delta(E)=\Delta(\tilde N)$ only, and $\Gamma(E)=\{E\}$. From $\tilde N'\subseteq\supp x$ in $\Delta_p$ one can move within $\Delta(\tilde N')=\Delta_p$ (back) or within $\Gamma(\tilde N')$ to $\tilde N$-blocks of $\Gamma(\tilde N')$, from which one can move within $\Gamma(\tilde N')$ (back) or within their $\Delta$-classes to $\N(y)$-blocks in $Y$, which are dead ends as before. From an $\N(y)$-block $E\subseteq Y$ in $\Delta_p$ one can move within $\Delta_p$ only. So the blocks of $c$ reachable from $C_p$ are: $C_p$; $\tilde N$-blocks (of $\Gamma_p$, or of the classes $\Gamma(\tilde N')$ for $\tilde N'\subseteq\supp x$ in $\Delta_p$); $\N(c)$-blocks inside $\supp x$ lying in $\Delta_p$; and $\N(y)$-blocks in $Y$. None of these is a $\T(c)$-block other than $C_p$, which proves \eqref{eq:Rp-goal}.

\emph{Step 8: counting.} Let $l\in\{1,\dots,t\}$ and $I_l=\{p:A_p\subseteq U_l\}$; by Step~3 every $A_p$ lies in some $U_l$. Consider the class of $U_l$ in $\varepsilon_u\vee\varepsilon_v$ ($\varepsilon_u=\varepsilon_y\vee\varepsilon_a$, $\varepsilon_v=\varepsilon_y\vee\varepsilon_c$); by $u\bowtie v$ (Step~1(iv)) it contains exactly one $\T(v)$-block, $V_{k(l)}$ say. Every $\Gamma_p$ is contained in a class of $\varepsilon_u\vee\varepsilon_v$ (its blocks of $a$ and of $c$ are contained in blocks of $u$ and of $v$), so $C_p\subseteq V_{k(l)}$ for $p\in I_l$, and conversely $C_q\subseteq V_{k(l)}$ implies $A_q\subseteq U_l$; thus $I_l=\{q:C_q\subseteq V_{k(l)}\}$. Now compare the two sides of $y\,ab=y\,acb$ (Step~1(vi)). The upper blocks of $ab$ are the $\N(a)$-blocks and the sets $D\cap\supp a$ for the classes $D$ of $\varepsilon_a\vee\varepsilon_b$ meeting $\supp a$, the latter being transversal iff $D$ meets $\supp b$, with lower part $D\cap\supp b$. By \eqref{eq:Rp-goal}, $D\cap\supp a$ is a single $\T(a)$-block; write $D(A_p)$ for the class of $A_p$. In the product $y\cdot(ab)$ the middle join is $\varepsilon_y\vee\varepsilon_{ab}$, where $\varepsilon_{ab}$ is the upper partition of $ab$; the class of $U_l$ in it is $U_l$ itself, because $U_l$ is a union of $\T(y)$-blocks and of the blocks $A_p$ ($p\in I_l$) of $\varepsilon_{ab}$, and no $\N(a)$-block or $\N(y)$-block meets it. So $y\,ab$ has the upper block $U_l$, with lower part $\bigcup_{p\in I_l}(D(A_p)\cap\supp b)$ (transversal iff this is nonempty). In $y\,acb=(ya)(cb)=u\,(cb)$ the middle join is $\varepsilon_u\vee\varepsilon_c$, whose class containing $U_l$ contains no other $\T(u)$-block and, being contained in the class of $U_l$ in $\varepsilon_u\vee\varepsilon_v$, contains exactly the $\T(c)$-blocks $C_p$ with $p\in I_l$ (all of which it does contain, as $\Gamma_p$ is a chain in $\varepsilon_u\vee\varepsilon_c$); so $u\,cb$ has the upper block $U_l$ with lower part $\bigcup_{p\in I_l}B_p$. Hence
\[
\bigcup_{p\in I_l}\bigl(D(A_p)\cap\supp b\bigr)=\bigcup_{p\in I_l}B_p .
\]
By \eqref{eq:Rp-goal}, $D(A_p)\cap\supp b\subseteq B_p$ for each $p$, and the $B_p$ are pairwise disjoint and nonempty; so $D(A_p)\cap\supp b=B_p$, i.e.\ $B_p\in D(A_p)$, for every $p\in I_l$ and every $l$, that is, for every $p$.

\emph{Conclusion.} Every class of $\varepsilon_a\vee\varepsilon_b$ meeting $\supp a$ is some $D(A_p)$, which contains exactly one $\T(a)$-block and exactly one $\T(b)$-block ($B_p$) by \eqref{eq:Rp-goal} and Step~8; and every class meeting $\supp b$ contains some $B_p$, hence is $D(A_p)$. By Lemma~\ref{lem:PX}(c) and its proof, $a\bowtie b$. The statements about $\mathfrak P_8^3$ and the algebra of order $9$ follow from Proposition~\ref{prop:census8}(a) and Remark~\ref{rem:census8}(3): both violate (R$'$), and all their proper subalgebras embed (for the algebra of order $9$, because its subalgebras of order $8$ satisfy \eqref{eq:Q1}, (S), (R), so they are among the $10\,828$ algebras of the census, and it contains none of the seven that do not embed).
\end{proof}

\begin{remark}\label{rem:Rp}
(1) In the language of Step~1, the hypotheses of (R$'$) say that $x\,ac=xc$ ($x$ does not distinguish the element $ac$ implementing $a\bowtie c$ from $c$) and $xb=yb$ ($x$ and $y$ act alike on $b$); the hypotheses of (R) say that $xb=x\,cb$. The conclusion actually obtained is stronger than $a\bowtie b$: by \eqref{eq:Rp-goal} and Step~8, the classes of $\varepsilon_a\vee\varepsilon_b\vee\varepsilon_c$ meeting $\supp a\cup\supp b\cup\supp c$ are exactly the classes containing the triples $A_i,B_i,C_i$, so that $ab=acb$; in $\P(\P_4)$ this identity, and its consequences $cab=cb$, $ba=bca$, hold in all instances of the hypotheses, whereas a triangle $a\bowtie b\bowtie c\bowtie a$ alone does not force $ab=acb$ (the three matchings need not compose to the identity). (2) Neither the counting of Step~8 nor its analogue for $x\,ab=x\,cb$ can replace \eqref{eq:Rp-goal}: both are blind to the merging of two $\T(a)$-blocks of the same $U_l$ into one class of $\varepsilon_a\vee\varepsilon_b$, which is exactly what happens in $\mathfrak P_8^3$ ($\theta_a(b)=u$, $\theta_b(a)=w$ there).
\end{remark}

\subsection{Proof of (P$_1$)}\label{subsec:P1}

We write the path of (P$_1$) as $a\bowtie c\bowtie b\bowtie d$ (so $p_0=a$, $p_1=c$, $p_2=b$, $p_3=d$); the hypotheses are then $xa=ax$, $xc=cx$, $\theta_a(x)=\theta_b(x)$ and $\theta_x(c)=\theta_x(d)$, and the conclusion is $a\bowtie d$.

\begin{theorem}\label{thm:P1}
(P$_1$) holds in every partition monoid $\P_X$: if $a,b,c,d,x$ are projections with $a\bowtie c$, $c\bowtie b$, $b\bowtie d$, $xa=ax$, $xc=cx$, $\theta_a(x)=\theta_b(x)$ and $\theta_x(c)=\theta_x(d)$, then $a\bowtie d$. Consequently $\mathfrak P_8^2$ embeds in no $\P(\P_X)$ and is a minimal non-embeddable projection algebra.
\end{theorem}

\begin{proof}
Write $u=xa=ax$ and $v=xc=cx$.

\smallskip
\emph{Step 1: $x$ commutes with $b$, and $xb=u$.} We have $\theta_b(x)=\theta_a(x)=axa=u$, so $u\le b$; as also $u\le x$, \eqref{eq:Q1} gives $xb=bx$, and then $\theta_b(x)=bxb=xb$, i.e.\ $xb=bx=u$. (In particular $u\le a,b$, so $ab=ba$ by \eqref{eq:Q1}; this is not needed below.)

\smallskip
\emph{Step 2: the blocks of $a$, $b$, $c$.} We use Lemma~\ref{lem:PX} and the notation of the proof of Theorem~\ref{thm:R}: $Y=X\setminus\supp x$; $A_1,\dots,A_s$ are the $\T(a)$-blocks, $C_i$ the $\T(c)$-block matched with $A_i$, $B_i$ the $\T(b)$-block matched with $C_i$, $D_i$ the $\T(d)$-block matched with $B_i$; $\Gamma_i$, $\Delta_i$, $\Lambda_i$ are the classes of $\varepsilon_a\vee\varepsilon_c$, $\varepsilon_c\vee\varepsilon_b$, $\varepsilon_b\vee\varepsilon_d$ containing $C_i$, $C_i$, $B_i$ respectively, so that $\Gamma_i$ contains exactly the $\T$-blocks $A_i,C_i$; $\Delta_i$ exactly $C_i,B_i$; $\Lambda_i$ exactly $B_i,D_i$ (Lemma~\ref{lem:PX}(c)). Since $x$ commutes with $a$, $b$ and $c$ (Step~1), every block of $a$, $b$ or $c$ lies inside $\supp x$ or is an $\N(x)$-block, and Step~1 of the proof of Theorem~\ref{thm:R}, applied to the matched pairs $a\bowtie c$ and $b\bowtie c$, shows that \emph{$a$, $b$ and $c$ have the same blocks inside $Y$, with the same labels}: for $i\notin I:=\{i:A_i\subseteq\supp x\}$ we have $A_i=B_i=C_i\subseteq Y$, and an $\N(x)$-block is an $\N(a)$-block iff it is an $\N(b)$-block iff it is an $\N(c)$-block. Inside $\supp x$ we compare the elements $xa$ and $xb$, which are equal by Step~1: by Lemma~\ref{lem:PX}(b) the upper blocks of $xa$ are the $\N(x)$-blocks, the $\T(u)$-blocks $U_1,\dots,U_t$ (the classes of type (i) of $\varepsilon_x\vee\varepsilon_a$, each a union of $\T(a)$-blocks and of $\T(x)$-blocks, with lower part $U_l'$) and the $\N(a)$-blocks inside $\supp x$ (type (ii)); likewise for $xb$. Hence \emph{the $\N(a)$-blocks inside $\supp x$ are exactly the $\N(b)$-blocks inside $\supp x$, and each $U_l$ is a union of $\T(a)$-blocks and also a union of $\T(b)$-blocks}. Altogether: $a$ and $b$ have the same $\N$-blocks, and $\supp a\cap\supp x=\supp b\cap\supp x=\supp u=U_1\cup\dots\cup U_t$. Finally, for $p\in I$ the block $C_p$ meets $\supp u$ only inside $A_p\cap B_p$: a point of $C_p\cap\supp u$ lies in some $A_q$ and some $B_r$, and $C_p$ meets no $\T(a)$-block other than $A_p$ and no $\T(b)$-block other than $B_p$.

\smallskip
\emph{Step 3: the blocks of $d$ inside $\supp x$.} By Lemma~\ref{lem:PX}(a) the blocks of $\theta_x(d)$ inside $\supp x$ are the sets $K\cap\supp x$ for the classes $K$ of $\varepsilon_x\vee\varepsilon_d$ meeting $\supp x$, transversal iff $K$ meets $\supp d$; the blocks of $\theta_x(c)=v$ inside $\supp x$ are the $\T(v)$-blocks $V_1,\dots,V_{t'}$ (unions of $\T(c)$-blocks; transversal) and the $\N(c)$-blocks inside $\supp x$ (non-transversal). Since $\theta_x(d)=\theta_x(c)$: \emph{every block $Z$ of $d$ meets $\supp x$ inside a single $V_k$, or inside a single $\N(c)$-block $\tilde N\subseteq\supp x$, or not at all; and in the second case $Z$ is an $\N(d)$-block}, because the class of $\varepsilon_x\vee\varepsilon_d$ whose trace is $\tilde N$ does not meet $\supp d$. We call $Z$ of \emph{$V$-type}, \emph{$\tilde N$-type} or \emph{outer} accordingly.

Now let $Z$ be a block of $d$ lying in $\Lambda_s$. A point of $Z$ in $\supp b$ lies in some $B_r$, and then $r=s$: \emph{a block of $\Lambda_s$ meets $\supp b$ only inside $B_s$}; in particular, if $Z$ meets a $\T(a)$-block $A_i\subseteq Y$, then $Z$ meets $B_i=A_i$ and $i=s$, and if $Z$ meets $\supp u$ then $s\in I$ and $Z$ meets no $\T(a)$-block inside $Y$. Suppose $Z$ meets $\supp u$ at a point $q$; then $q\in B_s$. If $Z$ is of $V$-type, $q\in V_k\cap\supp u\subseteq\supp c\cap\supp u$ lies in some $C_p$, hence in $A_p\cap B_p$ by Step~2, and $p=s$: \emph{a $V$-type block of $\Lambda_s$ meets $\supp a$ only inside $A_s$}. If $Z$ is of $\tilde N$-type, $q\in\tilde N\cap B_s$, so $\tilde N\in\Delta_s$ (a block of $c$ meeting $B_s$), and $\tilde N$ lies in a class $\Gamma(\tilde N)$ of $\varepsilon_a\vee\varepsilon_c$ which contains at most one $\T(a)$-block, $A_r$ say; then \emph{$Z$ meets $\supp a$ only inside $\tilde N\cap A_r$}, and $\tilde N\in\Gamma_r\cap\Delta_s$.

\smallskip
\emph{Step 4: the key invariant.} Call two blocks of $d$ \emph{linked} if they are joined by a chain of blocks of $d$ and $\N(a)$-blocks (consecutive blocks meeting); a linked family lies in a single class of $\varepsilon_b\vee\varepsilon_d$, since $\N(a)$-blocks are $\N(b)$-blocks (Step~2). We claim: \emph{a linked family $\mathcal F$ of blocks of $d$ meets at most one $\T(a)$-block; and if $\mathcal F$ contains a $\T(d)$-block $D_p$, it meets $\supp a$ only inside $A_p$.} If no block of $\mathcal F$ meets $\supp b$ there is nothing to prove (a $\T(a)$-block met by a block of $\mathcal F$ is then $A_i=B_i\subseteq Y$, and $\mathcal F\subseteq\Lambda_i$). Otherwise $\mathcal F\subseteq\Lambda_s$ for a unique $s$, and by Step~3 a block of $\mathcal F$ meets $\supp a$ only inside $A_s$ (if it is of $V$-type, or of $\tilde N$-type with $\tilde N\in\Gamma_s$), or inside $A_r$ for an $\N(c)$-block $\tilde N\subseteq\supp x$ with $\tilde N\in\Gamma_r\cap\Delta_s$, $r\ne s$ (if it is of $\tilde N$-type and meets $\supp u$), or inside a $\T(a)$-block $A_i\subseteq Y$, which forces $i=s$. So the claim follows if we show: \emph{if some $Z_0\in\mathcal F$ is of $\tilde N_0$-type with $\tilde N_0\in\Gamma_r\cap\Delta_s$ and $r\ne s$ (or with $\Gamma(\tilde N_0)$ containing no $\T(a)$-block), then every block of $\mathcal F$ meeting $\supp x$ is of $\tilde N$-type for some $\tilde N\in\Gamma_r\cap\Delta_s$ (resp.\ $\tilde N\in\Gamma(\tilde N_0)\cap\Delta_s$).} Indeed, such a family meets $\supp a$ only inside $A_r$ (resp.\ not at all) by Step~3, and it contains no $\T(d)$-block: a $\T(d)$-block is not of $\tilde N$-type, and if it were outer, the chain joining it to $Z_0$ would have a first block meeting $\supp x$, joined to it through $\N(a)$-blocks lying in $Y$, which are $\N(x)$-blocks (Step~1; the blocks of $d$ in $Y$ need not be), so that this first block would lie in a class of $\varepsilon_x\vee\varepsilon_d$ meeting $\supp d$ and could not be of $\tilde N$-type. We prove the statement by induction along a chain $Z_0,N_1,Z_1,N_2,\dots$ in $\mathcal F$. Let $Z_j$ be of $\tilde N_j$-type with $\tilde N_j\in\Gamma_r\cap\Delta_s$ and let $N_{j+1}$ be an $\N(a)$-block meeting $Z_j$. If $N_{j+1}\subseteq Y$, it is an $\N(x)$-block, so $Z_{j+1}$ lies in the same class of $\varepsilon_x\vee\varepsilon_d$ as $Z_j$, whose trace on $\supp x$ is $\tilde N_j$: $Z_{j+1}$ is of $\tilde N_j$-type or outer, and in the latter case the next $\N(a)$-block meeting it lies in $Y$ as well (an $\N(a)$-block inside $\supp x$ cannot meet an outer block), so the class of $\varepsilon_x\vee\varepsilon_d$ persists until the next block meeting $\supp x$, which is then of $\tilde N_j$-type. If $N_{j+1}\subseteq\supp x$, it meets $Z_j$ inside $\tilde N_j$; being a block of $a$ meeting the block $\tilde N_j$ of $c$, it lies in $\Gamma_r$, and being a block of $b$ (Step~2) meeting $\tilde N_j$, it lies in $\Delta_s$. Let $Z_{j+1}$ meet $N_{j+1}$ at a point $q\in\supp x$, and let $E$ be the block of $c$ containing $q$. Then $E$ meets $N_{j+1}$, so $E\in\Gamma_r\cap\Delta_s$; $E$ is not a $\T(c)$-block, since $C_r\in\Delta_r\ne\Delta_s$, $C_s\in\Gamma_s\ne\Gamma_r$ and $C_p\notin\Gamma_r$ for $p\ne r$ (in the case of a class without $\T(a)$-block, no $\T(c)$-block lies in it); so $E$ is an $\N(c)$-block, inside $\supp x$ as $q\in\supp x$, and by Step~3 $Z_{j+1}\cap\supp x\subseteq E$: $Z_{j+1}$ is of $E$-type with $E\in\Gamma_r\cap\Delta_s$. This completes the induction and the proof of the claim.

\smallskip
\emph{Step 5: the classes of $\varepsilon_a\vee\varepsilon_d$.} Let $\rho=\varepsilon_a\vee\varepsilon_d$ and let $D(A_i)$ be the $\rho$-class of $A_i$.

(a) \emph{$D(A_i)$ contains no $\T(a)$-block other than $A_i$.} Let $A_i,A_j$ lie in one class and take a chain in $\rho$ from $A_i$ to $A_j$; between two consecutive $\T(a)$-blocks on it, the blocks are blocks of $d$ and $\N(a)$-blocks, so the blocks of $d$ there are linked, and by Step~4 the two $\T(a)$-blocks coincide. Hence $A_i=A_j$.

(b) \emph{$D_i\in D(A_i)$.} Take a chain from $B_i$ to $D_i$ in $\Lambda_i$; its blocks of $b$ other than $B_i$ are $\N(b)$-blocks, i.e.\ $\N(a)$-blocks. If $i\notin I$ then $B_i=A_i$ and the chain is a chain in $\rho$. If $i\in I$, let $Z$ be the block of $d$ following $B_i$ on the chain; the blocks of $d$ on the chain are linked, and the family contains $D_i$, so by the last remark of Step~4 $Z$ meets $\supp a$ only inside $A_i$; but $Z$ meets $B_i\subseteq\supp u\subseteq\supp a$, so $Z$ meets $A_i$, and replacing $B_i$ by $A_i$ gives a chain in $\rho$ from $A_i$ to $D_i$.

(c) \emph{$D(A_i)$ contains no $\T(d)$-block other than $D_i$.} Let $D_p\in D(A_i)$ and take a chain in $\rho$ from $D_p$ to $A_i$; let $A_r$ be the first $\T(a)$-block on it. The blocks of $d$ before $A_r$ are linked to $D_p$, so by Step~4 the one meeting $A_r$ meets $\supp a$ only inside $A_p$, whence $r=p$; and by (a) applied along the rest of the chain, $A_i=A_r=A_p$.

\emph{Conclusion.} By (a)--(c) the classes $D(A_1),\dots,D(A_s)$ are distinct and $D(A_i)$ contains exactly the $\T$-blocks $A_i$ and $D_i$; every class of $\rho$ meeting $\supp a\cup\supp d$ contains some $A_i$ or some $D_i$, hence is one of them. By Lemma~\ref{lem:PX}(c) and its proof, $a\bowtie d$. The statement about $\mathfrak P_8^2$ follows from Proposition~\ref{prop:census8}(c) and the minimality noted in Section~\ref{subsec:census8}.
\end{proof}

\begin{remark}\label{rem:P1}
(1) The proof is shorter than that of (R$'$) because the hypotheses make $a$ and $b$ almost equal: they have the same blocks outside $\supp x$ and the same $\N$-blocks, and differ only in how they partition each $\T(u)$-block into $\T$-blocks (where $u=xa=xb$). Off the ``diagonal'' $A_p\cap B_p$, which is where the $\T(c)$-blocks meet $\supp u$, the overlaps $A_r\cap B_s$ lie in $\N(c)$-blocks, and Step~4 says that the blocks of $d$ reaching such an overlap form a closed family of $\N(d)$-blocks. (2) Unlike (R$'$), the conclusion is not accompanied by a separation statement for a coarser join: in $\P(\P_4)$ the identity $ad=acbd$ holds in all $9\,519$ instances of the hypotheses, but the classes of $\varepsilon_a\vee\varepsilon_d\vee\varepsilon_b$ may contain two $\T(a)$-blocks. Step~4 was found by testing candidate statements on those instances and on $41\,918$ instances in $\P(\P_5)$ (\texttt{pa8/p1/}); the statement actually used (a linked family of blocks of $d$ meets at most one $\T(a)$-block) holds in all of them.
\end{remark}

\subsection{Proof of (Q)}\label{subsec:Q}

We write the path of (Q) as $a\bowtie c\bowtie b\bowtie d$, with the bijections $A_i\mapsto C_i\mapsto B_i\mapsto D_i$ of $\T$-blocks and the classes $\Gamma_i\ni A_i,C_i$ of $\varepsilon_a\vee\varepsilon_c$, $\Delta_i\ni C_i,B_i$ of $\varepsilon_c\vee\varepsilon_b$, $\Lambda_i\ni B_i,D_i$ of $\varepsilon_b\vee\varepsilon_d$ as in Section~\ref{subsec:P1}; the hypotheses are $xa=ax$, $\theta_c(x)=\theta_d(x)$, $ya=ay$, $yb=by$, $\theta_y(c)=\theta_y(b)$, and the conclusion is $a\bowtie d$, i.e.\ $\theta_a(d)=a$, i.e.\ $\rk(ad)=\rk a$. This subsection proves (Q) (Theorem~\ref{thm:Q}). The statements of Steps~2--7 below were found and checked on all instances of the hypotheses in $\P(\P_4)$ and on $10^5$--$10^6$ sampled instances in $\P(\P_5)$ and $\P(\P_6)$ (\texttt{pa8/q2/}) before being proved.

\smallskip
\emph{Step 1: normalising the auxiliaries.} Put $u=xa=ax$. Then $u\le a$ and $\theta_c(u)=\theta_d(u)$: in the binary notation $u=a*x=x*a$, and the axioms give $(A*B)*(B*C)=(A*B)*C$, i.e.\ $\theta_C(q)=\theta_{\theta_C(B)}(q)$ for $q\le B$, so that $\theta_c(u)=(a*x)*c=(a*x)*(x*c)=(a*x)*(x*d)=\theta_d(u)$ (Prover9, \texttt{pa8/q2/p9/}). So we may and do assume $x=u\le a$. Since $\theta_c(u)=\theta_d(u)=:z$ is a common lower bound of $c$ and $d$, \eqref{eq:Q1} gives $cd=dc$. Further consequences derived by Prover9 in the language of regular $*$-semigroups (\texttt{pa8/q2/p9s/}): $yc=ybc$, $ycb=yb$ (from $ya=ay$, $yb=by$, $ycy=yby$ and $b\bowtie c$ only), and $ud=ucd=ucbd$; the first is $yc=(yc)(yc)^*(yc)=(ycy)(c)=(yby)c=byc$, the second follows by multiplying with $b$ on the right, and $ud=ucbd$ is the $u$-analogue of the identity $x\,ad=x\,acbd$ of Remark~\ref{rem:census8}(5). The element $u$ is determined by a partition $P$ of the index set $\{1,\dots,s\}$ into \emph{groups} $G$ and a label for each group: the blocks of $u$ are the $\N(a)$-blocks and the sets $A_G=\bigcup_{i\in G}A_i$, transversal iff $G$ is labelled $\T$. Changing the label of a group from $\T$ to $\N$ preserves the hypothesis $\theta_c(u)=\theta_d(u)$ (verified; it is also clear from Step~2), and conversely a non-hub group (Step~2) may be labelled $\T$; so the labels carry no information, but it is convenient to keep them. We write $C_G=\bigcup_{i\in G}C_i$, $B_G$, $D_G$ similarly. Not every $u\le a$ with singleton groups satisfies the hypothesis: in $\P(\P_5)$ there are instances in which the given $u$ has a group of size two and no $u'$ with smaller groups works (\texttt{u0.py}).

\smallskip
\emph{Step 2: the structure of $c$ and $d$.} By the transport lemma (Step~2 of the proof of Theorem~\ref{thm:wiring}) $z=\theta_c(u)$ has the blocks $\N(c)\cup\{C_G:G\in P\}$, $C_G$ transversal iff $G$ is. Put $\Pi_1=\varepsilon_u\vee\varepsilon_c$. Its classes are the sets $\Gamma_G=\bigcup_{i\in G}\Gamma_i$ and the \emph{plain} classes of $\varepsilon_a\vee\varepsilon_c$ (those without $\T$-blocks), and since $z\le c$, $z\le d$ and $z=\theta_d(u)$ we have $\Pi_1=\varepsilon_u\vee\varepsilon_d=\varepsilon_a\vee\varepsilon_z$, so that $\varepsilon_a\vee\varepsilon_d\le\Pi_1$. More precisely, since $\varepsilon_a\le\varepsilon_u$, $\Pi_1=\varepsilon_u\vee\varepsilon_d=\rho\vee\varepsilon_u$ with $\rho=\varepsilon_a\vee\varepsilon_d$, and $\rho\vee\varepsilon_u$ is obtained from $\rho$ by identifying the classes that meet the same block $A_G$; so a class of $\rho$ containing no $\T(a)$-block is a class of $\Pi_1$, and hence \emph{every class of $\rho$ contained in some $\Gamma_G$ contains a $\T(a)$-block $A_i$, $i\in G$}, and every class of $\rho$ not contained in any $\Gamma_G$ is a plain class. Comparing the two descriptions of $z$ (blocks $\N(c)\cup\{C_G\}$, and $\N(d)$ together with the traces $L\cap\supp d$ of the classes $L$ of $\Pi_1$) gives, for each group $G$, one of two cases:
\begin{itemize}[nosep]
\item[(nh)] $C_G\subseteq\supp d$ and $C_G=\Gamma_G\cap\supp d$ (a \emph{non-hub} group);
\item[(h)] $C_G$ is an $\N(d)$-block, and $\Gamma_G\cap\supp d$ is empty or a single $\N(c)$-block $E_G$ (a \emph{hub}, as in case (F2) of the proof of Theorem~\ref{thm:S});
\end{itemize}
every other $\N(c)$-block is an $\N(d)$-block and conversely, and a plain class $W$ meets $\supp d$ in an $\N(c)$-block $E_W$ or not at all. The same construction from the side of $b$, with $w=\theta_b(z)\le b$ (which satisfies $\theta_c(w)=\theta_d(w)=z$ by transport along $b\bowtie c$, $b\bowtie d$) and \emph{the same partition $P$} (because $\theta_c(w)$ has the blocks $\N(c)\cup\{C_{G'}\}$ for the groups $G'$ of $w$), gives $\Delta_G=\Lambda_G$ for every $G$, where $\Delta_G=\bigcup_{i\in G}\Delta_i$ and $\Lambda_G=\bigcup_{i\in G}\Lambda_i$, with $\Lambda_G\cap\supp d=D_G$ (a block of $z=\theta_d(w)$, by transport along $b\bowtie d$); hence in case (nh) $C_G=D_G$, and in case (h) $D_G$ is a single $\N(c)$-block $E'_G\subseteq\Lambda_G$. Comparing the blocks $\N(c)\cup\{C_G\}$ of $z=\theta_c(u)$ with the blocks $\N(d)\cup\{D_G\}$ of $z=\theta_d(w)$ (transport along $b\bowtie d$), and using $C_G=D_G$ for non-hubs, we see that the $\N(c)$-blocks which are not $\N(d)$-blocks are exactly the sets $D_{G'}$, $G'$ a hub. Each nonempty trace $E_G$ or $E_W$ is an $\N(c)$-block contained in $\supp d$, hence one of these; conversely each $D_{G'}$, $G'$ a hub, is an $\N(c)$-block contained in $\supp d$, hence contained in the trace of its class of $\Pi_1$, which is a hub class or a plain class (as $D_{G'}\cap\supp c=\emptyset$), and equals that trace. So the nonempty traces $E_G$ of the hubs and $E_W$ of the plain classes are exactly the sets $D_{G'}=E'_{G'}$, $G'$ a hub, each occurring once; write $E_G=D_{\pi(G)}$, $E_W=D_{\pi(W)}$. In particular the $\T(d)$-blocks contained in a non-hub class $\Gamma_G$ are exactly $D_k$, $k\in G$, and those contained in a hub class $\Gamma_G$ form $D_{\pi(G)}$ for a hub $\pi(G)$. Consequently $\varepsilon_a\vee\varepsilon_d\le\Pi_1$ has, inside a non-hub class, $|G|$ blocks of each kind, and a proof of $a\bowtie d$ has to show that these are matched by $\varepsilon_a\vee\varepsilon_d$, and that $|\pi(G)|=|G|$ for hubs, again with a matching. The hub $C_G$ lies in a single class $\Lambda_j$, $j\in G$, and $E_G$ in a single class $\Delta_{k'}$, $k'\in\pi(G)$. Hubs have $|G|\le2$, and a hub $G$ of size $2$ is \emph{cross-wired}: writing $G=\{k,l\}$ with $C_G\subseteq\Lambda_l$, we have $D_G\subseteq\Delta_k$, so that $C_k\subseteq\Delta_k\cap\Lambda_l$, $C_l\subseteq\Delta_l\cap\Lambda_l$, $D_k\subseteq\Delta_k\cap\Lambda_k$, $D_l\subseteq\Delta_k\cap\Lambda_l$, and in particular $C_k$ and $D_l$ do not meet $\supp b$ (the $\T(b)$-block $B_m$ lies in $\Delta_m\cap\Lambda_m$). Indeed, let $C_G\subseteq\Lambda_j$ and $D_G\subseteq\Delta_{j'}$ with $j,j'\in G$, and let $i\in G$. A path in $\Lambda_i$ from $D_i$ to $B_i$ alternates between blocks of $d$ and blocks of $b$; its interior blocks of $d$ are $\N(d)$-blocks, i.e.\ common blocks or hubs $C_{G''}$; every block of $b$, every common block (an $\N(c)$-block) and $D_i\subseteq D_G$ (contained in the $\N(c)$-block $D_G$) lies in a single class of $\varepsilon_c\vee\varepsilon_b$, while a hub $C_{G''}$ meets exactly the classes $\Delta_m$, $m\in G''$. So the class of $\varepsilon_c\vee\varepsilon_b$ changes along the path only at hubs $C_{G''}$, and only between indices in $G''$. The path starts in $\Delta_{j'}$ and ends in $\Delta_i$; if $i\ne j'$, the first change happens at a hub $C_{G''}$ with $j'\in G''$, i.e.\ $G''=G$, and since $C_G\subseteq\Lambda_j$ lies on a path in $\Lambda_i$, we get $i=j$. Hence $G\subseteq\{j,j'\}$, and if $|G|=2$ then $j\ne j'$; put $l=j$, $k=j'$. Further, $\pi$ need not be the identity: in $\P(\P_4)$, $a=(01|23)$, $c=(0|\N1|2|\N3)$, $b=(03|12)$, $d=(\N0|1|\N2|3)$ (writing only the upper parts, $\N$ marking $\N$-blocks) with $u=a$ relabelled $\N$ and $y$ the rank $0$ projection with the single block $X$ has two singleton hubs exchanged by $\pi$; here $ad\ne acbd$.

\smallskip
\emph{Step 3: the structure imposed by $y$.} Since $y$ commutes with $a$ and $b$, Lemma~\ref{lem:PX}(b) describes $\varepsilon_y\vee\varepsilon_a$ and $\varepsilon_y\vee\varepsilon_b$; in particular every $\N(y)$-block is an $\N(a)$-block or a union of $\T(a)$-blocks, and an $\N(b)$-block or a union of $\T(b)$-blocks (four kinds: \emph{$ab$-common} blocks $=$ $\N(a)$-blocks that are $\N(b)$-blocks, \emph{bridges} $=$ $\N(b)$-blocks that are unions of $\T(a)$-blocks, \emph{co-bridges}, and unions of $\T(a)$-blocks that are unions of $\T(b)$-blocks), and every block of $a$ or $b$ lies in $\supp y$ or in an $\N(y)$-block. The hypothesis $\theta_y(c)=\theta_y(b)$ says (Lemma~\ref{lem:PX}(a),(b)): the classes $K$ of $\varepsilon_y\vee\varepsilon_c$ meeting $\supp y$ have traces $K\cap\supp y$ equal to the classes of $\varepsilon_y\vee\varepsilon_b$ of type (i) (unions $Q$ of $\T(y)$-blocks which are unions of $\T(b)$-blocks; these $K$ meet $\supp c$; we call them classes \emph{of the first kind}) and of type (ii) (single $\N(b)$-blocks $F\subseteq\supp y$; these $K$ do not meet $\supp c$; \emph{of the second kind}). We shall use the following consequence repeatedly. \textup{($*$)} \emph{Let $K$ be a class of $\varepsilon_y\vee\varepsilon_c$ of the second kind, with trace an $\N(b)$-block $F\subseteq\supp y$, and let $P$ be a path in a class $\Delta_m$ of $\varepsilon_c\vee\varepsilon_b$ (blocks of $c$ and of $b$ alternately, consecutive blocks meeting, all distinct) or in a class $\Lambda_k$ of $\varepsilon_b\vee\varepsilon_d$ (blocks of $d$ and of $b$), whose first block is $F$ or a block of $c$ lying in $K$. Then every block of $c$ on $P$ or meeting a block of $b$ of $P$ lies in $K$; in particular $P$ contains no $\T(c)$-block; every block of $d$ on $P$ is a common block (an $\N(c)$-block which is an $\N(d)$-block; it lies in $K$), except possibly the $\T(d)$-block $D_k$, which is then not contained in $\supp c$, so that the $\N(c)$-block containing it lies in $K$.} Call a set \emph{covered} if each of its points lies in a block of $y$ or of $c$ belonging to $K$; $F$ is covered (its points lie in $\T(y)$-blocks, which are in $K$ as $F$ is the trace), and so is a block of $c$ in $K$. We show inductively that every block of $P$ is covered, which gives all assertions: let $W$ be covered and let $W'$ be the next block of $P$, meeting $W$ at $q$, so that $q$ lies in a block $Y$ of $y$ or of $c$ with $Y\in K$. If $W'$ is a block of $c$, then $W'=Y$ or $W'$ meets the $y$-block $Y$, so $W'\in K$; the same applies to any block of $c$ meeting $W$ if $W$ is a block of $b$, and shows that a block of $c$ meeting a covered set lies in $K$. If $W'$ is an $\N(b)$-block, it is an $\N(y)$-block, and then $W'\in K$ ($W'=Y$, or $W'$ meets the block $Y$ of $c$), or a union of $\T(y)$-blocks, and then the $\T(y)$-block through $q$ lies in $K$, hence in $F$, so $W'=F$; in both cases $W'$ is covered. If $W'$ is the $\T(b)$-block $B_m$ (resp.\ $B_k$), it is not contained in $\supp y$ (the $\T(y)$-block through $q$ would lie in $K$, hence in $F$, but $F\cap W'=\emptyset$), so it lies in an $\N(y)$-block $Q'$, which meets $Y$ and lies in $K$, and $W'$ is covered. If $W'$ is a block of $d$, it is a common block (a block of $c$, covered by the first case), or a hub, or $D_k$; a hub is a union of $\T(c)$-blocks, and the $\T(c)$-block through $q$ would lie in $K$, which does not meet $\supp c$, impossible; and for $D_k$ the same argument shows $D_k\not\subseteq\supp c$, so $D_k$ lies in an $\N(c)$-block (Step~2), which lies in $K$.

From this: (a) an $\N(b)$-block $F$ lying in the class $\Delta_m$ of a $\T(c)$-block is an $\N(y)$-block: otherwise it is a union of $\T(y)$-blocks, the class $K$ with trace $F$ is of the second kind, and a path in $\Delta_m$ from $F$ to $C_m$ contradicts ($*$); (a$'$) an $\N(b)$-block $F$ lying in the class $\Lambda_k$ of a $\T(d)$-block is an $\N(y)$-block: otherwise, with $K$ as before, a path in $\Lambda_k$ from $F$ to $D_k$ shows by ($*$) that $D_k$ lies in an $\N(c)$-block $E\in K$, which is $E_{G'}$ or $E_W$ for a hub $G'$ or a plain class $W$, hence equals $D_{G''}$ for a hub $G''$ (Step~2) and lies in a class $\Delta_{k''}$, $k''\in G''$ (as $D_{G''}\subseteq\Lambda_{G''}=\Delta_{G''}$), and a path in $\Delta_{k''}$ from $E$ to $C_{k''}$ contradicts ($*$); (b) a $\T(y)$-block meeting $\supp c$ lies in a class $K$ of the first kind, hence in a set $Q$ as above, hence in $\supp b$; so a $\T(y)$-block meets $C_i$ only inside $B_i\cap C_i$ (the only $\T(b)$-block meeting $C_i\subseteq\Delta_i$ is $B_i$); (b$'$) a $\T(y)$-block $Y$ meeting $D_k$ lies in $\supp b$, and meets $D_k$ only inside $B_k\cap D_k$: if $D_k\subseteq\supp c$ this is (b), and otherwise $D_k$ lies in an $\N(c)$-block $E\subseteq\Delta_{k''}$ as in (a$'$), and if the class $K$ of $Y$ were of the second kind, $E\in K$ and a path in $\Delta_{k''}$ from $E$ to $C_{k''}$ would contradict ($*$); (c) the $\T(c)$-blocks in the class with trace $Q$ are exactly the $C_i$ with $B_i\subseteq Q$ (this is $ycb=yb$); (d) every $\N(b)$-block meeting $\supp c\cup\supp d$ lies in some $\Delta_m$ or $\Lambda_k$, hence is an $\N(y)$-block by (a), (a$'$), hence an $\N(a)$-block or a bridge; (e) every block $Z$ of $a$ satisfies $Z\cap(\supp c\cup\supp d)\subseteq\supp b$, or $Z$ is contained in an $\N(b)$-block: if $Z\subseteq\supp y$ this is (b), (b$'$), and otherwise $Z$ lies in an $\N(y)$-block, which is an $\N(b)$-block or a union of $\T(b)$-blocks. (The identity $\theta_y(d)=\theta_y(b)$ itself holds in every instance, but it is not a consequence of the hypotheses in projection algebras, not even together with (S): Mace4 finds $6$-element countermodels; we do not need it.)

\smallskip
\emph{Step 4: plain classes contain no $\T(d)$-block.} Let $W$ be a plain class; then $W\cap\supp a=W\cap\supp c=\emptyset$, every block of $a$ or of $c$ meeting $W$ lies in $W$, and so does every block of $y$ meeting $W$: an $\N(y)$-block meeting $W$ is an $\N(a)$-block (it cannot be a union of $\T(a)$-blocks), and a $\T(y)$-block lies in a union of $\T(a)$-blocks or in an $\N(a)$-block (Lemma~\ref{lem:PX}(b), as $ya=ay$), the latter alternative being forced if it meets $W$. Hence \emph{$W$ is a union of classes of $\varepsilon_y\vee\varepsilon_c$}. Suppose now $D_j\subseteq E_W\subseteq W$. Take a path in $\Lambda_j$ from $D_j$ to $B_j$. Its $\N(b)$-blocks lie in $\Lambda_j$, so they are $\N(y)$-blocks (Step~3(a$'$)), hence $\N(a)$-blocks or bridges; a bridge is a union of $\T(a)$-blocks and cannot meet $W$, and an $\N(a)$-block meeting $W$ lies in $W$; the interior $\N(d)$-blocks of the path are common blocks (blocks of $c$, which lie in $W$ if they meet it) or hubs (unions of $\T(c)$-blocks, which cannot meet $W$). So inductively every block of the path except possibly $B_j$ lies in $W$, and $B_j$ meets $W$. If $B_j$ lies in an $\N(y)$-block $Q$, then $Q$ meets $W$, so $Q$ is an $\N(a)$-block contained in $W$ and $B_j\subseteq W$; a path in $\Delta_j$ from $B_j$ to $C_j$ then lies in $W$ by the same argument (its $\N(b)$-blocks are $\N(y)$-blocks by Step~3(a), and its blocks of $c$ lie in $W$ as soon as they meet it), so $C_j\subseteq W$, contradicting $W\cap\supp c=\emptyset$. If $B_j\subseteq\supp y$, take $q\in B_j\cap W$ and let $Y$ be the $\T(y)$-block of $q$ and $K$ its class in $\varepsilon_y\vee\varepsilon_c$; then $K\subseteq W$. The trace $K\cap\supp y$ is a class of $\varepsilon_y\vee\varepsilon_b$ containing $Y$, which meets $B_j$, so it is not a single $\N(b)$-block and contains $B_j$: $K$ is of the first kind, with trace $Q\supseteq B_j$, and $C_j\subseteq K\subseteq W$ by Step~3(c), again contradicting $W\cap\supp c=\emptyset$. Consequently every hub $G$ has $E_G\ne\emptyset$, and $\pi$ is a permutation of the hubs.

\smallskip
\emph{Step 5: types, index pairs and contacts.} Put $\rho=\varepsilon_a\vee\varepsilon_d$. By Step~2, $\rho\le\Pi_1$, every class of $\rho$ contained in a class $\Gamma_G$ contains some $A_i$, $i\in G$, and a class of $\rho$ not contained in any $\Gamma_G$ is a plain class, which contains no $\T$-block of $a$ (Step~2) or of $d$ (Step~4). By Lemma~\ref{lem:PX}(a) and (c) (the blocks of $\theta_a(d)$ are $\N(a)$ and the sets $U_K$, $K$ a class of $\rho$ meeting $\supp a$, transversal iff $K$ meets $\supp d$; and dually for $\theta_d(a)$), $a\bowtie d$ is therefore equivalent to: \emph{for every group $G$, every class of $\rho$ contained in $\Gamma_G$ contains exactly one $\T(a)$-block and exactly one $\T(d)$-block}. This is proved for non-hub groups in Step~6 and for hubs in Step~7. Inside $\Gamma_G$ the blocks of $a$ are $A_i$ ($i\in G$) and $\N(a)$-blocks, and the blocks of $d$ are the \emph{common} blocks (the $\N(c)$-blocks which are $\N(d)$-blocks; in a non-hub class these are all the $\N$-blocks of $c$ and of $d$ in $\Gamma_G$, and $C_G=D_G$ is the union of the $\T$-blocks of both), and, in a hub class, the hub $C_G$ and the $\T(d)$-blocks $D_m$, $m\in\pi(G)$, whose union is the $\N(c)$-block $E_G$.

\emph{Types of blocks.} By Step~3 every block $Z$ of $a$ either lies in $\supp y$, or lies in an $\N(y)$-block $Q$, which is then an $\N(b)$-block or a union of $\T(b)$-blocks. Call $Z$ of \emph{$\N$-type} if it lies in an $\N(y)$-block which is an $\N(b)$-block (if $Z$ is an $\N(a)$-block this means $Z$ itself is an $\N(b)$- and $\N(y)$-block; if $Z=A_i$ it means that $A_i$ lies in a bridge $Q$, an $\N(b)$- and $\N(y)$-block which is a union of $\T(a)$-blocks), and of \emph{$b$-type} otherwise; a $b$-type block lies in $\supp y$ or in an $\N(y)$-block contained in $\supp b$, so by Step~3(b), (b$'$) it meets each $C_m$ only inside $B_m\cap C_m$ and each $D_m$ only inside $B_m\cap D_m$. For an $\N(b)$-block $Q$ write $\delta(Q)=m$ if $Q\subseteq\Delta_m$ and $\lambda(Q)=k$ if $Q\subseteq\Lambda_k$ (undefined if the class of $Q$ in $\varepsilon_c\vee\varepsilon_b$, resp.\ $\varepsilon_b\vee\varepsilon_d$, contains no $\T$-block); an $\N$-type block $Z\subseteq Q$ meets $C_m$ only if $\delta(Q)=m$ and $D_k$ only if $\lambda(Q)=k$. A block of $a$ meeting $\supp c\cup\supp d$ at a point outside $\supp b$ is of $\N$-type (Step~3(e), (d)).

In a non-hub class every point of $C_G=D_G$ lies in a unique $\T(c)$-block $C_m$ and a unique $\T(d)$-block $D_k$ (its \emph{cell} $(m,k)$), and a point of $C_G$ lying in $\supp b$ lies in $B_m\cap C_m\cap D_m$ for some $m$: its cell is \emph{diagonal}. Hence a $b$-type block of $a$ meets $C_G$ only in diagonal cells, and since it meets at most one $\T(c)$-block (it lies in one class $\Gamma_m$), only in the cell $(m,m)$; while an $\N$-type block $Z\subseteq Q$ meets $C_G$ only in points outside $\supp b$, in the cell $(\delta(Q),\lambda(Q))$.

In a hub class, with $C_G\subseteq\Lambda_j$ ($j\in G$) and $E_G\subseteq\Delta_{k'}$ ($k'\in\pi(G)$) as in Step~2: a $b$-type block meets $C_G$ only in the sets $B_i\cap C_i$, $i\in G$, and $E_G$ only in the sets $B_m\cap D_m$, $m\in\pi(G)$; an $\N$-type block $Z\subseteq Q$ meets $C_G$ only if $\lambda(Q)=j$, and then inside $C_{\delta(Q)}$, and meets $E_G$ only if $\delta(Q)=k'$, and then inside $D_{\lambda(Q)}$.

\emph{Index pairs.} For a common block $E$ let $\lambda(E)$ be the index $k$ with $E\subseteq\Lambda_k$ (undefined if the class of $E$ in $\varepsilon_b\vee\varepsilon_d$ contains no $\T(d)$-block) and $\delta(E)$ the index $m$ with $E\subseteq\Delta_m$ (undefined likewise). (i) \emph{If $E$ meets a $b$-type block of $a$ at a point $q$, then either $q\in B_m$ for some $m$ and $(\lambda(E),\delta(E))=(m,m)$, or $q\notin\supp b$ and both indices are undefined.} Indeed, if $q\in B_m$ then $E$ meets $B_m$. Otherwise $q$ lies in $\supp y$ (a $b$-type block not contained in $\supp y$ lies in $\supp b$), in a $\T(y)$-block $Y$, which lies in an $\N(b)$-block $F\subseteq\supp y$ (Step~3), and the class $K$ of $\varepsilon_y\vee\varepsilon_c$ with trace $F$ is of the second kind and contains $E$ (a block of $c$ meeting $Y$). If $\delta(E)=m$ were defined, a path in $\Delta_m$ from $E$ to $C_m$ would contradict Step~3($*$). If $\lambda(E)=k$ were defined, a path in $\Lambda_k$ from $E$ to $D_k$ shows by ($*$) that $D_k$ lies in an $\N(c)$-block $E'\in K$, which lies in a class $\Delta_{k''}$ (Step~2, as in Step~3(a$'$)), and a path in $\Delta_{k''}$ from $E'$ to $C_{k''}$ contradicts ($*$). (ii) If $E$ meets an $\N$-type block $Z\subseteq Q$ then $(\lambda(E),\delta(E))=(\lambda(Q),\delta(Q))$ (as indices or as ``undefined''), as $E$ meets the $b$-block $Q$. (iii) Two common blocks meeting the same $\N$-type $\N(a)$-block have the same pair. Consequently, along a path $E_0,F_1,E_1,\dots,F_r,E_r$ alternating between common blocks and $\N(a)$-blocks (consecutive blocks meeting; no hub involved), the pair is preserved across $\N$-type $\N(a)$-blocks and, at a $b$-type $\N(a)$-block, the pairs of the two adjacent common blocks are each either diagonal ($(m,m)$) or undefined in both coordinates. In particular (\emph{rigidity}): \emph{if some common block of such a path, or some $\N$-type $\N(a)$-block of it, has a defined non-diagonal pair $(k,m)$, $k\ne m$, then every $\N(a)$-block of the path is of $\N$-type and every block of the path has the pair $(k,m)$; and a $\T(a)$-block met by a common block of the path is then of $\N$-type with the pair $(k,m)$ as well} (by (i) and (ii), which apply to $\T(a)$-blocks too).

\smallskip
\emph{Step 6: non-hub classes.}

\begin{proposition}\label{prop:Qnh}
In the situation of Steps~\textup{1--5}, let $G$ be a non-hub group. Then every class of $\rho$ contained in $\Gamma_G$ contains exactly one $\T(a)$-block and exactly one $\T(d)$-block.
\end{proposition}

\begin{proof}
Let $\mathcal K$ be the set of components of the graph on the $\N(a)$-blocks and the common blocks contained in $\Gamma_G$ (blocks joined when they meet); each $\kappa\in\mathcal K$ is contained in a single class $\Gamma_m$, so it touches at most one $\T(a)$-block, $A_m$, and its $\N(a)$-blocks meet $C_G$ only in $C_m$. Say that $D_k$ is \emph{attached} to $A_i$ if $A_i$ meets $D_k$, or if some $\kappa$ touching $A_i$ (a common block of $\kappa$ meets $A_i$) contains an $\N(a)$-block meeting $D_k$. The classes of $\rho$ inside $\Gamma_G$ are generated by these adjacencies, so it suffices to prove:
\begin{enumerate}[nosep,label=(\alph*)]
\item every $\kappa$ contains $\N(a)$-blocks meeting at most one $\T(d)$-block;
\item at most one $\T(d)$-block is attached to each $A_i$;
\item at least one $\T(d)$-block is attached to each $A_i$.
\end{enumerate}
Given (a)--(c): by (b), (c) the attachment defines a map $\tau\colon G\to G$; every $D_k$ lies in a class of $\rho$ containing some $A_i$ (Step~5), and a $\rho$-path from $D_k$ to $A_i$ starts with $D_k$, a block of $a$ meeting it, and, if that block is an $\N(a)$-block, continues inside its $\kappa$ until the first $\T$-block met, which by (a) is a $\T(a)$-block; so every $D_k$ is attached to some $A_i$, $\tau$ is surjective, hence bijective; and then the $\rho$-class of $A_i$ contains the $\T$-blocks $A_i$ and $D_{\tau(i)}$ only (a further $\T(a)$-block $A_j$ in it would be reached through some $D_k$, which would be attached to both, and a further $\T(d)$-block through some $\kappa$ or $A$, contradicting (a) or (b)), and every class of $\rho$ in $\Gamma_G$ is of this form.

(c): Take a path in $\Gamma_i$ from $A_i$ to $C_i$; its blocks of $c$ are common blocks and its blocks of $a$ are $\N(a)$-blocks, so they lie in one $\kappa$ touching $A_i$, and the last block of $a$ on it (or $A_i$ itself) meets $C_i\subseteq C_G=D_G$, hence some $D_k$.

(b): Let $D_k$ be attached to $A_i$. \emph{Case $A_i$ of $b$-type.} A direct contact $A_i\cap D_k\ne\emptyset$ is in the cell $(i,i)$, so $k=i$. Otherwise $\kappa$ touches $A_i$ at a common block $E$ whose pair is diagonal or undefined (i), and contains an $\N(a)$-block $F$ meeting $D_k$ at a point of $C_i$. If $F$ is of $b$-type, the contact is in the cell $(i,i)$ and $k=i$. If $F$ is of $\N$-type, $F$ is an $\N(b)$-block in $\Lambda_k\cap\Delta_i$, so the common block of $\kappa$ adjacent to $F$ on a path from $E$ has the pair $(k,i)$; but along the path from $E$ all pairs are diagonal or undefined, so $k=i$. \emph{Case $A_i\subseteq Q$ of $\N$-type.} A direct contact lies in the cell $(i,\lambda(Q))$, so $k=\lambda(Q)$ and $\delta(Q)=i$. Otherwise the common block $E$ of $\kappa$ meeting $A_i$ has the pair $(\lambda(Q),\delta(Q))$, and along the path from $E$ to the contact block $F$ the pair is preserved up to the first $b$-type $\N(a)$-block, if any. If there is none: $F$ is of $\N$-type with $(\lambda(F),\delta(F))=(k,i)$ equal to $(\lambda(Q),\delta(Q))$, or of $b$-type with $k=i$ and the pair $(\lambda(Q),\delta(Q))$ diagonal or undefined; if there is one, then $(\lambda(Q),\delta(Q))$ is diagonal or undefined and the contact gives $k=i$ (for an $\N$-type $F$ the adjacent pair $(k,i)$ is diagonal). So every attached $D_k$ is $D_{\lambda(Q)}$, and then $\delta(Q)=i$, or $D_i$, and then $\lambda(Q)=\delta(Q)$ or both are undefined; if both kinds occur then $\lambda(Q)=\delta(Q)=i$. Hence at most one $\T(d)$-block is attached.

(a): Let $\kappa$ contain $F$ meeting $D_k$ and $F'$ meeting $D_l$, in cells $(\gamma,k)$ and $(\gamma,l)$ where $\kappa\subseteq\Gamma_\gamma$. If $F$ is of $b$-type then $k=\gamma$, and likewise for $F'$. If $F'$ is of $\N$-type, it is an $\N(b)$-block in $\Lambda_l\cap\Delta_\gamma$, and the common block next to $F'$ on a path from $F'$ to $F$ has the pair $(l,\gamma)$; following the path, if it contains no $b$-type block before $F$ then either $F$ is of $\N$-type and $(k,\gamma)=(l,\gamma)$, or $F$ is of $b$-type and the pair $(l,\gamma)$ is diagonal, $l=\gamma=k$; if it contains a $b$-type block, the pair $(l,\gamma)$ is diagonal, $l=\gamma$, and by symmetry (the case $F$ of $\N$-type gives $k=\gamma$ in the same way) $k=\gamma$. In all cases $k=l$.
\end{proof}

\smallskip
\emph{Step 7: hub classes.}

\begin{proposition}\label{prop:Qh}
In the situation of Steps~\textup{1--5}, let $G$ be a hub. Then $|\pi(G)|=|G|$, and every class of $\rho$ contained in $\Gamma_G$ contains exactly one $\T(a)$-block and exactly one $\T(d)$-block.
\end{proposition}

\begin{proof}
Write $G'=\pi(G)$; by Steps~2 and~4, $E_G=D_{G'}$ is a nonempty $\N(c)$-block contained in $\Gamma_G$, hence in a single class $\Gamma_e$, $e\in G$, and in a single class $\Delta_{k'}$, $k'\in G'$; and $C_G\subseteq\Lambda_j$, $j\in G$. If $|G|=2$ we write $G=\{k,l\}$ with $l=j$, so that (cross wiring, Step~2) $C_k\subseteq\Delta_k\cap\Lambda_l$, $C_l\subseteq\Delta_l\cap\Lambda_l$ and $C_k\cap\supp b=\emptyset$. If $|G'|=2$ we write $G'=\{k',l'\}$, so that (cross wiring for $G'$) $D_{k'}\subseteq\Delta_{k'}\cap\Lambda_{k'}$, $D_{l'}\subseteq\Delta_{k'}\cap\Lambda_{l'}$ and $D_{l'}\cap\supp b=\emptyset$. Let $\mathcal N$ be the set of components $\nu$ of the graph on the $\N$-blocks of $a$ and $d$ inside $\Gamma_G$ (the $\N(a)$-blocks, the common blocks and $C_G$; joined when they meet), and say that $D_m$ ($m\in G'$) is \emph{attached} to $A_i$ ($i\in G$) if there is a path in the graph $\Gamma(a,d)$ from $A_i$ to $D_m$ all of whose interior blocks are $\N$-blocks, i.e.\ $A_i$ meets $D_m$, or some $\nu$ touching $A_i$ contains an $\N(a)$-block meeting $D_m$. All paths considered below are shortest paths between their endpoints in the graph in question, hence simple, and an attachment path from $A_i$ to $D_m$ is always taken shortest, so that its interior is a shortest path in a component $\nu$. A component $\nu$ not containing $C_G$ lies in a single class $\Gamma_i$; a shortest path in a component $\nu$ between two blocks lying in the same class $\Gamma_i$ (or a shortest attachment path with both ends in $\Gamma_i$) passes through $C_G$ at most once, and if it does, both neighbours of $C_G$ on it are $\N(a)$-blocks in $\Gamma_i$ meeting $C_G$ inside $C_i$, and the two parts of the path avoid $C_G$ (each part lies in a single class of $\varepsilon_a\vee\varepsilon_c$, being a chain of blocks of $a$ and of $c$); a path in $\nu$ between blocks in different classes $\Gamma_k$, $\Gamma_l$ passes through $C_G$. We write $G=\{k,l\}$, $G'=\{k',l'\}$ as above; when $G'=G$ we have $k'=k$ and $l'=l$, because $E_G=D_G$ lies in $\Delta_{k'}$ and in $\Delta_k$. By Step~5, inside $\Gamma_G$:
\begin{enumerate}[nosep,label=(\roman*)]
\item a block of $a$ meeting $C_G$ in $C_i$ is of $\N$-type, contained in an $\N(b)$-block $Q$ with $(\lambda(Q),\delta(Q))=(j,i)$, or it is of $b$-type, $i=j$ and the contact lies in $B_j\cap C_j$;
\item a block of $a$ meeting $D_{k'}$ is of $\N$-type with pair $(k',k')$, or of $b$-type with contact in $B_{k'}\cap D_{k'}$;
\item if $|G'|=2$, a block of $a$ meeting $D_{l'}$ is of $\N$-type with pair $(l',k')$.
\end{enumerate}

\emph{Claim 1: the $\N(a)$-blocks of a single component $\nu$ meet at most one $\T(d)$-block.} Suppose $Z,Z'\in\nu$ meet $D_m$, $D_{m'}$ with $m\ne m'$, so $|G'|=2$, say $m'=l'$, $m=k'$. Both lie in $\Gamma_e$; take a shortest path in $\nu$ from $Z'$ to $Z$. If it avoids $C_G$, then by (iii) and rigidity all its blocks have the pair $(l',k')$, contradicting (ii) for $Z$. If it passes through $C_G$, the part from $Z'$ to the neighbour $Z_1$ of $C_G$ avoids $C_G$, so $Z_1$ has the pair $(l',k')$ by rigidity, and $Z_1$ meets $C_G$ inside $C_e$, so $(\lambda,\delta)=(j,e)$ by (i); thus $k'=e\in G\cap G'$, so $G'=G$, $|G|=2$, $k'=k$, $e=k$, $l'=l=j$. Then the other neighbour $Z_2$ of $C_G$ meets $C_G$ inside $C_k$, so it is of $\N$-type with pair $(l,k)$ by (i) (as $C_k\cap\supp b=\emptyset$), and rigidity along the part from $Z_2$ to $Z$ gives $Z$ the pair $(l,k)$, contradicting (ii).

\emph{Claim 2: at most one $\T(d)$-block is attached to each $A_i$.} Suppose $D_{k'}$ and $D_{l'}$ are attached to $A_i$, so $|G'|=2$. By (iii) the contact block of $D_{l'}$ (which is $A_i$ itself or an $\N(a)$-block $Z'$) has the pair $(l',k')$. \emph{Case $G'\ne G$.} An attachment path from $A_i$ to $D_{l'}$ cannot pass through $C_G$: the neighbour of $C_G$ on the side of $D_{l'}$ would have the pair $(l',k')$ by rigidity and $\delta=i''\in G$ by (i), but $k'\notin G$. So by rigidity $A_i$ is of $\N$-type, $A_i\subseteq Q$ with $(\lambda(Q),\delta(Q))=(l',k')$. Now consider an attachment path from $A_i$ to $D_{k'}$. It does not start with a direct contact ($\lambda(Q)=l'\ne k'$), nor with $C_G$ ($\lambda(Q)=l'\ne j$, as $l'\notin G$), so it starts with a common block meeting $A_i$, of pair $(l',k')$, and by rigidity all its blocks have this pair until it reaches $C_G$ or the contact block of $D_{k'}$; the former is excluded as before ($k'\notin G$), the latter by (ii). \emph{Case $G'=G$}, so $|G|=2$, $k'=k$, $l'=l=j$. Suppose first $i=k$. If $A_k$ meets $D_l$, then $A_k$ is of $\N$-type with pair $(l,k)$ by (iii) and $e=k$. Otherwise take a shortest attachment path from $A_k$ to $D_l$; its contact block has the pair $(l,k)$. If the path passes through $C_G$, the neighbour of $C_G$ on the side of the contact block has the pair $(l,k)$ by rigidity, so by (i) it meets $C_G$ inside $C_k$ and lies in $\Gamma_k$, and so does the contact block, i.e.\ $e=k$; and the neighbour on the side of $A_k$ then also lies in $\Gamma_k$, meets $C_k$ and has the pair $(l,k)$ by (i) ($C_k\cap\supp b=\emptyset$). If the path avoids $C_G$, all its blocks lie in $\Gamma_k$, so again $e=k$. In both cases rigidity along the part of the path next to $A_k$ (or (i), if that part is empty and $A_k$ meets $C_G$) shows that $A_k$ is of $\N$-type with pair $(l,k)$. The attachment path from $A_k$ to $D_k$ then starts with the pair $(l,k)$ (a direct contact is excluded by $\lambda=l\ne k$; a common block meeting $A_k$ has the pair $(l,k)$; and $C_G$, if it comes first or later, has both neighbours in $\Gamma_k$, of pair $(l,k)$, as $e=k$), so by rigidity it never reaches a block of pair $(k,k)$ or of $b$-type, contradicting (ii). Suppose now $i=l$. If the attachment path from $A_l$ to $D_l$ avoids $C_G$, then $e=l$, $A_l$ is of $\N$-type with pair $(l,k)$ and, as $C_G$'s neighbours in $\Gamma_l$ have $\delta=l$ by (i), no path from $A_l$ with pair $(l,k)$ reaches $C_G$; so the attachment path from $A_l$ to $D_k$ keeps the pair $(l,k)$ by rigidity and cannot reach a contact block as in (ii); a direct contact is excluded by $\lambda(Q)=l\ne k$. If it passes through $C_G$, then $e=k$ (a shortest path between blocks of $\Gamma_l$ through $C_G$ would have both neighbours of $C_G$ in $\Gamma_l$, of $b$-type or with $\delta=l$ by (i), whereas rigidity from the contact block gives the neighbour on its side the pair $(l,k)$); a direct contact $A_l\cap D_k$ is then impossible ($A_l\subseteq\Gamma_l$, $D_k\subseteq\Gamma_k$), and an attachment path from $A_l\subseteq\Gamma_l$ to $D_k\subseteq\Gamma_k$ passes through $C_G$ with its neighbour on the side of $D_k$ in $\Gamma_k$, of pair $(l,k)$ by (i), which by rigidity cannot reach a contact block as in (ii).

\emph{Consequences.} Every $D_m$, $m\in G'$, is attached to some $A_i$: its class of $\rho$ contains some $A_i$ (Step~5), and on a shortest $\rho$-path from $D_m$ to $A_i$ the first $\T$-block after $D_m$ is reached through $\N$-blocks of one component $\nu$, which contains an $\N(a)$-block meeting $D_m$ (unless the block following $D_m$ is already a $\T(a)$-block), so by Claim~1 it is a $\T(a)$-block. Hence, by Claim~2, attachment is a partial map $\tau\colon G\rightharpoonup G'$ which is surjective, and $|\pi(G)|\le|G|$ for every hub $G$. Since $\pi$ is a permutation of the hubs (Step~4), $\sum_G|\pi(G)|=\sum_G|G|$, so $|\pi(G)|=|G|$ for every hub, and $\tau$ is a bijection $G\to G'$. In particular $|G'|=2$ when $|G|=2$.

\emph{Claim 3: every component $\nu$ touches at most one $\T(a)$-block.} Suppose $\nu$ touches $A_k$ and $A_l$, so $|G|=2=|G'|$. Then $\nu$ contains $C_G$, and a shortest path in $\nu$ from a block meeting $A_k$ to a block meeting $A_l$ passes through $C_G$; its neighbour of $C_G$ on the side of $A_k$ lies in $\Gamma_k$ and meets $C_k$, so it is of $\N$-type with pair $(l,k)$ by (i), and by rigidity $A_k$ is of $\N$-type with pair $(l,k)$ (this includes the case that $A_k$ meets $C_G$ directly, by (i)). \emph{Case $e=k$.} The block $D_{k'}$ is attached to some $A_i$, and every attachment path to $D_{k'}$ ends at a contact block as in (ii), of pair $(k',k')$ or of $b$-type, lying in $\Gamma_k$. An attachment path from $A_k$: a direct contact would need $\lambda=l=k'$, impossible ($k'=k$ if $G'=G$, $k'\notin G$ otherwise); otherwise it starts with the pair $(l,k)$, and if it passes through $C_G$ both neighbours lie in $\Gamma_k$ and have the pair $(l,k)$; so by rigidity it cannot reach the contact block. A direct contact $A_l\cap D_{k'}$ is impossible ($A_l\subseteq\Gamma_l$), and an attachment path from $A_l\subseteq\Gamma_l$ to $D_{k'}\subseteq\Gamma_k$ passes through $C_G$ with its neighbour on the side of $D_{k'}$ in $\Gamma_k$, of pair $(l,k)$, and rigidity gives the same contradiction. So $D_{k'}$ is attached to no $A_i$, contradicting the consequences of Claims~1 and~2. \emph{Case $e=l$.} The block $D_{l'}$ is attached to some $A_i$, through a contact block of pair $(l',k')$ in $\Gamma_l$ (iii). A direct contact $A_k\cap D_{l'}$ is impossible ($A_k\subseteq\Gamma_k$), and an attachment path from $A_k\subseteq\Gamma_k$ passes through $C_G$, whose neighbour on the side of $D_{l'}$ lies in $\Gamma_l$ and meets $C_l$, so by (i) it is of $b$-type or has the pair $(l,l)$, whereas rigidity from the contact block gives it the pair $(l',k')$, $l'\ne k'$: impossible. So $D_{l'}$ is attached to $A_l$, and by rigidity (the neighbours of $C_G$ in $\Gamma_l$ have $\delta=l$ by (i), and $k'=l$ is impossible: it would give $l\in G'$, $G'=G$, $k'=k$) the attachment path avoids $C_G$ and $A_l$ is of $\N$-type with pair $(l',k')$; a direct contact $A_l\cap D_{l'}\ne\emptyset$ gives the same. But then $A_l$ is not touched by $\nu$: a block of $\nu$ meeting $A_l$ is $C_G$ (excluded, as $\lambda=l'$ and $\lambda=j=l$ would give $l\in G'$, $G'=G$, and then $\delta=k'=k$, so the contact would lie in $C_k\subseteq\Gamma_k$) or a common block of pair $(l',k')$, from which, by rigidity, no path in $\nu$ reaches a neighbour of $C_G$.

\emph{Conclusion.} Let $R$ be a class of $\rho$ contained in $\Gamma_G$; it contains some $A_i$ (Step~5). We claim that its $\T$-blocks are exactly $A_i$ and $D_{\tau(i)}$. Otherwise take a shortest path in $\Gamma(a,d)$ from $\{A_i,D_{\tau(i)}\}$ to another $\T$-block; its interior consists of $\N$-blocks. It cannot lead from $A_i$ to a $\T(d)$-block (this would be attached to $A_i$, hence equal to $D_{\tau(i)}$), nor from $A_i$ to a $\T(a)$-block $A_{i'}$ (Claim~3), nor from $D_{\tau(i)}$ to a $\T(a)$-block $A_{i'}$ ($D_{\tau(i)}$ would be attached to $A_{i'}$, so $i'=i$ as $\tau$ is injective), nor from $D_{\tau(i)}$ to a $\T(d)$-block (Claim~1). So every class of $\rho$ in $\Gamma_G$ contains exactly one $\T(a)$-block and one $\T(d)$-block.
\end{proof}

\begin{theorem}\label{thm:Q}
The condition \textup{(Q)} holds in every $\P_X$: if $a\bowtie c\bowtie b\bowtie d$, $xa=ax$, $\theta_c(x)=\theta_d(x)$, $ya=ay$, $yb=by$ and $\theta_y(c)=\theta_y(b)$, then $a\bowtie d$.
\end{theorem}

\begin{proof}
By Step~1 we may assume $x=u\le a$. By Step~5 it suffices that every class of $\rho=\varepsilon_a\vee\varepsilon_d$ contained in a class $\Gamma_G$ contains exactly one $\T(a)$-block and one $\T(d)$-block, which is Proposition~\ref{prop:Qnh} for the non-hub groups and Proposition~\ref{prop:Qh} for the hubs.
\end{proof}

\begin{corollary}\label{cor:Q}
\textup{(P$_2$)} and \textup{(S)} hold in every $\P_X$, the latter with a computer-free proof of the hub case of Theorem~\ref{thm:S}; and $\mathfrak P_8^4$ and $\mathfrak P_8^5$ are minimal non-embeddable projection algebras.
\end{corollary}

\begin{proof}
(Q)$\Rightarrow$(P$_2$)$\Rightarrow$(S) in every projection algebra (Proposition~\ref{prop:indep}(a)); the hypotheses of (S) are those of (Q) with $x=y=\theta_a(cd)$. The two algebras violate (P$_2$), resp.\ (Q) (Section~\ref{subsec:census8}), and all their proper subalgebras embed.
\end{proof}

\begin{remark}\label{rem:Q}
(1) Of the consequences of the hypotheses in projection algebras derived in Step~1, the proof uses only the normalisation $x=u\le a$: Step~2 rests on the transport lemma and Lemma~\ref{lem:PX}(a) applied to $\theta_c(u)=\theta_d(u)$ and $\theta_c(w)=\theta_d(w)$, and Step~3 on Lemma~\ref{lem:PX}(a),(b) applied to $ya=ay$, $yb=by$ and $\theta_y(c)=\theta_y(b)$ (the identities $yc=ybc$, $ycb=yb$ are the semigroup form of Step~3(c), and $cd=dc$, $ud=ucbd$ are not needed). The identities $u\le\theta_a(d)$ ($u*(d*a)=u$), $z\le\theta_d(a)$ ($(u*c)*(a*d)=u*c$) and $\theta_{ya}(d)=ya$ ($d*(a*y)=a*y$), which also hold in every projection algebra satisfying the hypotheses (Prover9, \texttt{pa8/q2/p9/}), were used in earlier versions of the argument but are not needed: the fact that every class of $\rho$ inside $\Gamma_G$ contains a $\T(a)$-block is immediate from $\Pi_1=\varepsilon_u\vee\varepsilon_d$ (Step~2). Whether $ya\le\theta_a(d)$ follows as well is open (no countermodel of order $\le14$). The identity $\theta_y(d)=\theta_y(b)$ holds in every instance in $\P_X$ but is not a consequence of the hypotheses in projection algebras, not even together with (S) (Mace4 finds $6$-element countermodels); it is not used.
(2) The hub-free case ($\supp c=\supp d$, Proposition~\ref{prop:Qnh} alone) contains the case (F1) of the proof of Theorem~\ref{thm:S} (with $x=y=\theta_a(cd)$), for which Step~5 of that proof is the special case of (a), (b) in which all blocks of $a$ are of $b$-type; the hub case of Theorem~\ref{thm:S} is the case of Proposition~\ref{prop:Qh} in which $\pi$ is the identity and $A_G$ is a bridge or $A_G=B_G$. There the bound $|G|\le2$ was proved as the bound $p\le2$ by a different argument and the rest was a finite computer check (\texttt{hubtypes.py}); the cross wiring of Step~2 and the rigidity of Step~5 replace it.
(3) The hypothesis on $y$ enters only through the classification of the blocks of $a$ and $b$ into $\N(y)$-blocks and unions of $\T(y)$-blocks and through the description of the classes of $\varepsilon_y\vee\varepsilon_c$ meeting $\supp y$ (Step~3); the hypothesis on $x$ only through Step~2. Several further statements were verified on the instances in $\P(\P_m)$, $m\le6$, before the proof was found and are now consequences of it or of its steps: every block of $a$ meets at most one $\T(d)$-block, every block of $d$ (and every hub) meets at most one $\T(a)$-block; replacing $d$ by the projection $d_1$ with $\N$-blocks $\N(c)$ and $\T$-blocks $D_k$ ($k$ in non-hub groups) and $C_i$ ($i$ in hubs) gives instances $(a,c,b,d_1)$ and $(a,d_1,b,d)$ of the hypotheses (the second requires $\theta_y(d_1)=\theta_y(b)$, which we have not proved by hand); the hubs of size $2$ have $\pi(G)=G$, exactly one of $A_k,A_l$ meets $C_G$, namely the one of $\N$-type with pair $(l,k)$ found in Claim~3, and in $\P(\P_m)$, $m\le6$, the other one meets its own $D$. The example of Step~2 shows that $\pi$ can exchange singleton hubs.
\end{remark}

\subsection{Independence of the conditions}\label{subsec:indep}

We now have the following list of conditions, all of which hold in every $\P_X$ or are conjectured to: \eqref{eq:Q1}, the triangle condition (Tri), (S), (R), and the seven of Proposition~\ref{prop:census8}, (R$'$), (V), (P$_1$), (P$_2$), (Q), and the two of (e), which we call (E$_6$) (the one for $\mathfrak P_8^6$) and (E$_1$) (the one for $\mathfrak P_8^1$). Some of them are consequences of others.

\begin{proposition}\label{prop:indep}
\begin{enumerate}[label=\textup{(\alph*)},nosep]
\item In every projection algebra, (Q) implies (P$_2$), and (P$_2$) implies (S); hence (Q) implies (S). In every projection algebra, (R$'$) implies (R) (Proposition~\ref{prop:census8}(a)).
\item Each of the eight conditions \eqref{eq:Q1}, Tri, (R$'$), (V), (P$_1$), (Q), (E$_6$), (E$_1$) fails in a finite projection algebra satisfying the other ten conditions of the list: $Q(K_2,\mathrm{const})$ for \eqref{eq:Q1} (order $3$), $\P(Q_0)$ for Tri (order $11$), $\mathfrak P_8^3$ for (R$'$), $\mathfrak P_8^7$ for (V), $\mathfrak P_8^2$ for (P$_1$), $\mathfrak P_8^5$ for (Q), $\mathfrak P_8^6$ for (E$_6$), $\mathfrak P_8^1$ for (E$_1$). Consequently these eight conditions are independent: none of them follows from the conjunction of the other seven together with (S), (R), (P$_2$); in particular, for any two of them, neither implies the other even in the presence of \eqref{eq:Q1}, (S) and (R). The conditions (S), (R), (P$_2$) are redundant in the list, by (a).
\item The witnesses of the redundant conditions: $\mathfrak P$ violates (S), (P$_2$), (Q) and (E$_6$) and satisfies the rest; $\mathfrak P_9$ violates (S), (P$_2$), (Q); $\mathfrak P_7$ violates (R) and (R$'$); $\mathfrak P_8^4$ violates (P$_2$) and (Q). Mace4 finds a projection algebra of order $9$ (and none smaller) satisfying \eqref{eq:Q1} and (E$_6$) but violating (S), (P$_2$) and (Q); so (E$_6$) does not imply (S) either. Altogether, for every ordered pair $X\ne Y$ of conditions in the list, $X$ implies $Y$ (in the presence of \eqref{eq:Q1}) exactly in the cases of (a), (Q)$\Rightarrow$(P$_2$), (Q)$\Rightarrow$(S), (P$_2$)$\Rightarrow$(S), (R$'$)$\Rightarrow$(R); every other pair is separated by one of the algebras named in (b) and (c) or by this Mace4 model. The seven candidates of order $9$ (Remark~\ref{rem:census8}(3)) satisfy all eleven conditions.
\end{enumerate}
\end{proposition}

\begin{proof}
(a) Putting $y=x$ in (Q) gives exactly (P$_2$). For (P$_2$)$\Rightarrow$(S), let $a\bowtie c$, $b\bowtie c$, $b\bowtie d$, $cd=dc=:u$ and $\theta_a(u)=\theta_b(u)$ be the hypotheses of (S), and apply (P$_2$) to the path $p_0=d$, $p_1=b$, $p_2=c$, $p_3=a$ and $x=u$: $u\le d$ and $u\le c$, so $u$ commutes with $p_0$ and $p_2$; $\theta_u(b)=ubu=u\,cbc\,u=ucu=u=\theta_u(c)$ since $u=uc=cu$ and $cbc=c$; and $\theta_{p_1}(u)=\theta_b(u)=\theta_a(u)=\theta_{p_3}(u)$. The conclusion $p_0\bowtie p_3$ is $a\bowtie d$.  (b) The violated sets were computed with \texttt{pa8/indep.py} (which also contains the Mace4 inputs used in (c)); the last assertion is immediate: a condition of the eight that followed from the others would hold in the witness listed for it. (c) By computation.
\end{proof}

Thus, as far as is currently known, the obstructions to embeddability form the list \eqref{eq:Q1}, Tri, (R$'$), (V), (P$_1$), (Q), (E$_6$), (E$_1$), pairwise and jointly independent, with (S)$\subset$(P$_2$)$\subset$(Q) (that is, (Q) is the strongest of the three) and (R)$\subset$(R$'$). Note that (V) neither implies nor is implied by any of the others, despite its simple form; and that among the eight only \eqref{eq:Q1}, Tri, (R$'$), (V), (P$_1$) and (Q) have been proved to hold in every $\P_X$ (Proposition~\ref{prop:Q1}, Theorems~\ref{thm:hex}, \ref{thm:Rp}, \ref{thm:V}, \ref{thm:P1}, \ref{thm:Q}), together with the redundant (S), (P$_2$) and (R) (Theorems~\ref{thm:S}, \ref{thm:R}; (R) is now also a corollary of (R$'$), and (S), (P$_2$) of (Q)).

\subsection{The order-9 candidates: (R$''$) and the family of 2-path conditions}\label{subsec:nine}

This subsection records what was learnt (October 2026) from the seven candidate algebras of order $9$ (Remark~\ref{rem:census8}(3); \texttt{pa8/cand9.json}), none of which violates a condition of the list of Section~\ref{subsec:indep}. The evidence level of everything stated here is ``verified in $\P(\P_m)$ for $m\le6$ by SAT, variables not assumed distinct, and satisfied by all $92\,647$ embeddable census algebras of order $\le9$''; nothing in this subsection is proved for all $\P_X$. Code: \texttt{pa8/n4/}.

\smallskip
\emph{Length-$4$ path closing: negative.} The natural guess that the path-closing family continues with paths of length $4$ was tested first. On $16\,000$ sampled paths $p_0\bowtie\dots\bowtie p_4$ in $\P(\P_5)$ ($10^8$ negative instances, two auxiliaries, the $78$-atom language of \texttt{pathclose3.py} without the $z=\theta_p(q)$ atoms, coincidences of variables allowed; hitting sets by a C program, \texttt{hit4.c}), the $84$ minimal hypothesis sets of size $\le4$ that force $p_0\bowtie p_4$ on the sample all contain a chord $p_i\bowtie p_j$, except five, which fail in $\P(\P_4)$ as soon as variables may coincide. Of the seven candidates only one (candidate~$1$) has a chordless unclosed path of length $4$, and no two-auxiliary condition valid in $\P(\P_5)$ explains it. So the obstructions do not grow with the path length, at least up to order $9$ and two auxiliaries.

\smallskip
\emph{Unclosed 2-paths.} Six of the seven candidates have an unclosed path of length $2$, $a\bowtie c\bowtie b$ with $a\not\bowtie b$ (candidate $2$ has none; its extracted identities have an order conclusion $a\le b$, as (V)). Candidate $6$ is the graph semigroup $a-d-b$ over the $2\times2$ rectangular band $\{c,e,f,g\}$ ($a$ over $f$, $b$ over $e$, $d$ over $g$) with two further elements $h,i$: $hd=dh$, $ib=bi$, $\theta_h(a)=\theta_h(d)=g$, $\theta_h(b)=c=\theta_i(a)=\theta_i(d)$, $\theta_i(b)=e$. Its greedily minimised identity (eleven atoms, nine variables) is this configuration written out through restrictions, and it reduces to a condition of exactly the shape of (R$'$):
\begin{equation}\label{eq:Rpp}
\text{(R$''$)}\qquad a\bowtie c,\ b\bowtie c,\ xc=cx,\ \theta_x(a)=\theta_x(c),\ yb=by,\ \theta_x(b)=\theta_y(a)\ \Longrightarrow\ a\bowtie b .
\end{equation}
Compare (R$'$): $xc=cx$, $\theta_x(a)=\theta_x(c)$, $ya=ay$, $\theta_x(b)=\theta_y(b)$. (R$''$) is valid in $\P(\P_m)$ for $m\le6$, holds in every embeddable census algebra of order $\le9$, is violated by candidate~$6$ (with $x=h$, $y=i$) and by no other candidate of order $8$ or $9$; and it is independent of (R$'$): $\mathfrak P_8^3$ violates (R$'$) but not (R$''$), candidate~$6$ the other way round. The two remaining combinations of the same shape, $ya=ay$ with $\theta_x(b)=\theta_y(a)$ (equivalently $\theta_x(b)\le a$) and $yb=by$ with $\theta_x(b)=\theta_y(b)$ (equivalently $\theta_x(b)\le b$), are valid for $m\le5$ but have no witness of order $\le9$; the combinations with $y$ commuting with $c$ fail in $\P(\P_4)$. Why none of the earlier searches found (R$''$): the atom language of \texttt{pathclose3.py} (Section~\ref{subsec:pathclose}, Remark~\ref{rem:census8}(5)) contains $\theta_x(p_i)=\theta_y(p_i)$ but not the cross atoms $\theta_x(p_i)=\theta_y(p_j)$, $i\ne j$. With these added (\texttt{pathclose4b.py}), a sampled search in $\P(\P_5)$ at length $2$ finds (R$''$) and its mirror image as the only hypothesis sets of size $\le4$ explaining an order-$9$ algebra, among about forty cross-atom sets valid for $m\le5$ (whose status at $m=6$ is being checked, \texttt{verify6\_2.py}).

\smallskip
\emph{The other candidates.} For candidates $0,1,3,4,5$ every assignment of two auxiliaries in the extended language is realised by a negative instance in $\P(\P_5)$ (\texttt{explain2.py}), so they need three auxiliaries or a richer language. The right language turned out to be the equalities among the elements and all terms $\theta_p(q)$ over the path $a,c,b$ and three auxiliaries $x,y,z$: a hypothesis is the partition of this set of $36$ symbols induced by an assignment in the candidate, weakened greedily by splitting symbols off under the SAT test at $m=4$ (\texttt{explain3b.py}). The terms that matter are $u=\theta_a(b)=aba\le a$ and $v=\theta_b(a)=bab\le b$, the ``would-be'' matching restrictions along the unclosed path; the conclusion $a\bowtie b$ is $u=a$ and $v=b$. The smallest explanation found for each unclosed path (\texttt{catalog.log}; all valid for $m\le6$, none violated by an embeddable census algebra; ``$aya=a$'' is the half-matching $\theta_a(y)=a$):
\begin{center}\small
\begin{tabular}{@{}llll@{}}
\toprule
& hypotheses besides $a\bowtie c\bowtie b$ & conclusion & violated by\\
\midrule
K0 & $yc=cy$, $\theta_y(b)=v$, $bzb=b$, $\theta_z(a)=\theta_z(b)$, $\theta_y(z)=u$ & $a\bowtie b$ & cand.\ 0, $\mathfrak P_8^1$ (proved, Thm~\ref{thm:K0})\\
K1a & $aya=a$, $\theta_y(b)=\theta_z(x)$, $\theta_x(b)=u=v$, $\theta_c(z)=\theta_x(z)$ & $a\bowtie b$ & cand.\ 1, $\mathfrak P_8^2$\\
K1b & $aya=a$, $\theta_z(c)=b$, $\theta_y(a)=\theta_y(x)$, $\theta_x(c)=\theta_z(x)=u$, $\theta_y(z)=v$ & $a\bowtie b$ & cand.\ 1\\
K1c & $bxb=b$, $\theta_y(a)=\theta_z(y)$, $\theta_c(z)=\theta_x(a)=\theta_y(x)$, $\theta_z(x)=v$ & $a\bowtie b$ & cand.\ 1, $\mathfrak P_8^2$\\
K3a & $\theta_z(c)=a$, $b\le y$, $\theta_x(c)=\theta_x(y)$, $\theta_y(b)=\theta_y(x)$, $\theta_x(b)=u=v$, $yz=zy$ & $a\bowtie b$ & cand.\ 3\\
K3b & $\theta_x(c)=a$, $b\le y$, $b\le z$, $\theta_y(c)=\theta_y(b)$, $\theta_z(c)=\theta_z(b)$, $\theta_y(z)=\theta_z(y)=u$, $\theta_x(b)=v$ & $a\bowtie b$ & cand.\ 3\\
K3c & $bxb=b$, $\theta_x(a)=\theta_x(b)$, $\theta_x(y)=\theta_y(z)$, $\theta_z(y)=u=v$, $zc=cz$ & $a\bowtie b$ & cand.\ 3, $\mathfrak P_8^1$\\
K4a & $aya=a$, $b\le z$, $yc=cy$, $\theta_z(y)=v$, $\theta_z(a)=u$ & $a\bowtie b$ & cand.\ 4, $\mathfrak P_8^6$ (proved, Thm~\ref{thm:K4})\\
K4b & $aza=a$, $b\le y$, $zc=cz$, $\theta_z(y)=v$, $\theta_y(a)=u$ & $a\bowtie b$ & cand.\ 4\\
K5 & $\theta_x(c)=a$, $bxb=b$, $zc=cz$, $\theta_z(x)=u$, $\theta_z(b)=v$ & $a\bowtie b$ & cand.\ 5\\
K6 & $\theta_y(z)=\theta_z(a)$, $\theta_c(z)=\theta_y(a)$, $\theta_z(b)=v$ & $a\bowtie b$ & cand.\ 6\\
(R$''$) & as in \eqref{eq:Rpp} & $a\bowtie b$ & cand.\ 6\\
\bottomrule
\end{tabular}
\end{center}
Thus $\mathfrak P_8^1$ (the witness of (E$_1$)), $\mathfrak P_8^6$ (of (E$_6$)) and $\mathfrak P_8^2$ (of (P$_1$)) also violate 2-path conditions: the conditions (E$_1$), (E$_6$), whose shape looked unrelated to the path-closing ones, are path-closing conditions written in a language that cannot name $u$ and $v$. Whether $\mathfrak P_8^4$, $\mathfrak P_8^5$ and the order-$9$ witness of (R$'$) violate 2-path conditions as well has not been checked. The members of the table are different conditions (each explains a different witness), and no master condition in $a,b,c,u,v$ alone exists: the conjunction of all equalities among these five elements and their $\theta$-terms common to the six instances does not force $a\bowtie b$ even in $\P(\P_4)$.

\smallskip
\emph{The friendly relation.} The relation $\bowtie$ (``friendly'' in East--Azeef Muhammed) is reflexive and symmetric but not transitive, and every path-closing condition of the list is a statement of local transitivity: $a\bowtie c\bowtie b$ closes under side hypotheses. Some facts about it, the first two verified in $\P(\P_m)$ for $m\le5$ and the others valid in every regular $*$-semigroup:
\begin{enumerate}[nosep,label=(\alph*)]
\item $\bowtie$ is not compatible with the operation $\theta$ in either argument (so it is not a tolerance \emph{of the algebra} in the sense of Chajda and Zelinka): in $\P(\P_4)$, $p\bowtie p'$ fails to imply $\theta_q(p)\bowtie\theta_q(p')$ in $24\,816$ of $142\,128$ instances and $\theta_p(q)\bowtie\theta_{p'}(q)$ in $23\,304$.
\item Its transitive closure is ``same rank'', i.e.\ $\D$ restricted to projections ($\P_3$--$\P_5$), which is not a congruence of $\P(\P_X)$ either.
\item For $p\bowtie p'$ the maps $\theta_{p'}\colon[0,p]\to[0,p']$ and $\theta_p\colon[0,p']\to[0,p]$ are inverse isomorphisms (for $x\le p$, $p(p'xp')p=(pp'p)x(pp'p)=pxp=x$, and symmetrically), and they preserve $\bowtie$: if $x\le p$ then $x\bowtie p'xp'$. Indeed $x=px=xp$ gives $xp'x=x(pp'p)x=xpx=x$, hence $x(p'xp')x=(xp'x)p'x=x$ and $(p'xp')x(p'xp')=p'(xp'x)p'xp'=p'xp'$. The isomorphism $\theta_{p'}$ also carries $\bowtie$ inside $[0,p]$ to $\bowtie$ inside $[0,p']$: for $x,y\le p$ we have $xp'y=xpp'py=xy$, so $(p'xp')(p'yp')(p'xp')=p'(xyx)p'$, and $x\bowtie y$ gives $(p'xp')\bowtie(p'yp')$. This is part of the data of the \emph{chain groupoid} $\mathcal C(P)$ (paths of friendly pairs) through which East and Azeef Muhammed (arXiv:2301.04845) reconstruct a regular $*$-semigroup from its projection algebra and a \emph{chained projection groupoid}; we have not yet exploited their framework, in which the obstructions found here should be statements about which chained projection groupoids over a given $P$ arise inside $\P_X$.
\item The down-sets form a presheaf on $(\P(S),\le)$: for $w\le x\le p$ and $y\le p$, $w(xyx)w=wyw$, so $y\mapsto xyx$ is a restriction map $[0,p]\to[0,x]$ compatible with composition; and the transports of (c) commute with these restrictions (for $x,y\le p$ and $p\bowtie q$: $q(xyx)q=(qxq)(qyq)(qxq)$, using $xqy=xpqpy=xy$). In $\P(\P_X)$ the values of this presheaf are concrete: $[0,p]\cong\P(\P_{T(p)})$, a restriction of $p$ being a labelled partition of the set $T(p)$ of transversal blocks (Lemma~\ref{lem:PX}(a)), and the transport along $p\bowtie q$ is induced by the bijection $T(p)\to T(q)$ of Lemma~\ref{lem:PX}(c).
\item Pairwise matched projections of $\P_X$ generate a subsemigroup all of whose elements have the same rank: a product along a chain of matched projections keeps the upper blocks of its first factor and the lower blocks of its last; hence $tt^*=e$ for the first factor $e$, the subsemigroup is simple, and being finite it is completely simple. So the $\bowtie$-cliques play the role of the rectangular $\D$-classes of a band; but they overlap, and the obstructions are about how overlapping cliques may be glued.
\end{enumerate}
In this language the embedding problem is a descent problem: an abstract projection algebra comes with the presheaf of (d) and the transports of (c); the assignment $Y\mapsto\P(\P_Y)$ is a functor of finite sets (bijections act, and for arbitrary $p,q$ the map $\theta_q\colon[0,p]\to[0,q]$ is induced by the correspondence $T(p)\to(\text{classes of }\varepsilon_p\vee\varepsilon_q)\leftarrow T(q)$), and representability asks that the abstract data be induced through this functor from a system $p\mapsto T(p)$ of finite sets. The hypotheses of all our conditions are statements that certain transports of certain restrictions agree (the index pairs of the proof of Theorem~\ref{thm:Q} are these transports), and the conclusions are that an edge of $\bowtie$ exists; holonomy around closed paths of matchings is non-trivial in $\P_X$ itself (a permutation of $T(a)$; e.g.\ $ad\ne acbd$ in Remark~\ref{rem:census8}(5)), so the conditions are not flatness but agreement between the holonomy and the non-matched $\theta$'s. We know no literature in which projection algebras are set up this way.

\subsection{The commutation--identification language and the chain family (W$_n$)}\label{subsec:chain}

This subsection records the reconnaissance of (R$''$) made after the cold read of October~10 (the files are in \texttt{pa8/rpp/}). As in Section~\ref{subsec:nine}, nothing here is proved unless said so; ``verified'' means verified in $\P(\P_m)$ by SAT for the stated $m$, with variables not assumed distinct, and ``on all instances'' means on all assignments of the hypotheses in $\P(\P_4)$ ($89\,696$ for (R$''$)) and on $9\,482$ sampled instances in $\P(\P_5)$.

\smallskip
\emph{The family around (R$'$).} Fix the base hypotheses $a\bowtie c$, $b\bowtie c$, $xc=cx$, $\theta_x(a)=\theta_x(c)$ and write $r=\theta_x(b)$. Of the nine conditions ``$y$ commutes with $p$ and $r=\theta_y(q)$'' ($p,q\in\{a,b,c\}$) exactly the four with $p,q\in\{a,b\}$ imply $a\bowtie b$ (verified $m\le5$), and two of them are not new: $(p,q)=(b,b)$ says $r=\theta_y(b)=yb\le b$, i.e.\ $\theta_b(\theta_x(b))=\theta_x(b)$, which by the axiom $\theta_b(\theta_x(b))=\theta_b(x)$ is $xb=bx$, so this is (R) with $a$ and $b$ exchanged; and $(p,q)=(a,a)$ says $r=ya\le a$, and from $r\le a$ (together with \eqref{eq:Q1}) Prover9 derives $a=c$. The remaining two are (R$'$) ($(a,b)$) and (R$''$) ($(b,a)$). Likewise the order variants $r\le a$ (forces $a=c$) and $r\le b$ (is (R)) add nothing. So the family of the base hypotheses plus one further atom about $b$ is exactly $\{$(R), (R$'$), (R$''$)$\}$.

\smallskip
\emph{Algebraic consequences of (R$''$).} Put $v=xc=cx$, $w=yb=by$, $u=\theta_a(x)$, $m=xy$. On all instances, and verified for $m\le6$ (the first also for $m=7$):
\[
\theta_y(c)=\theta_y(a)=r,\qquad \theta_c(y)=\theta_c(x)=v,\qquad \theta_b(y)=\theta_b(x)=w,\qquad \theta_a(y)=\theta_a(x)=u,\qquad xy=yx .
\]
Prover9 (binary notation, the axioms and the clause \eqref{eq:Q1}) proves $\theta_a(y)=\theta_a(x)$ (via $\theta_a(r)=\theta_a(y)$ and $\theta_a(r)=a(xbx)a=a(vbv)a=ava$), $xy=yx$ (as $r\le x,y$), $\theta_y(v)=r$ and $\theta_x(w)=r$, but finds no proof of the first three identities in $300$\,s, with or without the others as extra assumptions; we expect them to need the geometry of $\P_X$. Further consequences: $u\le a$, $v\le c$, $w\le b$ with $w=\theta_b(v)$, $u=\theta_a(v)$, so $u\bowtie v\bowtie w$, and also $u\bowtie w$ (on all instances); $\theta_a(r)=u$, $\theta_c(r)=v$, $\theta_b(r)=w$; and $m$ sees $a$, $b$, $c$ alike: $\theta_m(a)=\theta_m(b)=\theta_m(c)=r$ (as $m\le x$, $m\le y$, $r\le m$). The special case $x=c$ of (R$''$) ($a\bowtie c\bowtie b$, $yb=by$, $\theta_y(a)=c$) is a theorem of projection algebras with \eqref{eq:Q1} (Prover9, $23$ steps). Since $v$ is a common lower bound of $x$ and $c$, any projection $x'\le y$ commuting with $c$ would make $y$ and $c$ commute by \eqref{eq:Q1} and reduce (R$''$) to (R); no such $x'$ is available ($v\not\le y$ in general).

\smallskip
\emph{The commutation--identification language.} The first three identities say that $y$, like $x$, \emph{identifies $a$ with $c$}, and suggest the language whose atoms are commutations $p\sim q$ ($pq=qp$) and identifications $\theta_p(q)=\theta_p(q')$, $q,q'\in\{a,b,c\}$, for auxiliary projections $p$. Enumerating all hypothesis sets with at most two auxiliaries $x,y$ ($13$ atoms, \texttt{ci2.py}) and testing validity at $m=4$ gives the following complete list of minimal valid sets (up to $x\leftrightarrow y$, and discarding the sets whose hypotheses force $a=c$ or $b=c$): the one-auxiliary sets $x\sim a$, $x\sim c$, $\theta_x(b)=\theta_x(c)$ (this is (R)) and $x\sim b$, $x\sim c$, $\theta_x(a)=\theta_x(c)$ (its mirror image (R)$'$); and the two-auxiliary sets
\[
c\sim x\sim y\sim e,\qquad \theta_x(q_1)=\theta_x(c),\qquad \theta_y(q_2)=\theta_y(c),
\]
with $e\in\{a,b\}$ and $(q_1,q_2)\in\{a,b\}^2$, except that for $e=b$ the pattern $(b,b)$ and for $e=a$ the pattern $(a,a)$ are degenerate. All of these are valid for $m\le5$. The pattern $e=b$, $(q_1,q_2)=(a,a)$ is
\begin{equation}\label{eq:W2}
\text{(W$_2$)}\qquad a\bowtie c,\ b\bowtie c,\ xc=cx,\ yb=by,\ xy=yx,\ \theta_x(a)=\theta_x(c),\ \theta_y(a)=\theta_y(c)\ \Longrightarrow\ a\bowtie b ,
\end{equation}
and by the identities above the hypotheses of (R$''$) imply those of (W$_2$) (in $\P(\P_m)$ for $m\le7$, since the identity $\theta_y(a)=\theta_y(c)$ used is verified, not proved; \emph{added later on 11 Oct:} it is a consequence of \eqref{eq:Q1} and (W$_2$), Theorem~\ref{thm:Rpp}, so (R$''$) is now proved in every $\P_X$); conversely (W$_2$) is violated by candidate~$6$ of order $9$ (with the same assignment as (R$''$)) and by no other non-embeddable algebra of order $\le9$ in the census, and Mace4 finds no projection algebra of order $\le8$ satisfying \eqref{eq:Q1} and the conditions of the status table that separates (W$_2$) from (R$''$) in either direction (Prover9 derives neither from the other in $900$\,s). We regard (W$_2$) as the right form of the order-$9$ condition: all its hypotheses are commutations and identifications, and (R)$'$ is the case $x=y$. Each hypothesis of (W$_2$) is needed (dropping any one of the five atoms gives a condition that fails at $m=4$).

\smallskip
\emph{Chains.} (W$_2$) generalises at once to chains of commuting auxiliaries:
\[
\text{(W$_n$)}\qquad a\bowtie c,\ b\bowtie c,\quad c\sim x_1\sim x_2\sim\dots\sim x_n\sim b,\quad \theta_{x_i}(a)=\theta_{x_i}(c)\ (1\le i\le n)\ \Longrightarrow\ a\bowtie b ,
\]
with (W$_1$) $=$ (R)$'$. \emph{(Added later on 11 Oct: (W$_n$) is proved for every $n$ in Section~\ref{subsec:count}, Theorem~\ref{thm:Wn}; the evidence recorded in this and the next subsections is superseded by that proof.)} (W$_2$), (W$_3$) and (W$_4$) are valid for $m\le6$ (\texttt{wn\_sat.py} for $m\le5$; the $m=6$ checks were run by the cold reader of 11~Oct with an independent harness, \texttt{review2/review-scripts/wsat.py}: $131$--$163$\,s, $360$--$392$\,s and $556$--$681$\,s for the two halves of the conclusion). Note that in $\P(\P_5)$ every chain end is within three links of $c$ (Section~\ref{subsec:joinstat}(g)), so that the $m\le5$ checks of (W$_4$) carry no evidence beyond (W$_3$); the $m=6$ check of (W$_4$) is the first one that does. More generally, for a word $q_1\cdots q_n\in\{a,b\}^n$ let $(q_1\cdots q_n)$ denote the chain condition in which $x_i$ identifies $q_i$ with $c$; at $m=4$ the valid words with $n\le4$ are exactly the words $a^ib^{n-i}$ with $1\le i\le n$ (the word $b^n$ is degenerate), $a^{n-2}ba$ and $b^{n-1}a$ (for $n=3$: $aaa$, $aab$, $abb$, $aba$, $bba$ valid, $baa$, $bab$ invalid; for $n=4$: $aaaa$, $aaab$, $aabb$, $abbb$, $aaba$, $bbba$ valid, the other nine invalid; the valid words with $n=3$ are also valid at $m=5$). We have no explanation of this pattern yet, and after Section~\ref{subsec:joinstat}(c),(f) a pattern established at $m=4$ in this language should not be trusted before it is confirmed at $m=6$. One direction between consecutive members is trivial: (W$_{n+1}$) implies (W$_n$) in every projection algebra, since an instance of (W$_n$) becomes an instance of (W$_{n+1}$) by repeating its last auxiliary ($x_{n+1}=x_n$); in particular every (W$_n$) implies (W$_1$) $=$ (R)$'$ (take all $x_i$ equal). The question that decides whether (W$_n$) is a genuinely infinite family is therefore whether the converse fails, i.e.\ whether some projection algebra satisfies (W$_n$) (and \eqref{eq:Q1} and the other known conditions) but not (W$_{n+1}$): Mace4 finds no algebra of order $\le8$ satisfying \eqref{eq:Q1}, the conditions of the table and (W$_2$) but violating (W$_3$) (order $8$ exhausted in a few minutes; order $9$ is running), and of the $263$ non-embeddable algebras of order $8$ and $9$ in the census only candidate~$6$ violates any chain condition (it violates (W$_2$), (W$_3$), $(aab)$ and $(aba)$, and satisfies $(ab)$, $(ba)$, $(abb)$, $(bba)$); so a witness, if it exists, has order $\ge10$, and the natural way to look for one is a targeted Mace4 search (\texttt{w3ind.in}) or a targeted extension of candidate~$6$. This is the first explicit candidate for the parametrised family conjectured in Section~\ref{sec:general}; note that the chain conditions, unlike (R$'$), (P$_1$), (Q), use no cross atoms $\theta_x(p)=\theta_y(q)$ and no terms, so that the language of Section~\ref{subsec:nine} would not have produced them directly.

\smallskip
\emph{Two further checks.} (1) The remaining witnesses of path-closing conditions are also 2-path phenomena: \texttt{explain3c.py} runs the three-auxiliary explanation search of Section~\ref{subsec:nine} on every unclosed 2-path of a given algebra. $\mathfrak P_8^4$ (the (P$_2$)-witness) has two unclosed 2-paths, with $14$ and $16$ minimal explanations valid at $m=5$; $\mathfrak P_8^5$ (the (Q)-witness) has two, with $18$ each (e.g.\ $\theta_x(y)=b$, $\theta_x(z)=\theta_z(y)=\theta_a(b)$, $\theta_c(z)=\theta_y(a)$ closes the path $a\bowtie c\bowtie b$ of $\mathfrak P_8^5$). By contrast $\mathfrak P_8^7$ (the (V)-witness) has no unclosed 2-path at all, and $\P(Q_0)$ (the triangle witness) has six, none of which has an explanation valid at $m=5$ (six minimal partitions each, all failing at $m=5$). So every known witness of a path-closing condition of order $\le9$ violates a 2-path condition with at most three auxiliaries, and the two witnesses outside the family do not. (2) All $42$ two-auxiliary 2-path conditions of Section~\ref{subsec:nine} that are valid for $m\le5$ remain valid at $m=6$ (\texttt{verify6\_2b.py}); the $m=7$ check of (R$''$) is running.

\smallskip
\emph{Towards a proof of (W$_2$).} The block-level reconnaissance of (R$''$) and of (W$_2$) (\texttt{harness*.py}, on all instances in $\P(\P_4)$ and samples in $\P(\P_5)$) shows the following structure, in the notation of the proof of Theorem~\ref{thm:Rp} (classes $\Gamma_i$, $\Delta_i$ of $C_i$ in $\varepsilon_a\vee\varepsilon_c$, $\varepsilon_b\vee\varepsilon_c$; $S_x=\supp x$, $S_y=\supp y$; $V_k$, $W_l$, $U_j$ the $\T$-blocks of $v$, $w$, $u$). (a) The $\N(c)$-blocks in $\Gamma_i$ are $\N(x)$-, $\N(y)$- and $\N(b)$-blocks (hence classes of $\varepsilon_b\vee\varepsilon_c$ by themselves). (b) The $\N(c)$-blocks in $\Delta_i$ do not meet $\supp a\cup\supp u\cup\supp v$, lie inside $S_x$ or inside $S_y$, and lie in classes of $\varepsilon_a\vee\varepsilon_c$ without $\T$-blocks. (c) An index $i$ is \emph{visible} if $C_i\subseteq S_x$; equivalently $A_i\subseteq\supp u$, $B_i\subseteq S_y$, $B_i\subseteq\supp w$; for invisible $i$ the classes $\Gamma_i$ and $\Delta_i$ do not meet $S_x\cup S_y$, and the $\N(x)$-block containing $C_i$ is an $\N(y)$-block, so that the class of $C_i$ in $\varepsilon_a\vee\varepsilon_b\vee\varepsilon_c$ is $\Gamma_i\cup\Delta_i$. (d) A class of $\varepsilon_a\vee\varepsilon_c$ without $\T$-blocks meets at most one $\Delta_i$ and at most one $W_l$; a class of $\varepsilon_b\vee\varepsilon_c$ without $\T$-blocks meets at most one $\Gamma_i$ and at most one $U_j$, and does not meet $\supp v\cup\supp w\cup\supp r$; the $W_l$ and $U_j$ so met are matched through $u\bowtie w$. (e) $\Gamma_i$ lies in a single class of $\varepsilon_y\vee\varepsilon_a$; every $\T(r)$-block meets at most one $V_k$, one $W_l$, one $U_j$, and conversely. The goal, as for (R$'$), is that no class of $\varepsilon_a\vee\varepsilon_b\vee\varepsilon_c$ contains two $\T(c)$-blocks (true on all instances; no coarser join $\varepsilon_c\vee\varepsilon_b\vee\varepsilon_p$ or $\varepsilon_c\vee\varepsilon_a\vee\varepsilon_p$, $p\in\{x,y,r,m\}$, has this property), followed by a counting argument. By (a)--(d) a shortest alternating chain of $c$-blocks from $C_p$ to $C_q$ leaves $\Delta_p$ into a $\T$-free class of $\varepsilon_a\vee\varepsilon_c$, then alternates between $\T$-free classes of the two joins, and enters $\Delta_q$ from a $\T$-free class of $\varepsilon_a\vee\varepsilon_c$; what is missing is an invariant carried along this chain (the labels $W_l$, $U_j$ of (d) are not always present), i.e.\ the analogue of Steps~4--6 of Theorem~\ref{thm:Rp}. The ``kinds and transitions'' abstraction (kinds of $c$-blocks by their relation to $x$, $y$, $a$, $b$ and the $\T$-content of their two classes; \texttt{kinds2.py}, \texttt{walk.py}) does not by itself exclude a chain of length four, so the invariant must involve more than these kinds.

Three further facts, all verified on instances (and the first also by SAT for $m\le6$, with the minimal hypotheses stated), point to the shape of the argument. (f) \emph{Lemma A (candidate; verified by SAT for $m\le7$).} If $a\bowtie c$, $xc=cx$, $\theta_x(a)=\theta_x(c)$, $xy=yx$ and $\theta_y(a)=\theta_y(c)$, then every block of $c$ lies in a single class of $\varepsilon_y\vee\varepsilon_a$; each of the four side hypotheses is needed (dropping any one gives counterexamples in $\P(\P_4)$ with $a,c,x,y$ all of positive rank), and $b$ plays no role (\texttt{lemA\_sat.py}). The statement is trivial when $x\le y$ (blocks of $c$ inside $\supp x$ are then unions of $\T(y)$-blocks) and when $y\le x$ (then $\varepsilon_x\le\varepsilon_y$ and one uses $\varepsilon_c\le\varepsilon_x\vee\varepsilon_a$); the content is in the incomparable case, where an $\N(y)$-block inside $\supp x$ meeting a block $E\subseteq\supp x$ of $c$ has to be joined to $E\cap\supp y$ through blocks of $a$ and $y$, and we do not yet see why. Together with $\varepsilon_c\le\varepsilon_x\vee\varepsilon_a$ (Step~5 of Theorem~\ref{thm:Rp}) it says that $a$ and $c$ are tied together with respect to both auxiliaries. (g) Along a chain in $\varepsilon_a\vee\varepsilon_b$, two consecutive blocks of $b$ meeting $\supp y$ (all intermediate blocks of $b$ lying outside $\supp y$) lie in the same class of $\varepsilon_b\vee\varepsilon_c$; the same holds for blocks of $b$ meeting $\supp w$, or $\supp s$ ($s=\theta_y(a)$), or $\supp m$, and for two consecutive blocks of $a$ meeting $\supp v$ with respect to $\varepsilon_a\vee\varepsilon_c$. These are the analogues of Step~6 of Theorem~\ref{thm:R}, and they imply the goal for the visible indices by the argument of Step~7 there. (h) Why (g) holds is partly visible: the intermediate part of such a chain lies in a single class $L$ of $\varepsilon_y\vee\varepsilon_a$ (the blocks of $b$ outside $\supp y$ lie in $\N(y)$-blocks), so the two end blocks are met at $\T(y)$-blocks $Y,Y'$ of the same class $L$, hence of the same class $L'$ of $\varepsilon_y\vee\varepsilon_c$, hence joined by a chain of blocks of $y$ and $c$. Converting this into a chain of blocks of $b$ and $c$ is possible as long as the blocks of $y$ on it lie inside single blocks of $b$, which fails exactly for the $\T(y)$-blocks inside a $\T(w)$-block $W_l$ (these may straddle several $\T(b)$-blocks $B_i\subseteq W_l$) and for $\N(y)$-blocks that are unions of several $\T(b)$-blocks. So the whole difficulty of (W$_2$) sits in the interleaving of the $\T(y)$- and $\T(b)$-blocks inside the $\T(w)$-blocks, i.e.\ in telling the $\T(b)$-blocks of one $W_l$ apart; no pairwise join of $a,b,c,x,y,m$ does this (two $B_i,B_j\subseteq W_l$ can share all of them), so the invariant that separates them is finer than any join, as the failure of the coarser joins for the goal already indicated.

\subsection{The witnesses $\mathfrak P_n^{\mathrm{ch}}$ of the chain conditions and a conditional non-finite-basis theorem}\label{subsec:chainwit}

\emph{(Added 11 Oct 2026.)} M.~Kinyon's Mace4 run on \texttt{w3ind.in} found a projection algebra of order $11$ satisfying \eqref{eq:Q1}, (S), (R), (R$'$), (V), (R$''$) and (W$_2$) but not (W$_3$); it also satisfies the triangle condition, (P$_1$), (P$_2$), (Q), (E$_1$), (E$_6$) and K0--K6, has no embedding into $\P(\P_m)$ for $m\le7$, and all six of its maximal proper subalgebras embed in $\P(\P_5)$ (\texttt{pa8/w3/}). Its structure---a chain $c\sim p\sim q\sim k\sim b$ of top elements whose consecutive members commute because they share a lower bound, over a rectangular band of bottom elements---generalises at once to every $n$.

\begin{definition}\label{def:chainwit}
For $n\ge1$ let $\mathfrak P^{\mathrm{ch}}_n$ be the set of $2n+5$ elements $a,c,b,x_1,\dots,x_n$ (\emph{top}) and $d,e_0,e_1,\dots,e_n$ (\emph{bottom}) with the operation $\theta$ defined by: $\theta_e\equiv e$ for every bottom element $e$; $\theta_a(c)=a$, $\theta_a(y)=d$ for $y\ne a,c$; $\theta_c(a)=\theta_c(b)=c$, $\theta_c(y)=e_0$ for $y\ne a,b,c$; $\theta_b(c)=b$, $\theta_b(y)=e_n$ for $y\ne b,c$; and, for $1\le i\le n$, $\theta_{x_i}(y)=e_{i-1}$ if $y\in\{a,c,d,e_0,\dots,e_{i-1},x_1,\dots,x_{i-1}\}$ (the elements to the \emph{left} of $x_i$) and $\theta_{x_i}(y)=e_i$ if $y\in\{e_i,\dots,e_n,x_{i+1},\dots,x_n,b\}$ (to the \emph{right}).
\end{definition}

So the bottom elements form a rectangular band, $d$ is the bottom of $a$, $e_0=\theta_c(x_1)=\theta_{x_1}(c)$, $e_i=\theta_{x_i}(x_{i+1})=\theta_{x_{i+1}}(x_i)$ and $e_n=\theta_{x_n}(b)=\theta_b(x_n)$ are the shared lower bounds along the chain $c\sim x_1\sim\dots\sim x_n\sim b$, $a\bowtie c\bowtie b$, every $x_i$ identifies $a$ with $c$ ($\theta_{x_i}(a)=\theta_{x_i}(c)=e_{i-1}$), and $a\not\bowtie b$ ($\theta_b(a)=e_n$). Hence $\mathfrak P^{\mathrm{ch}}_n$ violates (W$_n$), and (W$_k$) for every $k\ge n$ (repeat $x_1$). Computer checks (\texttt{pa8/wn/family.py}, \texttt{check.py}): for $n\le6$ (\texttt{family\_check.log}) $\mathfrak P^{\mathrm{ch}}_n$ is a projection algebra satisfying \eqref{eq:Q1}, the triangle condition and every condition of the status table other than the (W$_k$), $k\ge n$; it satisfies (W$_1$),\dots,(W$_{n-1}$); it has no embedding into $\P(\P_m)$ for $m\le7$ ($n\le3$; $m\le6$ for $n\le4$, the larger cases running); and each of its $n+3$ maximal proper subalgebras (omit one top element) embeds in some $\P(\P_m)$, $m\le7$ ($n\le3$; the larger cases running), so that it is a \emph{minimal} non-embeddable projection algebra as soon as (W$_n$) is valid in every $\P_X$. For $n=2$, $\mathfrak P^{\mathrm{ch}}_2$ is candidate~$6$ of order $9$ (the (R$''$)-witness) up to isomorphism, and $\mathfrak P^{\mathrm{ch}}_1$ is $\mathfrak P_7$ with $a$ and $b$ exchanged; for $n=3$ the Mace4 model of M.K.\ is not isomorphic to $\mathfrak P^{\mathrm{ch}}_3$ (it has $x_1\bowtie x_3$ and some other values $\theta_{x_i}(y)$; several choices of these values give witnesses, and the definition above is the uniform one, found by letting Mace4 complete the fixed part of the table, \texttt{pa8/wn/gen.py}). The minimality for all $n$ needs explicit embeddings of the $n+3$ maximal subalgebras, which are given in Section~\ref{subsec:maxsub}. The uniform verification of the axioms, of \eqref{eq:Q1}, of (W$_{n-1}$) and of all (R$'_j$) for every $n$ is short (it was supplied by the cold reader of 11~Oct 2026); we record it. Write $B$ for the bottoms ($\theta_e\equiv e$), $T$ for the tops, and for a top $t$ let $M_t=\{y:\theta_t(y)=t\}$ ($M_a=\{a,c\}$, $M_c=\{a,b,c\}$, $M_b=\{b,c\}$, $M_{x_i}=\{x_i\}$; note $y\in M_t\Leftrightarrow t\in M_y$) and $\beta_t(y)=\theta_t(y)\in B$ for $y\notin M_t$ ($\beta_a\equiv d$, $\beta_c\equiv e_0$, $\beta_b\equiv e_n$, $\beta_{x_i}=e_{i-1}$ on the left of $x_i$ and $e_i$ on the right). The one fact used repeatedly is ($*$): for a top $t\ne x_i$, $\beta_t$ maps every element to the same side of $x_i$ as $t$ itself ($a,c,d,e_0$ are left of every $x_i$, $b,e_n$ right, and for $x_j$ the value $e_j$ is left of $x_i$ when $j<i$ while $e_{j-1}$ is right when $j>i$), and $\beta_{x_i}$ maps left to $e_{i-1}$ (left) and right to $e_i$ (right). Then: $\theta_x(\theta_x(y))=\theta_x(y)$ is trivial for $x\in B$ or $\theta_x(y)=x$, and otherwise $\theta_x(\beta)=\beta$ for $\beta=\beta_x(y)$ since $d\notin M_a$, $e_0\notin M_c$, $e_n\notin M_b$ and $e_{i-1},e_i$ lie on the correct sides of $x_i$. $\theta_x(\theta_y(x))=\theta_x(y)$ is trivial unless $x,y$ are distinct tops with $y\notin M_x$, and then $\theta_x(\beta_y(x))=\beta_x(y)$ by ($*$). For $\theta_x(\theta_y(\theta_x(z)))=\theta_{\theta_x(y)}(z)$: if $x\in B$ both sides are $x$; if $y\in B$ both are $\beta_x(y)$; if $y\in M_x$ the right side is $\theta_x(z)$, and so is the left side (for $z\in M_x$ use $y\in M_x\Leftrightarrow x\in M_y$, otherwise $\beta_x(\beta_y(\beta_x(z)))=\beta_x(z)$ by ($*$)); if $y\notin M_x$ the right side is $\beta_x(y)$ and the left side is $\theta_x(\beta_y(x))=\beta_x(y)$ for $z\in M_x$ and $\beta_x(\beta_y(\beta_x(z)))=\beta_x(y)$ otherwise, by ($*$). For $\theta_y\theta_x\theta_y\theta_x(z)=\theta_y\theta_x(z)$ put $w=\theta_y(\theta_x(z))$: if $w=y$ then $\theta_x(z)\in M_y$ forces $\theta_x(z)=x\in M_y$ and $\theta_y\theta_x(y)=\theta_y(x)=y$; if $w\in B$ then $\beta_y(\beta_x(w))=w$ by ($*$), since $w=\beta_y(t)$ with $t\in\{x\}\cup B$ on the same side of $y$ as $x$. \eqref{eq:Q1}: $\theta_x(z)=\theta_y(z)=z$ with $x\ne y$ forces $z\in B$ below both; the only tops sharing a bottom are consecutive chain members ($e_0\le c,x_1$; $e_i\le x_i,x_{i+1}$; $e_n\le x_n,b$), and these commute. Finally, the identifiers of $(a,c)$ in $\mathfrak P^{\mathrm{ch}}_n$ are $c$, the $x_i$ and the bottoms, and the commuting graph on them is the path $c-x_1-\dots-x_n$ with $e_0$ adjacent to $c,x_1$, $e_i$ to $x_i,x_{i+1}$ and $e_n$ to $x_n$ (and $b$); the only elements commuting with $b$ are $b,x_n,e_n$, at distance $n$ or $n+1$ from $c$, so no chain of length $n-1$ closes the path, i.e.\ (W$_{n-1}$) holds; the other unclosed 2-path $(b,c,a)$ gives nothing (only $c$ and the bottoms identify $b$ with $c$, and none of them commutes with a partner of $a$); and every (R$'_j$) holds because the partners of $a$ are $a,d$ while $\theta_{x_j}(b)\in\{b,e_j\}$ is never $d$.

\begin{theorem}[conditional]\label{thm:nfb}
Suppose that \textup{(W$_n$)} holds in every partition monoid $\P_X$ for every $n\ge1$ and that every maximal proper subalgebra of $\mathfrak P^{\mathrm{ch}}_n$ embeds in some $\P(\P_X)$ \textup{(}both verified for small $n$ as stated above\textup{)}. Then for every finite set $\Sigma$ of quasi-identities valid in all $\P(\P_X)$ there is a finite projection algebra satisfying $\Sigma$ that embeds in no $\P(\P_X)$, namely $\mathfrak P^{\mathrm{ch}}_n$ for any $n$ such that every member of $\Sigma$ has at most $n+2$ variables. In particular the class of finite projection algebras embeddable in some $\P(\P_X)$ is not the class of finite models of any finite set of quasi-identities, the quasivariety generated by the $\P(\P_X)$ is not finitely based, and \textup{(W$_{n+1}$)} is not a consequence of \textup{(W$_1$)},\dots,\textup{(W$_n$)} together with any valid quasi-identities in at most $n+3$ variables.
\end{theorem}

\begin{proof}
The top elements of $\mathfrak P^{\mathrm{ch}}_n$ are irreducible: $\theta_x(y)$ is a top element only when it equals $x$. So a subalgebra generated by $k$ elements contains at most $k$ top elements, and for $k\le n+2$ it omits one of the $n+3$ top elements, hence lies in a maximal proper subalgebra, which embeds in some $\P(\P_X)$ by hypothesis. Let $\sigma\in\Sigma$ have $k\le n+2$ variables and suppose an assignment violates $\sigma$ in $\mathfrak P^{\mathrm{ch}}_n$; the assignment violates $\sigma$ in the subalgebra generated by its values, which embeds in some $\P(\P_X)$, contradicting the validity of $\sigma$ there. So $\mathfrak P^{\mathrm{ch}}_n\models\Sigma$, while $\mathfrak P^{\mathrm{ch}}_n$ violates \textup{(W$_n$)}, a quasi-identity valid in all $\P(\P_X)$ by hypothesis; so it embeds in none. The last assertion is the case $\Sigma=\{$(W$_1$),\dots,(W$_n$), and the given quasi-identities$\}$ applied to $\mathfrak P^{\mathrm{ch}}_{n+1}$, which has $n+4$ top elements.
\end{proof}

\emph{(Both hypotheses are now proved: Theorem~\ref{thm:Wn} and Proposition~\ref{prop:maxsub}; the unconditional statement is Theorem~\ref{thm:nfbW}.)} The theorem makes the two missing proofs precise: the validity of the chain conditions in $\P_X$ (the uniform proof of Section~\ref{subsec:chain}, now for all $n$ at once), and the embeddings of the maximal subalgebras of $\mathfrak P^{\mathrm{ch}}_n$. It also explains why the earlier searches found the family late: (W$_n$) has $n+3$ variables and its smallest witness has order $2n+5$, so the first new member beyond (R)$'$ needs order $9$ and five variables, the next order $11$ and six.

\subsection{How the conditions fit together (status on 11 Oct 2026)}\label{subsec:picture}

The conditions now on the list are \eqref{eq:Q1}; the triangle condition (Tri); (V); and the \emph{path-closing} conditions (S), (R), (R$'$), (R$''$), (P$_1$), (P$_2$), (Q), (E$_1$), (E$_6$), K0--K6, (W$_n$), (R$'_n$) (the last family is introduced below). Three facts organise them.

\emph{(a) Two conditions are not of path-closing type.} The triangle witness $\P(Q_0)$ satisfies every other condition, including (W$_n$) for $n\le4$, and its unclosed 2-paths have no explanation in the three-auxiliary equality language (Section~\ref{subsec:chain}); the (V)-witness $\mathfrak P_8^7$ has no unclosed 2-path at all. Everything else has a witness with an unclosed 2-path $a\bowtie c\bowtie b$ that is closed by a 2-path condition with at most three auxiliaries, valid at least for $m\le5$: this was shown in Section~\ref{subsec:nine} for the order-$9$ candidates and $\mathfrak P_8^1$, $\mathfrak P_8^2$, $\mathfrak P_8^6$, in Section~\ref{subsec:chain} for $\mathfrak P_8^4$, $\mathfrak P_8^5$, and now also for the (S)-witness $\mathfrak P$ itself (\texttt{pa8/impl/explain\_S.log}: six explanations for the path $a\bowtie c\bowtie b$ of $\mathfrak P$, e.g.\ $\theta_b(x)=b$, $\theta_y(z)=\theta_a(b)=\theta_b(a)$, $\theta_c(z)=\theta_x(a)=\theta_z(y)$, and six for $c\bowtie b\bowtie d$), and for $\mathfrak P_9$. So, conjecturally, the 3-path conditions (S), (P$_1$), (P$_2$), (Q) are consequences in $\P_X$ of 2-path conditions with three auxiliaries, and the whole list outside (Tri), (V) is ``2-paths closed by auxiliaries''.

\emph{(b) Implications and independence in projection algebras.} Besides Proposition~\ref{prop:indep} ((Q)$\Rightarrow$(P$_2$)$\Rightarrow$(S), (R$'$)$\Rightarrow$(R), the eight independent conditions), the new facts are: (W$_{n+1}$)$\Rightarrow$(W$_n$) (trivial), (W$_1$) is (R) with $a$ and $b$ exchanged (the hypotheses $a\bowtie c$, $b\bowtie c$ and the conclusion are symmetric in $a,b$, so (R)$'$ is (R) up to renaming; the Prover9 run \texttt{pa8/impl/r\_rmirror.in} was unnecessary), so that every (W$_n$) implies (R) (Prover9 finds the $17$-step derivation of (R) from (W$_2$) directly, as M.K.\ observed). The violation matrix of the $22$ known non-embeddable algebras against the $18$ labelled conditions (\texttt{pa8/impl/viol.py}, \texttt{viol.json}) shows that no other implication holds among the conditions of the list, with one exception: the triple (R$''$), (W$_2$), K6 is violated by the same algebras (candidate~$6$ $=\mathfrak P^{\mathrm{ch}}_2$, and $\mathfrak P_7$) and satisfied by all the others, so no known algebra separates any two of them, Mace4 finds no separating projection algebra of order $\le8$ with \eqref{eq:Q1} for any of the six ordered pairs, and Prover9 derives none from another in $120$\,s; likewise K0$\Rightarrow$(E$_1$) (\emph{added 13 Oct:} now proved from \eqref{eq:Q1} and (W$_2$), Theorem~\ref{thm:K0E1}), K4$\Rightarrow$(S) and K4$\Rightarrow$(E$_6$) are open in the same sense (every known violator of the right-hand side violates the left-hand side; no separating algebra of order $\le8$). In particular (W$_2$) does not imply (S), (R$'$), (V), (P$_1$), (Q), (E$_1$), (E$_6$), (Tri) or K0, K4, K5 (their witnesses satisfy (W$_2$)), and none of those implies (W$_2$) (candidate~$6$ satisfies them); and (W$_3$) is separated from (W$_2$) by $\mathfrak P^{\mathrm{ch}}_3$ and M.K.'s model, but implies nothing new beyond (W$_2$). Likewise (R$'_2$) is separated from (R$'_1$) and from every other listed condition by M.K.'s order-$10$ model $\mathfrak P^{\mathrm R'}_2$ (Section~\ref{subsec:rpfam}), and implies nothing new beyond (R$'$).

\emph{(c) Chains out of (R$'$): a correction.} The first attempt, ``$c\sim x_1\sim\dots\sim x_n\sim a$, $x_1$ identifies $a$ with $c$, $\theta_{x_i}(b)=\theta_{x_{i+1}}(b)$'', is not a new family: the agreements compose to $\theta_{x_1}(b)=\theta_{x_n}(b)$ and (R$'$) applied to $(x_1,x_n)$ gives the conclusion, the intermediate commutations being unused (M.K.: Prover9 proves the $n=3$ case from (R$'$) in $18$ steps). What (W$_n$) adds to (R)$'$ is that \emph{every} chain member identifies $a$ with $c$ and the chain carries the commutation from $c$ to $b$. The corresponding hybrid of (W$_n$) and (R$'$),
\[
\text{(R$'_n$)}\qquad a\bowtie c,\ b\bowtie c,\quad c\sim x_1\sim\dots\sim x_n,\quad \theta_{x_i}(a)=\theta_{x_i}(c)\ (1\le i\le n),\quad ya=ay,\quad \theta_{x_n}(b)=\theta_y(b)\ \Longrightarrow\ a\bowtie b,
\]
has (R$'_1$) $=$ (R$'$), (R$'_{n+1}$)$\Rightarrow$(R$'_n$) (repeat $x_n$), and is valid in $\P(\P_m)$ for $m\le5$ and $n\le3$ (\texttt{pa8/impl/rpchain2.py}); dropping the identification for $x_2,\dots,x_n$ makes it fail already at $m=4$, so here the chain is used. Whether (R$'_2$) is independent of (R$'_1$) (and of (W$_2$)) was the question (Mace4 input \texttt{rp2ind.in}); it is answered positively in Section~\ref{subsec:rpfam}. One expects the same hybrids for K0--K6, and the picture that emerges is a short list of chain families seeded by the order-$7$ to order-$9$ conditions, rather than a list of isolated conditions.

\subsection{The second chain family: (R$'_n$) and its witnesses $\mathfrak P^{\mathrm R'}_n$}\label{subsec:rpfam}

\emph{(Added 11 Oct 2026.)} M.~Kinyon's Mace4 run on \texttt{rp2ind.in} found a projection algebra of order $10$ satisfying \eqref{eq:Q1}, (S), (R), (R$'$), (V), (R$''$), (W$_2$) but not (R$'_2$), and showed that there is none of order $\le9$ (our own run exhausted the domains $\le8$). With the constants of the Mace4 output ($a=0$, $b=1$, $c=2$, $k=3$, $p=4$, $q=5$; \texttt{pa8/impl/rp2/mk10.json}) it has the top elements $a,c,b,p,q,k$ and the rectangular band $\{6,7,8,9\}$ of bottom elements, with $8\le a$, $6\le b$, $7\le c$, $7,9\le p$, $9\le q$, $p,q,7,8,9\le k$; the single violating instance of (R$'_2$) is $(a,c,b,x_1,x_2,y)=(a,c,b,p,q,k)$: $a\bowtie c\bowtie b$, $c\sim p\sim q$, $\theta_p(a)=\theta_p(c)=7$, $\theta_q(a)=\theta_q(c)=9$, $k\sim a$ (and $k\sim c$), $\theta_q(b)=\theta_k(b)=q$, and $\theta_b(a)=6$. Note that $b\bowtie q$ and $q=\theta_k(b)$. Checks (\texttt{pa8/impl/rp2/}): the algebra satisfies every condition of the status table (including (Tri), (W$_3$), (W$_4$), K0--K6); it has no embedding into $\P(\P_m)$ for $m\le7$; its five maximal proper subalgebras (omit $a$, $b$, $c$, $k$ or $p$; $q=\theta_k(b)$ is not irreducible) embed in $\P(\P_4)$ or $\P(\P_5)$; and among the $2^{11}$ subsets of the hypotheses it offers for the path $a\bowtie c\bowtie b$ ($b\bowtie q$, $p\sim c$, $p\sim q$, $p\le k$, $q\le k$, $k\sim a$, $k\sim c$, the two identifications, $\theta_k(b)=q$, $\theta_k(b)=\theta_q(b)$) exactly two are minimal valid conditions at $m=4$ (both valid at $m=5$; \texttt{minsets.py}): (R$'_2$) itself, and its variant with $\theta_k(b)=q$ in place of $\theta_k(b)=\theta_q(b)$. So nothing simpler than (R$'_2$) is responsible, and in particular $b\bowtie q$, $p,q\le k$ and $k\sim c$ play no r\^ole. (R$'_2$) is valid in $\P(\P_5)$ (SAT, $16$\,s; $m=6$ running), so, at the usual evidence level, the order-$10$ algebra is a new minimal non-embeddable projection algebra, the first of order $10$, and (R$'_2$) is independent of (R$'_1$) and of all other listed conditions.

The algebra is $\mathfrak P^{\mathrm{ch}}_2$ with the last link $x_2\sim b$ replaced by $x_2\bowtie b$ and a new top element $k$ above $x_1,x_2$ added, with $k\sim a$, $k\sim c$ and $\theta_k(b)=x_2$. This generalises to every $n$.

\begin{definition}\label{def:rpwit}
For $n\ge1$ let $\mathfrak P^{\mathrm R'}_n$ be the set of $2n+6$ elements $a,c,b,x_1,\dots,x_n,k$ (top) and $d,e_0,\dots,e_n$ (bottom) with $\theta$ defined by: $\theta_e\equiv e$ for every bottom element $e$; $\theta_a$ and $\theta_c$ as in Definition~\ref{def:chainwit}; $\theta_b(c)=\theta_b(k)=\theta_b(x_n)=b$ and $\theta_b(y)=e_n$ otherwise; for $1\le i<n$, $\theta_{x_i}(k)=x_i$ and $\theta_{x_i}$ as in Definition~\ref{def:chainwit} on the other elements; $\theta_{x_n}(b)=\theta_{x_n}(k)=x_n$ and $\theta_{x_n}(y)=e_{n-1}$ otherwise; $\theta_k(a)=d$, $\theta_k(c)=e_0$, $\theta_k(b)=x_n$, $\theta_k(e_n)=e_{n-1}$, and $\theta_k(y)=y$ for $y\in\{x_1,\dots,x_n,d,e_0,\dots,e_{n-1}\}$.
\end{definition}

So $x_1,\dots,x_n,d,e_0,\dots,e_{n-1}\le k$, $k\sim a$, $k\sim c$, $x_n\bowtie b$, $c\sim x_1\sim\dots\sim x_n$ with every $x_i$ identifying $a$ with $c$, $\theta_{x_n}(b)=x_n=\theta_k(b)$, and $\theta_b(a)=e_n$: the assignment $(x_1,\dots,x_n,y)=(x_1,\dots,x_n,k)$ violates (R$'_n$). Computer checks (\texttt{pa8/impl/rp2/rpfam.py}, \texttt{rpchk.py}, \texttt{rpchk.log}), for $n\le5$ resp.\ $n\le4$: $\mathfrak P^{\mathrm R'}_n$ is a projection algebra satisfying \eqref{eq:Q1}; it satisfies every condition of the status table, including all (W$_j$), $j\le n+1$, and (R$'_1$),\dots,(R$'_{n-1}$); it violates (R$'_n$) in exactly one instance (and (R$'_j$) for $j>n$ by repetition); it has no embedding into $\P(\P_m)$ for $m\le7$ ($n\le3$; $n=4$ running); and each of its $n+3$ maximal proper subalgebras (omit one of the irreducible elements $a,c,b,x_1,\dots,x_{n-1},k$) embeds in some $\P(\P_m)$, $m\le6$. Moreover $\mathfrak P^{\mathrm R'}_1\cong\mathfrak P_8^3$, the (R$'$)-witness of the order-$8$ census, and $\mathfrak P^{\mathrm R'}_2$ is M.K.'s order-$10$ model (\texttt{iso.py}). So the two chain families are exactly parallel: $\mathfrak P^{\mathrm{ch}}_n$ starts at $\mathfrak P_7$ (the (R)-witness) and $\mathfrak P^{\mathrm R'}_n$ at $\mathfrak P_8^3$ (the (R$'$)-witness), each adding two elements per chain link, and each new member satisfies everything previously known.

\begin{theorem}[conditional]\label{thm:nfb2}
Suppose that \textup{(R$'_n$)} holds in every $\P_X$ for every $n\ge1$ and that every maximal proper subalgebra of $\mathfrak P^{\mathrm R'}_n$ embeds in some $\P(\P_X)$. Then $\mathfrak P^{\mathrm R'}_n$ satisfies every quasi-identity valid in all $\P(\P_X)$ with at most $n+2$ variables, but violates \textup{(R$'_n$)}; hence the conclusions of Theorem~\ref{thm:nfb} hold with $\mathfrak P^{\mathrm R'}_n$ in place of $\mathfrak P^{\mathrm{ch}}_n$, and \textup{(R$'_{n+1}$)} is not a consequence of \textup{(R$'_1$)},\dots,\textup{(R$'_n$)}, of all \textup{(W$_j$)}, and of any valid quasi-identities in at most $n+3$ variables.
\end{theorem}

\begin{proof}
An element that is not in the subalgebra generated by the other elements (\emph{irreducible}) belongs to a subalgebra only as one of its generators. $\mathfrak P^{\mathrm R'}_n$ has the $n+3$ irreducible elements $a,c,b,x_1,\dots,x_{n-1},k$ (the top element $x_n=\theta_k(b)$ is not irreducible), so a subalgebra generated by $\le n+2$ elements omits one of them and lies in a maximal proper subalgebra, which embeds by hypothesis; the rest is as in the proof of Theorem~\ref{thm:nfb}. For the last assertion take $\mathfrak P^{\mathrm R'}_{m}$ with $m=n+1$: it has $n+4$ irreducible elements, so it satisfies every valid quasi-identity in $\le n+3$ variables; and it satisfies all (W$_j$) and (R$'_1$),\dots,(R$'_n$) by direct inspection (the variable count does not cover these: (R$'_n$) has $n+4$ variables), as follows (argument supplied by the cold reader of 11~Oct). In $\mathfrak P^{\mathrm R'}_m$ the identifiers of $(a,c)$ reachable from $c$ by commuting steps are $c,x_1,\dots,x_m,e_0,\dots,e_{m-1}$ ($k$ does not identify, as $\theta_k(a)=d\ne e_0=\theta_k(c)$; $e_m$ is adjacent only to $b$ and to no identifier), and the commuting graph on them is the path $c-x_1-\dots-x_m$ with $e_i$ pendant between $x_i$ and $x_{i+1}$, so $x_m$ is at distance $m$ from $c$. The only elements commuting with $b$ are $b$ and $e_m$, neither reachable, so no instance of any (W$_j$) exists (for the other unclosed paths $(c,b,x_m)$ and $(x_m,b,c)$ a chain stays in $\{b,e_m\}$ and would have to end at a partner of $x_m$, resp.\ of $c$, which is impossible). For (R$'_j$) with $j<m$: the partners of $a$ are $a,d,k$, with $\theta_y(b)\in\{d,x_m\}$, while a chain end $x_j$ has $\theta_{x_j}(b)\in\{b,e_j,x_m\}$; so an instance needs $x_j=x_m$ and $y=k$, i.e.\ a chain of length $\ge m$.
\end{proof}

The hypotheses of the two conditional theorems are the central open points: a uniform proof of (W$_n$) and of (R$'_n$) in $\P_X$ (Section~\ref{subsec:chain} made a start for (W$_2$)), and explicit embeddings of the maximal subalgebras of the two families; the latter are given in Section~\ref{subsec:maxsub}, so that only the validity of the conditions remains. Mace4 inputs for the next separations are \texttt{pa8/impl/rp2/rp3ind.in} ((R$'_2$) and all listed conditions, negated (R$'_3$); the expected smallest model is $\mathfrak P^{\mathrm R'}_3$ of order $12$, so this is a hard run) and the analogous \texttt{w4ind.in} (expected answer $\mathfrak P^{\mathrm{ch}}_4$, order $13$); the explicit families make these runs unnecessary for the existence of the witnesses, but a Mace4 exhaustion of the smaller domains would show that they are the smallest.

\subsection{The maximal subalgebras of $\mathfrak P^{\mathrm{ch}}_n$ and $\mathfrak P^{\mathrm R'}_n$ embed: explicit embeddings}\label{subsec:maxsub}

\emph{(Added 11 Oct 2026.)} This subsection removes the second hypothesis from Theorems~\ref{thm:nfb} and \ref{thm:nfb2}: the maximal proper subalgebras of the two families embed in $\P(\P_m)$ with $m\le2n+3$, by explicit projections. The constructions were found by hand from the SAT solutions for $n\le3$ and the following two observations, and checked by computer for $n\le7$ (\texttt{pa8/maxsub/chemb.py}, \texttt{rpemb.py}, \texttt{verify7.log}: every product $\theta_y(x)$ of the subalgebra is recomputed in $\P_m$ and compared with the table).

\emph{Rank-$0$ projections.} A projection $e$ of rank $0$ is a partition $\varepsilon_e$ of $X$ with no transversal block, and by Lemma~\ref{lem:PX}(a) $\theta_e(y)=e$ for every projection $y$, while for any projection $p$ the element $\theta_p(e)$ has rank $0$ and consists of the $\N(p)$-blocks together with the unions $U_K$ of the $\T(p)$-blocks lying in a common class $K$ of $\varepsilon_p\vee\varepsilon_e$. So all rank-$0$ projections are pairwise matched, as the bottom elements of the two families are, and the bottom elements will be realised by rank-$0$ projections; the values $\theta_p(e)$ then depend on $e$ only through the way $\varepsilon_e$ groups the $\T(p)$-blocks. More generally, for projections $p,q$, Lemma~\ref{lem:PX}(a) computes $\theta_p(q)$ from the classes of $\varepsilon_p\vee\varepsilon_q$: it has rank $0$ iff no class contains both a $\T(p)$- and a $\T(q)$-block; and $p\bowtie q$ iff the classes meeting $\supp p\cup\supp q$ contain exactly one $\T(p)$- and one $\T(q)$-block each (Lemma~\ref{lem:PX}(c), and the converse holds since then $pqp$ and $p$ have the same blocks).

\emph{The pairs.} Each chain element $x_l$ with two neighbours is realised with exactly two $\T$-blocks, the singletons $\{p_l\}$, $\{q_l\}$ of a \emph{pair} $X_l=\{p_l,q_l\}$ of points reserved for it; every other element has $X_l$ either as one $\N$-block (``$X_l$ merged'') or as two singleton $\N$-blocks (``$X_l$ separate''), and no other block meets $X_l$. Then for every element $y\ne x_l$ with $\theta_{x_l}(y)$ of rank $0$, $\theta_{x_l}(y)$ is $\N(x_l)$ together with $X_l$ merged if $X_l$ is merged in $y$, and with $X_l$ separate otherwise; this realises the dichotomy left/right of Definition~\ref{def:chainwit}, with ``left'' $=$ ``$X_l$ merged''. The bottom element $e_j$ has $X_l$ separate for $l\le j$ and merged for $l>j$; $d$, $a$, $c$ have all pairs merged, $b$ all pairs separate, and $x_l$ has $X_m$ separate for $m<l$ and merged for $m>l$. Writing $[\,\cdot\,]$ for a $\T$-block and $(\,\cdot\,)$ for an $\N$-block, the constructions are as follows (all pairs present are listed after the semicolon only through the rule just stated; $\alpha,\beta,\gamma,\delta,\delta',\zeta$ are further points).

\smallskip
\emph{$\mathfrak P^{\mathrm{ch}}_n$, omit $x_i$} ($2n+1$ points $\alpha,\beta,\gamma$ and the pairs $X_l$, $l\ne i$; $R=\{\alpha,\beta,\gamma\}$): $a=[\alpha](\beta\gamma)$, $c=[\alpha\beta](\gamma)$, $b=[\beta](\alpha)(\gamma)$, $d=(\alpha)(\beta\gamma)$; on $R$, every $x_l$ with $l<i$ and every $e_j$ with $j<i$ is $(\alpha\beta)(\gamma)$, every $x_l$ with $l>i$ and every $e_j$ with $j\ge i$ is $(\alpha)(\beta)(\gamma)$. \emph{Omit $b$}: the same with all $x_l$, $e_j$ equal to $(\alpha\beta)(\gamma)$ on $R$. \emph{Omit $c$}: $a=[\alpha](\beta\gamma)$, $b=[\beta](\alpha)(\gamma)$, $d=(\alpha)(\beta\gamma)$, all $x_l$, $e_j$ equal to $(\alpha)(\beta)(\gamma)$ on $R$. \emph{Omit $a$}: $c=[\alpha][\beta](\gamma)$ with all pairs merged, $b=[\alpha][\beta](\gamma)$ with all pairs separate, $d=(\alpha)(\beta\gamma)$, all $x_l$, $e_j$ equal to $(\alpha)(\beta)(\gamma)$ on $R$. (In the last three cases $2n+3$ points.)

\smallskip
\emph{$\mathfrak P^{\mathrm R'}_n$} ($x_n$ has one $\T$-block, $[\delta]$, and only the pairs $X_l$, $l<n$, are used; $k$ has the $\T$-blocks $[\delta]$, $[\delta']$, $[\zeta]$ (those present) and all $[p_l]$, $[q_l]$, so that $x_l\le k$). \emph{Omit $x_i$}, $i<n$ ($2n+2$ points; $R=\{\alpha,\beta,\gamma,\delta,\delta',\zeta\}$): $a=[\alpha\beta](\gamma)(\delta\zeta)(\delta')$, $c=[\alpha\beta](\gamma)(\delta\delta')(\zeta)$, $b=[\delta](\alpha)(\beta\delta')(\gamma)(\zeta)$, $x_n=[\delta](\alpha\beta)(\gamma)(\delta')(\zeta)$, $k=[\delta][\delta'][\zeta](\alpha\beta)(\gamma)$, $d=(\alpha\beta)(\gamma)(\delta\zeta)(\delta')$, $e_n=(\alpha)(\beta\delta')(\gamma)(\delta)(\zeta)$; on $R$, $x_l$ ($l<i$) and $e_j$ ($j<i$) are $(\alpha\beta)(\gamma)(\delta\delta')(\zeta)$, $x_l$ ($i<l<n$) and $e_j$ ($i\le j<n$) are $(\alpha\beta)(\gamma)(\delta)(\delta')(\zeta)$. \emph{Omit $a$} ($2n+3$ points, no $\delta'$): $c=[\alpha\beta](\gamma)(\delta)(\zeta)$, $b=[\beta\delta](\alpha)(\gamma)(\zeta)$, $x_n=[\delta](\alpha\beta)(\gamma)(\zeta)$, $k=[\delta][\zeta](\alpha\beta)(\gamma)$, $d=(\alpha\beta)(\gamma)(\delta\zeta)$, $e_n=(\alpha)(\gamma)(\zeta)(\beta\delta)$, all $x_l$ and $e_j$ ($j<n$) equal to $(\alpha\beta)(\gamma)(\delta)(\zeta)$ on $R$. \emph{Omit $b$}: as for omit $a$ with $a=[\alpha\beta](\gamma)(\delta\zeta)$ in place of $b$ and $e_n=(\alpha)(\beta)(\gamma)(\delta)(\zeta)$. \emph{Omit $c$}: $a=[\alpha](\beta)(\gamma)(\delta\zeta)$, $b=[\beta\delta](\alpha)(\gamma)(\zeta)$, $x_n=[\delta](\alpha)(\beta)(\gamma)(\zeta)$, $k=[\delta][\zeta](\alpha)(\beta)(\gamma)$, $d=(\alpha)(\beta)(\gamma)(\delta\zeta)$, $e_n=(\alpha)(\gamma)(\zeta)(\beta\delta)$, all $x_l$, $e_j$ ($j<n$) equal to $(\alpha)(\beta)(\gamma)(\delta)(\zeta)$ on $R$. \emph{Omit $k$} ($2n+2$ points, no $\delta',\zeta$): $a=[\alpha](\beta\gamma)(\delta)$, $c=[\alpha\beta](\gamma)(\delta)$, $b=[\beta\delta](\alpha)(\gamma)$, $x_n=[\delta](\alpha\beta)(\gamma)$, $d=(\alpha)(\beta\gamma)(\delta)$, $e_n=(\alpha)(\gamma)(\beta\delta)$, all $x_l$, $e_j$ ($j<n$) equal to $(\alpha\beta)(\gamma)(\delta)$ on $R$.

\begin{proposition}\label{prop:maxsub}
For every $n\ge1$, every maximal proper subalgebra of $\mathfrak P^{\mathrm{ch}}_n$ and of $\mathfrak P^{\mathrm R'}_n$ embeds in $\P(\P_m)$ for some $m\le2n+3$, by the maps listed above. Consequently Theorems~\ref{thm:nfb} and \ref{thm:nfb2} hold under the single hypothesis that \textup{(W$_n$)}, resp.\ \textup{(R$'_n$)}, is valid in every $\P_X$ for every $n$.
\end{proposition}

\begin{proof}
Each listed map is injective (distinct elements differ in a block of $R$ or in a pair), and each value $\theta_y(x)$ is computed by Lemma~\ref{lem:PX}(a) from the classes of $\varepsilon_y\vee\varepsilon_x$. These classes are read off from the lists: a pair $X_l$ is a class of $\varepsilon_y\vee\varepsilon_x$ by itself for every $x,y$, or the union of its two singletons, and the points of $R$ are connected only through the listed blocks on $R$. The products therefore fall into finitely many patterns independent of $n$---e.g.\ $\theta_{x_l}(y)$ for $y\ne x_l,k$: the classes of $\varepsilon_{x_l}\vee\varepsilon_y$ meeting $\{p_l,q_l\}$ are $X_l$ or its two singletons according as $X_l$ is merged or separate in $y$, no class contains a $\T(y)$-block together with $[p_l]$ or $[q_l]$ (the $\T(y)$-blocks lie in $R$ and the blocks on $R$ of $x_l$ are $\N$-blocks), so $\theta_{x_l}(y)$ is $\N(x_l)$ with $X_l$ merged or separate, i.e.\ $e_{l-1}$ or $e_l$, as $y$ is to the left or right of $x_l$; or $\theta_k(y)$ for $y\notin\{b,x_1,\dots,x_n\}$: no class joins a $\T(k)$-block, all of which are singletons outside $(\alpha\beta)$ and $(\gamma)$, to a $\T(y)$-block, so $\theta_k(y)$ is $\N(k)$ together with the groups of singletons $[\delta],[\delta'],[\zeta],[p_l],[q_l]$ connected through the blocks of $y$, which reproduces $y$ for $y\in\{d,e_0,\dots,e_{n-1}\}$ and $e_{n-1}$ for $y=e_n$ (there $(\beta\delta')$ or $(\beta\delta)$ leads from $\delta'$ or $\delta$ to $\beta$, whose block $(\alpha\beta)$ in $k$ leads to $\alpha$, a singleton of $e_n$). The remaining patterns ($a\bowtie c$, $c\bowtie b$, $x_n\bowtie b$ and $\theta_k(b)=x_n$ through the paths $\beta\to\delta'\to\delta$ resp.\ $\beta\to\delta$, the rank-$0$ values of $\theta_a,\theta_b,\theta_c$) are checked in the same way. That the computer check for $n\le7$ is a proof for all $n$ rests on two observations (made precise by the cold reader of 11~Oct). First, the point set is $R\sqcup X_1\sqcup\dots\sqcup X_n$ and no block of any listed projection meets two of these parts; hence $\varepsilon_y\vee\varepsilon_x$ is the disjoint union of the joins of the restrictions, $\theta_y(x)$ is computed part by part, and the maps are embeddings into $\P_R\times\prod_l\P_{X_l}\le\P_X$. Second, on a pair $X_l$ the restriction of every element is one of $[p_l][q_l]$, $(p_l\,q_l)$, $(p_l)(q_l)$, and the pair rule above is a three-line verification; on $R$ the restriction of every element takes one of finitely many values independent of $n$ (the ``$R$-types'', e.g.\ in ``omit $x_i$'' the left and right types of the $x_l$ and $e_j$), and the $R$-type of the table entry $\theta_y(x)$ is a function of the $R$-types of $x$ and $y$. So the whole proposition reduces to a finite table of combinations of $R$-types and pair types, and all combinations that occur for any $n$ already occur for $n=3$ (chain family) resp.\ $n=4$ ((R$'$) family); the computer check at $n\le7$ therefore covers every case.
\end{proof}

The $(W_n)$ and $(R'_n)$ witnesses are thus minimal non-embeddable projection algebras exactly when $(W_n)$, resp.\ $(R'_n)$, is valid in every $\P_X$; the two conditional theorems now depend on that validity alone.

\subsection{Towards the uniform proof: the two join statements and what the chain invariant is not}\label{subsec:joinstat}

\emph{(Added 11 Oct 2026; files in \texttt{pa8/red/}.)} This subsection records the attack on the obstruction of Section~\ref{subsec:chain}(h). Everything ``verified'' is verified by SAT in $\P(\P_m)$ for $m\le5$ (variables not assumed distinct) unless said otherwise.

\emph{(a) A uniform strong conclusion.} For a path $a\bowtie c\bowtie b$ consider the two statements
\[
\text{(G1)}\quad \varepsilon_c\le\varepsilon_a\vee\varepsilon_b,\qquad\qquad \text{(G2)}\quad\text{no class of }\varepsilon_a\vee\varepsilon_b\vee\varepsilon_c\text{ contains two }\T(c)\text{-blocks}
\]
((G2) is \eqref{eq:Rp-goal}). Both are verified consequences of the hypotheses of every 2-path condition on the list, proved or not: (R), (R)$'$, (R$'$), (R$''$), (W$_2$), (R$'_2$). Neither follows from $a\bowtie c\bowtie b$ together with the conclusion $a\bowtie b$ (counterexamples in $\P(\P_4)$), so they are genuinely stronger than $a\bowtie b$; but they imply it:

\begin{lemma}[closing lemma]\label{lem:closing}
If $a\bowtie c$, $c\bowtie b$, \textup{(G1)} and \textup{(G2)} hold in $\P_X$, then $a\bowtie b$; more precisely, every class of $\varepsilon_a\vee\varepsilon_b$ meeting $\supp a\cup\supp b\cup\supp c$ contains exactly one $\T(a)$-, one $\T(b)$- and one $\T(c)$-block, namely $A_i,B_i,C_i$ for one $i$.
\end{lemma}

\begin{proof}
By (G1), $\varepsilon_a\vee\varepsilon_b=\varepsilon_a\vee\varepsilon_b\vee\varepsilon_c\ge\varepsilon_a\vee\varepsilon_c,\ \varepsilon_b\vee\varepsilon_c$. Let $D$ be the class of $\varepsilon_a\vee\varepsilon_b$ containing $A_i$. Then $\Gamma_i\subseteq D$, so $C_i\subseteq D$, and $\Delta_i\subseteq D$, so $B_i\subseteq D$. If $A_j\subseteq D$ or $B_j\subseteq D$ for some $j\ne i$, then $C_j\subseteq D$ as well, contradicting (G2). The same applies to the class of any $B_i$ or $C_i$. So the classes meeting $\supp a\cup\supp b$ pair the $A$'s with the $B$'s bijectively, which characterises $a\bowtie b$ (Lemma~\ref{lem:PX}(c) and its proof).
\end{proof}

So for every condition of the list, and for the whole families (W$_n$), (R$'_n$), the uniform target is (G1)$\wedge$(G2), and the $a$--$b$ matching itself needs no counting argument (in the proof of (R$'$), Step~8 can be replaced by a proof of (G1)). Two qualifications (cold read of 11~Oct): the closing step is simplified uniformly, but the proof burden is not reduced, since (G1) is an extra statement to prove and is at present only verified (for (R$'$), Step~8 gives $a\bowtie b$ from $y\,ab=y\,acb$ without it); and the $\N(c)$-blocks in (G1) are essential: the weakening ``every $\T(c)$-block lies in one class of $\varepsilon_a\vee\varepsilon_b$'' together with (G2) does not imply $a\bowtie b$ ($1920$ counterexamples in $\P(\P_4)$), because $\Gamma_i\subseteq D$ needs the $\N(c)$-blocks of $\Gamma_i$ inside $D$. For (W$_2$), (G1) and (G2) are the statements to be proved from the chain $c\sim x\sim y\sim b$.

\emph{(b) (W$_2$) does not reduce to (W$_1$).} On $330$ of the $244\,469$ instances of the (W$_2$)-hypotheses in $\P(\P_4)$ there is no projection $z$ of $\P_4$, nor of $\P_5$ after adjoining a point, with $c\sim z$, $\theta_z(a)=\theta_z(c)$ and $z\sim b$ (\texttt{contract3.py}); a typical one is $a=[0](1\,2\,3)$, $c=[0](1)(2\,3)$, $b=[0\,1](2)(3)$, $x=(0)(1)[2][3]$, $y=[0][1](2)(3)$, where $x$ splits the $\N(c)$-block $\{2,3\}$ into two $\T$-blocks and $y$ turns them into two $\N$-blocks, so that $y$ refines $c$ without commuting with it. Hence no proof of (W$_2$) can proceed by exhibiting one auxiliary and quoting (W$_1$); the two links must be used together. Likewise, no substitution of terms of depth $\le1$ ($\theta_p(q)$ in the variables) into any proved condition ((R), (R$'$), (S), (P$_1$), (Q), either orientation) has hypotheses valid on all (W$_2$)- or (R$'_2$)-instances in $\P(\P_4)$ (\texttt{redsearch.py}; depth $2$ running), and the word surveys are empty: no identity of length $\le3$ in $a,b,c,x,y$, nor of length $\le5$ in $a,c,x,y$, holds on all (W$_2$)-instances beyond the consequences of the hypotheses and of the conclusion (\texttt{ids\_w2.log}, \texttt{ids\_w2\_long.log}); for (R$'_2$) the only non-trivial algebraic facts are $cy=yc$ (geometric: Prover9 finds no proof) and $bx_2=by$ (Prover9 proves it from $x_2\sim y$, $c\sim y$).

\emph{(c) Local statements along the chain.} In the two-point language (atoms: $u,v$ in the same block of $p$; $u\in\supp p$; $u,v$ joined in $\varepsilon_p\vee\varepsilon_q$; $u,v$ joined in $\varepsilon_a\vee\varepsilon_c\vee\varepsilon_p$; $p,q\in\{a,c,x\}$), the clauses with $\le4$ literals valid for $x$ with $c\sim x$, $\theta_x(a)=\theta_x(c)$, and still valid for the second and third element of a chain, are (modulo trivialities) a small join--meet calculus (\texttt{inv2.py}, \texttt{inv4.py}): Lemma~A, $\varepsilon_c\le\varepsilon_a\vee\varepsilon_x$; the meet condition $(\varepsilon_a\vee\varepsilon_c)\wedge(\varepsilon_c\vee\varepsilon_x)=\varepsilon_c$; two points of $\supp c$ (or of $\supp a$) in one class of $\varepsilon_a\vee\varepsilon_c\vee\varepsilon_x$ are joined in $\varepsilon_c\vee\varepsilon_x$ (resp.\ $\varepsilon_a\vee\varepsilon_x$); and $\supp a\cap\supp x\subseteq\supp c$. With five literals one finds moreover that every block of $x$ containing a point of $\supp a\setminus\supp c$ is an $\N(x)$-block contained in an $\N(c)$-block (\texttt{five.py}). All these hold at every chain position tested ($\le3$, $m\le5$). But they are not the invariant: in $\P(\P_4)$ the conjunction of all $67$ ($78$ with 3-joins) surviving clauses is preserved by one more link, while in $\P(\P_5)$ the propagation fixed point loses Lemma~A itself; and no set of such local statements about $y$, together with $y\sim b$, $b\bowtie c$, implies (G1) or (G2) ($m=5$). The reason is exhibited by the configuration
\[
a=[0\,2][1\,3](4),\quad c=(0)(1)[2][3\,4],\quad y=(0\,1)[2\,4][3],\quad b=(0\,1)[2](3)[4]
\]
in $\P_5$: it satisfies all two-point statements above for $y$, $y$ identifies $a$ with $c$, $y\sim b$, $a\bowtie c\bowtie b$, (G1) fails (the block $\{3,4\}$ of $c$ is split by $\varepsilon_a\vee\varepsilon_b$), and there is no projection $x$ with $c\sim x\sim y$, $\theta_x(a)=\theta_x(c)$ in $\P_5$, $\P_6$ or $\P_7$ (\texttt{findx.py}): the $\N(y)$-block $\{0,1\}$ is the union of two $\N(c)$-blocks containing the points $0\in A_1$, $1\in A_2$ of $\supp a\setminus\supp c$ from different classes $\Gamma_i$, which no commuting identifying $x$ can produce (a $\T(x)$-block $\{0\}$ would make $\theta_x(c)$ non-transversal and $\theta_x(a)$ transversal there). The $\T(y)$-block $\{2,4\}$ straddling $B_1=\{2\}$ and $B_2=\{4\}$ is exactly the obstruction of Section~\ref{subsec:chain}(h). The classes of $\varepsilon_y\vee\varepsilon_c$ at chain position $2$ are of all shapes (\texttt{classtypes.py}: classes meeting $\supp y$ and $\supp c$ in proper subsets, several $\N(c)$-blocks in one class, several $\Gamma_i$ per class), so there is no analogue of the four types of Lemma~\ref{lem:PX}(b) for $(y,c)$; the invariant carried along the chain must be a statement about chains of blocks, not about pairs of points.

\emph{(d) Consequence for the shape of the proof.} What the chain provides at its end $y$, and what is missing in the configuration above, is a statement of the kind of Section~\ref{subsec:chain}(g): \emph{two $\varepsilon_b$-blocks meeting $\supp y$ that are consecutive along a chain in $\varepsilon_a\vee\varepsilon_b$ (all intermediate $\varepsilon_b$-blocks outside $\supp y$) lie in one class of $\varepsilon_b\vee\varepsilon_c$}. Its proof in (h) was blocked at the conversion of a $y$--$c$ chain into a $b$--$c$ chain; the verified two-point consequence of the $x$-link ``$u,v$ joined in $\varepsilon_a\vee\varepsilon_b$ and in $\varepsilon_b\vee\varepsilon_c\vee\varepsilon_y$ $\Rightarrow$ joined in $\varepsilon_b\vee\varepsilon_c$'' (\texttt{inv6\_m5.log}) performs this conversion abstractly: the intermediate part of such a chain lies in one class of $\varepsilon_y\vee\varepsilon_a$ (its $\varepsilon_b$-blocks lie in $\N(y)$-blocks), hence the two end blocks are met at $\T(y)$-blocks of one class of $\varepsilon_y\vee\varepsilon_a$, hence of one class of $\varepsilon_y\vee\varepsilon_c$ ($y$ identifies), hence the end blocks are joined in $\varepsilon_b\vee\varepsilon_c\vee\varepsilon_y$ and, by the quoted statement, in $\varepsilon_b\vee\varepsilon_c$. So (g) reduces to that one statement, which in turn needs the chain.

\emph{(e) The statement (T$_n$).} The quoted statement is a meet condition, and it is the right chain statement: writing $p\wedge q$ for the meet of partitions,
\[
\text{(T$_n$)}\qquad a\bowtie c,\quad c\sim x_1\sim\dots\sim x_n,\quad \theta_{x_i}(a)=\theta_{x_i}(c),\quad x_n\sim b\ \Longrightarrow\ (\varepsilon_a\vee\varepsilon_b)\wedge(\varepsilon_b\vee\varepsilon_c\vee\varepsilon_{x_n})\le\varepsilon_b\vee\varepsilon_c
\]
is verified for $n\le3$ (\texttt{t1.py}; the cold reader of 11~Oct rechecked it exhaustively in $\P(\P_4)$ and on $10^7$ sampled instances in $\P(\P_5)$), with $b$ an \emph{arbitrary} projection commuting with $x_n$: unlike (G1), (G2) and $a\bowtie b$ it does not use $b\bowtie c$, and every other hypothesis is needed (\texttt{meets.py}). (An earlier version of these notes stated (T$_n$) with equality; the reverse inclusion $\varepsilon_b\vee\varepsilon_c\le\varepsilon_a\vee\varepsilon_b$ is (G1), which fails without $b\bowtie c$, e.g.\ for $a=[0](1\,2\,3)$, $c=[0\,1\,2\,3]$, $x=(0\,1\,2\,3)$, $b=[0][1\,2\,3]$ in $\P_4$.) With $b\bowtie c$ it gives (g) as explained in (d). Other verified meet identities under (W$_2$) are $(\varepsilon_a\vee\varepsilon_b)\wedge(\varepsilon_b\vee\varepsilon_y)=\varepsilon_b$ (this one uses $b\bowtie c$), $(\varepsilon_a\vee\varepsilon_c)\wedge(\varepsilon_b\vee\varepsilon_c\vee\varepsilon_y)=\varepsilon_c$, $(\varepsilon_a\vee\varepsilon_c\vee\varepsilon_y)\wedge(\varepsilon_b\vee\varepsilon_c\vee\varepsilon_y)=\varepsilon_c\vee\varepsilon_y$ and $(\varepsilon_a\vee\varepsilon_b\vee\varepsilon_y)\wedge(\varepsilon_a\vee\varepsilon_c\vee\varepsilon_y)=\varepsilon_a\vee\varepsilon_y$. Since (T$_n$) quantifies over all commuting partners $b$ of the chain end, it is a candidate for the invariant itself (``$x_n$ is a chain end'' $:=$ (T$_n$) for all $b\sim x_n$); but the induction step cannot be tested as such, and its finite approximations fail: (T) for $x_{n-1}$ with the partners $x_n$, $\theta_{x_{n-1}}(b)$ and $x_{n-1}x_n$, together with Lemma~A and the meet condition for both, does not imply (T) for $x_n$ with partner $b$ ($m\le5$, \texttt{tind.py}, \texttt{tind2.py}). Nor does (T$_2$) follow from the local statements about $y$ of (c). For $n=1$, (T$_1$) is the mirror image of Steps~4--6 of the proof of Theorem~\ref{thm:R} without the matching $b\bowtie c$, and its proof there rests on the fact that $b$ and $c$ have the same blocks outside $\supp x$ (both commute with $x$); in the chain case this is exactly what fails ($c\not\sim x_n$), and what replaces it is the five-literal invariant of (c), that the blocks of $x_n$ meeting $\supp a\setminus\supp c$ lie inside $\N(c)$-blocks, together with whatever the chain says about the $\N(x_n)$-blocks that are unions of several $\N(c)$-blocks (the configuration of (c)). A hand proof of (T$_1$) along these lines, then of (T$_2$), was the plan; what was found instead is recorded in (f)--(h).

\emph{(f) Local approximations, and why they fail.} \emph{(Superseded by Section~\ref{subsec:count}: (A), (pAC), (aCP) and (G2$'$) are proved there at every chain position, as consequences of the exact invariant.)} Call the following five statements about a projection $y$ relative to $(a,c)$ the \emph{end invariants}: (A) $\varepsilon_c\le\varepsilon_a\vee\varepsilon_y$ (Lemma~A); (pAC) $\varepsilon_y\wedge(\varepsilon_a\vee\varepsilon_c)\le\varepsilon_c$; (aCP) $\varepsilon_a\wedge(\varepsilon_c\vee\varepsilon_y)\le\varepsilon_c$; (N2) if an $\N(y)$-block contains a point of $\supp a\setminus\supp c$, then its points outside $\supp c$ lie in one $\N(c)$-block; (G2$'$) two points of $\supp c$ in one class of $\varepsilon_a\vee\varepsilon_c\vee\varepsilon_y$ are joined in $\varepsilon_c\vee\varepsilon_y$. All five hold at every chain position ($\le3$, $m\le5$), and in $\P(\P_5)$ they suffice at the end: if $y$ identifies, $y\sim b$, $b\bowtie c$ and $y$ satisfies the five, then (G1) and (G2) hold. Likewise in $\P(\P_5)$ Lemma~A for the second chain element follows from $x\sim y$, both identifying, and three local statements about $x$ ((A), the meet condition, and ``an $\N(x)$-block with a point outside $\supp c$ lies in one $c$-block''). \emph{Both implications fail in $\P(\P_6)$} (\texttt{modules.py}), exactly as the $\P(\P_4)$-invariants of (c) failed in $\P(\P_5)$. The reason is (g) below: no property of the chain end expressible by finitely many local statements characterises chain ends, and any such ``decomposition'' verified in $\P(\P_m)$ for a fixed $m$ is an artefact of $m$. Of the twelve chain-valid local statements, all but Lemma~A do propagate along a link at $m=5$ (\texttt{propA.py}), but after (g) this is not to be trusted either. The only statements with evidential value are those verified under the full chain hypotheses; the local language is useful for \emph{finding} candidate lemmas (as in (h)), not for decomposing the proof.

\emph{(g) The component of $c$ has unbounded radius.} Let $\mathcal I(a,c)$ be the graph whose vertices are the projections identifying $a$ with $c$, with $y\sim z$ as edges; the chain ends are the vertices of the component of $c$ (joined to $c$), and the chain length is the distance. In $\P_4$ and $\P_5$ every vertex of the component is at distance $\le3$ from $c$, so that (W$_n$) would follow from (W$_3$) there; in $\P_6$ a component of radius $4$ exists ($a=(0\,2\,3\,4)(1)[5]$, $c=(0\,3)(1)(2)(4)[5]$, $234$ vertices, \texttt{stab6.py}, \texttt{r4.py}; a geodesic is $c\sim(0\,3)[1][2](4)(5)\sim[0](1\,2)[3](4)(5)\sim(0)(1\,2)(3)[4][5]\sim(0)(1\,2)(3)(4\,5)$, where the transversality is carried across the set one step at a time). Likewise the normal form $c\le x_1\sim x_2\ge z$ (first step upwards, last step downwards) that every vertex of a component admits in $\P_4$ and $\P_5$ fails for eight vertices of that component (\texttt{normalform3.py}, \texttt{nf6.py}). After the cold read of 11~Oct the question was open: the radius is $3$ in $\P_4$ and $\P_5$ and $4$ in $\P_6$, and the cold reader found no radius $5$ in $\P_7$ (consistent with Lemma~\ref{lem:radius} below, which gives radius $\ge k+1$ in $\P_{2k+1}$, i.e.\ $5$ first in $\P_9$) on $1572$ random matched pairs with small supports and on eight hand-chosen generalisations of the $\P_6$ example (e.g.\ $a=(0\,2\,3\,4\,5)(1)[6]$, $c=(0\,3)(1)(2)(4)(5)[6]$: component of $328$ identifiers, radius $4$; \texttt{review2/review-scripts/radius3.py}). The question mattered: if the radius were bounded by some $R$ in every $\P_X$, then (W$_R$) would imply every (W$_n$) in $\P_X$ and a uniform proof would reduce to one condition. It is unbounded:

\begin{lemma}[radius]\label{lem:radius}
Let $k\ge1$, $X=\{p_1,\dots,p_k,q_1,\dots,q_k,t\}$, and let $a,c,z$ be the projections of $\P_X$ with the single transversal block $[t]$ and the nontransversal blocks $(p_1\cdots p_k)(q_1)\cdots(q_k)$ for $a$, all singletons for $c$, and $(p_1\,q_1)\cdots(p_k\,q_k)$ for $z$. Then $a\bowtie c$, $z$ identifies $a$ with $c$, and $z$ is at distance exactly $k+1$ from $c$ in the identification graph: there is a chain $c\sim y_1\sim\dots\sim y_k\sim z$ of identifiers and no shorter one. In particular the radius of $\mathrm{Comp}(c)$ is unbounded as $|X|\to\infty$.
\end{lemma}

\begin{proof}
$a\bowtie c$ since both have the single transversal block $[t]$. A projection $y$ with transversal blocks $Y_1,\dots,Y_r$ identifies $a$ with $c$ iff the partitions of $\supp y$ induced by $\varepsilon_y\vee\varepsilon_a$ and by $\varepsilon_y\vee\varepsilon_c$ coincide and a part is transversal for $a$ iff it is for $c$; here a part is transversal iff it contains $t$ (for both), so only the partitions matter. For $z$ both partitions of $\supp z=\{t\}$ are trivial, so $z$ identifies. Let $B=\{p_1,\dots,p_k\}$. Let $y_i$ ($1\le i\le k$) have the nontransversal blocks $(p_j\,q_j)$ for $j<i$, the transversal blocks $[p_i]$, $[q_i]$, $[t]$, and singletons elsewhere. Then $y_i$ identifies: the part of $p_i$ is $\{p_i\}$ for both joins (via $a$ its class is $B\cup\{q_j:j<i\}$, which meets $\supp y_i$ only in $p_i$; via $c$ it is $\{p_i\}$), and likewise for $q_i$ and $t$. Consecutive members commute by Lemma~\ref{lem:PX}(b): from $y_i$ to $y_{i+1}$ the two transversal blocks $[p_i],[q_i]$ merge into the nontransversal block $(p_i\,q_i)$ (type (ii)), the singletons $(p_{i+1}),(q_{i+1})$ become transversal (type (iii)) and all other blocks are common; the same for $c\sim y_1$ and $y_k\sim z$. So $\operatorname{dist}(c,z)\le k+1$.

For the lower bound, note first that an identifier $y$ has at most one transversal block meeting $B$: two such blocks lie in one class of $\varepsilon_y\vee\varepsilon_a$ (through the block $(p_1\cdots p_k)$ of $a$), but in different classes of $\varepsilon_y\vee\varepsilon_c$, since every block of $c$ is a singleton and the only blocks of $y$ meeting $\supp y$ are its own transversal blocks. Now let $c=y_0\sim y_1\sim\dots\sim y_n=z$ be any chain of identifiers. Fix $i$ and let $j$ be the largest index such that $(p_i\,q_i)$ is not a block of $y_j$ ($j$ exists, since it is not a block of $c$, and $j<n$). Then $(p_i\,q_i)$ is a nontransversal block of $y_{j+1}$ which is not a block of $y_j$, so by Lemma~\ref{lem:PX}(b) it is a union of transversal blocks of $y_j$; hence $y_j$ has a transversal block containing $p_i$ and contained in $\{p_i,q_i\}$. Call $j=j(i)$. Since $y_{j(i)}$ has exactly one transversal block meeting $B$, and that block determines $i$, the indices $j(1),\dots,j(k)$ are distinct; they lie in $\{1,\dots,n-1\}$ ($y_0=c$ has no transversal block meeting $B$). Hence $n-1\ge k$.
\end{proof}

(Computer check: $\operatorname{dist}(c,z)=2,3,4$ for $k=1,2,3$ in $\P_3,\P_5,\P_7$, \texttt{radfam.py}; the $k=2$ case is the radius-$3$ phenomenon of $\P_5$, and the radius-$4$ component of $\P_6$ above is a variant with the piece $(0\,3)$ split instead of merged.) So no bounded number of links reaches every chain end, and a uniform proof of the family has to handle arbitrary $n$ (with the proviso that a chain end $x_n$ together with its partner $b$ might still admit a shorter chain to some other partner; the lemma concerns the component, not the conditions). (An earlier version drew the conclusion ``the invariant of an inductive proof cannot be a property of the chain end alone'' from the growth of the radius; that inference is wrong: unbounded radius says that \emph{reachability from $c$} is not expressible in a bounded local language, but an inductive invariant need not characterise the component. The weakest invariant, $\Phi(y)$ $:=$ ``$y$ identifies $a$ with $c$ and $a\bowtie b$ for every $b\sim y$ with $b\bowtie c$'', is a property of the chain end, and the validity of all (W$_n$) is exactly the statement that $\Phi$ is closed under the two kinds of steps described next.) What the growth does show is that the local searches of (c) could succeed in $\P(\P_4)$ and fail in $\P(\P_5)$ for a structural reason.

\emph{The order formulation (cold read of 11~Oct).} Since in $\P_X$ two projections commute iff they have a common lower bound (\eqref{eq:Q1} and its converse), and since the set $\mathrm{Id}(a,c)=\{y:\theta_y(a)=\theta_y(c)\}$ of identifiers is a down-set ($\theta_z(a)=\theta_z(\theta_y(a))$ for $z\le y$), a chain $c\sim x_1\sim\dots\sim x_n$ is a zigzag $c\ge m_0\le x_1\ge m_1\le x_2\ge\dots$ in the poset $\mathrm{Id}(a,c)$, and the chain ends are the vertices of the connected component $\mathrm{Comp}(c)$ of $c$ in its comparability graph. Hence the validity of all (W$_n$) in $\P_X$ is the single statement
\[
\text{(W$_\infty$)}\qquad \text{for every } y\in\mathrm{Comp}(c) \text{ and every } b\bowtie c \text{ with } b\sim y:\ a\bowtie b,
\]
and an inductive proof splits into a \emph{down} step ($\Phi(y)$, $z\le y$ $\Rightarrow$ $\Phi(z)$) and an \emph{up} step ($\Phi(z)$, $z\le w\in\mathrm{Id}(a,c)$ $\Rightarrow$ $\Phi(w)$), of which only the up step has content; Lemma~A is the up step for the statement $\varepsilon_c\le\varepsilon_y\vee\varepsilon_a$ (the down step $z\le x$ is the easy case noted in Section~\ref{subsec:chain}(f)). For a finite projection algebra $P$, (W$_\infty$) is equivalent to (W$_{|P|}$) (a chain with a repeated element can be shortened), so the chain family is a decidable parametrised set of conditions, and the brute-force test of (W$_\infty$) for a matched pair $(a,c)$ in $\P(\P_m)$ (compute the component, test every $b$) checks all $n$ at once; the cold reader ran it for every matched pair of $\P(\P_4)$ and $\P(\P_5)$ and for $2863$ random pairs of $\P(\P_6)$ with no violation (\texttt{winf.py}). A chain-free consequence worth isolating: if (W$_\infty$) holds and $a\ne c$ then $a\notin\mathrm{Comp}(c)$ as soon as some $\Gamma_i$ contains an $\N(a)$-block (take $\varepsilon_b=\varepsilon_a$ with that block in place of $A_i$ among the transversal blocks: then $b\sim a$, $b\bowtie c$ by Lemma~\ref{lem:PX}(c), and $\theta_a(b)\ne a$); since $\mathrm{Id}(a,c)$ is symmetric in $a$ and $c$, the components of $a$ and of $c$ are then distinct. This holds for all ordered matched pairs $a\ne c$ of $\P(\P_4)$ and $\P(\P_5)$ (\texttt{acomp.py}) and is the simplest statement that any proof of the family must contain. (Proved in Corollary~\ref{cor:acomp}, for every matched pair $a\ne c$.)

\emph{(h) Decomposition of (T$_2$).} Removing from the pool of verified local statements those that imply (T$_2$) by monotonicity alone, (T$_2$) follows (with $b\bowtie c$, $m\le5$) from the two statements $\varepsilon_y\wedge(\varepsilon_a\vee\varepsilon_b\vee\varepsilon_c)\le\varepsilon_b$ (two points of one $y$-block joined in $\varepsilon_a\vee\varepsilon_b\vee\varepsilon_c$ lie in one $b$-block; in particular two $\T(b)$-blocks inside one $\T(w)$-block $W_l$, $w=yb$, are never joined in $\varepsilon_a\vee\varepsilon_b\vee\varepsilon_c$, which is the (h)-obstruction in its sharpest form) and $(\varepsilon_a\vee\varepsilon_c)\wedge(\varepsilon_b\vee\varepsilon_c\vee\varepsilon_y)\le\varepsilon_c$ (\texttt{pqdec4.py}); the first needs $b\bowtie c$, the second does not. Both are again chain statements. (G1) and (G2) themselves are verified under the full (W$_2$) hypotheses in $\P(\P_6)$ as well ($83$ and $94$\,s; \texttt{m6targets.py}), so the closing route via Lemma~\ref{lem:closing} stands; what is missing is a proof of (G1) and (G2) from the chain, and by (f)--(g) it has to be an argument about chains of blocks under the full hypotheses, in the manner of Steps~2--8 of Theorem~\ref{thm:Rp}, or a new idea.

\subsection{The counting invariant, and the proof of the chain family (W$_n$)}\label{subsec:count}

\emph{(Added 11 Oct 2026, evening; files \texttt{pa8/red/delta.py}, \texttt{deltaW.py}, \texttt{thm6.py}, \texttt{rpinf.py}, \texttt{rpinv.py}, \texttt{rpred.py}.)} The cold reader's item (d) asked for a hand proof of $a\notin\mathrm{Comp}(c)$. The invariant that separates the components of $\mathcal I(a,c)$ is a \emph{count}: it is exact (constant on every component, not only on that of $c$), it is global rather than local, which is why the searches of (c) and (f) could not find it, and it carries enough information to close the end step of the chain, so that the whole family (W$_n$) follows. For equivalence relations $\pi,\rho$ on $X$ write $|\pi|$ for the number of classes of $\pi$ and $|\pi|_Y$ for the number of classes of $\pi$ meeting a set $Y\subseteq X$. For projections $p,q,x$ put
\[
\delta_{pq}(x)=|\varepsilon_x\vee\varepsilon_p|-|\varepsilon_x\vee\varepsilon_q|,\qquad
\delta'_{pq}(x)=|\varepsilon_x\vee\varepsilon_p|_{\supp p}-|\varepsilon_x\vee\varepsilon_q|_{\supp q}.
\]
Both depend on $x$ only through the partition $\varepsilon_x$.

\begin{lemma}[invariance]\label{lem:count}
Let $p,q$ be projections of $\P_X$, let $w$ identify $p$ with $q$ ($\theta_w(p)=\theta_w(q)$) and let $z\le w$. Then $\delta_{pq}(z)=\delta_{pq}(w)$ and $\delta'_{pq}(z)=\delta'_{pq}(w)$. Consequently $\delta_{pq}$ and $\delta'_{pq}$ are constant on every component of the identification graph $\mathcal I(p,q)$: if $x\sim y$ and $y$ identifies $p$ with $q$, then $\delta_{pq}(xy)=\delta_{pq}(y)$ and $\delta'_{pq}(xy)=\delta'_{pq}(y)$.
\end{lemma}

\begin{proof}
Put $S=\supp w$, $\rho_p=\varepsilon_w\vee\varepsilon_p$ and $\rho_q=\varepsilon_w\vee\varepsilon_q$. By Lemma~\ref{lem:PX}(a) the blocks of $\theta_w(p)$ inside $S$ are the sets $K\cap S$ for the classes $K$ of $\rho_p$ meeting $S$, transversal iff $K$ meets $\supp p$; so $\theta_w(p)=\theta_w(q)$ says that $\rho_p$ and $\rho_q$ induce the same partition $\sigma$ of $S$, and that for every part $G$ of $\sigma$ the $\rho_p$-class containing $G$ meets $\supp p$ iff the $\rho_q$-class containing $G$ meets $\supp q$. By Lemma~\ref{lem:PX}(a) again, every block of $z\le w$ is an $\N(w)$-block or a union of $\T(w)$-blocks, so $\varepsilon_z=\varepsilon_w\vee\mu$ for an equivalence relation $\mu$ all of whose nontrivial classes lie inside $S$. Now $\varepsilon_z\vee\varepsilon_p=\rho_p\vee\mu$, and its classes are of two kinds. A class of $\rho_p$ disjoint from $S$ is untouched by $\mu$ and remains a class. The classes of $\rho_p$ meeting $S$ are merged according to $\sigma\vee\mu|_S$: a path alternating $\rho_p$-steps and nontrivial $\mu$-steps has all its $\mu$-steps inside $S$, so between two consecutive $\mu$-steps it joins a point of $S$ to a point of $S$ within one $\rho_p$-class, i.e.\ within one part of $\sigma$; hence two $\rho_p$-classes meeting $S$ lie in one class of $\rho_p\vee\mu$ iff their parts of $\sigma$ lie in one class of $\sigma\vee\mu|_S$. So the classes of the second kind correspond bijectively to the classes of $\sigma\vee\mu|_S$, and such a class meets $\supp p$ iff one of its constituent $\rho_p$-classes does. The same description holds for $q$, with the same $\sigma$ and the same $\mu|_S$, and corresponding constituents (same part $G$) meet $\supp p$, resp.\ $\supp q$, simultaneously. Therefore
\[
|\varepsilon_z\vee\varepsilon_p|-|\varepsilon_z\vee\varepsilon_q|=\bigl(|\rho_p|-|\sigma|\bigr)-\bigl(|\rho_q|-|\sigma|\bigr)=\delta_{pq}(w),
\]
and the same computation restricted to the classes meeting $\supp p$, resp.\ $\supp q$ (the classes of the first kind meeting $\supp p$ are the $\rho_p$-classes meeting $\supp p$ but not $S$, whose number is $|\rho_p|_{\supp p}$ minus the number of parts $G$ whose $\rho_p$-class meets $\supp p$, and the latter number is the same for $q$) gives $\delta'_{pq}(z)=\delta'_{pq}(w)$. For the last assertion, $xy=yx$ is a projection with $\theta_y(xy)=yxyy=xy$, i.e.\ $xy\le y$; so along a chain in $\mathcal I(p,q)$ every link $x\sim y$ is a zigzag $x\ge xy\le y$ with both upper elements identifying, and the values are transported.
\end{proof}

\begin{corollary}[separation]\label{cor:acomp}
Let $a\bowtie c$ in $\P(\P_X)$ and $\kappa=|\varepsilon_a\vee\varepsilon_c|$. Then $\delta_{ac}(c)=\kappa-|\varepsilon_c|\le0\le|\varepsilon_a|-\kappa=\delta_{ac}(a)$ and $\delta'_{ac}(c)=\delta'_{ac}(a)=0$; the two values of $\delta_{ac}$ coincide iff $\varepsilon_a=\varepsilon_c$, i.e.\ iff $a=c$. Hence for $a\ne c$ the projections $a$ and $c$ lie in different components of $\mathcal I(a,c)$, and more precisely every $y\in\mathrm{Comp}(c)$ has $\delta_{ac}(y)=\kappa-|\varepsilon_c|$ and $\delta'_{ac}(y)=0$.
\end{corollary}

\begin{proof}
$c$ identifies $a$ with $c$ ($\theta_c(a)=c=\theta_c(c)$) and so does $a$. $\varepsilon_c\vee\varepsilon_a$ has $\kappa$ classes and $\varepsilon_c\vee\varepsilon_c=\varepsilon_c$ has $|\varepsilon_c|\ge\kappa$; symmetrically for $a$. Both inequalities are equalities iff $\varepsilon_a\le\varepsilon_c\le\varepsilon_a$, and if $\varepsilon_a=\varepsilon_c$ then each block is a class of $\varepsilon_a\vee\varepsilon_c$ by itself and $a\bowtie c$ forces it to be transversal in $a$ iff in $c$ (Lemma~\ref{lem:PX}(c)), so $a=c$. For $\delta'$: the classes of $\varepsilon_a\vee\varepsilon_c$ meeting $\supp a$ are the classes $\Gamma_i$, as many as the $\T(c)$-blocks; the rest is Lemma~\ref{lem:count}.
\end{proof}

This settles the referee's item (d) for \emph{every} matched pair $a\ne c$ (the restriction ``some $\Gamma_i$ contains an $\N(a)$-block'' in (g) above was only needed to derive the statement from (W$_\infty$)), and it explains the failure of the local invariants: the separating invariant is a global count. The value of the count has structural content, which we extract with a lemma on contracting partitions.

\begin{lemma}[contraction]\label{lem:contr}
Let $\pi\le\pi'$ and $\rho$ be equivalence relations on $X$. Then $|\pi|-|\pi\vee\rho|\ge|\pi'|-|\pi'\vee\rho|$, and the inequality is strict as soon as two $\pi$-classes lying in a common $\pi'$-class lie in a common class of $\pi\vee\rho$.
\end{lemma}

\begin{proof}
Let $H$ be the graph on the $\pi$-classes in which two classes are adjacent iff they meet a common $\rho$-class; its components are the classes of $\pi\vee\rho$, so $|\pi|-|\pi\vee\rho|=|V(H)|-c(H)$. Identify, one pair at a time, the $\pi$-classes lying in a common $\pi'$-class. Each identification lowers $|V|$ by one and $c$ by at most one, by exactly one iff the two vertices lie in different components at that moment; at the end the vertices are the $\pi'$-classes and the components are the classes of $\pi'\vee\rho$. So $|V|-c$ never increases, and and if some identification joins two vertices of one component we may perform that identification first, so that $|V|-c$ drops by one at that step and never increases afterwards.
\end{proof}

\begin{lemma}[what the value means]\label{lem:value}
Let $a\bowtie c$, $\kappa=|\varepsilon_a\vee\varepsilon_c|$, and let $\varepsilon$ be an equivalence relation on $X$ with $|\varepsilon\vee\varepsilon_a|-|\varepsilon\vee\varepsilon_c|=\kappa-|\varepsilon_c|$ (for instance $\varepsilon=\varepsilon_y$ with $y\in\mathrm{Comp}(c)$). Then
\begin{itemize}[nosep]
\item[(A)] $\varepsilon_c\le\varepsilon\vee\varepsilon_a$ \textup{(}every block of $c$ lies in one class of $\varepsilon\vee\varepsilon_a$\textup{)};
\item[(M)] $(\varepsilon\vee\varepsilon_c)\wedge(\varepsilon_a\vee\varepsilon_c)=\varepsilon_c$ \textup{(}two blocks of $c$ in one class $\Gamma$ of $\varepsilon_a\vee\varepsilon_c$ lie in different classes of $\varepsilon\vee\varepsilon_c$\textup{)}.
\end{itemize}
If moreover $|\varepsilon\vee\varepsilon_a|_{\supp a}=|\varepsilon\vee\varepsilon_c|_{\supp c}$, then
\begin{itemize}[nosep]
\item[(G2$'$)] two $\T(c)$-blocks lying in one class of $\varepsilon\vee\varepsilon_a\vee\varepsilon_c$ lie in one class of $\varepsilon\vee\varepsilon_c$.
\end{itemize}
\end{lemma}

\begin{proof}
Put $g_a=|\varepsilon\vee\varepsilon_a|-|\varepsilon\vee\varepsilon_a\vee\varepsilon_c|\ge0$ and $g_c=|\varepsilon\vee\varepsilon_c|-|\varepsilon\vee\varepsilon_a\vee\varepsilon_c|\ge0$; by Lemma~\ref{lem:contr} with $\pi=\varepsilon_c\le\pi'=\varepsilon\vee\varepsilon_c$ and $\rho=\varepsilon_a$, $g_c\le|\varepsilon_c|-\kappa$. The hypothesis reads $g_a-g_c=\kappa-|\varepsilon_c|$, so $g_a=g_c-(|\varepsilon_c|-\kappa)\le0$, whence $g_a=0$ and $g_c=|\varepsilon_c|-\kappa$. The first equation is (A); the second is the equality case of Lemma~\ref{lem:contr}, which is (M). For (G2$'$): by (A), $\varepsilon\vee\varepsilon_a=\varepsilon\vee\varepsilon_a\vee\varepsilon_c$, whose classes are unions of classes of $\varepsilon_a\vee\varepsilon_c$; such a class meets $\supp a$ iff it contains some $\Gamma_i$ iff it meets $\supp c$. So $|\varepsilon\vee\varepsilon_a\vee\varepsilon_c|_{\supp c}=|\varepsilon\vee\varepsilon_a|_{\supp a}=|\varepsilon\vee\varepsilon_c|_{\supp c}$. Every class of $\varepsilon\vee\varepsilon_a\vee\varepsilon_c$ meeting $\supp c$ contains at least one class of the finer relation $\varepsilon\vee\varepsilon_c$ meeting $\supp c$; equal counts force exactly one.
\end{proof}

In particular Lemma~A of (f) holds at \emph{every} position of \emph{every} chain, together with (pAC) and (aCP) (both are contained in (M), as $\varepsilon_y\le\varepsilon_y\vee\varepsilon_c$ and $\varepsilon_a\le\varepsilon_a\vee\varepsilon_c$) and (G2$'$); the obstacle described in (f)--(g) was that these statements are not closed under the up step by themselves, while the pair $(\delta,\delta')$ is. The end step now needs only the value of the two counts at the partition $\varepsilon_{x_nb}$.

\begin{theorem}[end step]\label{thm:endstep}
Let $a\bowtie c$ and $b\bowtie c$ in $\P(\P_X)$, let $\mu$ be an equivalence relation on $X$ whose nontrivial classes are unions of $\T(b)$-blocks, and put $\varepsilon=\varepsilon_b\vee\mu$. If
\[
|\varepsilon\vee\varepsilon_a|-|\varepsilon\vee\varepsilon_c|=|\varepsilon_a\vee\varepsilon_c|-|\varepsilon_c|
\qquad\text{and}\qquad
|\varepsilon\vee\varepsilon_a|_{\supp a}=|\varepsilon\vee\varepsilon_c|_{\supp c},
\]
then \textup{(G1)} and \textup{(G2)} hold, and hence $a\bowtie b$ by Lemma~\ref{lem:closing}.
\end{theorem}

\begin{proof}
Let $B_1,\dots,B_s$ be the $\T(b)$-blocks, $C_i$ the $\T(c)$-block matched with $B_i$ and $\Delta_i$ the class of $\varepsilon_b\vee\varepsilon_c$ containing them; by Lemma~\ref{lem:PX}(c), $\Delta_i\cap\supp b=B_i$, $\Delta_i\cap\supp c=C_i$, and the other classes of $\varepsilon_b\vee\varepsilon_c$ meet neither $\supp b$ nor $\supp c$. Write $\kappa=|\varepsilon_a\vee\varepsilon_c|$, $\kappa_{bc}=|\varepsilon_b\vee\varepsilon_c|$, $\pi=\varepsilon_a\vee\varepsilon_b\vee\varepsilon_c$ and $N=|\pi\vee\mu|$. Call the nontrivial classes of $\mu$ the \emph{groups}; let there be $g$ groups and let $s'$ be the number of $\T(b)$-blocks lying in groups, so that $\mu$ identifies $s'-g$ pairs of $\T(b)$-blocks ``independently''. Precisely, choose one point $p_i$ in each $B_i$ and a set $\Lambda$ of $s'-g$ pairs $(p_i,p_j)$ forming a spanning tree on the $\T(b)$-blocks of each group; then $\varepsilon_b\vee\mu=\varepsilon_b\vee\nu$ for the equivalence $\nu$ generated by $\Lambda$, and for every equivalence relation $\tau\ge\varepsilon_b$ the relation $\tau\vee\mu=\tau\vee\nu$ is obtained from $\tau$ by adding the pairs of $\Lambda$ one at a time, each addition lowering the number of classes by at most one. Hence
\begin{itemize}[nosep]
\item[(i)] $|\tau|-|\tau\vee\mu|\le s'-g$ for every $\tau\ge\varepsilon_b$, with equality only if no pair of $\Lambda$ joins two points already in one $\tau$-class; since $\Lambda$ may be chosen to contain any given pair $(p_i,p_j)$ with $B_i,B_j$ in one group, equality implies that two $\T(b)$-blocks of one group lie in different $\tau$-classes;
\item[(ii)] $|\varepsilon\vee\varepsilon_c|=|\varepsilon_b\vee\varepsilon_c\vee\nu|=\kappa_{bc}-(s'-g)$: the pairs of $\Lambda$ join the distinct classes $\Delta_i$, and a tree edge always joins two different pieces.
\end{itemize}
By Lemma~\ref{lem:value}(A), $\varepsilon\vee\varepsilon_a=\varepsilon\vee\varepsilon_a\vee\varepsilon_c=\pi\vee\mu$, so the first hypothesis and (ii) give
\begin{equation}\label{eq:N}
N=|\varepsilon\vee\varepsilon_a|=|\varepsilon\vee\varepsilon_c|+\kappa-|\varepsilon_c|=\kappa_{bc}-(s'-g)+\kappa-|\varepsilon_c|.
\end{equation}
Put $r_1=\kappa-|\pi|\ge0$ and $r_2=|\pi|-N\ge0$, so that $r_1+r_2=\kappa-N=(|\varepsilon_c|-\kappa_{bc})+(s'-g)$. Lemma~\ref{lem:contr} with $\pi_0=\varepsilon_c\le\varepsilon_a\vee\varepsilon_c$ and $\rho=\varepsilon_b$ gives $|\varepsilon_c|-\kappa_{bc}\ge\kappa-|\pi|=r_1$, and (i) with $\tau=\pi$ gives $r_2\le s'-g$. Hence both are equalities, $r_1=|\varepsilon_c|-\kappa_{bc}$ and $r_2=s'-g$, and the equality clause of (i) yields
\begin{equation}\label{eq:star}
\text{two $\T(b)$-blocks of one group lie in different classes of }\pi=\varepsilon_a\vee\varepsilon_b\vee\varepsilon_c.
\end{equation}

\emph{(G1).} Let $r_c=|\varepsilon_a\vee\varepsilon_b|-|\pi|\ge0$. By \eqref{eq:N} and (i) with $\tau=\varepsilon_a\vee\varepsilon_b$, $N=|\varepsilon_a\vee\varepsilon_b\vee\mu|\ge|\varepsilon_a\vee\varepsilon_b|-(s'-g)$, while $N=|\pi|-r_2=|\varepsilon_a\vee\varepsilon_b|-r_c-(s'-g)$. So $r_c\le0$, i.e.\ $|\varepsilon_a\vee\varepsilon_b|=|\varepsilon_a\vee\varepsilon_b\vee\varepsilon_c|$, which says $\varepsilon_c\le\varepsilon_a\vee\varepsilon_b$.

\emph{(G2).} Suppose $C_i$ and $C_j$, $i\ne j$, lie in one class of $\pi$. Since $\Delta_i$ and $\Delta_j$ lie inside classes of $\pi$, so do $B_i$ and $B_j$, in that same class. If $B_i,B_j$ lie in one group this contradicts \eqref{eq:star}. Otherwise $\Delta_i$ and $\Delta_j$ lie in different classes of $\varepsilon\vee\varepsilon_c=\varepsilon_b\vee\varepsilon_c\vee\mu$ (the classes of this relation meeting $\supp c$ are the unions of the $\Delta_i$ over the $\T(b)$-blocks of a group, together with the $\Delta_i$ of the $\T(b)$-blocks in no group, and these are distinct because $\Delta_i\cap\supp b=B_i$), but in one class of $\pi\le\pi\vee\mu=\varepsilon\vee\varepsilon_a\vee\varepsilon_c$; this contradicts (G2$'$) of Lemma~\ref{lem:value}, which applies by the second hypothesis.
\end{proof}

\begin{theorem}[the chain family]\label{thm:Wn}
\textup{(W$_n$)} holds in every partition monoid $\P_X$ for every $n\ge1$. More generally \textup{(W$_\infty$)} holds: if $a\bowtie c$, $b\bowtie c$, $y$ identifies $a$ with $c$, $y$ lies in the component of $c$ in $\mathcal I(a,c)$---or merely $\delta_{ac}(y)=\delta_{ac}(c)$ and $\delta'_{ac}(y)=0$---and $yb=by$, then $a\bowtie b$.
\end{theorem}

\begin{proof}
By Corollary~\ref{cor:acomp} every $y\in\mathrm{Comp}(c)$ has $\delta_{ac}(y)=\kappa-|\varepsilon_c|$ and $\delta'_{ac}(y)=0$; the end $x_n$ of a chain $c\sim x_1\sim\dots\sim x_n$ of identifiers is such a $y$. Let $m=yb=by$. Then $m\le y$, so by Lemma~\ref{lem:count} $\delta_{ac}(m)=\delta_{ac}(y)$ and $\delta'_{ac}(m)=\delta'_{ac}(y)$. By Lemma~\ref{lem:PX}(b) the classes of $\varepsilon_m=\varepsilon_y\vee\varepsilon_b$ are common $\N$-blocks, $\N(y)$-blocks which are unions of $\T(b)$-blocks, $\N(b)$-blocks which are unions of $\T(y)$-blocks, and classes inside $\supp y\cap\supp b$, which are unions of $\T(b)$-blocks; so $\varepsilon_m=\varepsilon_b\vee\mu$, where $\mu$ identifies the $\T(b)$-blocks lying in a common class. Theorem~\ref{thm:endstep} applies and gives $a\bowtie b$.
\end{proof}

\begin{theorem}[non-finite basis]\label{thm:nfbW}
The quasivariety generated by the projection algebras $\P(\P_X)$ of the finite partition monoids is not finitely based, and the class of finite projection algebras embeddable in some $\P(\P_X)$ is not the class of finite models of any finite set of quasi-identities: for every $n$, $\mathfrak P^{\mathrm{ch}}_n$ (Definition~\ref{def:chainwit}) is a minimal non-embeddable projection algebra of order $2n+5$ which satisfies every quasi-identity in at most $n+2$ variables valid in all $\P(\P_X)$.
\end{theorem}

\begin{proof}
Theorem~\ref{thm:nfb} with its two hypotheses discharged: the validity of all (W$_n$) is Theorem~\ref{thm:Wn}, and the embeddings of the maximal subalgebras are Proposition~\ref{prop:maxsub}. Minimality of $\mathfrak P^{\mathrm{ch}}_n$: it violates (W$_n$), so it embeds in no $\P(\P_X)$, and every proper subalgebra lies in a maximal one, which embeds.
\end{proof}

\emph{Computer checks.} (1) \texttt{delta.py}: for all $1418$ ordered matched pairs $a\ne c$ of $\P(\P_4)$, $\delta_{ac}$ and $\delta'_{ac}$ are constant on every component of $\mathcal I(a,c)$ (as proved); the graph has between $2$ and $15$ components, and in $334$ pairs some component other than $\mathrm{Comp}(c)$ has the same values, so the pair $(\delta,\delta')$ separates $a$ from $c$ but does not characterise $\mathrm{Comp}(c)$ (e.g.\ $a=[0\,1\,2\,3]$, $c=[0](1\,2\,3)$: $\mathrm{Comp}(c)=\{c,(0)(1\,2\,3)\}$, while $(0)(1)(2\,3)$ and $[0](1)(2\,3)$ form a component with the same values). This does not matter for the chain conditions, by Theorem~\ref{thm:Wn}. (2) \texttt{deltaW.py}: for all matched pairs $(a,c)$ of $\P(\P_4)$ and of $\P(\P_5)$, all $y$ identifying $a$ with $c$ with $\delta_{ac}(y)=\delta_{ac}(c)$, and all $b\bowtie c$ commuting with $y$ ($53\,239$, resp.\ $2\,635\,385$ instances): $a\bowtie b$ fails in $1\,992$, resp.\ $196\,670$ instances, in none of which $\delta'_{ac}(y)=0$. So the second count is needed, and the theorem is sharp in this sense. (3) \texttt{thm6.py}: for all $a\bowtie c\bowtie b$ of $\P(\P_4)$ and $\P(\P_5)$ and all groupings $\mu$ of the $\T(b)$-blocks with the first count right ($8\,383$, resp.\ $297\,879$ partitions $\varepsilon$), (G1) holds in every case, and with the second count also right ($7\,891$, resp.\ $261\,669$) (G2) and $a\bowtie b$ hold in every case, as proved. (4) \texttt{delta6.py}: in $\P(\P_6)$, Lemma~\ref{lem:count} on $2.45\cdot10^6$ random pairs $z\le w$ with $w$ identifying, and Theorem~\ref{thm:Wn} on $71\,057$ random instances, with no failure.

\emph{(R$''$) and K6 are consequences of (W$_2$).} The implications are purely algebraic: M.~Kinyon's Prover9 runs on \texttt{rpp\_vs\_chain.in} and \texttt{k6\_vs\_chain.in} (with (S), (R), (R$'$), (V) removed, as they only clog the search; $99$ and $41$ steps, found with hints, \texttt{pa8/p9/rpp\_vs\_chain.pf}, \texttt{k6\_vs\_chain.pf}) show:

\begin{theorem}\label{thm:Rpp}
In every projection algebra satisfying \eqref{eq:Q1} and \textup{(W$_2$)}, \textup{(R$''$)} and \textup{K6} hold. Hence both hold in every partition monoid $\P_X$, and candidate~$6=\mathfrak P^{\mathrm{ch}}_2$ is explained by \textup{(W$_2$)} alone.
\end{theorem}

The shape of the derivation is the one anticipated in Section~\ref{subsec:nine}: under the hypotheses of (R$''$) put $r=\theta_x(b)=\theta_y(a)$; then $r\le x$ and $r\le y$, so $xy=yx$ by \eqref{eq:Q1} (the only use of \eqref{eq:Q1}), and $\theta_a(y)=\theta_a(x)$; the work is to show $\theta_y(c)=\theta_x(b)$ (step 4572 of the proof), i.e.\ that $y$ identifies $a$ with $c$, which the prover does from (W$_2$) applied in intermediate instantiations; then $c\sim x\sim y\sim b$ with both $x$ and $y$ identifying is an instance of (W$_2$), which gives $a\bowtie b$. In Section~\ref{subsec:nine} the identity $\theta_y(a)=\theta_y(c)$ had been verified for $m\le7$ and was expected to need the geometry of $\P_X$; it needs only (W$_2$). For K6 the derivation is shorter and entirely transparent: with $u=\theta_a(b)$, $v=\theta_b(a)$ the hypotheses $\theta_y(z)=\theta_z(a)$, $\theta_c(z)=\theta_y(a)$, $\theta_z(b)=v$ give $\theta_y(a)\le c$, hence $yc=cy$ by \eqref{eq:Q1}, then $\theta_z(a)=\theta_z(c)$ and $\theta_y(a)=\theta_y(c)$, and $yz=zy$, $zb=bz$; so $c\sim y\sim z\sim b$ with $y,z$ identifying $a$ with $c$, and (W$_2$) closes. Of the cluster (R$''$), (W$_2$), K6 of Section~\ref{subsec:picture}(b) this settles (W$_2$)$\Rightarrow$(R$''$) and (W$_2$)$\Rightarrow$K6; the converses in projection algebras remain open (no separating algebra of order $\le8$). The same question for K0, K4, K5 against all proved conditions is set up in \texttt{pa8/p9/k0\_vs\_proved.in} etc. For K0 the answer is negative and known: Mace4 (M.K.) finds a model of order $8$ of \eqref{eq:Q1}, (S), (R), (R$'$), (V), (P$_1$), (P$_2$), (Q), (W$_2$)--(W$_4$) violating K0, and it is $\mathfrak P_8^1$ (isomorphism checked, \texttt{pa8/impl/mk8\_k0.json}); it also satisfies (Tri), (E$_6$), (R$''$), (R$'_2$), (R$'_3$), K4--K6 and violates exactly (E$_1$) and K0 (\texttt{evalcond2.py}). So K0 and (E$_1$) are not consequences of the proved conditions together with the chain family, and likewise for K4: M.K.'s Mace4 model of order $8$ violating K4 is $\mathfrak P_8^6$ (\texttt{mk8\_k4.json}), which satisfies (Tri) and everything else on the list and violates exactly (E$_6$) and K4. So $\mathfrak P_8^1$ and $\mathfrak P_8^6$ remain the order-$8$ candidates whose non-embeddability needs a proof of (E$_1$) or K0, resp.\ (E$_6$) or K4.

\emph{(E$_1$) is a consequence of K0 (added 13 Oct 2026).} The two conditions killing $\mathfrak P_8^1$ are not independent: with semantic guidance and hints, and two days of work, M.~Kinyon obtained a Prover9 proof of

\begin{theorem}\label{thm:K0E1}
In every projection algebra satisfying \eqref{eq:Q1}, \textup{(W$_2$)} and \textup{K0}, \textup{(E$_1$)} holds. Consequently \textup{(E$_1$)} holds in every $\P_X$ in which \textup{K0} holds, and a proof of \textup{K0} (five variables) makes $\mathfrak P_8^1$ a proved minimal non-embeddable algebra; \textup{(E$_1$)} (seven variables) need not be attacked separately. \textup{(K0 is proved in Theorem~\ref{thm:K0} below.)}
\end{theorem}

The main proof ($227$ steps, $701$\,s; \texttt{pa8/p9/k0\_implies\_e1.pf}) uses, besides \eqref{eq:Q1}, (W$_2$) and K0, five lemmas proved separately (\texttt{k0\_lemmas.pf}; notation $x*y=\theta_y(x)$, and the hypotheses of (E$_1$) as constants: $z\le b$, $\theta_a(b)=z$, $d\le b$, $\theta_x(c)=z$, $\theta_a(c)=a$, $\theta_b(c)=b$, $e\le x$, $e\le d$, $\theta_a(d)=a$, conclusion $z=a$): (L1) $\theta_{\theta_b(\theta_d(a))}(a)=z$ and (L3) $\theta_{\theta_d(a)}(z)=\theta_x(d)$ from \eqref{eq:Q1} alone; (L2) $\theta_d(\theta_c(a))=\theta_d(a)$ from \eqref{eq:Q1} and (R); (L4) $\theta_d(\theta_b(\theta_c(d)))=d$ and (L5) $\theta_d(c)=d$ from (W$_2$), (L1), (L3). Since (R) is a consequence of (W$_2$), the whole derivation lives in \eqref{eq:Q1}$+$(W$_2$)$+$K0. The chain condition is needed: \eqref{eq:Q1}$+$K0 alone has a model of order $10$ violating (E$_1$) (\texttt{pa8/impl/mk10\_k0\_e1.json}), which violates (R); with (R) added, Mace4 finds no such model of order $\le18$. \emph{(Added 13 Oct, evening.)} The same holds for the other order-$8$ pair, and more simply: M.~Kinyon's Prover9 proof \texttt{k4\_implies\_e6.pf} ($72$ steps, $1700$\,s) derives (E$_6$) from \eqref{eq:Q1} and K4 alone, no chain condition needed. So the order-$8$ candidates $\mathfrak P_8^1$, $\mathfrak P_8^6$ are settled by K0, resp.\ K4, and the hand proofs of K0 and K4 in $\P_X$ are the remaining targets at order $8$.

\emph{The hybrid family.} (R$'_n$) is not covered: its end is $\theta_{x_n}(b)=\theta_y(b)$ with $ya=ay$, not $x_n\sim b$. In the order formulation it reads (R$'_\infty$): $x\in\mathrm{Comp}(c)$, $ya=ay$, $\theta_x(b)=\theta_y(b)$, $b\bowtie c$ $\Rightarrow$ $a\bowtie b$, which holds for all $146\,233$ instances in $\P(\P_4)$ (\texttt{rpinf.py}). As in the proof of Theorem~\ref{thm:Rp}, one may assume $x\le y$ (replace $x$ by $xy$, which lies in $\mathrm{Comp}(c)$ since $xy\le x$; here $x\sim y$ by \eqref{eq:Q1}, both lying above $\theta_x(b)=\theta_y(b)$). Two reductions do \emph{not} work: there need be no $z\in\mathrm{Comp}(c)$ with $z\sim b$ (so (W$_\infty$) cannot be invoked directly; $1\,528$ instances in $\P(\P_4)$, e.g.\ $a=[0\,1\,2\,3]$, $c=[0](1\,2\,3)$, $b=[0](1)(2\,3)$, $x=c$, $y=[0][1\,2\,3]$), and there need be no $z$ in the component of $c$ in $\mathcal I(b,c)$ with $z\sim a$ ($5\,130$ instances, \texttt{rpred.py}). What does work, in $\P(\P_4)$: the projection $\theta_b(x)=\theta_b(y)$ (the two coincide because $\varepsilon_x\vee\varepsilon_b=\varepsilon_y\vee\varepsilon_b$ and a class meeting $\supp y$ meets $\supp x$, Step~4 of Theorem~\ref{thm:Rp}) lies below $b$, its partition is $\varepsilon_b\vee\mu$ for the grouping $\mu$ of the $\T(b)$-blocks by the classes of $\varepsilon_x\vee\varepsilon_b$, and it has $\delta_{ac}=\delta_{ac}(c)$ and $\delta'_{ac}=0$ in all $86\,221$ instances of $\P(\P_4)$ (\texttt{rpinv.py}); by Theorem~\ref{thm:endstep} this would give $a\bowtie b$. So (R$'_n$) for all $n$ reduces to: \emph{if $x\in\mathrm{Comp}(c)$, $x\le y$, $ya=ay$, $\theta_x(b)=\theta_y(b)$ and $b\bowtie c$, then the two counts at $\varepsilon_{\theta_b(x)}$ take the values they take at $c$.} Lemma~\ref{lem:count} gives this when $x\sim b$ (then $\theta_b(x)=xb\le x$); the general case must use $ya=ay$, and is the next target.

\emph{K4: reductions, and the case $z=b$ (15 Oct).} Write the hypotheses of K4 as $a\bowtie c\bowtie b$, $aya=a$, $b\le z$, $yc=cy$, $\theta_z(y)=\theta_b(a)=:v$, $\theta_z(a)=\theta_a(b)=:u$, and put $r=\theta_y(a)$, $r'=\theta_y(b)$, $m=\rk a$, $n=\rk u$. In every projection algebra with \eqref{eq:Q1}: $u\le a$ and $u\le z$, so $az=za=u$; $v\le b\le z$, so $\theta_b(y)=\theta_b(\theta_z(y))=\theta_b(v)=v$; $\theta_u(y)=u$; $a\bowtie r$, $v\bowtie r'$ and $u\bowtie v$ (each of the form $\theta_p(q)\bowtie\theta_q(p)$); $\theta_z(c)=b$ (M.~Kinyon, Prover9, from \eqref{eq:Q1} alone, \texttt{pa8/p9/k4.pf1}); and, since $r\le y$, $\theta_r(b)=\theta_r(r')=\theta_r(z)=\theta_y(\theta_a(r'))$. In $\P_X$ the conclusion $a\bowtie b$ is equivalent to $n=m$, because $u\le a$ and $v\le b$ have rank $n$, and a projection below a projection of the same rank equals it. \emph{Rank lemma: if $\theta_r(b)=r$ then $a\bowtie b$.} Indeed $\theta_r(r')=r$ says (Lemma~\ref{lem:PX}(a)) that every class of $\varepsilon_r\vee\varepsilon_{r'}$ meeting $\supp r$ contains exactly one $\T(r)$-block and meets $\supp r'$, so $\rk r'\ge\rk r=m$, while $\rk r'=\rk v=n\le m$. So K4 follows from: \emph{the hypotheses imply $\theta_a(\theta_y(b))=a$}, i.e.\ $aybya=a$ (then $\theta_r(b)=\theta_y(\theta_a(r'))=\theta_y(a)=r$). In fact the rank lemma is not needed: M.~Kinyon's Prover9 proof \texttt{pa8/p9/k4.pf2} ($30$ steps) derives $a\bowtie b$ from the projection-algebra axioms and the seven equations $\theta_c(a)=c$, $aya=a$, $b\le z$, $\theta_z(y)=\theta_b(a)$, $\theta_z(a)=\theta_a(b)$, $\theta_z(c)=b$, $\theta_a(\theta_y(b))=a$ alone, so that in every projection algebra with \eqref{eq:Q1} the condition K4 is \emph{equivalent} to: its hypotheses imply $\theta_a(\theta_y(b))=a$. This identity fails in the K4-instance of $\mathfrak P_8^6$, so it is not an algebraic consequence of the hypotheses and the proved conditions, and a proof must use the geometry; candidate intermediate statements, each verified in $\P(\P_4)$ but not distinguishable there from $a\bowtie b$ itself, are: every class of $\varepsilon_a\vee\varepsilon_b$ meeting $\supp a$ contains one $\T(a)$-block; $\varepsilon_z\le\varepsilon_a\vee\varepsilon_c$; no $\T(z)$-block meets two $\T(a)$-blocks; every $\T(b)$-block $B_i$ lies in the class $E_i$ of $\varepsilon_a\vee\varepsilon_y$ containing $A_i$; no $y$-block meets two $\T(b)$-blocks (\texttt{red/k4rec*.py}).

What is proved about the geometry (the proofs are in Theorem~\ref{thm:K4}; notation of Theorem~\ref{thm:R}; $D_j$ the classes of $\varepsilon_a\vee\varepsilon_b$ meeting $\supp a$, $K_j=D_j\cap\supp a$, $S_j=D_j\cap\supp b$, $I_j=\{i:A_i\subseteq K_j\}$): $a\sim z$ and $\supp a\subseteq\supp z$; the $\N(z)$-blocks are exactly the $\N(a)$-blocks and the $\N(b)$-blocks outside $\supp z$ (``outer blocks''), every $\N(a)$-block inside $\supp z$ is a union of $\T(z)$-blocks, and the $K_j$ are the classes of $\varepsilon_a\vee\varepsilon_z$ that are unions of $\T(a)$-blocks (so each $D_j\subseteq\supp z$); $\varepsilon_z\vee\varepsilon_c=\varepsilon_b\vee\varepsilon_c$ and $\Delta_i\cap\supp z=B_i$; $\varepsilon_z\vee\varepsilon_y=\varepsilon_b\vee\varepsilon_y$, and a class of it meeting $\supp z$ meets it in some $S_j$ (and then meets $\supp y$) or in a single $\N(b)$-block (and then not); the classes of $\varepsilon_a\vee\varepsilon_y$ meeting $\supp a$ are $E_1,\dots,E_m$ with $E_i\cap\supp a=A_i$, and the class of $\varepsilon_a\vee\varepsilon_y\vee\varepsilon_z$ containing $K_j$ is exactly $\bigcup_{i\in I_j}E_i$ (Lemma~\ref{lem:count} for $u\le a$, in $\mathcal I(a,c)$ and in $\mathcal I(y,c)$, gives $|\varepsilon_a\vee\varepsilon_c|-|\varepsilon_a\vee\varepsilon_b\vee\varepsilon_c|=m-n=|\varepsilon_a\vee\varepsilon_y|-|\varepsilon_a\vee\varepsilon_y\vee\varepsilon_z|$, and $\theta_u(y)=u$ leaves no room for other merges); and, with $w=yc\le c$ in the component of $c$ of both $\mathcal I(a,c)$ and $\mathcal I(b,c)$, the classes of $\varepsilon_a\vee\varepsilon_c\vee\varepsilon_y$ meeting $\supp a$ and those of $\varepsilon_b\vee\varepsilon_c\vee\varepsilon_y$ meeting $\supp b$ are both in bijection with the classes of $\varepsilon_y\vee\varepsilon_c$ containing $\T(c)$-blocks (the $\delta'$ half of Lemma~\ref{lem:count}). If $n<m$, some $I_j$ contains two indices, say $1,2$; then $S_j\supseteq B_1\cup B_2$ is connected in $\varepsilon_z\vee\varepsilon_y$, so some $y$-block $Y$ meets $B_1$ and $B_2$, and by Lemma~\ref{lem:PX}(b) for $y\sim c$ the $c$-blocks at these points are transversal: $Y\cap B_1\subseteq C_1$, $Y\cap B_2\subseteq C_2$. \emph{The case $z=b$} (the shape of the K4-instance of $\mathfrak P_8^6$) is now short, granted Step~5 of Theorem~\ref{thm:K4} (which excludes a $\T(a)$-block inside an $\N(b)$-block, a possibility this argument does not see): then $\theta_a(b)=\theta_b(a)$, so $a\sim b$, $K_j$ is a union of $\T(b)$-blocks and $B_1\cup B_2\subseteq K_j\subseteq\supp a$; a point of $C_1\cap\supp a$ lies in $A_1$ ($\Gamma_1$ contains no other $\T(a)$-block), so $Y$ meets $A_1$ and $A_2$ and $E_1=E_2$, contradicting $aya=a$. So K4 holds in every $\P_X$ when $z=b$. In general $B_1\not\subseteq\supp a$ and the link from $Y\cap B_1\subseteq C_1$ to $A_1$ runs through $\N$-blocks of $\Gamma_1$; the general case is Theorem~\ref{thm:K4} in the next subsection, proved on 16 Oct by tracking three labels of each point along the connection of $B_1$ with $B_2$. (SAT had confirmed K4 in $\P(\P_m)$ for $m\le7$: kissat, M.~Kinyon, 15 Oct, the instance denying $\theta_a(b)=a$ at $m=7$ is unsatisfiable, $12\,600$ variables, $3$ hours; the other half $\theta_b(a)=b$ is equivalent to it in $\P_X$ since $u\bowtie v$ gives $\rk u=\rk v$ and $\rk a=\rk b$.)

\subsection{K4 holds in every partition monoid}\label{subsec:K4}

\begin{theorem}\label{thm:K4}
Let $a,b,c,y,z$ be projections of $\P_X$ with $a\bowtie c$, $c\bowtie b$, $aya=a$, $b\le z$, $yc=cy$, $\theta_z(y)=\theta_b(a)$ and $\theta_z(a)=\theta_a(b)$. Then $a\bowtie b$. Consequently \textup{(E$_6$)} holds in every $\P_X$ (Section~\ref{subsec:count}, \texttt{k4\_implies\_e6.pf}), $\mathfrak P_8^6$ embeds in no $\P(\P_X)$, and $\mathfrak P_8^6$ is a minimal non-embeddable projection algebra (every projection algebra of order $\le7$ satisfying \eqref{eq:Q1} and \textup{(S)} other than $\mathfrak P_7$ embeds, Section~\ref{sec:comp}; the subalgebras of $\mathfrak P_8^6$ inherit \eqref{eq:Q1}, \textup{(S)} and \textup{(R)}, and $\mathfrak P_7$ violates \textup{(R)}).
\end{theorem}

\begin{proof}
We use the notation of Theorem~\ref{thm:R} ($A_i,C_i,B_i,\Gamma_i,\Delta_i$) and of the paragraph ``K4: reductions'' of Section~\ref{subsec:count}: $u=\theta_a(b)=\theta_z(a)$, $v=\theta_b(a)=\theta_z(y)$, $m=\rk a$, $n=\rk u$.

\smallskip\emph{Step 1: algebraic consequences} (valid in every projection algebra with \eqref{eq:Q1}). $u\le a$ and $u\le z$, so $az=za=u$ by \eqref{eq:Q1}; $v\le b\le z$; $\theta_b(y)=\theta_b(\theta_z(y))=\theta_b(v)=v$; $u\bowtie v$ (as $\theta_a(b)\bowtie\theta_b(a)$); $\theta_u(y)=\theta_u(\theta_a(y))=\theta_u(a)=u$ and $\theta_u(c)=u$ likewise; $\theta_z(c)=b$ (M.~Kinyon, Prover9, \texttt{pa8/p9/k4.pf1}, from \eqref{eq:Q1}, $b\bowtie c$, $b\le z$, $yc=cy$ and $\theta_z(y)=\theta_b(a)$). Finally $\theta_y(u)=\theta_y(z)$: the identity $\theta_p(\theta_{\theta_q(p)}(x))=\theta_p(\theta_q(x))$, valid in every projection algebra (it is the axiom $\theta_p\theta_q\theta_p\theta_q=\theta_p\theta_q$ with the inner $\theta_q\theta_p\theta_q$ rewritten by the axiom $\theta_q\theta_p\theta_q=\theta_{\theta_q(p)}$), gives, with $(p,q)=(y,b)$ and $(p,q)=(y,z)$, $\theta_y(\theta_b(x))=\theta_y(\theta_v(x))=\theta_y(\theta_z(x))$ for every $x$ (both $\theta_b(y)$ and $\theta_z(y)$ equal $v$); with $x=a$ this is $\theta_y(v)=\theta_y(u)$, and $\theta_y(v)=\theta_y(\theta_z(y))=\theta_y(z)$ by the axiom $\theta_p(\theta_q(p))=\theta_p(q)$.

\smallskip\emph{Step 2: the blocks of $a,b,z$.} Since $a\sim z$, the classes of $\varepsilon_a\vee\varepsilon_z$ are of the four types of Lemma~\ref{lem:PX}(b) and $\varepsilon_u=\varepsilon_a\vee\varepsilon_z$: the $\T(u)$-blocks are the classes $K_1,\dots,K_n$ of type (i), each a union of $\T(a)$-blocks and of $\T(z)$-blocks; the $\N(u)$-blocks are the $\N(a)$-blocks inside $\supp z$ (type (iii), unions of $\T(z)$-blocks), the $\N(z)$-blocks $Z^{(1)},\dots,Z^{(t)}$ that are unions of $\T(a)$-blocks (type (ii)), and the common $\N$-blocks of $a$ and $z$ (type (iv)). We call $K_1,\dots,K_n,Z^{(1)},\dots,Z^{(t)}$ the \emph{generalised $K$'s}; every $\T(a)$-block lies in exactly one of them, and $\rk u=n$. Every block of $a$ lies inside one block of $z$ (the $\N(a)$-blocks outside $\supp z$, which we call \emph{outer}, and the $\T(a)$-blocks inside some $Z^{(k)}$) or inside $\supp z$. Since $b\le z$, every block of $b$ is an $\N(z)$-block or a union of $\T(z)$-blocks, $\varepsilon_z\le\varepsilon_b$, and $B_i\subseteq\supp z$. By $u\bowtie v$, every class of $\varepsilon_u\vee\varepsilon_v=\varepsilon_a\vee\varepsilon_b$ meeting $\supp u\cup\supp v$ contains exactly one $K_j$ and exactly one $\T(v)$-block; write $D_j$ for the class containing $K_j$, $S_j=D_j\cap\supp b$ for its $\T(v)$-block (a union of $\T(b)$-blocks), and $I_j=\{i:A_i\subseteq K_j\}$. By Lemma~\ref{lem:PX}(a) for $v=\theta_b(a)$, the blocks of $v$ inside $\supp b$ are the sets $D\cap\supp b$ for the classes $D$ of $\varepsilon_a\vee\varepsilon_b$ meeting $\supp b$: the $S_j$, and the $\N(v)$-blocks $D\cap\supp b$ for the classes $D$ meeting $\supp b$ but not $\supp a$; the other $\N(v)$-blocks are the $\N(b)$-blocks. In particular every block of $v$ inside $\supp z$ lies in a single class of $\varepsilon_a\vee\varepsilon_b$, which we denote $D(\cdot)$; and a $Z^{(k)}$ is a block of $b$ which is a union of blocks of $a$, hence a class of $\varepsilon_a\vee\varepsilon_b$ by itself.

\smallskip\emph{Step 3: the blocks of $c$ and $y$ relative to $z$.} By $\theta_z(c)=b$ and Lemma~\ref{lem:PX}(a), the blocks of $b$ inside $\supp z$ are the sets $K\cap\supp z$ for the classes $K$ of $\varepsilon_z\vee\varepsilon_c$ meeting $\supp z$; hence $\varepsilon_b\le\varepsilon_z\vee\varepsilon_c$, so $\varepsilon_z\vee\varepsilon_c=\varepsilon_b\vee\varepsilon_c$ and $\Delta_i\cap\supp z=B_i$: every block of $c$ meets $\supp z$ inside a single block of $b$, and the only $\T(z)$-blocks in $\Delta_i$ are those of $B_i$. By $\theta_z(y)=v$ and Lemma~\ref{lem:PX}(a), every class $K$ of $\varepsilon_z\vee\varepsilon_y$ meeting $\supp z$ meets it in exactly one block of $v$, which is a $\T(v)$-block $S_j$ if $K$ meets $\supp y$ and an $\N(v)$-block otherwise. We call a class of $\varepsilon_z\vee\varepsilon_y$ \emph{bright} if it contains a $\T(y)$-block and \emph{dark} otherwise; for a dark class $K$ meeting $\supp z$ put $D(K)=D(K\cap\supp z)$. Since every block of $b$ inside $\supp z$ lies in a block of $v$, $\varepsilon_b\le\varepsilon_z\vee\varepsilon_y$ and $\varepsilon_z\vee\varepsilon_y=\varepsilon_b\vee\varepsilon_y$. By $yc=cy$ and Lemma~\ref{lem:PX}(b), a block of $c$ of type (iv) is a block of $y$ and one of type (iii) lies in an $\N(y)$-block, so these lie in a single class of $\varepsilon_z\vee\varepsilon_y$; a block of $c$ of type (ii) is a union of $\T(y)$-blocks and one of type (i) lies in $\supp y$, so these lie in bright classes. Hence \emph{a block of $c$ meeting a dark class lies inside it}. Also, every block of $y$ meeting an $\N(c)$-block lies inside it (types (ii), (iv)), and a block of $y$ meeting no $\N(c)$-block is contained in $\bigcup_iC_i$ (types (i), (iii)).

\smallskip\emph{Step 4: two counts.} Put $\Omega=\varepsilon_u\vee\varepsilon_c=\varepsilon_a\vee\varepsilon_b\vee\varepsilon_c$ (by Step~3 and $\varepsilon_z\le\varepsilon_b$). Lemma~\ref{lem:count} for the pair $(a,c)$, which $a$ identifies, and $u\le a$ gives, in its $\delta'$ half, $|\varepsilon_u|_{\supp a}=|\Omega|_{\supp c}$; the left side is $n+t$. By $\theta_u(c)=u$ every class of $\Omega$ meeting $\supp u$ contains exactly one $K_j$, and it meets $\supp c$; so exactly $t$ classes of $\Omega$ meet $\supp c$ without meeting $\supp u$, and each of them contains some $C_i$, hence $A_i$, hence a $Z^{(k)}$ (as $A_i\not\subseteq\supp u$). Since there are $t$ blocks $Z^{(k)}$, \emph{every class of $\Omega$ meeting $\supp c$ contains exactly one generalised $K$, and a class containing some $Z^{(k)}$ contains no $K_j$.} Next, by $aya=a$ the classes of $\varepsilon_a\vee\varepsilon_y$ meeting $\supp a$ are $E_1,\dots,E_m$ with $E_i\cap\supp a=A_i$ and each $E_i$ meets $\supp y$. Lemma~\ref{lem:count} for $u\le a$ with the pairs $(a,c)$ and $(y,c)$ (both identified by $a$, as $\theta_a(y)=a=\theta_a(c)$) gives $|\varepsilon_u|-|\Omega|=|\varepsilon_a|-|\varepsilon_a\vee\varepsilon_c|$ and $|M|-|\Omega|=|\varepsilon_a\vee\varepsilon_y|-|\varepsilon_a\vee\varepsilon_c|$, where $M:=\varepsilon_u\vee\varepsilon_y=\varepsilon_a\vee\varepsilon_z\vee\varepsilon_y$; hence $|\varepsilon_a\vee\varepsilon_y|-|M|=|\varepsilon_a|-|\varepsilon_u|=m-n-t$. The classes $E_i$ whose $\T(a)$-block lies in a common generalised $K$ lie in a common class of $M$ (a generalised $K$ is a block of $z$ or a union of blocks of $a$ and $z$ in one class of $\varepsilon_a\vee\varepsilon_z$), so the $m$ classes $E_i$ fall into at most $n+t$ classes of $M$, which already accounts for a drop of at least $m-n-t$; hence they fall into exactly $n+t$ classes and no other class is merged: writing $M_j$ for the class of $M$ containing $K_j$ and $M(Z^{(k)})$ for the one containing $Z^{(k)}$, these are $n+t$ distinct classes, $M_j=\bigcup_{i\in I_j}E_i$ and $M(Z^{(k)})=\bigcup_{A_i\subseteq Z^{(k)}}E_i$. Since $\varepsilon_b\le\varepsilon_z\vee\varepsilon_y$, $D_j\subseteq M_j$; in particular a block of $a$ inside $\supp z$ meeting some $S_j$ lies in $D_j\subseteq M_j$.

\smallskip\emph{Step 5: $t=0$, and every $B_i$ lies in some $S_j$.} Suppose that $A_1\subseteq Z:=Z^{(1)}$ (case I), or that $t=0$ and $B_1$ lies in an $\N(v)$-block (case II; then $A_1\subseteq K_j$ for some $j$). Let $O$ be the class of $\Omega$ containing $C_1$; it contains $\Gamma_1$ and $\Delta_1$, and by Step~4 it contains $Z$ and no $K_{j'}$ in case I, and $K_j$ and no other generalised $K$ in case II.

(A) \emph{$D(B_1)$ does not meet $\supp a$, the class $K'$ of $\varepsilon_z\vee\varepsilon_y$ containing $B_1$ is dark, $D(K')=D(B_1)$, and $\Delta_1\subseteq K'$.} In case II the first assertion is the definition of an $\N(v)$-block. In case I, if $D(B_1)$ contained some $A_i$, then either $A_i\subseteq K_{j'}$, so $D(B_1)=D_{j'}$, $D_{j'}\subseteq O$ and $K_{j'}\subseteq O$, which is impossible, or $A_i\subseteq Z^{(k)}$, so $D(B_1)=Z^{(k)}$ (Step~2), which does not contain the $\T(b)$-block $B_1$. So $B_1$ lies in the $\N(v)$-block $D(B_1)\cap\supp b$, the class $K'$ meets $\supp z$ in that $\N(v)$-block, hence is dark (Step~3), and $D(K')=D(B_1)$. Finally $\Delta_1\subseteq K'$: every block of $b$ lies in one class of $\varepsilon_z\vee\varepsilon_y$ and every block of $c$ meeting the dark class $K'$ lies inside it (Step~3), so a chain from $B_1$ in $\Delta_1$ never leaves $K'$. In particular $C_1\subseteq K'$.

(B) \emph{In case I, the class of $\varepsilon_z\vee\varepsilon_y$ containing $Z$ is bright.} $E_1$ contains a $\T(y)$-block $H$. Follow a chain $A_1=P_0,Y_1,P_1,\dots,Y_r=H$ in $E_1$ ($P$'s blocks of $a$, $Y$'s blocks of $y$) point by point and track the class of $\varepsilon_z\vee\varepsilon_y$ of the current point. A step inside a block of $y$, or inside a block of $a$ lying in one block of $z$, keeps the class. The remaining blocks of $a$ on the chain are $\N(a)$-blocks $P\subseteq\supp z$ ($E_1\cap\supp a=A_1\not\subseteq\supp z$); a step inside such a $P$ joins two points of $\supp z$, neither of which lies in a bright class: a bright class meets $\supp z$ in some $S_{j'}$, and then $P\subseteq D_{j'}\subseteq M_{j'}$ (Step~4), while $P\subseteq E_1\subseteq M(Z)\ne M_{j'}$. So the class changes only from a dark class to a dark class; since the chain ends in a bright class (at $H$) it starts in one, and the start is the class of $Z\supseteq A_1$.

(C) \emph{The contradiction.} Take a chain $A_1=P_0,F_1,P_1,\dots,F_\ell,P_\ell$ in $\Gamma_1$ with $F_\ell=C_1$ ($P$'s blocks of $a$, $F$'s blocks of $c$, all inside $O$) and follow it point by point, tracking the class $K(x)$ of $\varepsilon_z\vee\varepsilon_y$ of the current point $x$. We claim that throughout, \emph{$K(x)$ is bright, or dark and disjoint from $\supp z$, or dark, meeting $\supp z$, with $D(K(x))=D_j$} (the third alternative is void in case I). At the start this holds by (B) in case I, and in case II because $x\in A_1\subseteq K_j\subseteq D_j\cap\supp z$, so if $K(x)$ is dark then $K(x)\cap\supp z$ is the block of $v$ containing $x$ and $D(K(x))=D_j$. A step inside a block of $c$ preserves the claim by Step~3 (the class does not change, or both points lie in bright classes). A step inside a block of $a$ lying in one block of $z$ keeps the class. For a step inside a block of $a$, $P\subseteq\supp z$: the current point $x$ lies in $\supp z$, so $K(x)$ is not of the second kind; if $K(x)$ is bright, $x$ lies in some $S_{j'}$, so $P\subseteq D_{j'}\subseteq O$ and $K_{j'}\subseteq O$, impossible in case I and forcing $j'=j$ in case II; if $K(x)$ is dark, $x\in K(x)\cap\supp z\subseteq D(K(x))=D_j$. In both cases $P\subseteq D_j$, so the new point $x'$ lies in $D_j\cap\supp z$ and $K(x')$ is bright or satisfies $D(K(x'))=D_j$. This proves the claim. But the chain ends in $C_1\subseteq K'$, a dark class meeting $\supp z$ (in $B_1$) with $D(K')=D(B_1)$, which does not meet $\supp a$ (by (A)) while $D_j$ does; a contradiction in both cases.

\emph{Consequences.} $t=0$, so $\supp a\subseteq\supp z$ and $\supp u=\supp a$; the $\N(z)$-blocks are exactly the common $\N$-blocks of $a$, $b$ and $z$ (``outer blocks''); every class of $\varepsilon_a\vee\varepsilon_b$ meeting $\supp a$ or $\supp b$ is some $D_j$; an outer block is a class of $\varepsilon_a\vee\varepsilon_b$ by itself, so $D_j\subseteq\supp z$; every class of $\varepsilon_z\vee\varepsilon_y$ meeting $\supp z$ meets it in some $S_j$ (bright; write $K'_j$ for this class) or in an $\N(b)$-block inside $\supp z$ (dark); $A_i\subseteq K_j$ iff $B_i\subseteq S_j$ (both $D(A_i)$ and $D(B_i)$ lie in the class of $\Omega$ containing $C_i$, which contains one $K_j$); the classes of $M$ meeting $\supp a$ are $M_1,\dots,M_n$ with $M_j=\bigcup_{i\in I_j}E_i$ and $M_j\cap\supp a=K_j$; and $K'_j\subseteq M_j$. Finally $\theta_y(u)=\theta_y(z)$ (Step~1) says, by Lemma~\ref{lem:PX}(a), that the sets $M_j\cap\supp y$ and the sets $K'_j\cap\supp y$ are the same family of blocks; as $K'_j\subseteq M_j$, we get $K'_j\cap\supp y=M_j\cap\supp y$: \emph{every $\T(y)$-block inside $M_j$ lies in $K'_j$.}

\smallskip\emph{Step 6: labels.} Suppose, for a contradiction, that $a\not\bowtie b$. As $u\le a$ has rank $n$, this means $n<m$, so some $I_j$ has $|I_j|\ge2$; fix such a $j$. Every point $x$ of $K'_j$ lies in $M_j$, so its block of $a$ lies in a unique $E_i$, $i\in I_j$; put $\sigma(x)=i$ (the \emph{side} of $x$; the block of $y$ containing $x$ lies in the same $E_i$). Let $\gamma(x)\in[m]\cup\{\emptyset\}$ be the index $i$ such that the block of $c$ containing $x$ lies in $\Gamma_i$ ($\emptyset$ if it lies in a class of $\varepsilon_a\vee\varepsilon_c$ without $\T$-blocks); the block of $a$ containing $x$ lies in the same class, so $\gamma$ is constant on blocks of $a$ and of $c$. Let $\lambda(x)$ be the index $k$ such that the block of $c$ containing $x$ lies in $\Delta_k$ ($\emptyset$ if in a $\T$-less class of $\varepsilon_b\vee\varepsilon_c$); the block of $b$ containing $x$ lies in the same class, so $\lambda$ is constant on blocks of $b$ and of $c$, and $\lambda(x)=k$ for $x\in B_k$. The class $K'_j$ is connected by chains alternating blocks of $z$ and blocks of $y$; we record how the label $(\sigma,\gamma,\lambda)$ can vary inside a block of $z$ or of $y$ lying in $K'_j$. An outer block is a block of $a$ and of $b$, so all three labels are constant on it. A $\T(z)$-block $Z\subseteq S_j$ lies in one block of $b$, hence $\lambda$ is constant; it lies in a class of $\varepsilon_a\vee\varepsilon_z$ of type (i) or (iii); in type (iii) it lies in one $\N(a)$-block and $\sigma,\gamma$ are constant; in type (i) it lies in $K_j$, its points lie in the $\T(a)$-blocks $A_i$, $i\in I_j$, and a point of $A_i$ has $\sigma=\gamma=i$: so the labels on $Z$ are of the form $(i,i,k)$ with $k$ fixed. A block $Y$ of $y$ has constant $\sigma$; if $Y$ lies in an $\N(c)$-block then $\gamma$ and $\lambda$ are constant on $Y$; otherwise $Y\subseteq\bigcup C_i$, and a point of $C_i$ has $\gamma=\lambda=i$: the labels on $Y$ are of the form $(s,i,i)$ with $s$ fixed. Hence the labels occurring in $K'_j$ form a connected set in the graph $\mathcal L$ on $[m]\times([m]\cup\{\emptyset\})^2$ whose edges join $(i,i,k)$ to $(i',i',k)$ and $(s,i,i)$ to $(s,i',i')$. The connected components of $\mathcal L$ are the sets $\mathcal C_s=\{(s,i,i):i\in[m]\}\cup\{(i,i,s):i\in[m]\}$, $s\in[m]$, the set $\{(i,i,\emptyset):i\in[m]\}$, and single vertices. Since $K'_j$ contains $B_i$ and $B_{i'}$ for distinct $i,i'\in I_j$, with labels $(\cdot,\cdot,i)$ and $(\cdot,\cdot,i')$, its labels lie in one component $\mathcal C_s$, for a unique $s$; as $s$ is the side of some point, $s\in I_j$; and no point of $K'_j$ has $\gamma=\emptyset$ or $\lambda=\emptyset$.

\smallskip\emph{Step 7: the contradiction.} Take $i\in I_j$, $i\ne s$. A point $x\in B_i$ has $\lambda(x)=i\ne s$, so its label in $\mathcal C_s$ is $(s,i,i)$: \emph{every point of $B_i$ has side $s$ and its block of $c$ lies in $\Gamma_i$}. Choose $x\in B_i$ and let $P$ be its block of $a$; then $P\in\Gamma_i$, $P\subseteq\supp z$ and $P\subseteq E_s$. For a block $Q$ of $a$ write $\hat\sigma(Q)$ for its class of $\varepsilon_a\vee\varepsilon_y$ (an $E_k$ or a class without $\T(a)$-blocks). Take a chain $A_i=P_0,F_1,P_1,\dots,F_\ell,P_\ell=P$ in $\Gamma_i$ ($P_k$ blocks of $a$, $F_k$ blocks of $c$, consecutive ones meeting). We show by induction on $k$ that either $\hat\sigma(P_k)=E_i$, or $P_k$ is an outer block contained in $\Delta_s$; since $P_\ell=P$ lies in $\supp z$ and in $E_s\ne E_i$, this is the desired contradiction. First let $\hat\sigma(P_{k-1})=E_i$ and $x'\in P_{k-1}\cap F_k$, $x''\in P_k\cap F_k$. If $F_k$ is a single block of $y$ (type (iv)), $\hat\sigma(P_k)=E_i$. If $F_k$ is an $\N(c)$-block of type (ii), $x'$ lies in a $\T(y)$-block $H\subseteq F_k$ with $H\subseteq E_i\subseteq M_j$, so $H\subseteq K'_j$ (Step~5) and $x'$ has label $(i,i,\lambda(F_k))\in\mathcal C_s$; as $i\ne s$ this forces $\lambda(F_k)=s$. If $\hat\sigma(P_k)=E_{i''}$ with $i''\in I_j$, the same argument at $x''$ gives the label $(i'',i,s)\in\mathcal C_s$, which forces $i''=i$. Otherwise $P_k\not\subseteq M_j$; but $F_k\cap\supp z\subseteq B_s\subseteq K'_j\subseteq M_j$, so $P_k\cap\supp z=\emptyset$, i.e.\ $P_k$ is an outer block (Step~5), hence a block of $b$ meeting $F_k\subseteq\Delta_s$, so $P_k\subseteq\Delta_s$. If $F_k=C_i$ then, as $C_i\subseteq\Delta_i$ and $i\ne s$, the point $x'$ cannot lie in a $\T(y)$-block of $K'_j$ (its label would be $(i,i,i)\notin\mathcal C_s$), so $C_i$ lies in an $\N(y)$-block (type (iii)), which is a single block of $y$, and $\hat\sigma(P_k)=E_i$. Now let $P_{k-1}$ be an outer block contained in $\Delta_s$; then $F_k\subseteq\Delta_s$. If $P_k$ is outer, it lies in $\Delta_s$. Otherwise $P_k\subseteq\supp z$ meets $F_k\cap\supp z\subseteq B_s\subseteq K'_j$ at a point with label $(\hat\sigma(P_k),i,s)\in\mathcal C_s$ (its $\gamma$-label is $i$ because $F_k\in\Gamma_i$), which forces $\hat\sigma(P_k)=E_i$. This completes the induction. Hence $|I_j|=1$ for every $j$, $n=m$, $u=a$, $v=b$ and $a\bowtie b$.
\end{proof}

The proof uses the $y$-conditions exactly where the computer said they are needed (dropping any one hypothesis produces counterexamples in $\P(\P_4)$, \texttt{red/k4drop.py}): $yc=cy$ confines the blocks of $y$ to single blocks of $c$ or to $\bigcup C_i$ (Step~3), $aya=a$ gives the sides, $\theta_z(y)=\theta_b(a)$ forces the connection of $B_i$ with $B_{i'}$ inside $K'_j$ and the bright/dark dichotomy of Step~3, and $\theta_z(a)=\theta_a(b)$ gives $\theta_y(u)=\theta_y(z)$, i.e.\ that the connection carries all $\T(y)$-blocks of $M_j$; Step~5, which excludes $\T(a)$-blocks outside $\supp z$ and $\T(b)$-blocks outside $\supp v$, needs $aya=a$ as well (through the $\T(y)$-block of $E_1$ in (B)). A first version of the proof (16 Oct, morning) derived Step~5 from a count that is in fact a tautology; the gap was found by the cold reader of 16 Oct, who also proposed the present Step~5 and the \emph{dropped-hypothesis test} that exposed it: for each intermediate statement of a proof and each hypothesis $h$, test the statement on instances satisfying all hypotheses but $h$ (\texttt{review3/review-scripts/dropmap.py}); a statement that fails when $h$ is dropped cannot be proved by an argument that does not use $h$. The special role of the side $s$ is curious: the connection of the $B_i$ through $y$ and $z$ can only be made from one side, and that side cannot see its own $C$'s.

\subsection{K0 holds in every partition monoid}\label{subsec:K0}

\emph{(16 Oct 2026, evening; files \texttt{pa8/k0/}.)} The second order-$8$ condition is shorter to prove than K4, because its $z$-conditions make $z$ a common matched partner of $a$ and $b$, which pins every block of $z$ to a ``cell'', and the $y$-conditions then force the connection of two $\T(a)$-blocks to run through the diagonal cells.

\begin{theorem}\label{thm:K0}
Let $a,b,c,y,z$ be projections of $\P_X$ with $a\bowtie c$, $c\bowtie b$, $yc=cy$, $\theta_y(b)=\theta_b(a)$, $\theta_b(z)=b$, $\theta_z(a)=\theta_z(b)$ and $\theta_y(z)=\theta_a(b)$. Then $a\bowtie b$. Consequently \textup{(E$_1$)} holds in every $\P_X$ (Theorem~\ref{thm:K0E1}, with (W$_2$) from Theorem~\ref{thm:Wn}), $\mathfrak P_8^1$ embeds in no $\P(\P_X)$, and $\mathfrak P_8^1$ is a minimal non-embeddable projection algebra (as in Theorem~\ref{thm:K4}: every proper subalgebra embeds, Section~\ref{sec:comp}). With Theorem~\ref{thm:K4}, all seven candidates of order $8$ are proved minimal non-embeddable.
\end{theorem}

\begin{proof}
We use the notation of Theorem~\ref{thm:R}: $A_i,C_i,B_i$ ($1\le i\le m$) are the $\T$-blocks of $a,c,b$, indexed so that $\Gamma_i\ni A_i,C_i$ are the classes of $\varepsilon_a\vee\varepsilon_c$ meeting $\supp a\cup\supp c$ and $\Delta_i\ni C_i,B_i$ those of $\varepsilon_b\vee\varepsilon_c$ (Lemma~\ref{lem:PX}(c)); in particular every block of $c$ meeting $A_i$ lies in $\Gamma_i$, every block of $b$ meeting $C_i$ lies in $\Delta_i$, and no $\T$-block of one of $a,b,c$ is an $\N$-block of a matched partner. Put $u=\theta_a(b)$, $v=\theta_b(a)$, $n=\rk u$. The blocks of $u$ are the $\N(a)$-blocks and the sets $M\cap\supp a$ for the classes $M$ of $\varepsilon_a\vee\varepsilon_b$ meeting $\supp a$, transversal iff $M$ meets $\supp b$ (Lemma~\ref{lem:PX}(a)); so $a\bowtie b$ iff $n=m$ and every class of $\varepsilon_a\vee\varepsilon_b$ meeting $\supp a$ meets $\supp b$ (Lemma~\ref{lem:PX}(c) and its proof).

\smallskip\emph{Step 1: algebra, and the cells.} $u\le a$, $v\le b$, and $\theta_y(u)=\theta_y(\theta_y(z))=\theta_y(z)=u$, $\theta_y(v)=\theta_y(\theta_y(b))=v$ give $u\le y$, $v\le y$; by \eqref{eq:Q1}, $ay=ya$ and $by=yb$, and then $\theta_y(a)=ya=u$, $\theta_y(b)=yb=v$. Put $p=\theta_z(a)=\theta_z(b)$. Then $b=\theta_b(z)\bowtie\theta_z(b)=p$ (Lemma~\ref{lem:friends}), so $\rk p=m$. The $\T$-blocks of $p=\theta_z(a)$ are the sets $L\cap\supp z$ for the classes $L$ of $\varepsilon_a\vee\varepsilon_z$ meeting both $\supp a$ and $\supp z$; there are $m$ of them, each contains a $\T(a)$-block, and there are $m$ $\T(a)$-blocks; so every class of $\varepsilon_a\vee\varepsilon_z$ meeting $\supp a$ meets $\supp z$ and contains exactly one $\T(a)$-block, which by Lemma~\ref{lem:PX}(a) says $\theta_a(z)=a$ (hence also $a\bowtie p$, by Lemma~\ref{lem:friends}, though we shall only use $\theta_a(z)=a$ and $\theta_b(z)=b$). Let $L_i$ be the class of $\varepsilon_a\vee\varepsilon_z$ containing $A_i$ and $L'_i$ the class of $\varepsilon_b\vee\varepsilon_z$ containing $B_i$; the $L_i$ are distinct, and so are the $L'_i$ (this is $\theta_b(z)=b$). A \emph{cell} is a nonempty set $L\cap L'$ with $L$ a class of $\varepsilon_a\vee\varepsilon_z$ and $L'$ a class of $\varepsilon_b\vee\varepsilon_z$; the cells partition $X$, every block of $z$ lies in a cell, every block of $a$ lies in one class of $\varepsilon_a\vee\varepsilon_z$, and every block of $b$ in one class of $\varepsilon_b\vee\varepsilon_z$. The \emph{diagonal} cells are $D_i=L_i\cap L'_i$ (possibly empty).

\smallskip\emph{Step 2: $y$ against $a,b,c$.} Lemma~\ref{lem:PX}(b) applies to the commuting pairs $(y,a)$, $(y,b)$, $(y,c)$. We claim that for $(y,a)$ and $(y,b)$ there is no class of type (iii): \emph{every $\T(a)$-block and every $\T(b)$-block meets $\supp y$.} Let $N$ be an $\N(y)$-block which is a union of $\T(a)$-blocks, $N\supseteq A_i$. As an $\N(y)$-block, $N$ is a nontransversal block of $\theta_y(b)=v=\theta_b(a)$, i.e.\ an $\N(b)$-block or a set $M\cap\supp b$ for a class $M$ of $\varepsilon_a\vee\varepsilon_b$ missing $\supp a$; the latter is impossible as $N\supseteq A_i$, so $N$ is an $\N(b)$-block. For $(y,c)$, $N$ is of type (iii) or (iv). If $N$ is an $\N(c)$-block, it meets $A_i$, so $N\subseteq\Gamma_i$, and $N\subseteq\supp a$ gives $N=A_i$, a $\T(a)$-block that is an $\N(c)$-block, impossible. So $N$ is a union of $\T(c)$-blocks, and the $\N(b)$-block $N$ contains some $C_k$, impossible. Next let $N$ be an $\N(y)$-block which is a union of $\T(b)$-blocks, $N\supseteq B_k$. It is a nontransversal block of $v=\theta_b(a)$ which is not an $\N(b)$-block, so $N=M\cap\supp b$ for a class $M$ of $\varepsilon_a\vee\varepsilon_b$ missing $\supp a$; thus $N$ contains no $\T(a)$-block, and by what was just proved (no type (iii) for $(y,a)$) $N$ is an $\N(a)$-block. For $(y,c)$, if $N$ were an $\N(c)$-block then $N\subseteq\Delta_k$ and $N=B_k$, impossible; so $N$ is a union of $\T(c)$-blocks. Every $\T(b)$- and every $\T(c)$-block lies in some $\Delta_l$, so $N=\bigsqcup_l(N\cap\Delta_l)$, where $N\cap\Delta_l$ is both a union of $\T(b)$-blocks (hence $B_l$ or empty) and a union of $\T(c)$-blocks (hence $C_l$ or empty): $N$ is a disjoint union of blocks $B_l=C_l$. Being a block of $a$, $N$ lies in one $\Gamma_l$, so $N=B_k=C_k$ is an $\N(a)$-block which is a $\T(c)$-block, impossible.

Consequences. Every $\N(y)$-block is an $\N(a)$-block and an $\N(b)$-block (type (iv) for $(y,a)$ and $(y,b)$); every $\T(y)$-block $Y$ lies in a class of type (i) or (ii), so $Y$ lies in a union of $\T(a)$-blocks or in an $\N(a)$-block, in a union of $\T(b)$-blocks or in an $\N(b)$-block, and in a union of $\T(c)$-blocks or in an $\N(c)$-block; and every $\N$-block of $a$, $b$ or $c$ meeting a $\T(y)$-block contains it (types (ii) and (iv)). The $\T$-blocks of $u=ya$ are the classes $U_1,\dots,U_n$ of type (i) of $\varepsilon_y\vee\varepsilon_a$, each a union of $\T(a)$-blocks and of $\T(y)$-blocks, and its $\N$-blocks are the $\N(a)$-blocks; comparing with the description of $u=\theta_a(b)$ above, the $U_j$ are the sets $M_j\cap\supp a$ for the classes $M_1,\dots,M_n$ of $\varepsilon_a\vee\varepsilon_b$ meeting $\supp a$, and every $M_j$ meets $\supp b$. Put $I_j=\{i:A_i\subseteq U_j\}$; the $I_j$ partition $[m]$, and $a\bowtie b$ iff $n=m$ iff every $|I_j|=1$.

\smallskip\emph{Step 3: the classes of $\varepsilon_y\vee\varepsilon_z$.} By Lemma~\ref{lem:PX}(a) the blocks of $\theta_y(z)$ inside $\supp y$ are the sets $K\cap\supp y$ for the classes $K$ of $\varepsilon_y\vee\varepsilon_z$ meeting $\supp y$, transversal iff $K$ meets $\supp z$; the blocks of $u=ya$ inside $\supp y$ are the $U_j$ (transversal) and the $\N(a)$-blocks contained in $\supp y$ (nontransversal; type (ii)). As $\theta_y(z)=u$, the classes of $\varepsilon_y\vee\varepsilon_z$ meeting $\supp y$ are $K_1,\dots,K_n$ with $K_j\cap\supp y=U_j$ and $K_j\cap\supp z\ne\emptyset$, and, for each $\N(a)$-block $F\subseteq\supp y$, a class $K_F$ with $K_F\cap\supp y=F$. In particular each $A_i$ lies in $K_j$ for $i\in I_j$, and each $\N(a)$-block is an $\N(y)$-block or lies in its $K_F$.

\smallskip\emph{Step 4: blocks of $y$ and cells.} Call a $\T(y)$-block \emph{crossing} if it meets two $\T(a)$-blocks.

(a) \emph{A crossing block $Y$ satisfies $Y\subseteq\bigcup_i(A_i\cap B_i\cap C_i)$, the union over the $i$ with $Y\cap A_i\ne\emptyset$; in particular $Y$ lies in the union of the diagonal cells.} $Y$ meets $\supp a$, so it lies in a type (i) class $U_j$ and $Y\subseteq\supp a$. $Y$ is not inside an $\N(c)$-block $F'$: such an $F'$ would meet $A_i$ and $A_{i'}$ for two $i\ne i'$, so $F'\subseteq\Gamma_i\cap\Gamma_{i'}$. So $Y\subseteq\bigcup_kC_k$; a point of $Y\cap A_i$ lies in some $C_k\subseteq\Gamma_k$ and in $A_i\subseteq\Gamma_i$, so $k=i$: $Y\cap A_i\subseteq C_i$. A point $x\in Y\cap C_i$ lies in a block of $b$ inside $\Delta_i$, namely $B_i$ or an $\N(b)$-block $G\subseteq\Delta_i$; in the second case $G\supseteq Y$, but $Y\cap A_{i'}\subseteq C_{i'}\subseteq\Delta_{i'}$ for a second index $i'$. So $Y\cap A_i\subseteq A_i\cap B_i\cap C_i$, and a point of $A_i\cap B_i$ lies in $L_i\cap L'_i=D_i$.

(b) \emph{A non-crossing $\T(y)$-block $Y$ lies in one cell.} If $Y$ meets no $\T(a)$-block then $Y$ lies in an $\N(a)$-block, and if it meets exactly one, $A_i$, then it lies in $U_j\subseteq\supp a$ and $Y\subseteq A_i$; either way $Y$ lies in one block of $a$, hence in one class of $\varepsilon_a\vee\varepsilon_z$ and in one $\Gamma$-class or $\T$-less class of $\varepsilon_a\vee\varepsilon_c$. For $\varepsilon_b\vee\varepsilon_z$: if $Y$ lies in an $\N(b)$-block we are done; otherwise $Y\subseteq\bigcup_kB_k$, and we claim $Y$ meets only one $B_k$. Suppose $Y$ meets $B_k$ and $B_l$, $k\ne l$. $Y$ is not inside an $\N(c)$-block (it would lie in $\Delta_k\cap\Delta_l$), so $Y\subseteq\bigcup C$'s, and as in (a) $Y\cap B_k\subseteq C_k\subseteq\Gamma_k$, $Y\cap B_l\subseteq C_l\subseteq\Gamma_l$, contradicting the fact that $Y$ lies in one class of $\varepsilon_a\vee\varepsilon_c$. So $Y\subseteq B_k\subseteq L'_k$.

(c) \emph{Every $\N(y)$-block and every block of $z$ lies in one cell}: an $\N(y)$-block is a block of $a$ and of $b$ (Step~2).

(d) \emph{Every class $K$ of $\varepsilon_y\vee\varepsilon_z$ containing no crossing block lies in one cell, and every class containing a crossing block lies in $\bigcup_iD_i$.} $K$ is connected by chains of blocks of $y$ and $z$ in which consecutive blocks meet. By (b), (c) a non-crossing block lies in one cell, so it does not leave the cell of the block before it; by (a) a crossing block lies in $\bigcup D_i$, and a non-crossing block meeting it lies in some $D_i$. Starting the chains at a crossing block, resp.\ noting that without crossing blocks no chain changes cell, gives the claim.

\smallskip\emph{Step 5: $\varepsilon_a\le\varepsilon_b\vee\varepsilon_z$.} Let $A_i\subseteq U_j\subseteq K_j$. If $K_j$ contains no crossing block, $A_i$ lies in one cell. Otherwise $K_j\subseteq\bigcup_{i'}D_{i'}$, and $A_i\subseteq L_i$ with $D_{i'}\subseteq L_{i'}$ and the $L$'s distinct, so $A_i\subseteq D_i\subseteq L'_i$. An $\N(a)$-block $F\subseteq\supp y$ lies in $K_F$, which meets $\supp a$ in nothing (its intersection with $\supp y$ is $F$) and so contains no crossing block (crossing blocks lie in $\supp a$); so $F$ lies in one cell. An $\N(a)$-block not inside $\supp y$ is an $\N(y)$-block, in one cell by Step~4(c). So every block of $a$ lies in one class of $\varepsilon_b\vee\varepsilon_z$, and therefore $\varepsilon_a\vee\varepsilon_b\le\varepsilon_b\vee\varepsilon_z$.

\smallskip\emph{Step 6: the contradiction.} Suppose $n<m$, so that some $I_j$ contains two indices, say $1,2\in I_j$. The class $U_j$ of type (i) is connected by chains alternating $\T(a)$- and $\T(y)$-blocks, so some $\T(y)$-block of $U_j$ meets $A_1$ and $A_2$: $K_j$ contains a crossing block, and by Step~5, $A_1\subseteq L'_1$ and $A_2\subseteq L'_2$, two distinct classes of $\varepsilon_b\vee\varepsilon_z$. But $A_1,A_2\subseteq M_j$, a single class of $\varepsilon_a\vee\varepsilon_b$, which by Step~5 lies in one class of $\varepsilon_b\vee\varepsilon_z$. Hence $n=m$, every class of $\varepsilon_a\vee\varepsilon_b$ meeting $\supp a$ meets $\supp b$ (Step~2), and $a\bowtie b$.
\end{proof}

Every step was tested by the dropped-hypothesis map before being written down (\texttt{pa8/k0/k0claims.py}: all instances of $\P(\P_4)$, and a sample of $\P(\P_5)$; the statements are: $\theta_a(z)=a$; $\N(y)$-blocks are blocks of $a,b,c$; crossing blocks lie in $\bigcup(A_i\cap B_i\cap C_i)$; non-crossing $\T(y)$-blocks and $\N(y)$-blocks lie in one cell; classes of $\varepsilon_y\vee\varepsilon_z$ lie in one cell or in the diagonal; $\varepsilon_a\le\varepsilon_b\vee\varepsilon_z$). The map shows where each hypothesis enters: $\theta_b(z)=b$ and $\theta_z(a)=\theta_z(b)$ are used only through Step~1, i.e.\ through $\theta_a(z)=a$ and $\theta_b(z)=b$ (the $L_i$ distinct, the $L'_i$ distinct) --- so the proof gives also the condition with $\theta_z(a)=\theta_z(b)$ replaced by $\theta_a(z)=a$; $\theta_y(z)=\theta_a(b)$ gives $y\sim a$, $y\sim b$ (Step~1) and the classes $K_j$, $K_F$ (Step~3), without which Step~4(a) and Step~5 fail; $yc=cy$ and $\theta_y(b)=\theta_b(a)$ carry Step~2 and the confinement of the blocks of $y$ in Step~4 (without either, non-crossing blocks leave their cells). The proof is shorter than that of K4 because here $z$ is matched with both $a$ and $b$ (Step~1), so that its blocks are confined to cells from the start, whereas in K4 the roles of $y$ and $z$ are mixed ($z$ sees $y$ as $v$ and $a$ as $u$) and the confinement has to be established along chains (Step~7 there). Remarkably, the whole argument never uses the $\T$-blocks of $z$: only the classes of $\varepsilon_a\vee\varepsilon_z$ and $\varepsilon_b\vee\varepsilon_z$ enter, through the cells.

\section{Computations}\label{sec:comp}

All regular $*$-semigroups of order $n\le8$ were enumerated up to $*$-isomorphism by a backtracking search on multiplication tables with propagation of associativity, $(xy)^*=y^*x^*$ and $xx^*x=x$; for $n=7,8$ only the semigroups whose involution is not the identity were generated, since $x^*=x$ for all $x$ forces commutativity and hence an inverse (Clifford) semigroup, which embeds by Wagner--Preston. For each non-inverse semigroup satisfying \eqref{eq:Q1}, a $*$-embedding into $\P_n$ was searched for, $n\le5$, by assigning images to projections first and forcing the images of all products.

\begin{center}
\begin{tabular}{@{}rrrrrr@{}}
\toprule
order & classes & inverse & non-inverse, \eqref{eq:Q1} & non-inverse, not \eqref{eq:Q1} & of these, in $\P_5$\\
\midrule
4 & 17 & 16 & 1 & 0 & 1\\
5 & 55 & 52 & 2 & 1 & 2\\
6 & 222 & 208 & 9 & 5 & 9\\
7 & $259^{\dagger}$ & 196 & 35 & 28 & 35\\
8 & $1310^{\dagger}$ & 970 & 178 & 162 & 178\\
\bottomrule
\end{tabular}

\smallskip
{\small $^\dagger$ nontrivial involution only. The last column counts the non-inverse semigroups satisfying \eqref{eq:Q1} that embed in $\P_5$.}
\end{center}

Thus, up to order $8$, a regular $*$-semigroup embeds in a partition monoid if and only if it satisfies \eqref{eq:Q1} (the ones failing \eqref{eq:Q1} do not embed, the others embed in $\P_5$, inverse ones by Wagner--Preston). The smallest known counterexample to sufficiency of \eqref{eq:Q1} is the $61$-element $Q_0$ of Theorem~\ref{thm:witness}. The smallest known counterexample to sufficiency of \eqref{eq:Q1} together with the triangle condition is the $19$-element $\mathfrak T_{19}$ of Theorem~\ref{thm:R} (before it, the $20$-element $\mathfrak Q$ of Definition~\ref{def:Qpath}); so the smallest counterexample has order between $9$ and $19$ (the order-$9$ enumeration, some $10^7$ tables, was not completed).

The ten $*$-congruences of $\P_3$ have quotients of orders $203$, $179$, $103$, $79$, $67$, $43$, $7$, $3$, $2$, $1$. The quotients of orders $179,79,43$ have a zero and are not inverse, so they fail \eqref{eq:Q1}; those of orders $7,3,2,1$ are inverse; the two remaining proper quotients, of orders $103$ and $67$, satisfy \eqref{eq:Q1} and are the semigroups of Theorem~\ref{thm:witness}. For $\P_2$ all proper quotients satisfying \eqref{eq:Q1} embed. For $\P_4$ the classification of East, Mitchell, Ru\v skuc and Torpey (Adv.\ Math.\ 2018) leaves only two proper $*$-congruences whose quotients have no zero (a quotient with a zero satisfying \eqref{eq:Q1} is inverse, as every idempotent lies above the zero); they are $\mu$, which identifies each partition of rank $\le1$ with the rank-$0$ partition obtained by cutting its transversal block, and $\mu\vee\lambda_{S_2}$, which moreover identifies the two elements of each rank-$2$ $\Hh$-class. We verified directly that these are $*$-congruences of $\P_4$ (with $2771$ and $1810$ classes) and that they are the only ones contained in the second. Both quotients satisfy \eqref{eq:Q1} and fail the triangle condition, as they must: they contain the corresponding quotients of $\P_3$ (add the singletons $\{4\},\{4'\}$). So the $*$-congruences of $\P_4$ give nothing new.

\emph{Projection algebras.} Mace4 with isofilter enumerated all projection algebras (axioms (P1)$'$--(P5)$'$) satisfying \eqref{eq:Q1} and the triangle clause, up to isomorphism: there are $2,5,16,62$ of orders $2,3,4,5$. Each was tested for embeddability into $\P(\P_m)$, $m\le5$, by a backtracking search which forces the image of every product of already assigned elements; all embed, in $\P(\P_2)$, $\P(\P_3)$ or $\P(\P_4)$ ($3$, $39$, $20$ of the $62$ of order $5$). Of order $6$ there are exactly $296$ (from $79332$ Mace4 models before isofilter). All but one embed in $\P(\P_m)$ with $m\le5$: one in $\P(\P_2)$, $90$ in $\P(\P_3)$, $173$ in $\P(\P_4)$ and $31$ in $\P(\P_5)$; the exception is $\mathfrak P$, which is also the only one of the $296$ violating the wiring condition (S). So up to order $6$ the projection algebras that embed in some $\P(\P_X)$ are exactly those satisfying \eqref{eq:Q1}, the triangle condition and (S). Of order $9$, $\mathfrak P_9$ (Example~\ref{ex:P9}) satisfies \eqref{eq:Q1}, the triangle condition and the quasi-identity of Theorem~\ref{thm:wiring} but not (S), and does not embed in $\P(\P_m)$ for $m\le8$ (SAT). ($\mathfrak P$ does not embed in $\P(\P_m)$ for $m\le6$ by exhaustive search, nor for $m\le8$ by a SAT encoding of the join and product structure solved by CaDiCaL, $m=8$ taking about $25$ minutes, in agreement with Theorem~\ref{thm:wiring}.) Of order $7$ there are $1\,663$ projection algebras satisfying \eqref{eq:Q1} and (S) (Section~\ref{subsec:fourth}; $295$ of order $6$); all satisfy the triangle condition, and all but $\mathfrak P_7$ embed in $\P(\P_m)$ with $m\le6$ ($161$ in $\P(\P_3)$, $924$ in $\P(\P_4)$, $531$ in $\P(\P_5)$, $46$ in $\P(\P_6)$), while $\mathfrak P_7$ does not embed at all (Theorem~\ref{thm:R}; the SAT search confirms this for $m\le8$). Of order $8$ there are $10\,828$ projection algebras satisfying \eqref{eq:Q1}, (S) and (R) (Lemma~\ref{lem:extension}); all satisfy the triangle condition, $10\,821$ embed in $\P(\P_m)$ with $m\le7$, and the seven of Section~\ref{subsec:census8} do not embed for $m\le8$ (whether they embed for larger $m$ is open). Of order $9$ there are $80\,039$, of which seven embed in no $\P(\P_m)$ with $m\le8$ and contain none of the seven of order $8$ (Remark~\ref{rem:census8}(3)). The graph semigroups $Q(G,\bot)$ with $|V|\le3$ and $\bot$ a proper colouring all embed in $\P_4$ (and Theorems~\ref{thm:graphsg}, \ref{thm:layered} settle all graph and layered graph semigroups).

Further checks: several completely simple semigroups with non-flat sandwich matrices (an $\mathcal M[\Z_2;3,3;P]$ embeds in $\P_4$; an $\mathcal M[\Z_3;3,3;P]$ and the $48$-element $\mathcal M[\Z_3;4,4;P]$ of the previous paper do not embed in $\P_5$ but embed in $\P_{27}$ and $\P_{48}$ by Theorem~\ref{thm:cs}); strong semilattices of Rees matrix semigroups with injective-on-projections structure maps (all embedded); $H\cup(A\times A)$ with various actions (all embedded).

\section{General remarks and questions}\label{sec:general}

\begin{proposition}\label{prop:finite}
A finite regular $*$-semigroup embeds in some partition monoid if and only if it satisfies every quasi-identity (in the language $\{\cdot,{}^*\}$) that holds in all partition monoids.
\end{proposition}

\begin{proof}
Necessity is clear. Conversely, a finite $S$ satisfying all quasi-identities of the class $\{\P_X\}$ lies in the quasivariety it generates, so $S$ embeds in a product $\prod_j A_j$ of ultraproducts $A_j$ of partition monoids $\P_{X_{ij}}$, $i\in I_j$, over ultrafilters $\mathcal U_j$. As $S$ is finite, finitely many factors suffice, and by \L o\'s's theorem each $*$-homomorphism $S\to A_j$ is, for $\mathcal U_j$-almost all $i$, induced by a $*$-homomorphism $S\to\P_{X_{ij}}$ with the same kernel. Taking one such $i$ for each of the finitely many $j$ and using $\P_X\times\P_Y\le\P_{X\sqcup Y}$ gives an embedding of $S$ into a partition monoid.
\end{proof}

So for finite semigroups the programme of finding quasi-identities is complete in principle; whether finitely many suffice is unclear. We have three pairwise independent ones: \eqref{eq:Q1}, the triangle condition and the wiring condition (S) of Theorem~\ref{thm:S} (independence: Section~\ref{sec:wiring}), and a fourth one, the condition (R) of Theorem~\ref{thm:R}, independent of the other three; by Theorem~\ref{thm:master} all the other matching-pattern conditions of Section~\ref{sec:hex} follow from the triangle condition, so they add nothing; up to order $6$ the first three describe the projection algebras embeddable in some $\P(\P_X)$ exactly, and up to order $7$ the four do (Section~\ref{sec:comp}). Whether the class of all (possibly infinite) embeddable regular $*$-semigroups is closed under ultraproducts, i.e.\ is a universal class, I do not know; the natural map from an ultraproduct of partition monoids to a partition monoid is not a homomorphism, because connectivity through the middle layer is not a first-order property.

\begin{proposition}\label{prop:relations}
For every set $X$, the map $\alpha\mapsto R_\alpha$, where
\[
R_\alpha=\{(A,B)\in 2^X\times2^X: A\sqcup B'\text{ is a union of blocks of }\alpha\},
\]
is an injective homomorphism of involuted semigroups from $\P_X$ into the semigroup of binary relations on $2^X$ with converse as involution.
\end{proposition}

\begin{proof}
$(A,C)\in R_\alpha\circ R_\beta$ iff there is $B$ which is a union of lower blocks of $\alpha$ and of upper blocks of $\beta$, i.e.\ a union of classes of the middle join, with $A\sqcup B'$ and $B\sqcup C'$ unions of blocks; this is exactly the condition for $A\sqcup C'$ to be a union of blocks of $\alpha\beta$. Injectivity: the minimal nonempty $A$ with $(A,\emptyset)\in R_\alpha$ are the upper nontransversal blocks, and the pairs $(A,B)$ with $A,B$ minimal and nonempty are the transversal blocks.
\end{proof}

Consequently every $*$-subsemigroup of a partition monoid is representable by binary relations with converse, so Schein's quasi-identities for representable involuted semigroups are necessary conditions. They seem, however, to carry no information here: every regular $*$-semigroup of order $\le5$ is representable by relations on at most $5$ points (checked by computer), including the $5$-element band with zero which fails \eqref{eq:Q1} (represent the $2\times2$ band by $A\times A$, $B\times B$ with $A\cap B\ne\emptyset$ and the zero by the empty relation). The point of $\P_X$ is that its relations $R_\alpha$ are \emph{Boolean subalgebras} of $2^X\times2^X$ containing $(X,X)$; conversely every such subalgebra comes from a partition. Thus $S$ embeds in $\P_X$ iff $S$ has a faithful representation by relations on the Boolean algebra $2^X$ each of which is a subalgebra of $2^X\times 2^X$.

\begin{question}
\begin{enumerate}[label=\textup{(\arabic*)}]
\item Is embeddability of a finite regular $*$-semigroup determined by its projection algebra? All known obstructions, \eqref{eq:Q1}, the triangle condition, (S), (R) and those of Proposition~\ref{prop:census8}, are projection-algebra quasi-identities, and Theorem~\ref{thm:cs} shows that groups and sandwich matrices impose nothing; but a regular $*$-semigroup is far from determined by its projection algebra.
\item \eqref{eq:Q1} and the triangle condition do \emph{not} suffice, even together with all matching-pattern conditions (Corollary~\ref{cor:wiring}). What is the smallest counterexample (order $9$ to $19$)? Is there a common generalisation of (R$'$) and (V), or of (R$'$) and the remaining path-closing conditions (P$_1$), (Q), with a uniform proof? (The hub case of Theorem~\ref{thm:S} now has a computer-free proof, Corollary~\ref{cor:Q}.) Are there further wiring-type obstructions, e.g.\ with longer paths of matchings or with two commuting pairs? For graph semigroups and for layered graph semigroups the answer is complete (Theorems~\ref{thm:graphsg} and \ref{thm:layered}): \eqref{eq:Q1} and (S) suffice, and the triangle condition is automatic; so new obstructions have to be sought in families whose projections have non-chain down-sets or large $\T$-blocks (Remark~\ref{rem:layered}(2)). More generally, is the quasivariety generated by the projection algebras $\P(\P_X)$ finitely based? Its members of order $\le6$ are exactly the algebras satisfying \eqref{eq:Q1}, the triangle condition and (S), and (S) is needed only once, for $\mathfrak P$; of order $7$ one further condition, (R), is needed, once ($\mathfrak P_7$); of order $8$ seven further candidates appear (Section~\ref{subsec:census8}), and the quasi-identities they violate, valid in $\P(\P_m)$ for $m\le7$, are of at least three different shapes, so the list of obstructions is not closing up quickly. (V), (R$'$), (P$_1$) and (Q) are proved (Theorems~\ref{thm:V}, \ref{thm:Rp}, \ref{thm:P1}, \ref{thm:Q}); is (E$_1$) of Proposition~\ref{prop:census8} valid in every $\P_X$ ((E$_6$) is, Theorem~\ref{thm:K4})? (They are independent of each other and of the rest, Proposition~\ref{prop:indep}.)
\item Is there a conceptual description of the quasi-identities of Section~\ref{sec:hex}, e.g.\ in terms of ``two commuting families with a common meet-preserving matching must have the same meet''? Theorem~\ref{prop:path} says that for finite $X$ the pattern of matchings is irrelevant as long as it is not a star. The proof of Theorem~\ref{thm:hex} suggests a general principle: in $\P_X$, a projection is determined by its behaviour outside its support together with its $\T$-blocks, and matched pairs cannot share a set which is transversal for one and nontransversal for the other.
\item Is the condition K5 of Section~\ref{subsec:nine} valid in every $\P_X$? ((W$_2$), (R$''$), K6, K4 and K0 are: Theorems~\ref{thm:Wn}, \ref{thm:Rpp}, \ref{thm:K4}, \ref{thm:K0}; K4 gives (E$_6$) and K0 gives (E$_1$), Theorem~\ref{thm:K0E1}, so all candidates of order $8$ are settled.) The chain families (W$_n$) and (R$'_n$) of Sections~\ref{subsec:chain}, \ref{subsec:rpfam} are genuinely infinite, with the explicit witnesses $\mathfrak P^{\mathrm{ch}}_n$, $\mathfrak P^{\mathrm R'}_n$ (Sections~\ref{subsec:chainwit}, \ref{subsec:rpfam}), provided (R$'_n$) is valid in every $\P_X$ ((W$_n$) is, Theorem~\ref{thm:Wn}; (R$'_n$) is verified for small $m$ only; the maximal subalgebras of the witnesses do embed, Proposition~\ref{prop:maxsub}): is it? What explains the pattern of valid identification words along a chain? Their witnesses would be the first minimal non-embeddable projection algebras of order $9$ beyond the (R$'$)-witness. More importantly: is there a combinatorial criterion deciding, for a configuration of auxiliaries attached to an unclosed path $a\bowtie c\bowtie b$ (a ``wiring diagram'': commutations, half-matchings, equalities of transports of restrictions), whether it forces $a\bowtie b$, with one uniform proof? The known proofs ((R), (R$'$), (P$_1$), (Q)) all run by transport along the matchings and propagation of an invariant along a join, which suggests that the criterion exists; the growth of the list (one auxiliary at order $7$, two at order $8$, two or three plus the terms $u,v$ at order $9$) suggests that the family is infinite, so that no finite basis exists and a parametrised one is the best possible. Related: is embeddability of a finite projection algebra in some $\P(\P_X)$ decidable? A recursive basis of quasi-identities makes non-embeddability semi-decidable; embeddability itself is semi-decidable provided a \emph{compression lemma} holds: \emph{if a finite $P$ embeds in some $\P(\P_X)$, then it embeds in $\P(\P_m)$ for some finite $m$ (say $m\le f(|P|)$)}. We have not proved this; it is not obvious, since restriction to a finite subset of $X$ does not preserve products (connectivity through the middle layer can leave the subset), and $\P_X$ for infinite $X$ is not an ultraproduct of the finite $\P_m$. The census data ($m\le n-1$ suffices for all embeddable algebras of order $n\le9$) suggest that it is true, and a bound $f$ would turn the ``verified for $m\le7$'' evidence of the status table into proofs. Hand proofs are unaffected: Lemma~\ref{lem:PX} holds for every $X$.
\item A cheap necessary condition for a counterexample to Question~(1) (a regular $*$-semigroup $S$ with $P(S)$ embeddable but $S$ not): if $p\bowtie q$ with $p\ne q$ then $pq$ and $qp$ are distinct non-projections determined by $\{p,q\}$ (as $(pq)(pq)^*=pqp=p$ and $(pq)^*(pq)=q$; if $pq$ were a projection then $pq=qp$ and $p=pqp=qp\cdot p=qp=pq$, whence $p=q$), so $|S|\ge|P(S)|+2\cdot\#\{\{p,q\}:p\bowtie q,\ p\ne q\}$. For $\mathfrak P$ ($4$ matched pairs) this gives $14$ ($|\mathfrak Q|=20$), for $\mathfrak P_7$ ($5$ pairs) $17$ ($|\mathfrak T_{19}|=19$), for the order-$8$ witnesses at least $20$ and for the order-$9$ candidates at least $21$ (the bottom rectangular bands contribute many matched pairs). Since matched pairs of a subalgebra are matched pairs of the algebra, every regular $*$-semigroup whose projection algebra contains one of the known minimal non-embeddable algebras has order at least $14$. So a counterexample in Question~(2) of order $\le13$ would have an embeddable projection algebra, i.e.\ would answer Question~(1) negatively; this makes the orders $9$ to $13$ (regular $*$-semigroups satisfying \eqref{eq:Q1} with embeddable projection algebra) the natural range for an experiment on Question~(1).
\item The infinite case. By Theorem~\ref{thm:master} the triangle condition implies every non-star condition in every projection algebra, and by Theorem~\ref{thm:hinf} the hexagon condition does not imply the triangle condition; so for infinite regular $*$-semigroups the triangle condition is the matching-pattern condition to use; but by Section~\ref{sec:wiring} the wiring conditions are needed as well, and their infinite-case behaviour is unexplored. Does the hexagon condition imply the conditions on $k\ge4$ indices in general (Remark~\ref{rem:hinf}(1))? Which pattern conditions are equivalent to the triangle condition in the variety of projection algebras (the path and the $2$-cycle are; the hexagon is not)?
\end{enumerate}
\end{question}

\section*{Files}
The directory \texttt{embed/} contains the code: \texttt{pmon.py} (partition monoids), \texttt{enum\_rss.py} (enumeration of regular $*$-semigroups), \texttt{embed.py} (embedding search), \texttt{analyse.py}, \texttt{congr.py} (congruences and quotients), \texttt{qsearch.py}, \texttt{hexsearch*.py}, \texttt{hexmin*.py}, \texttt{hexqi.py}, \texttt{qi.py} (searches for the hexagon configuration and quasi-identity tests), \texttt{hexagon.in}, \texttt{triangle.in}, \texttt{octagon.in} (Prover9/Mace4 inputs; the projection algebra is written with the binary operation $X*Y=\theta_Y(X)$, a device for Prover9 and Mace4, rather than with the unary operations $\theta_Y$ of East and Azeef Muhammed, arXiv:2301.04845; see Section~2), \texttt{digraph.py} (clauses and $\P_n$-checks for the matching-pattern conditions), \texttt{oct14.py}, \texttt{p4mu.py}, \texttt{octmin.py}, \texttt{octsub6.py} (the octagon experiments of Remark~\ref{rem:2k}), \texttt{invariant.in}, \texttt{stab.in}, \texttt{period.in}, \texttt{uniform.in}, \texttt{derived.in} (Prover9/Mace4 inputs behind Section~\ref{subsec:equiv}), \texttt{mace4parse.py}, \texttt{walk\_hex.py}, \texttt{uniform\_hex.py} (analysis of Mace4 models), \texttt{path.in}, \texttt{path2.in}, \texttt{square.in}, \texttt{square2.in}, \texttt{disjoint.in}, \texttt{disjoint2.in}, \texttt{disjoint3.in}, \texttt{sqfilter2.py}, \texttt{sqstar.py}, \texttt{digraph3.py}, \texttt{digraph4.py}, \texttt{nofull.py}, \texttt{nofull3.py}, \texttt{p4mu\_disjoint.py} (Theorem~\ref{prop:path} and Remark~\ref{rem:path}); \texttt{pa\_enum.py}, \texttt{pa\_embed2.py}, \texttt{pa\_sat.py} (enumeration of projection algebras with \eqref{eq:Q1} and the triangle clause, embedding into $\P(\P_m)$ by backtracking and by SAT), \texttt{graphsg.py}, \texttt{graphsg\_survey.py}, \texttt{graphsg\_thm.py} (the semigroups $Q(G,\bot)$, the survey of Remark~\ref{rem:graphthm}(2) and the verification of the construction of Theorem~\ref{thm:graphsg}), \texttt{layered.py}, \texttt{layered\_search.py}, \texttt{layered\_search2.py}, \texttt{layered\_thm.py} (layered graph semigroups, Section~\ref{subsec:layered}), \texttt{pa7\_q1S.in}, \texttt{pa7\_driver.sh}, \texttt{pa7\_post.py}, \texttt{fourth.py}, \texttt{fourth\_pre.py} (the order-$7$ census, $\mathfrak P_7$, $\mathfrak T_{19}$, $\mathfrak T_{26}$ and the tests of (R), Section~\ref{subsec:fourth}), \texttt{pathclose.py}, \texttt{pathclose\_np.py}, \texttt{pathclose4.py}, \texttt{pathclose5.py}, \texttt{pathclose\_check.py}, \texttt{pc\_prover.py}, \texttt{pc\_mace\_in.py}, \texttt{pc\_mace.sh}, \texttt{pathclose2.py}, \texttt{pathclose2r.py}, \texttt{pc2\_mace.py} (the search for path-closing conditions, Section~\ref{subsec:pathclose}), \texttt{pa8/extend.py}, \texttt{pa8/post8.py}, \texttt{pa8/cand8.py}, \texttt{pa8/cand8b.py}, \texttt{pa8/qi\_extract.py}, \texttt{pa8/qi\_all.py}, \texttt{pathclose3.py}, \texttt{pa8/pc3\_hit.py}, \texttt{pa8/pc3\_verify.py}, \texttt{pa8/pc1\_classify.py}, \texttt{pa8/verify6.py}, \texttt{pa8/verify7.py}, \texttt{pa8/sat8.py}, \texttt{pa8/newconds.py}, \texttt{pa8/pass2\_9.py}, \texttt{pa8/indep.py}, \texttt{pa8/run\_indep.py}, \texttt{pa8/vexp.py} (the experiments behind the proof of Theorem~\ref{thm:V}), \texttt{pa8/p9/} (Prover9 inputs; \texttt{pa8/p9/conditions.in} collects the projection algebra axioms and all conditions of Section~\ref{subsec:indep} with labels), \texttt{pa8/p1/} (instances of (P$_1$) in $\P(\P_4)$ and $\P(\P_5)$ and the block-level tests behind Theorem~\ref{thm:P1}), \texttt{pa8/q/}, \texttt{pa8/gen/} (the reconnaissance of (Q) and the search for a common generalisation, Remark~\ref{rem:census8}(5)), \texttt{pa8/rp/} (the experiments behind the proof of Theorem~\ref{thm:Rp}: \texttt{rp\_inst.py} enumerates the instances of the hypotheses of (R$'$) in $\P(\P_n)$, \texttt{rp\_words.py}, \texttt{rp\_min.py}, \texttt{rp\_rel.py} the identities and relations they satisfy, \texttt{harness*.py} the block-level statements, \texttt{pa\_*.in}, \texttt{g\_*.in} the Prover9 inputs for Step~1), and the census files \texttt{pa8/ord\emph{n}.json}, \texttt{pa8/post8.log} (the census by extension, the seven algebras of order $8$ and the quasi-identities of Section~\ref{subsec:census8}), \texttt{m172*.py}, \texttt{qi2.py}, \texttt{wiringcheck.py} (tests of the wiring conditions in $\P_n$ and in given projection algebras), \texttt{hubtypes.py}, \texttt{hubcheck\_rand.py} (the finite verification in the proof of Theorem~\ref{thm:S}), \texttt{wiring\_*.in} (Prover9/Mace4 inputs relating (S), Theorem~\ref{thm:wiring}, \eqref{eq:Q1} and the triangle clause; \texttt{mace\_WnotS9.out} is $\mathfrak P_9$), \texttt{p4cong.py}, \texttt{p4cong2.py} (the $*$-congruences of $\P_4$) (Section~\ref{sec:wiring}; the Mace4 models are the files \texttt{mace\_*.out}), \texttt{hinf2.py}, \texttt{hinf2\_run.py}, \texttt{hinf2\_pat.py}, \texttt{chain.py} (the unrolled hexagon of Theorem~\ref{thm:hinf} and the chain models), \texttt{periodic.py}, \texttt{periodic3.py} (periodic projections of $\P_{\Z\times[m]}$ and the pattern searches of Remark~\ref{rem:hinf}(3)), \texttt{cs\_embed.py} (the construction of Theorem~\ref{thm:cs}), \texttt{normalband*.py}, \texttt{strongsl.py}, \texttt{examples*.py} (further test families), \texttt{relrep.py} (representations by binary relations), and the pickled enumerations \texttt{rss\_*.pkl}. The directory \texttt{pa8/n4/} contains the code of Section~\ref{subsec:nine} (\texttt{README.md} describes it: \texttt{pathclose4b.py}, \texttt{hit4.c}, \texttt{verify\_hits.py}, \texttt{qi.py}, \texttt{explain2.py}, \texttt{explain3b.py}, \texttt{catalog.py}, \texttt{checkK.py}, \texttt{min9.py}, \texttt{mace\_gen.py}; logs \texttt{catalog.log}, \texttt{explain3b\_all.log}). The directory \texttt{pa8/rpp/} contains the reconnaissance of (R$''$) and (W$_n$) of Section~\ref{subsec:chain} (\texttt{rpp\_inst.py}, \texttt{rpp\_inst5.py}, \texttt{words.py}, \texttt{rel.py}, \texttt{harness*.py}, \texttt{kinds2.py}, \texttt{walk.py}, \texttt{ci2.py}, \texttt{chain3.py}, \texttt{wn\_sat.py}, the Prover9 inputs \texttt{h*.in}, \texttt{sc1.in}, \texttt{v3p.in}, \texttt{v4p.in}, \texttt{red.in} and the Mace4 inputs \texttt{w3ind.in}, \texttt{w2\_rpp.in}, \texttt{rpp\_w2.in}); \texttt{pa8/n4/explain3c.py} and \texttt{explain3c.log} are the 2-path check of the remaining witnesses. The directory \texttt{pa8/q2/} contains the instance generators, the harness \texttt{hq.py}, the claim files \texttt{t*.py} (earlier reconnaissance) and \texttt{h1.py}--\texttt{h4.py} (the statements of Steps~2--7 of Section~\ref{subsec:Q}: cross wiring, the claims of Proposition~\ref{prop:Qh}, $\Pi_1=\varepsilon_u\vee\varepsilon_d$, and the contact rules), and \texttt{NOTES\_Q.md}.

\end{document}
